\chapter{Complexity — Catalog}
%%% generated by the blueprint pipeline from TCSlib.Complexity.TuringMachine.Build.Catalog
\section{Overview}

Declarations of the Lean module \texttt{TCSlib.Complexity.TuringMachine.Build.Catalog}.

\section{Declarations}

\begin{definition}[Sweep Phases]
\lean{Turing.SweepPhase}
\label{Turing.SweepPhase}
\leanok
A sweep phase is one of three states: a forward sweep over the stored word, a rewind
back to the origin, or a live exit state.
\end{definition}

\begin{definition}[Flagged Computation Phases]
\lean{Turing.FlagPhase}
\label{Turing.FlagPhase}
\leanok
A flagged computation phase is either a forward working pass, a return pass carrying a
Boolean flag, or a live exit state carrying a Boolean flag.
\end{definition}

\begin{definition}[Tape Transfer Machine]
\lean{Turing.transferTM}
\label{Turing.transferTM}
\leanok
\uses{Turing.MultiTapeTM, Turing.SweepPhase}
Given a natural number $k$ and two work-tape indices
$\mathrm{src},\mathrm{dst}\in\{0,\ldots,k-1\}$, define a machine over any symbol
alphabet with sweep, rewind, and done states. In the sweep state, if the source head
scans a symbol, the machine writes that symbol at the destination head and moves both
heads right; if the source scans blank, it moves both heads left and enters rewind. In
rewind, if the source scans a symbol, it erases that cell and moves both heads left; if
the source scans blank, it moves both heads right and enters done. In done, it stays in
that state without moving or writing.
\end{definition}

\begin{definition}[Tape Copying Machine]
\lean{Turing.copyTM}
\label{Turing.copyTM}
\leanok
\uses{Turing.MultiTapeTM, Turing.SweepPhase}
Given a natural number $k$ and two work-tape indices
$\mathrm{src},\mathrm{dst}\in\{0,\ldots,k-1\}$, define a machine over any symbol
alphabet with sweep, rewind, and done states. In the sweep state, if the source head
scans a symbol, the machine writes it at the destination head and moves both heads
right; if the source scans blank, it moves both heads left and enters rewind. In rewind,
it moves both heads left while the source scans a symbol, and when the source scans
blank it moves both heads right and enters done. In done, it stays in that state without
moving or writing.
\end{definition}

\begin{definition}[Tape Clearing Machine]
\lean{Turing.clearTM}
\label{Turing.clearTM}
\leanok
\uses{Turing.MultiTapeTM, Turing.SweepPhase}
Given a natural number $k$, a work-tape index $i<k$, and any symbol alphabet, define a
machine with sweep, rewind, and done states. In the sweep state, it moves the head on
tape $i$ right while that head scans a symbol; when it scans blank, it moves that head
left and enters rewind. In rewind, it erases the scanned cell and moves left while the
head scans a symbol; when it scans blank, it moves right and enters done. All other
work-tape heads remain stationary, and in done the machine makes no changes or
movements.
\end{definition}

\begin{definition}[Tape Comparison Machine]
\lean{Turing.compareTM}
\label{Turing.compareTM}
\leanok
\uses{Turing.FlagPhase, Turing.MultiTapeTM}
Given a natural number $k$, two work-tape indices $\mathrm{fst},\mathrm{snd}<k$, and a
symbol alphabet with decidable equality, define a machine with run, flagged rewind, and
flagged done states. In the run state, if both heads scan symbols, it moves both right
when the symbols agree and otherwise moves both left and enters rewind with flag false;
if both scan blank, it moves both left and enters rewind with flag true; if exactly one
scans blank, it moves both left and enters rewind with flag false. In rewind it moves
both heads left while the first scans a symbol; when the first scans blank, it moves
both right and enters done carrying the flag. In done it stays in that state without
moving or writing.
\end{definition}

\begin{definition}[Binary increment machine]
\lean{Turing.incrementTM}
\label{Turing.incrementTM}
\leanok
\uses{Turing.FlagPhase, Turing.MultiTapeTM}
For a natural number $k$ and a selected work tape $i$ among the $k$ tapes, define a
machine with Boolean tape symbols and phases for running, rewinding with a Boolean flag,
and finishing with that flag. It starts in the running phase: when the selected cell
contains $1$, it writes $0$, moves that tape head right, and continues; when it contains
$0$, it writes $1$, moves the head left, and begins rewinding with flag true; when it is
blank, it moves left and begins rewinding with flag false. During rewinding it moves
left over nonblank cells and, at the first blank, moves right and enters the finished
phase; in the finished phase it remains there. Every transition leaves the input head
and output unchanged and leaves all unselected work tapes unchanged.
\end{definition}

\begin{definition}[Catalog configuration]
\lean{Turing.catalogCfg}
\label{Turing.catalogCfg}
\leanok
\uses{Turing.Cfg, Turing.Cfg.ofWords}
Given a natural number $k$, a state type $S$, a Boolean input string $x$, a state $q\in
S$, a list $w$ of Boolean words indexed by the $k$ work tapes, and integer head
positions for those tapes, define a configuration on input $x$ using state $q$ and the
words $w$, with the specified head positions. Its other configuration data are those
supplied by the configuration formed from $q$ and $w$.
\end{definition}

\begin{definition}[Catalog trace]
\lean{Turing.catalogTrace}
\label{Turing.catalogTrace}
\leanok
\uses{Turing.Cfg}
Let $F$ and $R$ be sequences of configurations for a fixed input, and let $D$ be another
configuration; for natural numbers $L$ and $t$, define the configuration at time $t$ to
be $F(t)$ if $t\le L$, to be $R(2L+1-t)$ if $L<t\le 2L+1$, and to be $D$ if $2L+1<t$.
\end{definition}

\begin{theorem}[Run Follows a Prescribed Trace]
\lean{Turing.catalog_trace_run}
\label{Turing.catalog_trace_run}
\leanok
\uses{Turing.Cfg, Turing.MultiTapeTM, Turing.MultiTapeTM.runFrom, Turing.MultiTapeTM.runFrom_succ_eq_step', Turing.MultiTapeTM.step, Turing.catalogTrace}
\difficulty{5}
Let $M$ be a $k$-tape machine, let $F$ and $R$ be sequences of configurations, let $D$
be a configuration, and let $L$ be a natural number. Suppose each step from $F(r)$ gives
$F(r+1)$ for $r<L$, the step from $F(L)$ gives $R(L)$, each step from $R(r+1)$ gives
$R(r)$ for $r<L$, the step from $R(0)$ gives $D$, and $D$ is fixed by a step. Then at
every time $t$, the run from $F(0)$ is $F(t)$ if $t\le L$, is $R(2L+1-t)$ if $L<t\le
2L+1$, and is $D$ if $t>2L+1$.
\end{theorem}

\begin{lemma}[Bound on Visited Tape Cells]
\lean{Turing.catalog_space_bound}
\label{Turing.catalog_space_bound}
\leanok
\uses{Turing.Cfg, Turing.MultiTapeTM, Turing.MultiTapeTM.runFrom, Turing.MultiTapeTM.spaceUsedByTape, Turing.MultiTapeTM.spaceUsedByTape_le_card_Icc}
\difficulty{3}
Let $M$ be a machine with $k$ work tapes, let $c$ be an initial configuration, and let
$L,t$ be natural numbers. If, at every time $u$, the position of work-tape head $i$ lies
between $-1$ and $L$, inclusive, then the number of distinct positions visited by that
head from time $0$ through time $t$ is at most $L+2$.
\end{lemma}

\begin{lemma}[Space Used by a Stationary Tape Head]
\lean{Turing.catalog_space_one}
\label{Turing.catalog_space_one}
\leanok
\uses{Turing.Cfg, Turing.MultiTapeTM, Turing.MultiTapeTM.runFrom, Turing.MultiTapeTM.spaceUsedByTape, Turing.MultiTapeTM.visitedByTapeHead}
\difficulty{3}
Let $M$ be a machine with $k$ work tapes, let $c$ be an initial configuration, let $t$
be a natural number, and let $i$ be a work-tape index. If head $i$ is at position $0$ at
every time $u$, then it visits exactly one distinct cell through time $t$.
\end{lemma}

\begin{definition}[Catalog tape contents]
\lean{Turing.catalogTape}
\label{Turing.catalogTape}
\leanok
\uses{Turing.FinTM.bufferTape}
Let $f$ be a tape-content function from integer positions to optional symbols, let $a$
be an integer, let $w$ be a finite word, and let $q$ be an integer position. Define the
tape content at $q$ by the buffer representation of $w$ at offset $q-a$ when $a\le
q<a+|w|$, and by $f(q)$ otherwise.
\end{definition}

\begin{theorem}[Extension of a tape prefix]
\lean{Turing.catalog_write_take}
\label{Turing.catalog_write_take}
\leanok
\uses{Turing.FinTM.bufferTape}
\difficulty{4}
Let $f\colon\mathbb{Z}\to\mathrm{Option}(S)$ be a tape content function, let
$a\in\mathbb{Z}$, and let $w$ be a finite word over an alphabet $S$. For a natural
number $r<|w|$, changing the tape that contains the prefix of $w$ in cells $a$ through
$a+r-1$ so that cell $a+r$ contains the symbol $w_r$ gives exactly the tape that
contains the prefix of $w$ in cells $a$ through $a+r$, leaving all other cells as in
$f$.
\end{theorem}

\begin{definition}[Forward catalog clearing configuration]
\lean{Turing.catalogClearF}
\label{Turing.catalogClearF}
\leanok
\uses{Turing.Cfg, Turing.SweepPhase}
Let $i$ be one of the $k$ work tapes, let $d$ be a configuration with a phase in the
three-phase sweep system, and let $r$ be a natural number. Define a new configuration by
setting its phase to the forward sweep and moving the head on tape $i$ right by $r$
positions; all other fields, including all tape contents and other head positions, are
retained.
\end{definition}

\begin{definition}[Rewind catalog clearing configuration]
\lean{Turing.catalogClearR}
\label{Turing.catalogClearR}
\leanok
\uses{Turing.Cfg, Turing.SweepPhase}
Let $i$ be one of the $k$ work tapes, let $d$ be a configuration with a sweep phase, and
let $n$ and $r$ be natural numbers. Define a configuration in the rewind phase by
clearing tape $i$ on the integer interval from $d$'s head position on that tape plus $r$
(inclusive) to that position plus $n$ (exclusive). Move that tape's head by $r-1$
positions, and retain the other fields and all other tape contents and head positions.
\end{definition}

\begin{lemma}[Trace of the clearing machine]
\lean{Turing.catalog_clear_trace}
\label{Turing.catalog_clear_trace}
\leanok
\uses{Turing.Action.apply, Turing.Cfg, Turing.Cfg.workTapeSymbols, Turing.FinTM.bufferTape, Turing.MultiTapeTM.runFrom, Turing.MultiTapeTM.step, Turing.SweepPhase, Turing.catalogClearF, Turing.catalogClearR, Turing.catalogTrace, Turing.catalog_trace_run, Turing.clearTM}
\difficulty{6}
Let $i$ be a work-tape index, $w$ a finite word, and $d$ a configuration whose state is
the forward-sweep phase. Suppose that, at every offset $p$ from $-1$ through the length
of $w$, tape $i$ at position equal to its head position plus $p$ contains the symbol of
$w$ at offset $p$ when that offset is a word position, and is blank otherwise. Then, for
every time $t$, the clearing machine started from $d$ has configuration given by the
following trace: its forward-sweep configuration at time $t$ when $t\le |w|$, its rewind
configuration indexed by $2|w|+1-t$ when $|w|<t\le 2|w|+1$, and its specified completed
configuration thereafter. The completed configuration has the same data as $d$ except
that the state is the done phase and tape $i$ is blank on the interval from its initial
head position inclusive to that position plus $|w|$ exclusive.
\end{lemma}

\begin{definition}[Forward catalog copy configuration]
\lean{Turing.catalogCopyF}
\label{Turing.catalogCopyF}
\leanok
\uses{Turing.Cfg, Turing.FinTM.bufferTape, Turing.SweepPhase}
Let $src$ and $dst$ be indices among the $k$ work tapes, let $w$ be a finite word, let
$d$ be a configuration with a sweep phase, and let $r$ be a natural number. Define a
configuration in the forward sweep phase by replacing the contents of tape $dst$ on the
interval from its head position in $d$ (inclusive) to that position plus $r$ (exclusive)
with the buffer representation of $w$ at the corresponding offsets. Move the heads on
$src$ and $dst$ right by $r$ positions and retain all other fields and tape contents.
\end{definition}

\begin{definition}[Rewind catalog copy configuration]
\lean{Turing.catalogCopyR}
\label{Turing.catalogCopyR}
\leanok
\uses{Turing.Cfg, Turing.SweepPhase, Turing.catalogCopyF}
Let $src$ and $dst$ be indices among the $k$ work tapes, let $w$ be a finite word, and
let $d$ be a configuration with a sweep phase. Define a configuration by first applying
the forward catalog copy operation to $d$ using a segment of length $|w|$, then setting
its phase to rewind and setting the heads on $src$ and $dst$ to their positions in $d$
plus $|w|-1$. All other fields are those of the forward copy configuration.
\end{definition}

\begin{definition}[Rightward Catalog Transfer]
\lean{Turing.catalogTransferR}
\label{Turing.catalogTransferR}
\leanok
\uses{Turing.Cfg, Turing.SweepPhase, Turing.catalogCopyR}
Given a positive number $k$ of work tapes, a source tape and a destination tape, a
finite word $w$, a configuration in the sweep, rewind, or done phase, and a natural
number $r$, the resulting configuration is obtained from the catalog copy operation and
has every source-tape cell with index from the source head position plus $r$ up to, but
not including, the source head position plus the length of $w$ left blank.
\end{definition}

\begin{lemma}[Forward step of the copying machine]
\lean{Turing.catalog_copy_forward}
\label{Turing.catalog_copy_forward}
\leanok
\uses{Turing.Action.apply, Turing.Cfg, Turing.Cfg.workTapeSymbols, Turing.FinTM.bufferTape, Turing.MultiTapeTM.step, Turing.SweepPhase, Turing.catalogCopyF, Turing.catalog_write_take, Turing.copyTM}
\difficulty{5}
Let $M$ be the copying machine, with distinct source and destination tape indices $s$
and $d$. Suppose a word $w$ is stored on the source tape of a configuration $c$: at
every offset $p$ from $-1$ through $|w|$, the cell at the source head position plus $p$
agrees with the word buffer at offset $p$. If $r<|w|$, then one step of the copying
machine from its prescribed forward configuration at progress $r$ gives the forward
configuration at progress $r+1$.
\end{lemma}

\begin{lemma}[Trace of the copying machine]
\lean{Turing.catalog_copy_trace}
\label{Turing.catalog_copy_trace}
\leanok
\uses{Turing.Action.apply, Turing.Cfg, Turing.Cfg.workTapeSymbols, Turing.FinTM.bufferTape, Turing.MultiTapeTM.runFrom, Turing.MultiTapeTM.step, Turing.SweepPhase, Turing.catalogCopyF, Turing.catalogCopyR, Turing.catalogTrace, Turing.catalog_copy_forward, Turing.catalog_trace_run, Turing.copyTM}
\difficulty{6}
Let $M$ be the copying machine, with distinct source and destination tape indices $s,d$.
Let $w$ be a word and let $c$ be a configuration in the sweep phase. Suppose the source
tape contains $w$ at every offset from $-1$ through $|w|$, relative to its head
position, with blank cells outside the word. Then for every time $t$, the run from $c$
follows the piecewise trace consisting of the forward configuration at time $t$ for
$t\le |w|$, the rewind configuration indexed by $2|w|+1-t$ for $|w|<t\le2|w|+1$, and the
completed configuration thereafter. In the completed configuration the destination tape
has the copied word as specified by the forward configuration at progress $|w|$.
\end{lemma}

\begin{lemma}[Trace of a Tape Transfer]
\lean{Turing.catalog_transfer_trace}
\label{Turing.catalog_transfer_trace}
\leanok
\uses{Turing.Action.apply, Turing.Cfg, Turing.Cfg.workTapeSymbols, Turing.FinTM.bufferTape, Turing.MultiTapeTM.runFrom, Turing.MultiTapeTM.step, Turing.SweepPhase, Turing.catalogCopyF, Turing.catalogCopyR, Turing.catalogTrace, Turing.catalogTransferR, Turing.catalog_copy_forward, Turing.catalog_trace_run, Turing.transferTM}
\difficulty{6}
Let $k$ be a natural number, let $x$ be an input word, let $S$ be a symbol alphabet, and
let $w$ be a finite word over $S$. Choose distinct work tapes $s$ and $d$, and let $c$
be a configuration in the sweep phase. Suppose that, for every integer $p$ with $-1\le
p\le |w|$, tape $s$ at position $	ext{head}_s(c)+p$ contains exactly the symbol or blank
specified by storing $w$ from offset $p$ (with positions outside the word blank). Then
the run of the transfer machine from $c$ follows this trace: during the forward pass,
both selected heads advance by one at each step and the destination tape is filled with
the corresponding symbols of $w$; during the return pass, both heads move back across
the word while the source cells are erased; and after the passes the machine remains in
the done phase. The trace configurations use the prescribed copied tape contents and
head positions, and its final configuration has the done phase and the tape contents at
the start of the return pass.
\end{lemma}

\begin{lemma}[First stopping position for word comparison]
\lean{Turing.catalog_compare_stop}
\label{Turing.catalog_compare_stop}
\leanok
\difficulty{4}
For any two finite words $u$ and $v$ over a type with decidable equality, there is an
index $d\le \min(|u|,|v|)$ such that every index $r<d$ contains the same symbol in both
words, while at index $d$ there is no symbol present in both. Moreover, the
symbols-or-absence at index $d$ are equal if and only if $u=v$.
\end{lemma}

\begin{definition}[Forward Comparison Configuration]
\lean{Turing.catalogCompareF}
\label{Turing.catalogCompareF}
\leanok
\uses{Turing.Cfg, Turing.FlagPhase}
Given a positive number $k$ of work tapes, two designated work tapes, a configuration in
the run, rewind, or done phase, and a natural number $r$, the resulting configuration
has state run, and the head positions on the two designated tapes are each increased by
$r$; all other head positions and configuration fields are unchanged.
\end{definition}

\begin{definition}[Rewind Comparison Configuration]
\lean{Turing.catalogCompareR}
\label{Turing.catalogCompareR}
\leanok
\uses{Turing.Cfg, Turing.FlagPhase}
Given a positive number $k$ of work tapes, two designated work tapes, a configuration in
the run, rewind, or done phase, a Boolean verdict $v$, and a natural number $r$, the
resulting configuration has state rewind carrying $v$. The head positions on the two
designated tapes are each increased by $r-1$, and all other head positions and
configuration fields are unchanged.
\end{definition}

\begin{lemma}[Trace of the word-comparison machine]
\lean{Turing.catalog_compare_trace}
\label{Turing.catalog_compare_trace}
\leanok
\uses{Turing.Action.apply, Turing.Cfg, Turing.Cfg.inputSymbol, Turing.Cfg.workTapeSymbols, Turing.FinTM.bufferTape, Turing.FlagPhase, Turing.MultiTapeTM.runFrom, Turing.MultiTapeTM.step, Turing.catalogCompareF, Turing.catalogCompareR, Turing.catalogTrace, Turing.catalog_trace_run, Turing.compareTM}
\difficulty{6}
Let $M$ be the comparison machine on $k$ work tapes, with distinct tape indices $f$ and
$s$. Let $u,v$ be finite words over a type with decidable equality, and let $d$ be a
configuration in the run phase. Suppose tapes $f$ and $s$, respectively, contain $u$ and
$v$ at every offset from $-1$ through their word lengths, with blank cells outside each
word, relative to their current head positions. Suppose also that an index
$\delta\le\min(|u|,|v|)$ has matching symbols at every earlier index, has no symbol
present in both words at index $\delta$, and the two words are equal exactly when their
symbols-or-absence at $\delta$ agree. Then for every time $t$, the run of the comparison
machine from $d$ is the corresponding piecewise trace: its forward configuration at time
$t$ for $t\le\delta$, its rewind configuration with verdict indicating whether $u=v$ at
index $2\delta+1-t$ for $\delta<t\le2\delta+1$, and thereafter its completed
configuration with that same verdict.
\end{lemma}

\begin{lemma}[Prefix and optional separator decomposition]
\lean{Turing.catalog_increment_split}
\label{Turing.catalog_increment_split}
\leanok
\difficulty{4}
Every finite Boolean word $w$ can be written as a string of $p$ true bits, followed by
either no separator or one false bit, followed by a remaining Boolean word; that is,
there exist $p\in\mathbb{N}$ and an optional word $z$ such that $w=\mathrm{true}^p$
followed by the optional false separator and then $z$.
\end{lemma}

\begin{lemma}[Incrementing a Run of Ones]
\lean{Turing.catalog_increment_value}
\label{Turing.catalog_increment_value}
\leanok
\uses{Turing.incFixed, Turing.incFixed_replicate_true_append}
\difficulty{3}
For every natural number $p$ and optional Boolean list $T$, incrementing the word formed
by $p$ true bits followed by a false bit and then the list in $T$ yields, if $T$ is
present, the word of $p$ false bits followed by a true bit and that list; if $T$ is
absent, the increment result is absent.
\end{lemma}

\begin{lemma}[Updating an Interior Tape Symbol]
\lean{Turing.catalog_write_middle}
\label{Turing.catalog_write_middle}
\leanok
\uses{Turing.FinTM.bufferTape, Turing.catalogTape}
\difficulty{4}
Let $f:br	o	ext{Option}(S)$ be a tape, let $z$ be an integer, let $	ext{pre}$ and
$	ext{rest}$ be finite words over a symbol set $S$, and let $a,b\in S$. If the tape is
formed by storing $	ext{pre}\,a\,	ext{rest}$ starting at position $z$, then updating the
cell at position $z+|	ext{pre}|$ to contain $b$ gives exactly the tape formed by storing
$	ext{pre}\,b\,	ext{rest}$ starting at $z$, with the original tape retained outside the
stored segment.
\end{lemma}

\begin{definition}[Forward Increment Configuration]
\lean{Turing.catalogIncF}
\label{Turing.catalogIncF}
\leanok
\uses{Turing.Cfg, Turing.FlagPhase, Turing.catalogTape}
Given a positive number $k$ of work tapes, a designated work tape, a configuration in
the run, rewind, or done phase, natural numbers $p$ and $r$, and an optional finite
Boolean word $t$, the resulting configuration has state run. On the designated tape it
applies the catalog tape operation at the current head position to the word consisting
of $r$ zeros, followed by $p-r$ ones, followed by a zero and $t$ if $t$ is present (and
no further symbols otherwise); its head position is increased by $r$. The other tapes
and configuration fields are unchanged.
\end{definition}

\begin{definition}[Rewind Increment Configuration]
\lean{Turing.catalogIncR}
\label{Turing.catalogIncR}
\leanok
\uses{Turing.Cfg, Turing.FlagPhase, Turing.catalogTape}
Given a positive number $k$ of work tapes, a designated work tape, a configuration in
the run, rewind, or done phase, natural numbers $p$ and $r$, and an optional finite
Boolean word $t$, the resulting configuration has state rewind carrying true exactly
when $t$ is present. On the designated tape it applies the catalog tape operation at the
current head position to the word consisting of $p$ zeros, followed by a one and $t$ if
$t$ is present (and no further symbols otherwise); its head position is increased by
$r-1$. The other tapes and configuration fields are unchanged.
\end{definition}

\begin{lemma}[Trace of the increment machine]
\lean{Turing.catalog_increment_trace}
\label{Turing.catalog_increment_trace}
\leanok
\uses{Turing.Action.apply, Turing.Cfg, Turing.Cfg.workTapeSymbols, Turing.FinTM.bufferTape, Turing.FlagPhase, Turing.MultiTapeTM.runFrom, Turing.MultiTapeTM.step, Turing.catalogIncF, Turing.catalogIncR, Turing.catalogTape, Turing.catalogTrace, Turing.catalog_trace_run, Turing.catalog_write_middle, Turing.incrementTM}
\difficulty{6}
Let $M$ be the increment machine on tape $i$, and let $w$ be a Boolean word stored on
tape $i$ of a configuration $c$ in the run phase: at every offset from $-1$ through
$|w|$, the tape agrees with the word buffer relative to its head position. Suppose $w$
consists of $p$ true bits followed by either no separator or one false bit and then a
remaining word $z$. Then for every time $t$, the run from $c$ is the piecewise trace
whose forward configuration at time $t$ applies for $t\le p$, whose rewind configuration
indexed by $2p+1-t$ applies for $p<t\le2p+1$, and whose completed configuration applies
thereafter; the verdict in the rewind and completed phases records whether the optional
tail $z$ is present.
\end{lemma}

\begin{theorem}[Behavior of the transfer machine]
\lean{Turing.transferTM_run_ofCfg}
\label{Turing.transferTM_run_ofCfg}
\leanok
\uses{Turing.Cfg, Turing.FinTM.bufferTape, Turing.MultiTapeTM.runFrom, Turing.SweepPhase, Turing.catalogCopyF, Turing.catalogCopyR, Turing.catalogTrace, Turing.catalogTransferR, Turing.catalog_transfer_trace, Turing.transferTM}
\difficulty{4}
Let $k$ be a natural number, let $x$ be a finite input word over an alphabet of symbols,
and let $w$ be another finite word over that alphabet. Choose distinct indices
$\mathrm{src},\mathrm{dst}\in\{0,\ldots,k-1\}$ and a configuration $d$ on input $x$
whose state is the sweep phase. Suppose that, for every integer offset $p$ with $-1\le
p\le |w|$, the source tape of $d$ at the cell $p$ steps from its initial head position
agrees with the buffer storing $w$ at offset $p$. After $2|w|+2$ steps, the transfer
machine is in the done phase; it has erased the source tape and written the buffered
word on the destination tape at offsets $0,\ldots,|w|-1$, leaving every other tape cell
unchanged. It does not enter the done phase at any earlier step. At every step up to and
including step $2|w|+2$, the source and destination heads lie between one cell to the
left of their respective initial positions and $|w|$ cells to their right, inclusive,
while every other work-tape head remains at its initial position.
\end{theorem}

\begin{theorem}[Transferring a Word Between Tapes]
\lean{Turing.Cfg.ofWords_transfer_endpoint}
\label{Turing.Cfg.ofWords_transfer_endpoint}
\leanok
\uses{Turing.Cfg, Turing.Cfg.ofWords, Turing.FinTM.bufferTape, Turing.FinTM.bufferTape_nil}
\difficulty{4}
Let $k$ be a nonnegative integer, let $x$ be a Boolean input word, let $q,q'$ be states,
and let $w$ assign a Boolean word to each of the $k$ work tapes. Choose distinct source
and destination tape indices, and suppose the destination word is empty. Starting from
the configuration determined by $q$ and $w$, replace the state by $q'$, erase the source
word, and copy it to the destination tape in cells starting at position $0$, leaving all
other tape cells and all other configuration data unchanged. The resulting configuration
is the one determined by $q'$ and the family obtained by replacing the source word with
the empty word and the destination word with the original source word.
\end{theorem}

\begin{theorem}[Copying from a general configuration]
\lean{Turing.copyTM_run_ofCfg}
\label{Turing.copyTM_run_ofCfg}
\leanok
\uses{Turing.Cfg, Turing.FinTM.bufferTape, Turing.MultiTapeTM.runFrom, Turing.SweepPhase, Turing.catalogCopyF, Turing.catalogCopyR, Turing.catalogTrace, Turing.catalog_copy_trace, Turing.copyTM}
\difficulty{4}
Let $k\in\mathbb{N}$, let $x$ be an input word over an alphabet $\Sigma$, and let
$\mathrm{src},\mathrm{dst}$ be distinct indices among the $k$ work tapes. Let $w$ be a
word over $\Sigma$ and $d$ a configuration of the copying machine on input $x$. Suppose
that $d$ is in the sweep state and that, at every integer offset $p$ from $-1$ through
$|w|$ relative to the source tape's head, the source tape contains the symbol specified
by the word-buffer representation of $w$ at offset $p$. After $2|w|+2$ steps, the
machine is in the done state; its destination tape contains the word-buffer
representation of $w$ beginning at its original head position, and every other tape cell
retains its original content. At all earlier times it is not in the done state.
Throughout these steps, the source and destination head positions lie between their
respective initial positions minus one and their initial positions plus $|w|$,
inclusive, while every other tape head remains at its initial position.
\end{theorem}

\begin{theorem}[Copying a Word to an Empty Tape]
\lean{Turing.Cfg.ofWords_copy_endpoint}
\label{Turing.Cfg.ofWords_copy_endpoint}
\leanok
\uses{Turing.Cfg, Turing.Cfg.ofWords, Turing.FinTM.bufferTape, Turing.FinTM.bufferTape_nil}
\difficulty{3}
Let $k$ be a nonnegative integer, let $x$ be a Boolean input word, let $q,q'$ be states,
and let $w$ assign a Boolean word to each of the $k$ work tapes. Choose source and
destination tape indices, and suppose the destination word is empty. Starting from the
configuration determined by $q$ and $w$, replace the state by $q'$ and replace the
destination tape's cells from position $0$ through position $|w_{
m src}|-1$ with the
source word, leaving all other tape cells and all other configuration data unchanged.
The resulting configuration is the one determined by $q'$ and the family of words
obtained by replacing the destination word with the source word.
\end{theorem}

\begin{theorem}[Clearing a tape from a configuration]
\lean{Turing.clearTM_run_ofCfg}
\label{Turing.clearTM_run_ofCfg}
\leanok
\uses{Turing.Cfg, Turing.FinTM.bufferTape, Turing.MultiTapeTM.runFrom, Turing.SweepPhase, Turing.catalogClearF, Turing.catalogClearR, Turing.catalogTrace, Turing.catalog_clear_trace, Turing.clearTM}
\difficulty{4}
Let $i<k$, let $w$ be a word over an alphabet, and let $d$ be a configuration in the
sweep state. Suppose that for every integer $p$ with $-1\le p\le |w|$, tape $i$ at
position $d_i+p$ contains the symbol of $w$ at position $p$ when that position is in the
word, and is blank otherwise. After $2|w|+2$ steps, the clearing machine is in the done
state, with tape $i$ blank on $[d_i,d_i+|w|)$ and all other tape cells unchanged; it was
not in the done state at any earlier time. Throughout those steps, the head on tape $i$
stays in $[d_i-1,d_i+|w|]$, while every other tape head stays at its initial position.
\end{theorem}

\begin{theorem}[Successful increment from an arbitrary configuration]
\lean{Turing.incrementTM_run_succ_ofCfg}
\label{Turing.incrementTM_run_succ_ofCfg}
\leanok
\uses{Turing.Cfg, Turing.FinTM.bufferTape, Turing.FlagPhase, Turing.MultiTapeTM.runFrom, Turing.catalogIncF, Turing.catalogIncR, Turing.catalogTape, Turing.catalogTrace, Turing.catalog_increment_split, Turing.catalog_increment_trace, Turing.catalog_increment_value, Turing.incFixed, Turing.incrementTM}
\difficulty{5}
Let $k$ be a natural number, let $x,w,v$ be Boolean words, choose a work tape
$i\in\mathrm{Fin}(k)$, and let $d$ be a configuration on input $x$. Suppose incrementing
$w$ yields $v$, the state of $d$ is the run phase, and on tape $i$ the cells from the
head position minus one through the head position plus $|w|$ agree with the buffer tape
storing $w$ at the corresponding relative positions. After exactly $2m+2$ steps, where
$m$ is the length of the initial run of true bits in $w$, the state is the done phase
with true verdict, and the tape contents are those of $d$ except that the chosen tape
stores $v$ in the interval of $|w|$ cells starting at its initial head position; no
earlier time is in either done phase. At every time up to and including that bound, the
chosen tape's head lies between one cell left of its initial position and $m$ cells
right of it, inclusive, while every other tape head stays fixed.
\end{theorem}

\begin{theorem}[Overflow run from an arbitrary configuration]
\lean{Turing.incrementTM_run_overflow_ofCfg}
\label{Turing.incrementTM_run_overflow_ofCfg}
\leanok
\uses{Turing.Cfg, Turing.FinTM.bufferTape, Turing.FlagPhase, Turing.MultiTapeTM.runFrom, Turing.catalogIncF, Turing.catalogIncR, Turing.catalogTape, Turing.catalogTrace, Turing.catalog_increment_split, Turing.catalog_increment_trace, Turing.catalog_increment_value, Turing.incFixed, Turing.incrementTM}
\difficulty{5}
Let $k$ be a natural number, let $x$ and $w$ be Boolean words, choose a work tape
$i\in\mathrm{Fin}(k)$, and let $d$ be a configuration on input $x$. Suppose incrementing
$w$ overflows, the state of $d$ is the run phase, and on tape $i$ the cells from the
head position minus one through the head position plus $|w|$ agree with the buffer tape
storing $w$ at the corresponding relative positions. After exactly $2|w|+2$ steps, the
state is the done phase with false verdict, and the tape contents are those of $d$
except that the chosen tape's cells from its head position up to but excluding the
position $|w|$ cells to its right contain the all-false word; no earlier time is in
either done phase. At every time up to and including that bound, the chosen tape's head
lies between one cell left of its initial position and $|w|$ cells right of it,
inclusive, while every other tape head stays fixed.
\end{theorem}

\begin{theorem}[Bounded Tape Transfer Run]
\lean{Turing.transferTM_run}
\label{Turing.transferTM_run}
\leanok
\uses{Turing.Cfg.ofWords, Turing.FinTM.bufferTape, Turing.FinTM.bufferTape_nil, Turing.MultiTapeTM.runFrom, Turing.transferTM, Turing.transferTM_run_ofCfg}
\difficulty{5}
Let $k\in\mathbb{N}$, let $\mathrm{src}$ and $\mathrm{dst}$ be distinct indices among
the $k$ work tapes, let $x$ be a Boolean input, and assign a Boolean word $w_i$ to each
tape $i$, with $w_{\mathrm{dst}}$ empty. There is a natural number $T\leq
3|w_{\mathrm{src}}|+3$ such that the transfer machine, started in its sweep phase with
these tape words, has not reached its done phase at any time $t<T$. At time $T$ it is in
the done phase, with the source tape emptied and the original source word placed on the
destination tape; all other tape words are unchanged.
\end{theorem}

\begin{theorem}[Space used by the transfer machine]
\lean{Turing.transferTM_spaceUsedByTape}
\label{Turing.transferTM_spaceUsedByTape}
\leanok
\uses{Turing.Cfg.ofWords, Turing.MultiTapeTM.spaceUsedByTape, Turing.catalogCfg, Turing.catalogCopyF, Turing.catalogCopyR, Turing.catalogTrace, Turing.catalogTransferR, Turing.catalog_space_bound, Turing.catalog_space_one, Turing.catalog_transfer_trace, Turing.transferTM}
\difficulty{4}
Let $k$ be a natural number, let $x$ be a Boolean input word, and let $w$ assign a
Boolean word to each of the $k$ work tapes. Choose distinct tape indices $\mathrm{src}$
and $\mathrm{dst}$, and suppose the destination tape's word is empty. Starting from the
configuration in the sweep phase with these words stored from cell zero, for every time
$t$, the number of cells visited by the source head through step $t$ is at most
$|w(\mathrm{src})|+2$, and the same bound holds for the destination head. Every other
work-tape head visits exactly one distinct cell through step $t$.
\end{theorem}

\begin{theorem}[Running the copying machine]
\lean{Turing.copyTM_run}
\label{Turing.copyTM_run}
\leanok
\uses{Turing.Cfg.ofWords, Turing.FinTM.bufferTape, Turing.FinTM.bufferTape_nil, Turing.MultiTapeTM.runFrom, Turing.copyTM, Turing.copyTM_run_ofCfg}
\difficulty{4}
Let $k\in\mathbb{N}$, let $x$ be a Boolean input word, and let
$\mathrm{src},\mathrm{dst}$ be distinct indices among the $k$ work tapes. For any family
$w$ of Boolean words on those tapes with $w_{\mathrm{dst}}$ empty, there is a time $T\le
3|w_{\mathrm{src}}|+3$ such that the copying machine, started in its sweep state with
tape contents $w$, is not in its done state at any time $t<T$, and at time $T$ it is in
its done state with tape contents given by replacing the destination word by
$w_{\mathrm{src}}$.
\end{theorem}

\begin{theorem}[Space used during copying]
\lean{Turing.copyTM_spaceUsedByTape}
\label{Turing.copyTM_spaceUsedByTape}
\leanok
\uses{Turing.Cfg.ofWords, Turing.MultiTapeTM.spaceUsedByTape, Turing.catalogCfg, Turing.catalogCopyF, Turing.catalogCopyR, Turing.catalogTrace, Turing.catalog_copy_trace, Turing.catalog_space_bound, Turing.catalog_space_one, Turing.copyTM}
\difficulty{5}
Let $k\in\mathbb{N}$, let $x$ be a Boolean input word, and let
$\mathrm{src},\mathrm{dst}$ be distinct indices among the $k$ work tapes. For any family
$w$ of Boolean words on those tapes with empty destination word, and for every time
$t\in\mathbb{N}$, the copying machine started in its sweep state with tape contents $w$
has visited at most $|w_{\mathrm{src}}|+2$ distinct cells on each of the source and
destination tapes by time $t$. Every other tape has visited exactly one distinct cell by
time $t$.
\end{theorem}

\begin{theorem}[Running time and outcome of clearing]
\lean{Turing.clearTM_run}
\label{Turing.clearTM_run}
\leanok
\uses{Turing.Cfg.ofWords, Turing.FinTM.bufferTape, Turing.FinTM.bufferTape_nil, Turing.MultiTapeTM.runFrom, Turing.clearTM, Turing.clearTM_run_ofCfg}
\difficulty{4}
Let $k$ be the number of work tapes, let $x$ be a Boolean input word, choose a tape
$i<k$, and let $w$ assign a Boolean word to each tape. There is a time $T\le 2|w_i|+2$
such that the clearing machine has not reached its done state at any time $t<T$, and at
time $T$ its configuration is the initial configuration with state done and tape $i$
replaced by the empty word.
\end{theorem}

\begin{theorem}[Space used by the clearing machine]
\lean{Turing.clearTM_spaceUsedByTape}
\label{Turing.clearTM_spaceUsedByTape}
\leanok
\uses{Turing.Cfg.ofWords, Turing.MultiTapeTM.spaceUsedByTape, Turing.catalogCfg, Turing.catalogClearF, Turing.catalogClearR, Turing.catalogTrace, Turing.catalog_clear_trace, Turing.catalog_space_bound, Turing.catalog_space_one, Turing.clearTM}
\difficulty{4}
Let $k$ be the number of work tapes, let $x$ be a Boolean input word, choose a tape
$i<k$, and let $w$ assign a Boolean word to each tape. At every time $t$, the clearing
machine started with these words visits at most $|w_i|+2$ distinct cells on tape $i$,
and it visits exactly one distinct cell on each tape $j\ne i$.
\end{theorem}

\begin{theorem}[Runtime and outcome of the comparison machine]
\lean{Turing.compareTM_run}
\label{Turing.compareTM_run}
\leanok
\uses{Turing.Cfg.ofWords, Turing.MultiTapeTM.runFrom, Turing.catalogCompareF, Turing.catalogCompareR, Turing.catalogTrace, Turing.catalog_compare_stop, Turing.catalog_compare_trace, Turing.compareTM}
\difficulty{4}
Let $k$ be a natural number, let $x$ be a Boolean input string, let $fst,sndin k$ be
work-tape indices, and let $w$ assign a Boolean string to each work tape. There exists a
natural number $T\le 2
oot{	ext{min}(|w(fst)|,|w(snd)|)}+2$ such that at every time
$t<T$ the machine has not entered either completed state, and at time $T$ its
configuration is the initial word configuration with state set to completed with the
truth value of whether $w(fst)=w(snd)$.
\end{theorem}

\begin{theorem}[Space used by the comparison machine]
\lean{Turing.compareTM_spaceUsedByTape}
\label{Turing.compareTM_spaceUsedByTape}
\leanok
\uses{Turing.Cfg.ofWords, Turing.MultiTapeTM.spaceUsedByTape, Turing.catalogCompareF, Turing.catalogCompareR, Turing.catalogTrace, Turing.catalog_compare_stop, Turing.catalog_compare_trace, Turing.catalog_space_bound, Turing.catalog_space_one, Turing.compareTM}
\difficulty{5}
Let $k$ be a natural number, choose work-tape indices $fst,sndin k$, let $w$ assign a
Boolean word to each of the $k$ work tapes, and let $t$ be a natural number. Starting
from the configuration containing those words, the comparison machine visits at most
$min(|w(fst)|,|w(snd)|)+2$ distinct cells on each of tapes $fst$ and $snd$ by time $t$.
Every tape whose index differs from both $fst$ and $snd$ visits exactly one distinct
cell by that time.
\end{theorem}

\begin{theorem}[Running time for a successful increment]
\lean{Turing.incrementTM_run_succ}
\label{Turing.incrementTM_run_succ}
\leanok
\uses{Turing.Cfg.ofWords, Turing.FinTM.bufferTape, Turing.MultiTapeTM.runFrom, Turing.catalog_increment_split, Turing.catalog_increment_value, Turing.incFixed, Turing.incrementTM, Turing.incrementTM_run_succ_ofCfg}
\difficulty{5}
Let $k$ be a natural number, choose one of $k$ work tapes, and let $x$ be a Boolean
input word and $w$ a family of Boolean words on the work tapes. If incrementing the
chosen tape's word yields a Boolean word $v$, then the increment machine started in its
run phase from the configuration storing $w$ reaches the done phase with true verdict
after some $T\le 2|w_i|+2$ steps. No earlier time $t<T$ is in either done phase, and at
time $T$ the configuration stores the updated family obtained by replacing $w_i$ by $v$.
\end{theorem}

\begin{theorem}[Running time on overflow]
\lean{Turing.incrementTM_run_overflow}
\label{Turing.incrementTM_run_overflow}
\leanok
\uses{Turing.Cfg.ofWords, Turing.FinTM.bufferTape, Turing.MultiTapeTM.runFrom, Turing.incFixed, Turing.incrementTM, Turing.incrementTM_run_overflow_ofCfg}
\difficulty{6}
Let $k$ be a natural number, choose one of $k$ work tapes, and let $x$ be a Boolean
input word and $w$ a family of Boolean words on the work tapes. If fixed-width increment
of the chosen tape's word overflows, then the increment machine started in its run phase
from the configuration storing $w$ reaches the done phase with false verdict after some
$T\le 2|w_i|+2$ steps. No earlier time $t<T$ is in either done phase, and at time $T$
the configuration is the one storing the same words except that the chosen tape contains
$|w_i|$ false bits.
\end{theorem}

\begin{theorem}[Space used by the increment machine]
\lean{Turing.incrementTM_spaceUsedByTape}
\label{Turing.incrementTM_spaceUsedByTape}
\leanok
\uses{Turing.Cfg.ofWords, Turing.MultiTapeTM.spaceUsedByTape, Turing.catalogCfg, Turing.catalogIncF, Turing.catalogIncR, Turing.catalogTrace, Turing.catalog_increment_split, Turing.catalog_increment_trace, Turing.catalog_space_bound, Turing.catalog_space_one, Turing.incrementTM}
\difficulty{4}
Let $k\ge 0$, choose a work-tape index $i<k$, a family of Boolean words $w=(w_j)_{j<k}$,
and a time $t\in\mathbb N$. For the increment machine on any Boolean input, the number
of cells visited on tape $i$ by time $t$ is at most $|w_i|+2$, while every other tape
$j\ne i$ has visited exactly one cell.
\end{theorem}

\begin{theorem}[Capture tape visitation and space]
\lean{Turing.capture_visitedByTapeHead}
\label{Turing.capture_visitedByTapeHead}
\leanok
\uses{Turing.Cfg, Turing.Cfg.Halted, Turing.MultiTapeTM, Turing.MultiTapeTM.output_prefix, Turing.MultiTapeTM.runFrom, Turing.MultiTapeTM.runFrom_zero, Turing.MultiTapeTM.spaceUsedByTape, Turing.MultiTapeTM.visitedByTapeHead, Turing.captureAction, Turing.captureCfg, Turing.capture_run}
\difficulty{5}
Let $M$ be a deterministic machine with $k$ work tapes and Boolean output, and let $H$
be a machine with $k+1$ work tapes and a Boolean output. Suppose a map from the states
of $M$ to those of $H$ and a return state are chosen, and at every mapped state the
transition of $H$ agrees with the capture transform of the transition of $M$. For any
Boolean input, initial configuration $c_0$, Boolean words $p,o$, and time $t$, suppose
$M$ is not halted at any time strictly less than $t$. Then each of the first $k$ work
tapes of $H$, started from the capture image of $c_0$, visits exactly the same head
positions through time $t$ as the corresponding tape of $M$, and uses the same amount of
space; the additional tape of $H$ uses at most the output length of $M$ at time $t$
minus its initial output length, plus one.
\end{theorem}

\begin{definition}[Binary Identity Machine]
\lean{Turing.FinTM.f2_idTM}
\label{Turing.FinTM.f2_idTM}
\leanok
\uses{Turing.FinTM}
There is a finite Turing machine over the Boolean alphabet with no work tapes and a
single state. On each input bit it moves the input head one position to the right,
outputs that bit, and remains in its state; when it scans a blank, it stays in place,
outputs nothing, and halts.
\end{definition}

\begin{lemma}[Identity machine run]
\lean{Turing.FinTM.f2_idTM_run}
\label{Turing.FinTM.f2_idTM_run}
\leanok
\uses{Turing.Action.apply, Turing.Cfg.workTapeSymbols, Turing.FinTM.f2_idTM, Turing.MultiTapeTM.initCfg, Turing.MultiTapeTM.runFrom, Turing.MultiTapeTM.runFrom_succ_eq_step', Turing.MultiTapeTM.step, Turing.inputSymbolInner, Turing.moveInputPos_pos_of_ne_right}
\difficulty{4}
For every Boolean word $x$ and every natural number $t\le |x|$, the identity machine
started on $x$ is still in its initial state after $t$ steps, its input head is at
position $t+1$, and its output is the prefix of $x$ of length $t$.
\end{lemma}

\begin{definition}[Binary Constant-Output Machine]
\lean{Turing.FinTM.f2_constTM}
\label{Turing.FinTM.f2_constTM}
\leanok
\uses{Turing.FinTM, Turing.FinTM.emitAction}
For each finite Boolean word $w$, there is a finite Turing machine over the Boolean
alphabet with no work tapes and $|w|+1$ states, whose initial state is $0$ and whose
transition in state $i$ is the action obtained from the constant-output construction for
$w$, the identity state map, and $i$.
\end{definition}

\begin{definition}[Binary Prefix-Then-Copy Machine]
\lean{Turing.FinTM.f2_catalogPrefixTM}
\label{Turing.FinTM.f2_catalogPrefixTM}
\leanok
\uses{Turing.FinTM}
For each finite Boolean word $w$, there is a finite Turing machine over the Boolean
alphabet with no work tapes and $|w|+1$ states, initially in state $0$. In any state $q$
with $q<|w|$, it stays on the input cell, outputs the $q$th bit of $w$, and advances to
state $q+1$. In any state $q$ with $q\ge |w|$, it moves right and outputs the scanned
input bit while that bit is present, remaining in state $q$; on a blank cell it stays in
place, outputs nothing, and halts.
\end{definition}

\begin{definition}[Catalog prefix configuration]
\lean{Turing.FinTM.f2_catalogPrefixCfg}
\label{Turing.FinTM.f2_catalogPrefixCfg}
\leanok
\uses{Turing.Cfg}
Given Boolean lists $w$ and $x$, a state $q\in\mathrm{Option}(\mathrm{Fin}(|w|+1))$, an
input-head position $p\in\mathrm{Fin}(|x|+2)$, and an output list $o$, the associated
configuration on input $x$ has state $q$, head position $p$, no work tapes, and output
$o$.
\end{definition}

\begin{lemma}[Emitted catalog prefix]
\lean{Turing.FinTM.f2_catalogPrefixTM_emit}
\label{Turing.FinTM.f2_catalogPrefixTM_emit}
\leanok
\uses{Turing.Action.apply, Turing.Cfg.ext_zero_tapes, Turing.FinTM.f2_catalogPrefixCfg, Turing.FinTM.f2_catalogPrefixTM, Turing.MultiTapeTM.initCfg, Turing.MultiTapeTM.runFrom, Turing.MultiTapeTM.runFrom_succ_eq_step', Turing.MultiTapeTM.step}
\difficulty{4}
For Boolean words $w$ and $x$, and every $i \le |w|$, after $i$ steps the catalog-prefix
machine started in its initial configuration on input $x$ is in state $i$, has
input-head position $1$, and has output equal to the first $i$ symbols of $w$.
\end{lemma}

\begin{lemma}[Copied catalog prefix]
\lean{Turing.FinTM.f2_catalogPrefixTM_copy}
\label{Turing.FinTM.f2_catalogPrefixTM_copy}
\leanok
\uses{Turing.Action.apply, Turing.Cfg.ext_zero_tapes, Turing.Cfg.inputSymbol, Turing.FinTM.f2_catalogPrefixCfg, Turing.FinTM.f2_catalogPrefixTM, Turing.MultiTapeTM.runFrom, Turing.MultiTapeTM.runFrom_succ_eq_step', Turing.inputSymbolInner, Turing.moveInputPos, Turing.moveInputPos_pos_of_ne_right}
\difficulty{4}
For Boolean words $w$ and $x$, and every $i \le |x|$, after $i$ steps the catalog-prefix
machine initialized with state $|w|$, input-head position $1$, and output $w$ is in
state $|w|$, has input-head position $i+1$, and has output $w$ followed by the first $i$
symbols of $x$.
\end{lemma}

\begin{lemma}[Prefix machine computes concatenation]
\lean{Turing.FinTM.f2_catalogPrefixTM_computes}
\label{Turing.FinTM.f2_catalogPrefixTM_computes}
\leanok
\uses{Turing.Action.apply, Turing.Cfg.inputSymbol, Turing.FinTM.ComputesFunInTime, Turing.FinTM.computesInTime_iff, Turing.FinTM.f2_catalogPrefixCfg, Turing.FinTM.f2_catalogPrefixTM, Turing.FinTM.f2_catalogPrefixTM_copy, Turing.FinTM.f2_catalogPrefixTM_emit, Turing.MultiTapeTM.runFrom_add, Turing.MultiTapeTM.runFrom_succ_eq_step', Turing.MultiTapeTM.step}
\difficulty{3}
For every Boolean word $w$, the finite Turing machine determined by $w$ computes the
function sending each Boolean input word $x$ to the concatenation $w \mathbin{++} x$ in
time $|w|+|x|+1$.
\end{lemma}

\begin{definition}[Scan configuration]
\lean{Turing.FinTM.f2_scanCfg}
\label{Turing.FinTM.f2_scanCfg}
\leanok
\uses{Turing.Cfg}
Given a Boolean input list $x$, a state $q$, an integer $i$ with $0\le i\le |x|$, and an
output list $o$, the associated configuration has state $q$, input-head position $i+1$
in the shifted indexing of the input, no work tapes, and output $o$.
\end{definition}

\begin{lemma}[Input symbol in a scan configuration]
\lean{Turing.FinTM.f2_scanCfg_read}
\label{Turing.FinTM.f2_scanCfg_read}
\leanok
\uses{Turing.Cfg.inputSymbol, Turing.FinTM.f2_scanCfg, Turing.inputSymbolInner}
\difficulty{4}
For every Boolean word $x$, state option $q$, output word $o$, and position $i\le |x|$,
the symbol read by the input head in the scan configuration at position $i+1$ is the
$i$th entry of $x$, interpreted as an optional symbol.
\end{lemma}

\begin{lemma}[Output during a partial input scan]
\lean{Turing.FinTM.f2_scanCopy_run}
\label{Turing.FinTM.f2_scanCopy_run}
\leanok
\uses{Turing.Action.apply, Turing.Cfg.ext_zero_tapes, Turing.Cfg.inputSymbol, Turing.FinTM.f2_scanCfg, Turing.FinTM.f2_scanCfg_read, Turing.MultiTapeTM, Turing.MultiTapeTM.runFrom, Turing.MultiTapeTM.runFrom_succ_eq_step', Turing.MultiTapeTM.step, Turing.moveInputPos_pos_of_ne_right}
\difficulty{5}
Let $M$ be a Turing machine with no work tapes, let $q$ be a state, and suppose that in
state $q$ each input bit is copied to the output while the input head moves right and
the state remains $q$, and a blank input symbol causes the machine to emit nothing and
halt. For every Boolean word $x$, output word $o$, and $j\le |x|$, after $j$ steps from
the scan configuration at position zero in state $q$ with output $o$, the machine is in
state $q$ at position $j$ and its output is $o$ followed by the first $j$ entries of
$x$.
\end{lemma}

\begin{lemma}[Completion of a full input scan]
\lean{Turing.FinTM.f2_scanCopy_finish}
\label{Turing.FinTM.f2_scanCopy_finish}
\leanok
\uses{Turing.Action.apply, Turing.Cfg.ext_zero_tapes, Turing.Cfg.inputSymbol, Turing.FinTM.f2_scanCfg, Turing.FinTM.f2_scanCfg_read, Turing.FinTM.f2_scanCopy_run, Turing.MultiTapeTM, Turing.MultiTapeTM.runFrom, Turing.MultiTapeTM.runFrom_succ_eq_step', Turing.MultiTapeTM.step}
\difficulty{3}
Let $M$ be a Turing machine with no work tapes, let $q$ be a state, and suppose that in
state $q$ its transition on any input symbol and work-tape reading copies the scanned
input bit to the output, moves the input head right, and remains in state $q$, while on
a blank input symbol it emits nothing and halts. Then for every Boolean word $x$ and
output word $o$, after $|x|+1$ steps from the scan configuration at position zero in
state $q$ with output $o$, the machine is in the halted scan configuration at position
$|x|$, with output $o$ followed by $x$.
\end{lemma}

\begin{definition}[Pair-duplication machine]
\lean{Turing.FinTM.f2_pairDupTM}
\label{Turing.FinTM.f2_pairDupTM}
\leanok
\uses{Turing.FinTM, Turing.FinTM.controlAction}
There is a deterministic Boolean Turing machine with no work tapes and five states,
numbered $0$ through $4$, whose initial state is $0$. Its transition rules are as
follows: in state $0$, on a Boolean input symbol $b$, it emits $b$ and enters state $1$
without moving the input head; on blank it moves left, emits false, and enters state
$2$. In state $1$ it moves right, emits the scanned symbol if present, and returns to
state $0$. In state $2$, it moves left while reading symbols, and on blank moves right
into state $3$. In state $3$ it emits true and enters state $4$ without moving; in state
$4$ it moves right and emits each scanned symbol, halting on blank.
\end{definition}

\begin{lemma}[Pair duplication run invariant]
\lean{Turing.FinTM.f2_pairDup_double}
\label{Turing.FinTM.f2_pairDup_double}
\leanok
\uses{Turing.Action.apply, Turing.Cfg.ext_zero_tapes, Turing.FinTM.f2_pairDupTM, Turing.FinTM.f2_scanCfg, Turing.FinTM.f2_scanCfg_read, Turing.MultiTapeTM.initCfg, Turing.MultiTapeTM.runFrom, Turing.MultiTapeTM.runFrom_succ_eq_step', Turing.MultiTapeTM.step, Turing.moveInputPos_pos_of_ne_right}
\difficulty{5}
For every finite bit string $x$, every natural number $j$ with $j \\le |x|$, and every
such proof of the inequality, after $2j$ steps the constructed machine is in state $0$
of its five-state scanning phase, with input head at position $j+1$, no work tapes, and
output equal to the first $j$ bits of $x$, each doubled in place.
\end{lemma}

\begin{lemma}[Pair duplication computation]
\lean{Turing.FinTM.f2_pairDup_computes}
\label{Turing.FinTM.f2_pairDup_computes}
\leanok
\uses{Turing.Action.apply, Turing.Cfg, Turing.Cfg.ext_zero_tapes, Turing.FinTM.ComputesInTime, Turing.FinTM.computesInTime_iff, Turing.FinTM.f2_pairDupTM, Turing.FinTM.f2_pairDup_double, Turing.FinTM.f2_scanCfg, Turing.FinTM.f2_scanCfg_read, Turing.FinTM.f2_scanCopy_finish, Turing.FinTM.rewind_scan, Turing.MultiTapeTM.initCfg, Turing.MultiTapeTM.runFrom, Turing.MultiTapeTM.runFrom_add, Turing.MultiTapeTM.runFrom_succ_eq_step', Turing.MultiTapeTM.step, Turing.moveInputPos_neg_of_ne_left, Turing.pairEncode}
\difficulty{5}
For every finite bit string $x$, the constructed machine computes the encoding of the
pair $(x,x)$—the string obtained by doubling each bit of $x$, then appending the
separator $01$ and a further copy of $x$—in exactly $4(|x|+1)$ steps.
\end{lemma}

\begin{lemma}[Scanning a suffix of the input]
\lean{Turing.FinTM.f2_scanCopy_suffix}
\label{Turing.FinTM.f2_scanCopy_suffix}
\leanok
\uses{Turing.Action.apply, Turing.Cfg.ext_zero_tapes, Turing.FinTM.f2_scanCfg, Turing.FinTM.f2_scanCfg_read, Turing.MultiTapeTM, Turing.MultiTapeTM.runFrom, Turing.MultiTapeTM.runFrom_succ_eq_step, Turing.MultiTapeTM.runFrom_zero, Turing.MultiTapeTM.step, Turing.moveInputPos_pos_of_ne_right}
\difficulty{5}
Let $M$ be a Turing machine with no work tapes, let $q$ be a state, and suppose that in
state $q$ each input bit is copied to the output while the input head moves right and
the state remains $q$, and a blank input symbol causes the machine to emit nothing and
halt. If Boolean words $x$, $\mathrm{pre}$, and $\mathrm{rest}$ satisfy
$x=\mathrm{pre}\,\mathrm{rest}$, then for every output word $o$, after
$|\mathrm{rest}|+1$ steps from the scan configuration on $x$ at position
$|\mathrm{pre}|$ in state $q$ with output $o$, the machine is halted at position $|x|$
with output $o$ followed by $\mathrm{rest}$.
\end{lemma}

\begin{lemma}[Scanning a true prefix]
\lean{Turing.FinTM.f2_scanTrues_run}
\label{Turing.FinTM.f2_scanTrues_run}
\leanok
\uses{Turing.Action.apply, Turing.Cfg.ext_zero_tapes, Turing.FinTM.f2_scanCfg, Turing.FinTM.f2_scanCfg_read, Turing.MultiTapeTM, Turing.MultiTapeTM.runFrom, Turing.MultiTapeTM.runFrom_succ_eq_step', Turing.MultiTapeTM.step, Turing.moveInputPos_pos_of_ne_right}
\difficulty{5}
Let $M$ be a Turing machine with no work tapes, let $q$ be a state, and let $e$ be a
Boolean bit. Suppose that whenever the input symbol is true, every work-tape reading
gives a transition that moves the input head right, emits false if $e$ is true and emits
nothing if $e$ is false, and remains in state $q$. For every Boolean word $x$ and every
$j\le |x|$ such that the first $j$ entries of $x$ are all true, after $j$ steps from the
scan configuration at position zero in state $q$ with empty output, the machine is in
state $q$ at position $j$; its output is the word of $j$ false entries if $e$ is true,
and is empty if $e$ is false.
\end{lemma}

\begin{lemma}[Cases for fixed-width increment]
\lean{Turing.FinTM.f2_incFixed_cases}
\label{Turing.FinTM.f2_incFixed_cases}
\leanok
\uses{Turing.incFixed}
\difficulty{4}
For every Boolean word $x$, either $x$ consists entirely of true bits and its
fixed-width increment is undefined, or there are a natural number $j$ and a Boolean word
$rest$ such that $x$ consists of $j$ true bits followed by false and then $rest$, and
its increment consists of $j$ false bits followed by true and then $rest$.
\end{lemma}

\begin{definition}[Fixed-increment machine]
\lean{Turing.FinTM.f2_incFixedTM}
\label{Turing.FinTM.f2_incFixedTM}
\leanok
\uses{Turing.FinTM, Turing.FinTM.controlAction}
There is a deterministic Boolean Turing machine with no work tapes and four states,
numbered $0$ through $3$, whose initial state is $0$. In state $0$, it moves right
without output on true and remains in state $0$; on false it moves left into state $1$;
on blank it halts without moving. In state $1$, it moves left while reading symbols, and
on blank moves right into state $2$. In state $2$, it moves right, changing true to
false and entering state $2$, or changing false to true and entering state $3$; on blank
it halts. In state $3$, it moves right while preserving each scanned symbol and halts on
blank.
\end{definition}

\begin{lemma}[Running time of fixed-width increment]
\lean{Turing.FinTM.f2_incFixed_computes}
\label{Turing.FinTM.f2_incFixed_computes}
\leanok
\uses{Turing.Action.apply, Turing.Cfg, Turing.Cfg.ext_zero_tapes, Turing.FinTM.ComputesInTime, Turing.FinTM.ComputesInTime.mono, Turing.FinTM.computesInTime_iff, Turing.FinTM.controlAction, Turing.FinTM.f2_incFixedTM, Turing.FinTM.f2_incFixed_cases, Turing.FinTM.f2_scanCfg, Turing.FinTM.f2_scanCfg_read, Turing.FinTM.f2_scanCopy_suffix, Turing.FinTM.f2_scanTrues_run, Turing.FinTM.rewind_scan, Turing.MultiTapeTM.initCfg, Turing.MultiTapeTM.runFrom, Turing.MultiTapeTM.runFrom_add, Turing.MultiTapeTM.runFrom_succ_eq_step', Turing.MultiTapeTM.step, Turing.incFixed, Turing.moveInputPos_neg_of_ne_left, Turing.moveInputPos_pos_of_ne_right}
\difficulty{6}
For every Boolean word $x$, the fixed-width increment machine computes the word given by
the fixed-width increment of $x$, using the empty word when that increment is undefined,
within $3(|x|+1)$ steps.
\end{lemma}

\begin{lemma}[One rightward scan step]
\lean{Turing.FinTM.f2_scanStep_right}
\label{Turing.FinTM.f2_scanStep_right}
\leanok
\uses{Turing.Action.apply, Turing.Cfg.ext_zero_tapes, Turing.FinTM.f2_scanCfg, Turing.FinTM.f2_scanCfg_read, Turing.MultiTapeTM, Turing.MultiTapeTM.step, Turing.moveInputPos_pos_of_ne_right}
\difficulty{3}
Let $M$ be a Turing machine with no work tapes, and let $x$ be a Boolean word. Suppose
that for every work-tape reading, the transition from state $q$ on the $i$th bit of $x$
moves the input head right, leaves the work tapes unchanged, emits the optional bit $e$,
and enters optional state $q'$. If $i<|x|$, then one step from the scan configuration at
position $i$ in state $q$ with output $o$ gives the scan configuration at position $i+1$
in state $q'$ with output $o$ followed by $e$ when present.
\end{lemma}

\begin{definition}[Pair-validity machine]
\lean{Turing.FinTM.f2_pairValidTM}
\label{Turing.FinTM.f2_pairValidTM}
\leanok
\uses{Turing.FinTM}
There is a deterministic Boolean Turing machine with no work tapes and state set
$\mathrm{Option}(\{\mathrm{false},\mathrm{true}\})$, initially in the state representing
no stored Boolean. From that state, on input $b$ it moves right, emits nothing, and
stores $b$. When it has stored $b$ and reads $c$, it moves right: if $b=c$, it emits
nothing and returns to the state with no stored Boolean; if $b\ne c$, it emits $\neg
b\land c$ and halts. On blank in either kind of state, it emits false and halts without
moving.
\end{definition}

\begin{lemma}[Pair validity on one block]
\lean{Turing.FinTM.f2_pairValid_block}
\label{Turing.FinTM.f2_pairValid_block}
\leanok
\uses{Turing.FinTM.f2_pairValidTM, Turing.FinTM.f2_scanCfg, Turing.FinTM.f2_scanStep_right, Turing.MultiTapeTM.runFrom, Turing.MultiTapeTM.step}
\difficulty{4}
Let $x$ be a finite bit string, with $x$ written as a prefix followed by bits $b,c$ and
a suffix. Starting in the parser's initial scanning state with its input head at the
position immediately after the prefix and empty output, the constructed machine after
two steps has its input head two positions farther along. If $b=c$, it remains in the
initial scanning state with empty output; otherwise it halts and outputs the single bit
that is true exactly when $b=0$ and $c=1$.
\end{lemma}

\begin{lemma}[Pair validity run bound]
\lean{Turing.FinTM.f2_pairValid_run}
\label{Turing.FinTM.f2_pairValid_run}
\leanok
\uses{Turing.Action.apply, Turing.FinTM.f2_pairValidTM, Turing.FinTM.f2_pairValid_block, Turing.FinTM.f2_scanCfg, Turing.FinTM.f2_scanCfg_read, Turing.FinTM.f2_scanStep_right, Turing.MultiTapeTM.runFrom, Turing.MultiTapeTM.runFrom_add, Turing.MultiTapeTM.runFrom_succ_eq_step, Turing.MultiTapeTM.runFrom_zero, Turing.MultiTapeTM.step, Turing.pairDecode}
\difficulty{6}
Let $x$ and $r$ be finite bit strings. For every prefix $p$ such that $x$ is the
concatenation of $p$ and $r$, there is a time $t \\le |r|+1$ such that the constructed
machine, started in its initial scanning state with input head immediately after $p$ and
empty output, has halted by time $t$ and output the one-bit indicator of whether the
pair parser succeeds on $r$.
\end{lemma}

\begin{lemma}[Running time of pair validation]
\lean{Turing.FinTM.f2_pairValid_computes}
\label{Turing.FinTM.f2_pairValid_computes}
\leanok
\uses{Turing.Cfg.ext_zero_tapes, Turing.FinTM.ComputesInTime, Turing.FinTM.ComputesInTime.mono, Turing.FinTM.computesInTime_iff, Turing.FinTM.f2_pairValidTM, Turing.FinTM.f2_pairValid_run, Turing.FinTM.f2_scanCfg, Turing.MultiTapeTM.initCfg, Turing.pairDecode}
\difficulty{4}
For every finite bit string $x$, the fixed pair-validation Turing machine halts within
$|x|+1$ steps and outputs the one-symbol word containing true exactly when the aligned
pair parser is defined on $x$, and false otherwise.
\end{lemma}

\begin{definition}[Pair-extraction machine]
\lean{Turing.FinTM.f2_pairExtractTM}
\label{Turing.FinTM.f2_pairExtractTM}
\leanok
\uses{Turing.FinTM}
For each pair of Boolean parameters $(f,s)$, there is a deterministic Boolean Turing
machine with one work tape and state set
$\mathrm{Option}(\mathrm{Option}(\{\mathrm{false},\mathrm{true}\})\oplus\mathrm{Fin}(3))$,
initially in the state representing no stored Boolean. Its transition function is the
following: in the initial state it moves right through the input, storing the first
symbol in its state, and halts if it encounters blank. With a stored Boolean $b$, it
moves right on the next input symbol $c$: if $b=c$, it writes $b$ on the work tape and
enters the initial state; if $b\ne c$ and $b$ is true, it halts; if $b\ne c$ and $b$ is
false, it moves the work-tape head left and enters phase $0$. In phase $0$, it moves the
work-tape head left over occupied cells, then right into phase $1$ at a blank cell. In
phase $1$, it moves the work-tape head right over occupied cells, writing the scanned
symbol to the output exactly when $f$ is true; at a blank cell it enters phase $2$. In
phase $2$, if $s$ is true, it moves right along the input, emitting each scanned symbol,
and halts on blank; if $s$ is false, it halts immediately. All transitions not specified
as emitting a symbol emit nothing, and omitted successor states halt.
\end{definition}

\begin{definition}[Extraction configuration]
\lean{Turing.FinTM.f2_extractCfg}
\label{Turing.FinTM.f2_extractCfg}
\leanok
\uses{Turing.Cfg, Turing.FinTM.bufferTape}
Given a Boolean input list $x$, a state
$q\in\mathrm{Option}(\mathrm{Option}(\{\mathrm{false},\mathrm{true}\})\oplus\mathrm{Fin}(3))$,
a natural number $i\le |x|$, a Boolean list $a$, an integer $z$, and an output list $o$,
the associated one-work-tape configuration has input-head position $i+1$, work-tape
contents given by the tape encoding of $a$, work-tape head position $z$, and output $o$.
\end{definition}

\begin{lemma}[Input Symbol in an Extractor Configuration]
\lean{Turing.FinTM.f2_extractCfg_read}
\label{Turing.FinTM.f2_extractCfg_read}
\leanok
\uses{Turing.Cfg.inputSymbol, Turing.FinTM.f2_extractCfg, Turing.FinTM.f2_scanCfg_read}
\difficulty{1}
Let $x$ and $a$ be Boolean words, let $q$ be an optional state in
$\operatorname{Option}(\mathbb{B}\oplus\mathrm{Fin}(3))$, let $i\in\mathbb{N}$ satisfy
$i\le |x|$, and let $z\in\mathbb{Z}$ and $out$ be a Boolean word. In the configuration
with input head at position $i+1$, work tape storing $a$, work-tape head at $z$, state
$q$, and output $out$, the scanned input symbol is the $i$th symbol of $x$, or no symbol
if $i$ is beyond its end.
\end{lemma}

\begin{lemma}[First Step of Pair Extraction]
\lean{Turing.FinTM.f2_extract_first}
\label{Turing.FinTM.f2_extract_first}
\leanok
\uses{Turing.Action.apply, Turing.FinTM.f2_extractCfg, Turing.FinTM.f2_extractCfg_read, Turing.FinTM.f2_pairExtractTM, Turing.MultiTapeTM.step, Turing.moveInputPos_pos_of_ne_right}
\difficulty{4}
Let $first,second$ be Boolean flags and let $x,pre,rest,a$ be Boolean words with
$x=pre\mathbin{++}[b]\mathbin{++}rest$. From the extractor configuration at position
$|pre|$, in its initial reading state, with tape word $a$, head at $|a|$, and empty
output, one step reaches the configuration at position $|pre|+1$ with the stored bit $b$
in the state and all other tape and output data unchanged.
\end{lemma}

\begin{lemma}[Two-Step Extraction of a Pair]
\lean{Turing.FinTM.f2_extract_block}
\label{Turing.FinTM.f2_extract_block}
\leanok
\uses{Turing.Action.apply, Turing.FinTM.bufferTape_append, Turing.FinTM.f2_extractCfg, Turing.FinTM.f2_extractCfg_read, Turing.FinTM.f2_extract_first, Turing.FinTM.f2_pairExtractTM, Turing.MultiTapeTM.runFrom, Turing.MultiTapeTM.step, Turing.moveInputPos, Turing.moveInputPos_pos_of_ne_right}
\difficulty{5}
Let $first,second\in\{\mathrm{false},\mathrm{true}\}$ and let Boolean words satisfy
$x=pre\mathbin{++}[b,c]\mathbin{++}rest$. Starting from the extractor configuration at
position $|pre|$ with state indicating that it is ready to read the first symbol, tape
word $a$, head at $|a|$, and empty output, after two steps the configuration is at
position $|pre|+2$ with tape word and head position unchanged and empty output if $b=c$;
if $b\ne c$ and $b$ is true, it is halted with those same tape and output contents; if
$b$ is false, it is in state $0$ with tape word $a$, head position $|a|-1$ (truncated at
zero), and empty output.
\end{lemma}

\begin{lemma}[Rewind to the Replay Start]
\lean{Turing.FinTM.f2_extract_rewind}
\label{Turing.FinTM.f2_extract_rewind}
\leanok
\uses{Turing.Action.apply, Turing.Cfg.workTapeSymbols, Turing.FinTM.bufferTape_left, Turing.FinTM.bufferTape_nat, Turing.FinTM.f2_extractCfg, Turing.FinTM.f2_pairExtractTM, Turing.MultiTapeTM.runFrom, Turing.MultiTapeTM.runFrom_succ_eq_step, Turing.MultiTapeTM.runFrom_zero, Turing.MultiTapeTM.step, Turing.moveInputPos_zero}
\difficulty{5}
Let $first,second$ be Boolean flags, let $x,a$ be Boolean words, and let $i\le |x|$. For
every $j\le |a|$, running the extractor machine for $j+1$ steps from state $0$ at input
position $i$, with tape word $a$, head at integer position $j-1$, and empty output,
reaches the replay state at the same input position, with tape head at the origin and
empty output.
\end{lemma}

\begin{lemma}[Replay Phase Configuration]
\lean{Turing.FinTM.f2_extract_replay}
\label{Turing.FinTM.f2_extract_replay}
\leanok
\uses{Turing.Action.apply, Turing.Cfg.workTapeSymbols, Turing.FinTM.bufferTape_nat, Turing.FinTM.f2_extractCfg, Turing.FinTM.f2_pairExtractTM, Turing.MultiTapeTM.runFrom, Turing.MultiTapeTM.runFrom_succ_eq_step', Turing.MultiTapeTM.step, Turing.moveInputPos_zero}
\difficulty{4}
Let $first,second$ be Boolean flags, let $x,a$ be Boolean words, and let $i\le |x|$. For
each $j\le |a|$, running the extractor machine for $j$ steps from the replay state at
input position $i$, with tape word $a$, head at the origin, and empty output, reaches
the same replay state and input position, with tape head at $j$ and output equal to the
first $j$ symbols of $a$ when $first$ is true and empty otherwise.
\end{lemma}

\begin{lemma}[End of Replay Phase]
\lean{Turing.FinTM.f2_extract_replay_finish}
\label{Turing.FinTM.f2_extract_replay_finish}
\leanok
\uses{Turing.Action.apply, Turing.Cfg.workTapeSymbols, Turing.FinTM.bufferTape_nat, Turing.FinTM.f2_extractCfg, Turing.FinTM.f2_extract_replay, Turing.FinTM.f2_pairExtractTM, Turing.MultiTapeTM.runFrom, Turing.MultiTapeTM.runFrom_succ_eq_step', Turing.MultiTapeTM.step, Turing.moveInputPos_zero}
\difficulty{3}
Let $first,second$ be Boolean flags, let $x,a$ be Boolean words, and let $i\le |x|$.
After $|a|+1$ steps from the replay state at input position $i$, with tape word $a$,
head at the origin, and empty output, the extractor is in state $2$ at the same input
position, with tape head at $|a|$ and output equal to $a$ if $first$ is true and empty
otherwise.
\end{lemma}

\begin{lemma}[Suffix extraction by a Turing machine]
\lean{Turing.FinTM.f2_extract_suffix}
\label{Turing.FinTM.f2_extract_suffix}
\leanok
\uses{Turing.Action.apply, Turing.FinTM.f2_extractCfg, Turing.FinTM.f2_extractCfg_read, Turing.FinTM.f2_pairExtractTM, Turing.MultiTapeTM.runFrom, Turing.MultiTapeTM.runFrom_succ_eq_step, Turing.MultiTapeTM.runFrom_zero, Turing.MultiTapeTM.step, Turing.moveInputPos_pos_of_ne_right}
\difficulty{5}
Let $first$ be a Boolean value, and let $x$, $rest$, and $a$ be finite Boolean words.
For every pair of Boolean words $pre$ and $out$ such that $x$ is the concatenation of
$pre$ and $rest$, running the pair-extraction machine with its second extraction enabled
for $|rest|+1$ steps from the configuration whose state is the second-extraction state,
whose input head is just after $pre$, whose work tape stores $a$ with its head at
position $|a|$, and whose output is $out$, yields a halted configuration with input head
just after $x$, the same work tape and head position, and output equal to the
concatenation of $out$ and $rest$.
\end{lemma}

\begin{lemma}[Completion of Pair Extraction]
\lean{Turing.FinTM.f2_extract_finish}
\label{Turing.FinTM.f2_extract_finish}
\leanok
\uses{Turing.Action.apply, Turing.FinTM.f2_extractCfg, Turing.FinTM.f2_extract_replay_finish, Turing.FinTM.f2_extract_rewind, Turing.FinTM.f2_extract_suffix, Turing.FinTM.f2_pairExtractTM, Turing.MultiTapeTM.runFrom, Turing.MultiTapeTM.runFrom_add, Turing.MultiTapeTM.runFrom_succ_eq_step', Turing.MultiTapeTM.step}
\difficulty{4}
Let $first,second$ be Boolean flags, let $x,a,pre,rest$ be Boolean words, and suppose
$x=pre\mathbin{++}rest$. Starting from the extractor configuration in state $0$, at
input position $|pre|$, with tape word $a$, tape head at $|a|-1$ (truncated at zero),
and empty output, there is a time $t\le 2|a|+|rest|+3$ at which the machine is halted
and its output is $(a\text{ if }first\text{, otherwise }[])\mathbin{++}(rest\text{ if
}second\text{, otherwise }[])$.
\end{lemma}

\begin{lemma}[Halting of the pair extraction machine]
\lean{Turing.FinTM.f2_extract_run}
\label{Turing.FinTM.f2_extract_run}
\leanok
\uses{Turing.Action.apply, Turing.FinTM.f2_extractCfg, Turing.FinTM.f2_extractCfg_read, Turing.FinTM.f2_extract_block, Turing.FinTM.f2_extract_finish, Turing.FinTM.f2_extract_first, Turing.FinTM.f2_pairExtractTM, Turing.MultiTapeTM.runFrom, Turing.MultiTapeTM.runFrom_add, Turing.MultiTapeTM.runFrom_succ_eq_step, Turing.MultiTapeTM.runFrom_zero, Turing.MultiTapeTM.step, Turing.pairDecode}
\difficulty{6}
Let $first$ and $second$ be Boolean flags, and let $x$ and $rest$ be Boolean words. For
every pair of Boolean words $pre$ and $a$ satisfying $x=pre\mathbin{++}rest$, the pair
extraction machine started with input head just after $pre$, work tape containing $a$
with its head at position $|a|$, and empty output reaches a halting configuration after
some number $t$ of steps with $t\le 3|rest|+2|a|+5$. Its output is $(a\mathbin{++}b$ if
$first$ is true, and the empty word otherwise$)\mathbin{++}(c$ if $second$ is true, and
the empty word otherwise$)$ when pair decoding of $rest$ yields $(b,c)$, and is empty
when that decoding is undefined.
\end{lemma}

\begin{lemma}[Pair extraction computation]
\lean{Turing.FinTM.f2_pairExtract_computes}
\label{Turing.FinTM.f2_pairExtract_computes}
\leanok
\uses{Turing.Cfg.init, Turing.FinTM.ComputesInTime, Turing.FinTM.ComputesInTime.mono, Turing.FinTM.computesInTime_iff, Turing.FinTM.f2_extractCfg, Turing.FinTM.f2_extract_run, Turing.FinTM.f2_pairExtractTM, Turing.MultiTapeTM.initCfg, Turing.pairDecode}
\difficulty{3}
Let $x$ be a finite bit string and let $a,b$ be Boolean flags. The constructed machine
outputs the first component of the parsed pair if $a$ is true, followed by the second
component if $b$ is true; components whose flags are false are omitted, and the output
is empty if parsing fails. It computes this output within $5(|x|+1)$ steps.
\end{lemma}

\begin{definition}[Catalog polynomial control states]
\lean{Turing.FinTM.f2_CatalogPolyControl}
\label{Turing.FinTM.f2_CatalogPolyControl}
\leanok
For natural numbers $c$ and $C$, define a finite type of control states with two
distinguished states, copy and setup; loop and rewind states indexed by
$i\in\{0,\ldots,c\}$; advance states indexed by $i\in\{0,\ldots,c+1\}$; and emit states
indexed by $j\in\{0,\ldots,C\}$.
\end{definition}

\begin{definition}[Bounded Integer Indicator]
\lean{Turing.FinTM.f2_catalogPolyTape}
\label{Turing.FinTM.f2_catalogPolyTape}
\leanok
For each natural number $q$, define a function on the integers that returns the optional
Boolean value $	ext{some}(	ext{true})$ at $z$ when $0 \le z < q$, and returns no value
otherwise.
\end{definition}

\begin{definition}[Single-tape head movement]
\lean{Turing.FinTM.f2_catalogPolyMove}
\label{Turing.FinTM.f2_catalogPolyMove}
\leanok
\uses{Turing.Action, Turing.FinTM.f2_CatalogPolyControl}
For natural numbers $c,C$, an index $i\in\{0,\ldots,c\}$, a direction $d\in\{-1,0,1\}$,
and a control state $s$ from the finite control set with copy, setup, loop, rewind,
advance, and emit states, define an action on $c+1$ work tapes that leaves the input
head stationary, writes nothing on every work tape, moves work-tape head $i$ by $d$ and
leaves all other work-tape heads stationary, emits no output, and transitions to state
$s$.
\end{definition}

\begin{definition}[Unary catalog machine]
\lean{Turing.FinTM.f2_catalogPolyUnaryTM}
\label{Turing.FinTM.f2_catalogPolyUnaryTM}
\leanok
\uses{Turing.FinTM, Turing.FinTM.f2_CatalogPolyControl, Turing.FinTM.f2_catalogPolyMove}
For natural numbers $c,C$, define a finite Turing machine over the Boolean alphabet with
$c+1$ work tapes and control states indexed by $c,C$. It starts in the copy state; its
transition function copies each input symbol onto every work tape while moving all work
heads right, then enters setup at the end of input. In setup it moves the first work
head right until it scans a blank, then starts a loop at index $c$; the loop and rewind
states inspect the indexed tape cell and move that head to rewind or advance the indexed
counter, while advance either resumes the loop at the corresponding index or halts. At
the loop's zero index it enters an emission state indexed by $C$, which emits one true
symbol per positive index before returning through advance; at index zero it emits
nothing and resumes advance.
\end{definition}

\begin{definition}[Catalog machine configuration]
\lean{Turing.FinTM.f2_catalogPolyCfg}
\label{Turing.FinTM.f2_catalogPolyCfg}
\leanok
\uses{Turing.Cfg, Turing.FinTM.f2_CatalogPolyControl, Turing.FinTM.f2_catalogPolyTape}
For natural numbers $c,C,q$, a Boolean input list $x$, a control state $s$ indexed by
$c,C$, a function $h$ assigning an integer position to each of the $c+1$ work tapes, and
a Boolean output list $o$, define the configuration with state $s$, input-head position
$|x|+1$ in the shifted indexing, each work tape containing true exactly at integer
positions $z$ with $0\le z<q$ and blank elsewhere, work-head positions given by $h$, and
output $o$.
\end{definition}

\begin{lemma}[Work tape head movement]
\lean{Turing.FinTM.f2_catalogPolyMove_apply}
\label{Turing.FinTM.f2_catalogPolyMove_apply}
\leanok
\uses{Turing.Action.apply, Turing.FinTM.f2_CatalogPolyControl, Turing.FinTM.f2_catalogPolyCfg, Turing.FinTM.f2_catalogPolyMove, Turing.moveInputPos_zero}
\difficulty{4}
Let $c,C$ be natural numbers, $x$ a finite Boolean input, $q$ a natural number, and
$s,s'$ control states. Given work-tape head positions $h$ and output word $o$, choose a
tape index $i$ and a movement $d\in\{-1,0,1\}$. Applying the action that keeps the input
head fixed, leaves all tape contents and the output unchanged, moves only head $i$ by
$d$, and sets the control state to $s'$ yields the configuration with state $s'$, head
positions updated by replacing $h(i)$ with $h(i)+d$, and all other data unchanged.
\end{lemma}

\begin{lemma}[Output phase emits a unary block]
\lean{Turing.FinTM.f2_catalogPoly_emit}
\label{Turing.FinTM.f2_catalogPoly_emit}
\leanok
\uses{Turing.Action.apply, Turing.FinTM.f2_catalogPolyCfg, Turing.FinTM.f2_catalogPolyUnaryTM, Turing.MultiTapeTM.runFrom, Turing.MultiTapeTM.runFrom_succ_eq_step, Turing.MultiTapeTM.runFrom_zero, Turing.MultiTapeTM.step}
\difficulty{4}
For all natural numbers $c,C,q$, every Boolean word $x$, every work-head position
function $h\colon\{0,\ldots,c\}\to\mathbb Z$, every Boolean output word $o$, and all
$j\le C$, starting in state $\mathrm{emit}(j)$ with input-head position $|x|+1$, tape
contents determined by $q$, and head positions $h$, the machine reaches after $j+1$
steps state $\mathrm{advance}(0)$ with the same head positions and output $o$ followed
by $j$ copies of true.
\end{lemma}

\begin{lemma}[Rewind phase restores a work head]
\lean{Turing.FinTM.f2_catalogPoly_rewind}
\label{Turing.FinTM.f2_catalogPoly_rewind}
\leanok
\uses{Turing.Action, Turing.Action.apply, Turing.Cfg.workTapeSymbols, Turing.FinTM.f2_CatalogPolyControl, Turing.FinTM.f2_catalogPolyCfg, Turing.FinTM.f2_catalogPolyMove_apply, Turing.FinTM.f2_catalogPolyTape, Turing.FinTM.f2_catalogPolyUnaryTM, Turing.MultiTapeTM.runFrom, Turing.MultiTapeTM.runFrom_succ_eq_step, Turing.MultiTapeTM.runFrom_zero, Turing.MultiTapeTM.step}
\difficulty{5}
For all natural numbers $c,C,q$, every Boolean word $x$, every work-head position
function $h\colon\{0,\ldots,c\}\to\mathbb Z$, every Boolean output word $o$, every tape
index $i\in\{0,\ldots,c\}$, and every $j\le q$, starting in state $\mathrm{rewind}(i)$
with head positions obtained from $h$ by setting position $i$ to $j-1$ leads after $j+1$
steps to state $\mathrm{advance}(i+1)$, with position $i$ reset to $0$ and all other
head positions as in $h$, and with output still $o$.
\end{lemma}

\begin{lemma}[Advance transition]
\lean{Turing.FinTM.f2_catalogPoly_advance}
\label{Turing.FinTM.f2_catalogPoly_advance}
\leanok
\uses{Turing.FinTM.f2_catalogPolyCfg, Turing.FinTM.f2_catalogPolyMove_apply, Turing.FinTM.f2_catalogPolyUnaryTM, Turing.MultiTapeTM.step}
\difficulty{3}
Let $c,C,q\in\mathbb{N}$, let $x$ and $o$ be finite Boolean words, let $h$ assign an
integer head position to each of the $c+1$ work tapes, and let $i$ be one of those
tapes. One step of the specified unary machine from the configuration with control state
$\mathrm{advance}(i)$ leaves the input-head position, work-tape contents, and output
word unchanged, moves work-tape head $i$ one cell to the right, leaves the other
work-tape head positions unchanged, and enters control state $\mathrm{loop}(i)$.
\end{lemma}

\begin{definition}[Recursive catalog cost]
\lean{Turing.FinTM.f2_catalogPolyCost}
\label{Turing.FinTM.f2_catalogPolyCost}
\leanok
For natural numbers $q,C$, define a function $T:\mathbb{N}\to\mathbb{N}$ recursively by
$T(0)=C+1$ and $T(r+1)=q(T(r)+2)+q+2$ for every $r\in\mathbb{N}$.
\end{definition}

\begin{lemma}[Loop phase output and head advance]
\lean{Turing.FinTM.f2_catalogPoly_loop}
\label{Turing.FinTM.f2_catalogPoly_loop}
\leanok
\uses{Turing.Action.apply, Turing.Cfg.workTapeSymbols, Turing.FinTM.f2_catalogPolyCfg, Turing.FinTM.f2_catalogPolyCost, Turing.FinTM.f2_catalogPolyMove_apply, Turing.FinTM.f2_catalogPolyTape, Turing.FinTM.f2_catalogPolyUnaryTM, Turing.FinTM.f2_catalogPoly_advance, Turing.FinTM.f2_catalogPoly_emit, Turing.FinTM.f2_catalogPoly_rewind, Turing.MultiTapeTM.runFrom, Turing.MultiTapeTM.runFrom_add, Turing.MultiTapeTM.runFrom_succ_eq_step, Turing.MultiTapeTM.runFrom_zero, Turing.MultiTapeTM.step}
\difficulty{7}
Let $c,C,q$ be natural numbers with $q>0$, let $x$ be a Boolean word, and let $i<c+1$.
Suppose $h\colon\{0,\ldots,c\}\to\mathbb Z$ satisfies $h(k)=0$ whenever $k\le i$. For
natural numbers $r,j$ with $j+r=q$ and any Boolean output word $o$, start in state
$\mathrm{loop}(i)$ with tape-head positions obtained from $h$ by setting position $i$ to
$j$. After $r(\mathrm{cost}(q,C,i)+2)+q+2$ steps, the machine is in state
$\mathrm{advance}(i+1)$ with head positions $h$ and output $o$ followed by $rCq^i$
copies of true.
\end{lemma}

\begin{lemma}[Extending the marked tape interval]
\lean{Turing.FinTM.f2_catalogPolyTape_write}
\label{Turing.FinTM.f2_catalogPolyTape_write}
\leanok
\uses{Turing.FinTM.f2_catalogPolyTape}
\difficulty{4}
For every natural number $q$, writing true at cell $q$ of the tape that contains true
exactly at the integer cells $z$ with $0\le z<q$ produces the tape that contains true
exactly at the cells $z$ with $0\le z<q+1$.
\end{lemma}

\begin{lemma}[Polynomial cost bound]
\lean{Turing.FinTM.f2_catalogPolyCost_le}
\label{Turing.FinTM.f2_catalogPolyCost_le}
\leanok
\uses{Turing.FinTM.f2_catalogPolyCost}
\difficulty{4}
For natural numbers $q,C$ with $q>0$, and every natural number $r$, the recursively
defined cost satisfies $\mathrm{cost}(q,C,r)\le (C+1+5r)q^r$, where
$\mathrm{cost}(q,C,0)=C+1$ and $\mathrm{cost}(q,C,r+1)=q(\mathrm{cost}(q,C,r)+2)+q+2$.
\end{lemma}

\begin{definition}[Copy-stage configuration]
\lean{Turing.FinTM.f2_catalogPolyCopyCfg}
\label{Turing.FinTM.f2_catalogPolyCopyCfg}
\leanok
\uses{Turing.Cfg, Turing.FinTM.f2_CatalogPolyControl, Turing.FinTM.f2_catalogPolyTape}
For natural numbers $c,C$, a Boolean input list $x$, and a natural number $i$ with $i\le
|x|$, define the configuration in the copy state with input-head position $i+1$ in the
shifted indexing, each of the $c+1$ work tapes containing true exactly at integer
positions $z$ with $0\le z<i$ and blank elsewhere, each work-head at position $i$, and
empty output.
\end{definition}

\begin{lemma}[Copy phase of the unary machine]
\lean{Turing.FinTM.f2_catalogPoly_copy}
\label{Turing.FinTM.f2_catalogPoly_copy}
\leanok
\uses{Turing.Action.apply, Turing.Cfg.init, Turing.Cfg.inputSymbol, Turing.FinTM.f2_catalogPolyCopyCfg, Turing.FinTM.f2_catalogPolyTape, Turing.FinTM.f2_catalogPolyTape_write, Turing.FinTM.f2_catalogPolyUnaryTM, Turing.MultiTapeTM.initCfg, Turing.MultiTapeTM.runFrom, Turing.MultiTapeTM.runFrom_succ_eq_step', Turing.MultiTapeTM.step, Turing.inputSymbolInner, Turing.moveInputPos, Turing.moveInputPos_pos_of_ne_right}
\difficulty{5}
For all natural numbers $c,C$, every Boolean word $x$, and every natural number $i\le
|x|$, after $i$ steps the unary machine with parameters $c,C$, started on $x$, is in its
copy state with empty output, input-head position $i+1$, every one of its $c+1$ work
heads at position $i$, and each work tape containing true exactly at the integer
positions from $0$ through $i-1$.
\end{lemma}

\begin{lemma}[Setup phase initializes the loop]
\lean{Turing.FinTM.f2_catalogPoly_setup}
\label{Turing.FinTM.f2_catalogPoly_setup}
\leanok
\uses{Turing.Action.apply, Turing.Cfg.workTapeSymbols, Turing.FinTM.f2_catalogPolyCfg, Turing.FinTM.f2_catalogPolyTape, Turing.FinTM.f2_catalogPolyUnaryTM, Turing.MultiTapeTM.runFrom, Turing.MultiTapeTM.runFrom_succ_eq_step, Turing.MultiTapeTM.runFrom_zero, Turing.MultiTapeTM.step}
\difficulty{4}
For all natural numbers $c,C,q$, every Boolean word $x$, and every $j\le q$, start in
state $\mathrm{setup}$ with every work-head position equal to $j-1$, tape contents
determined by $q$, and empty output. After $j+1$ steps the machine is in state
$\mathrm{loop}(c)$, with every work head at position $0$ and output still empty.
\end{lemma}

\begin{lemma}[Initial copy and setup phases]
\lean{Turing.FinTM.f2_catalogPoly_start}
\label{Turing.FinTM.f2_catalogPoly_start}
\leanok
\uses{Turing.Action.apply, Turing.Cfg.inputSymbol, Turing.FinTM.f2_catalogPolyCfg, Turing.FinTM.f2_catalogPolyCopyCfg, Turing.FinTM.f2_catalogPolyTape_write, Turing.FinTM.f2_catalogPolyUnaryTM, Turing.FinTM.f2_catalogPoly_copy, Turing.FinTM.f2_catalogPoly_setup, Turing.MultiTapeTM.initCfg, Turing.MultiTapeTM.runFrom, Turing.MultiTapeTM.runFrom_add, Turing.MultiTapeTM.runFrom_succ_eq_step', Turing.MultiTapeTM.step, Turing.moveInputPos_zero}
\difficulty{5}
For all natural numbers $c,C$ and every Boolean word $x$, the unary machine with
parameters $c,C$, started in its initial configuration on $x$, reaches after $2(|x|+1)$
steps state $\mathrm{loop}(c)$ with every work-head position equal to $0$, each work
tape containing true exactly at positions $0$ through $|x|$, input-head position
$|x|+1$, and empty output.
\end{lemma}

\begin{lemma}[Unary output computation]
\lean{Turing.FinTM.f2_catalogPoly_unary_computes}
\label{Turing.FinTM.f2_catalogPoly_unary_computes}
\leanok
\uses{Turing.Action.apply, Turing.FinTM.ComputesFunInTime, Turing.FinTM.ComputesInTime, Turing.FinTM.ComputesInTime.mono, Turing.FinTM.computesInTime_iff, Turing.FinTM.f2_catalogPolyCfg, Turing.FinTM.f2_catalogPolyCost, Turing.FinTM.f2_catalogPolyCost_le, Turing.FinTM.f2_catalogPolyUnaryTM, Turing.FinTM.f2_catalogPoly_loop, Turing.FinTM.f2_catalogPoly_start, Turing.MultiTapeTM.runFrom, Turing.MultiTapeTM.runFrom_add, Turing.MultiTapeTM.runFrom_succ_eq_step', Turing.MultiTapeTM.step}
\difficulty{5}
For all $c,C\in\mathbb N$ and every finite Boolean word $x$, the machine specified by
$c$ and $C$ halts after exactly $(C+5(c+1)+4)(|x|+1)^{c+1}$ steps and outputs the word
consisting of exactly $C(|x|+1)^{c+1}$ copies of true.
\end{lemma}

\begin{definition}[Head-position bounds]
\lean{Turing.FinTM.f2_polyHeads}
\label{Turing.FinTM.f2_polyHeads}
\leanok
\uses{Turing.FinTM.f2_CatalogPolyControl}
For natural numbers $d,C,q$, an optional control state indexed by $d,C$, and integer
work-head positions $h_j$ for $j=0,\ldots,d$, define the predicate that constrains these
positions as follows. If the state is absent, require $-1\le h_j\le q$ for every $j$. In
loop state $i$, require $0\le h_j\le q$ for every $j$, and $h_j<q$ whenever $j\ne i$; in
rewind state $i$, require $-1\le h_j<q$ for every $j$, and $h_j\ge0$ whenever $j\ne i$.
In advance or emit states, require $0\le h_j<q$ for every $j$; in all other states the
predicate is false.
\end{definition}

\begin{lemma}[Bounds on polynomial machine heads]
\lean{Turing.FinTM.f2_polyHeads_bounds}
\label{Turing.FinTM.f2_polyHeads_bounds}
\leanok
\uses{Turing.FinTM.f2_CatalogPolyControl, Turing.FinTM.f2_polyHeads}
\difficulty{4}
Let $d,C,q$ be natural numbers, let $s$ be either absent or a control stage (copy,
setup, loop indexed by $i<d+1$, rewind indexed by $i<d+1$, advance indexed by $i<d+2$,
or emit indexed by $j<C+1$), and let $h$ assign an integer head position to each
$j<d+1$. If the positions satisfy the bounds prescribed by the stage—when $s$ is absent,
$-1 \\le h(j) \\le q$; in a loop stage indexed by $i$, $0 \\le h(j) \\le q$ and $h(j)<q$
for $j \\ne i$; in a rewind stage indexed by $i$, $-1 \\le h(j)<q$ and $h(j)\\ge0$ for
$j \\ne i$; or in an advance or emit stage, $0 \\le h(j)<q$—then every head position
satisfies $-1 \\le h(j) \\le q$.
\end{lemma}

\begin{lemma}[Preservation of the tape and head bounds]
\lean{Turing.FinTM.f2_poly_step}
\label{Turing.FinTM.f2_poly_step}
\leanok
\uses{Turing.Action.apply, Turing.Cfg, Turing.Cfg.workTapeSymbols, Turing.FinTM.f2_CatalogPolyControl, Turing.FinTM.f2_catalogPolyMove, Turing.FinTM.f2_catalogPolyTape, Turing.FinTM.f2_catalogPolyUnaryTM, Turing.FinTM.f2_polyHeads, Turing.MultiTapeTM.step}
\difficulty{6}
Let $d,C,q$ be natural numbers with $q>0$, let $x$ be a Boolean word, and let $c$ be a
configuration of the unary catalog machine with parameters $d,C$ on input $x$. Suppose
every work tape equals the tape containing true exactly at integer positions $z$ with
$0\le z<q$, and suppose the work-head positions satisfy the bounds specified by the
machine's control state: in the halted state they lie between $-1$ and $q$ inclusive; in
loop state $i$ they lie between $0$ and $q$ inclusive, with every head other than $i$
strictly below $q$; in rewind state $i$ they lie between $-1$ inclusive and $q$
exclusive, with every head other than $i$ nonnegative; in advance or emit states they
lie between $0$ inclusive and $q$ exclusive; and in copy or setup states the condition
is false. After one step, every work tape still has that content and the resulting state
and head positions satisfy the corresponding bounds.
\end{lemma}

\begin{lemma}[Work head displacement bound]
\lean{Turing.FinTM.f2_head_steps}
\label{Turing.FinTM.f2_head_steps}
\leanok
\uses{Turing.Cfg, Turing.MultiTapeTM, Turing.MultiTapeTM.runFrom, Turing.MultiTapeTM.runFrom_succ_eq_step', Turing.MultiTapeTM.workTapePos_step_le}
\difficulty{3}
Let $M$ be a $k$-work-tape Turing machine over the Boolean alphabet, let $c$ be a
configuration of $M$ on a Boolean input $x$, let $t\in\mathbb N$, and let $i$ be a
work-tape index. After $t$ steps, the position of work-tape head $i$ lies between its
initial position minus $t$ and its initial position plus $t$, inclusive.
\end{lemma}

\begin{lemma}[Space bound for the unary catalog machine]
\lean{Turing.FinTM.f2_poly_space}
\label{Turing.FinTM.f2_poly_space}
\leanok
\uses{Turing.FinTM.f2_catalogPolyCfg, Turing.FinTM.f2_catalogPolyTape, Turing.FinTM.f2_catalogPolyUnaryTM, Turing.FinTM.f2_catalogPoly_start, Turing.FinTM.f2_head_steps, Turing.FinTM.f2_polyHeads, Turing.FinTM.f2_polyHeads_bounds, Turing.FinTM.f2_poly_step, Turing.MultiTapeTM.initCfg, Turing.MultiTapeTM.runFrom, Turing.MultiTapeTM.runFrom_add, Turing.MultiTapeTM.runFrom_succ_eq_step', Turing.MultiTapeTM.spaceUsed, Turing.MultiTapeTM.spaceUsedByTape, Turing.MultiTapeTM.visitedByTapeHead}
\difficulty{6}
For all natural numbers $d,C,t$ and every Boolean word $x$, the total number of
work-tape cells visited by the unary catalog machine with parameters $d,C$, started on
$x$ and run for $t$ steps, is at most $5(d+1)(|x|+1)$.
\end{lemma}

\begin{definition}[Boolean counter increment]
\lean{Turing.FinTM.f2_counterInc}
\label{Turing.FinTM.f2_counterInc}
\leanok
Define a function on finite Boolean lists recursively: the empty list maps to the
singleton list containing true; a list beginning with false maps to the list beginning
with true and having the same remaining tail; and a list beginning with true maps to the
list beginning with false followed by the increment of its tail.
\end{definition}

\begin{definition}[Initial run of true bits]
\lean{Turing.FinTM.f2_counterCarry}
\label{Turing.FinTM.f2_counterCarry}
\leanok
For a Boolean list $b$, define its initial run length of true bits recursively: it is
one plus the initial run length of the tail when $b$ begins with true, and zero
otherwise.
\end{definition}

\begin{lemma}[Counter increment potential]
\lean{Turing.FinTM.f2_counterInc_potential}
\label{Turing.FinTM.f2_counterInc_potential}
\leanok
\uses{Turing.FinTM.f2_counterCarry, Turing.FinTM.f2_counterInc}
\difficulty{4}
For every finite Boolean list $b$, the number of true entries after incrementing the
counter, plus the number of leading true entries before incrementing, equals the
original number of true entries plus one.
\end{lemma}

\begin{lemma}[Binary increment on bit lists]
\lean{Turing.FinTM.f2_counterInc_bits}
\label{Turing.FinTM.f2_counterInc_bits}
\leanok
\uses{Turing.FinTM.f2_counterInc}
\difficulty{5}
For every natural number $n$, incrementing the list of binary digits of $n$ by the
recursive Boolean-list operation gives the binary digit list of $n+1$.
\end{lemma}

\begin{lemma}[Length and carry bounds]
\lean{Turing.FinTM.f2_counterInc_length}
\label{Turing.FinTM.f2_counterInc_length}
\leanok
\uses{Turing.FinTM.f2_counterCarry, Turing.FinTM.f2_counterInc}
\difficulty{4}
For every finite Boolean list $b$, incrementing the binary counter produces a list of
length at most $|b|+1$, and the number of leading true bits of $b$ is at most the length
of the incremented list.
\end{lemma}

\begin{definition}[Counter tape update action]
\lean{Turing.FinTM.f2_counterBump}
\label{Turing.FinTM.f2_counterBump}
\leanok
\uses{Turing.Action}
For a head movement $d\in\{-1,0,+1\}$ and a scanned tape value
$w\in\{\text{blank},\text{false},\text{true}\}$, define a one-step action for a
one-work-tape Boolean machine with states $0,1,2,3$. If $w=\text{true}$, the action
moves the input head by $d$, writes false on the work tape, moves its head right, emits
no output, and enters state $1$; otherwise it moves the input head by $d$, writes true,
moves the work-tape head left, emits no output, and enters state $2$.
\end{definition}

\begin{definition}[Counter Turing machine]
\lean{Turing.FinTM.f2_counterTM}
\label{Turing.FinTM.f2_counterTM}
\leanok
\uses{Turing.FinTM, Turing.FinTM.f2_counterBump}
Define a one-work-tape deterministic Turing machine over the Boolean alphabet with
states $0,1,2,3$ and initial state $0$. In state $0$, on a blank input cell it stays
put, does not write or move its work-tape head, emits nothing, and enters state $3$; on
a nonblank input cell it uses the counter tape update action with input-head movement
right. In state $1$ it uses that update action with input-head movement zero. In state
$2$, it stays on the input and, if the work cell is blank, moves the work head right
without writing and enters state $0$; otherwise it moves the work head left without
writing and remains in state $2$. In state $3$, if the work cell is blank it halts
without moving or writing; otherwise it moves the work head right, writes the scanned
bit to the output, and remains in state $3$.
\end{definition}

\begin{definition}[Finite Boolean tape lookup]
\lean{Turing.FinTM.f2_counterTape}
\label{Turing.FinTM.f2_counterTape}
\leanok
For a Boolean list $b$ and an integer position $z$, define the tape value to be blank
when $z<0$; otherwise it is the list entry at index $z$, or blank if that index is
beyond the list.
\end{definition}

\begin{definition}[Counter machine configuration]
\lean{Turing.FinTM.f2_counterCfg}
\label{Turing.FinTM.f2_counterCfg}
\leanok
\uses{Turing.Cfg, Turing.FinTM.f2_counterTape}
Given a Boolean input list $x$, a state $q\in\{0,1,2,3\}$, a valid input-head position
$p\in\{0,\ldots,|x|+1\}$, an integer work-head position $z$, and Boolean lists $b$ and
$o$, define a configuration in state $q$ with input-head position $p$, work-tape
contents given by the finite Boolean tape determined by $b$, work-head position $z$, and
output $o$.
\end{definition}

\begin{lemma}[Reading the counter tape]
\lean{Turing.FinTM.f2_counterTape_read}
\label{Turing.FinTM.f2_counterTape_read}
\leanok
\uses{Turing.FinTM.f2_counterTape}
\difficulty{4}
For finite Boolean lists $u$ and $b$, the counter tape formed from their concatenation,
read at the nonnegative position $|u|$, contains the first entry of $b$, or no entry if
$b$ is empty.
\end{lemma}

\begin{lemma}[Writing to the counter tape]
\lean{Turing.FinTM.f2_counterTape_write}
\label{Turing.FinTM.f2_counterTape_write}
\leanok
\uses{Turing.FinTM.f2_counterTape, Turing.FinTM.f2_counterTape_read}
\difficulty{5}
For Boolean lists $u$ and $v$, define their counter tape on integer positions by
assigning no value at negative positions and the Boolean entry at each nonnegative
position, when that entry exists. For any Boolean value $c$, changing the tape of $u$
followed by $v$ at position $|u|$ to the value $c$ gives the tape of $u$ followed by $c$
and then the tail of $v$.
\end{lemma}

\begin{lemma}[One counter carry step]
\lean{Turing.FinTM.f2_counter_carry_step}
\label{Turing.FinTM.f2_counter_carry_step}
\leanok
\uses{Turing.Action.apply, Turing.FinTM.f2_counterBump, Turing.FinTM.f2_counterCfg, Turing.FinTM.f2_counterTM, Turing.FinTM.f2_counterTape, Turing.FinTM.f2_counterTape_read, Turing.FinTM.f2_counterTape_write, Turing.MultiTapeTM.step, Turing.moveInputPos_zero}
\difficulty{4}
Let $x$, $u$, and $b$ be finite Boolean lists, and let $p$ be an input-head position
among the $|x|+2$ positions for input $x$. Starting in state $1$ with the input head at
$p$, the work-tape head at position $|u|$, tape contents given by the counter tape for
$u$ concatenated with $b$, and empty output, one counter-machine step yields the
following configuration: if the first entry of $b$ is true, the state remains $1$, the
work-tape head moves to $|u|+1$, and the tape becomes the counter tape for $u$ followed
by false and then the tail of $b$; otherwise the state becomes $2$, the work-tape head
moves to $|u|-1$, and the tape becomes the counter tape for $u$ followed by true and
then the tail of $b$. In both cases the input head remains at $p$ and the output remains
empty.
\end{lemma}

\begin{lemma}[F2 Counter Carry]
\lean{Turing.FinTM.f2_counter_carry}
\label{Turing.FinTM.f2_counter_carry}
\leanok
\uses{Turing.FinTM.f2_counterCarry, Turing.FinTM.f2_counterCfg, Turing.FinTM.f2_counterInc, Turing.FinTM.f2_counterTM, Turing.FinTM.f2_counter_carry_step, Turing.MultiTapeTM.runFrom, Turing.MultiTapeTM.runFrom_succ_eq_step, Turing.MultiTapeTM.runFrom_zero}
\difficulty{5}
Let $x$ and $b$ be finite Boolean lists, and let $p$ be an input-head position among the
$|x|+2$ positions associated with input $x$. For every finite Boolean list $u$, starting
the counter machine in state $1$ with input-head position $p$, work-tape head at $|u|$,
tape contents specified by the counter tape for $u$ concatenated with $b$, and empty
output, after $c(b)+1$ steps it is in state $2$, with input-head position $p$, work-tape
head at $|u|+c(b)-1$, tape contents specified by the counter tape for $u$ concatenated
with the increment of $b$, and empty output, where $c(b)$ is the number of leading true
entries of $b$.
\end{lemma}

\begin{lemma}[F2 Counter Rewind]
\lean{Turing.FinTM.f2_counter_rewind}
\label{Turing.FinTM.f2_counter_rewind}
\leanok
\uses{Turing.Action.apply, Turing.Cfg.workTapeSymbols, Turing.FinTM.f2_counterCfg, Turing.FinTM.f2_counterTM, Turing.FinTM.f2_counterTape, Turing.MultiTapeTM.runFrom, Turing.MultiTapeTM.runFrom_succ_eq_step, Turing.MultiTapeTM.runFrom_zero, Turing.MultiTapeTM.step, Turing.moveInputPos_zero}
\difficulty{6}
For every Boolean list $x$, valid input-head position $p$ for $x$, and Boolean list
$bs$, for each natural number $j\le |bs|$, running the counter machine for $j+1$ steps
from state $2$, input-head position $p$, work-tape head position $j-1$, tape containing
$bs$, and empty output yields state $0$, the same input-head position $p$, work-tape
head position $0$, tape containing $bs$, and empty output.
\end{lemma}

\begin{lemma}[F2 Counter Initial Step Agreement]
\lean{Turing.FinTM.f2_counter_start}
\label{Turing.FinTM.f2_counter_start}
\leanok
\uses{Turing.Action.apply, Turing.Cfg.inputSymbol, Turing.FinTM.f2_counterBump, Turing.FinTM.f2_counterCfg, Turing.FinTM.f2_counterTM, Turing.FinTM.f2_counterTape, Turing.MultiTapeTM.step, Turing.inputSymbolInner, Turing.moveInputPos, Turing.moveInputPos_pos_of_ne_right, Turing.moveInputPos_zero}
\difficulty{4}
For every Boolean list $x$, natural number $i<|x|$, and Boolean list $bs$, the one-step
successors of the configurations with state $0$ and input-head position $i+1$ and with
state $1$ and input-head position $i+2$ are equal, when both have work-tape head at
position $0$, tape containing $bs$, and empty output.
\end{lemma}

\begin{lemma}[F2 Counter Increment]
\lean{Turing.FinTM.f2_counter_increment}
\label{Turing.FinTM.f2_counter_increment}
\leanok
\uses{Turing.FinTM.f2_counterCarry, Turing.FinTM.f2_counterCfg, Turing.FinTM.f2_counterInc, Turing.FinTM.f2_counterInc_length, Turing.FinTM.f2_counterTM, Turing.FinTM.f2_counter_carry, Turing.FinTM.f2_counter_rewind, Turing.FinTM.f2_counter_start, Turing.MultiTapeTM.runFrom, Turing.MultiTapeTM.runFrom_add, Turing.MultiTapeTM.runFrom_succ_eq_step}
\difficulty{4}
For every Boolean list $x$, natural number $i<|x|$, and Boolean list $bs$, running the
counter machine for $2\,\mathrm{carry}(bs)+2$ steps from state $0$, input-head position
$i+1$, work-tape head position $0$, tape containing $bs$, and empty output yields state
$0$, input-head position $i+2$, work-tape head position $0$, tape containing the binary
increment of $bs$, and empty output.
\end{lemma}

\begin{lemma}[Counter configuration after counting]
\lean{Turing.FinTM.f2_counter_count}
\label{Turing.FinTM.f2_counter_count}
\leanok
\uses{Turing.FinTM.f2_counterCarry, Turing.FinTM.f2_counterCfg, Turing.FinTM.f2_counterInc_bits, Turing.FinTM.f2_counterInc_potential, Turing.FinTM.f2_counterTM, Turing.FinTM.f2_counterTape, Turing.FinTM.f2_counter_increment, Turing.MultiTapeTM.initCfg, Turing.MultiTapeTM.runFrom, Turing.MultiTapeTM.runFrom_add}
\difficulty{4}
For every finite Boolean list $x$ and every natural number $i\le |x|$, there exists a
time $t$ such that $t+2$ times the number of true bits in the binary representation of
$i$ is at most $4i$, and after $t$ steps from the machine's initial configuration on
$x$, the machine is in state $0$, its input head is at position $i+1$, its work-tape
head is at position $0$, its tape contains the binary representation of $i$, and its
output is empty.
\end{lemma}

\begin{lemma}[F2 Counter Output at Each Step]
\lean{Turing.FinTM.f2_counter_emit_run}
\label{Turing.FinTM.f2_counter_emit_run}
\leanok
\uses{Turing.Action.apply, Turing.Cfg.workTapeSymbols, Turing.FinTM.f2_counterCfg, Turing.FinTM.f2_counterTM, Turing.FinTM.f2_counterTape, Turing.MultiTapeTM.runFrom, Turing.MultiTapeTM.runFrom_succ_eq_step', Turing.MultiTapeTM.step, Turing.moveInputPos_zero}
\difficulty{5}
For every Boolean list $x$, every valid input-head position $p$ for $x$, and every
Boolean list $bs$, at each natural time $i \le |bs|$, the counter machine run for $i$
steps from state $3$, input-head position $p$, work-tape head position $0$, tape
containing $bs$, and empty output is in state $3$, with the same input-head and
work-tape head positions and tape, and with output equal to the first $i$ entries of
$bs$.
\end{lemma}

\begin{lemma}[F2 Counter Halting Output]
\lean{Turing.FinTM.f2_counter_emit}
\label{Turing.FinTM.f2_counter_emit}
\leanok
\uses{Turing.Action.apply, Turing.Cfg.workTapeSymbols, Turing.FinTM.f2_counterCfg, Turing.FinTM.f2_counterTM, Turing.FinTM.f2_counterTape, Turing.FinTM.f2_counter_emit_run, Turing.MultiTapeTM.runFrom, Turing.MultiTapeTM.runFrom_succ_eq_step', Turing.MultiTapeTM.step}
\difficulty{3}
For every Boolean list $x$, every valid input-head position $p$ for input $x$, and every
Boolean list $bs$, after $|bs|+1$ steps from state $3$, input-head position $p$,
work-tape head position $0$, tape containing $bs$, and empty output, the counter machine
is halted and its output is $bs$.
\end{lemma}

\begin{lemma}[Running time for binary length]
\lean{Turing.FinTM.f2_counter_computes}
\label{Turing.FinTM.f2_counter_computes}
\leanok
\uses{Turing.Action.apply, Turing.Cfg.inputSymbol, Turing.FinTM.ComputesFunInTime, Turing.FinTM.ComputesInTime, Turing.FinTM.ComputesInTime.mono, Turing.FinTM.f2_counterCfg, Turing.FinTM.f2_counterTM, Turing.FinTM.f2_counter_count, Turing.FinTM.f2_counter_emit, Turing.MultiTapeTM.initCfg, Turing.MultiTapeTM.runFrom, Turing.MultiTapeTM.runFrom_add, Turing.MultiTapeTM.runFrom_succ_eq_step', Turing.MultiTapeTM.step, Turing.length_bits_le_self}
\difficulty{4}
The finite Turing machine computes, on every finite input word $x$ over its symbol
alphabet, the binary expansion of the length of $x$ in time $5(|x|+1)$, where the binary
expansion is represented by the list of binary digits returned by the natural-number bit
function.
\end{lemma}

\begin{lemma}[Space bounds during counter execution]
\lean{Turing.FinTM.f2_counter_count_space}
\label{Turing.FinTM.f2_counter_count_space}
\leanok
\uses{Turing.Cfg.init, Turing.FinTM.f2_counterCarry, Turing.FinTM.f2_counterCfg, Turing.FinTM.f2_counterInc_bits, Turing.FinTM.f2_counterInc_length, Turing.FinTM.f2_counterInc_potential, Turing.FinTM.f2_counterTM, Turing.FinTM.f2_counterTape, Turing.FinTM.f2_counter_increment, Turing.FinTM.f2_head_steps, Turing.MultiTapeTM.initCfg, Turing.MultiTapeTM.runFrom, Turing.MultiTapeTM.runFrom_add}
\difficulty{5}
For every finite Boolean word $x$ and every natural number $i \le |x|$, there is a time
$t$ with $t + 2\,\mathrm{count}_{\mathrm{true}}(i.bits) \le 4i$ such that, after $t$
steps, the counter machine started on $x$ is in state $0$, its input head is at position
$i+1$, its work-tape head is at position $0$, its work tape contains the binary
expansion of $i$, and its output is empty. At every time $u \le t$, the work-tape head
position lies between $-2(\mathrm{size}(|x|)+1)$ and $2(\mathrm{size}(|x|)+1)$,
inclusive.
\end{lemma}

\begin{lemma}[F2 Counter Work-Tape Head Bound]
\lean{Turing.FinTM.f2_counter_heads}
\label{Turing.FinTM.f2_counter_heads}
\leanok
\uses{Turing.Action.apply, Turing.Cfg.inputSymbol, Turing.FinTM.f2_counterCfg, Turing.FinTM.f2_counterTM, Turing.FinTM.f2_counter_count_space, Turing.FinTM.f2_counter_emit, Turing.FinTM.f2_head_steps, Turing.MultiTapeTM.initCfg, Turing.MultiTapeTM.runFrom, Turing.MultiTapeTM.runFrom_add, Turing.MultiTapeTM.runFrom_of_halt, Turing.MultiTapeTM.runFrom_succ_eq_step', Turing.MultiTapeTM.step}
\difficulty{5}
For every Boolean list $x$, every natural time $u$, and the sole work tape, the head
position after running the counter machine on $x$ for $u$ steps lies between
$-2(\mathrm{size}(|x|)+1)$ and $2(\mathrm{size}(|x|)+1)$, inclusive.
\end{lemma}

\begin{lemma}[F2 Counter Space Bound]
\lean{Turing.FinTM.f2_counter_space}
\label{Turing.FinTM.f2_counter_space}
\leanok
\uses{Turing.FinTM.f2_counterTM, Turing.FinTM.f2_counter_heads, Turing.MultiTapeTM.initCfg, Turing.MultiTapeTM.spaceUsed, Turing.MultiTapeTM.spaceUsedByTape, Turing.MultiTapeTM.visitedByTapeHead}
\difficulty{4}
For every Boolean list $x$ and natural time $t$, the total number of distinct work-tape
cells visited by the counter machine on input $x$ through time $t$ is at most
$5(\mathrm{size}(|x|)+1)$.
\end{lemma}

\begin{definition}[Length-machine action]
\lean{Turing.FinTM.f2_lenAction}
\label{Turing.FinTM.f2_lenAction}
\leanok
\uses{Turing.Action, Turing.FinTM}
Let $M$ be a finite Boolean Turing machine with $k$ work tapes. For input-head movement
$m$, work-head movement $d$, optional Boolean output $b$, and optional successor state
$q$ of the enlarged machine, define an action with $k+1$ work tapes: it moves the input
head by $m$, writes nothing on every work tape, leaves the first $k$ work heads
stationary, moves the last work head by $d$, emits $b$ when present, and passes to $q$,
halting when $q$ is absent.
\end{definition}

\begin{definition}[Pair-counting machine construction]
\lean{Turing.FinTM.f2_pairCountTM}
\label{Turing.FinTM.f2_pairCountTM}
\leanok
\uses{Turing.FinTM, Turing.FinTM.controlAction, Turing.FinTM.f2_lenAction, Turing.captureAction}
Let $M$ be a finite Turing machine over the Boolean alphabet. The construction defines a
machine with one more work tape than $M$ and state set
$S\oplus(\mathrm{Fin}(4)\oplus\mathrm{Option}(\mathrm{Bool}))$, where $S$ is the state
set of $M$. Its initial state is $M$’s initial state in the left summand; there it uses
$M$’s transition on the original work tapes and the specified capture action, while
transitions in the added states are given by the case distinction on the state, input
symbol, and extra work-tape symbol in the construction.
\end{definition}

\begin{definition}[Length-machine configuration]
\lean{Turing.FinTM.f2_lenCfg}
\label{Turing.FinTM.f2_lenCfg}
\leanok
\uses{Turing.Cfg, Turing.FinTM, Turing.FinTM.f2_pairCountTM, Turing.captureCfg}
Let $M$ be a finite Boolean Turing machine, let $x$ be a Boolean input list, and let $c$
be a configuration of $M$ on $x$. For an optional state $q$ of the enlarged machine, an
index $i$ with $i\le |x|$, and a natural number $r$, define a configuration of the
enlarged machine on $x$ whose state is $q$, input-head position is $i+1$, and positions
on the original $k$ work tapes are those inherited from $c$; the additional work-head
position is $r-1$. Its tape contents and output are obtained from $c$ by the capture
construction.
\end{definition}

\begin{lemma}[Input symbol at a prescribed position]
\lean{Turing.FinTM.f2_lenCfg_read}
\label{Turing.FinTM.f2_lenCfg_read}
\leanok
\uses{Turing.Cfg, Turing.Cfg.inputSymbol, Turing.FinTM, Turing.FinTM.f2_lenCfg, Turing.FinTM.f2_pairCountTM, Turing.FinTM.inputSymbol_at}
\difficulty{1}
For a Boolean finite Turing machine $M$, a configuration $c$ on input $x$, a state $q$
of the augmented machine or no state, and natural numbers $i,r$ with $i\le |x|$, the
configuration formed from $c$ with input head at position $i+1$ scans exactly the $i$-th
symbol of $x$, or no symbol if $i\ge |x|$.
\end{lemma}

\begin{lemma}[Effect of a length-counting action]
\lean{Turing.FinTM.f2_lenAction_apply}
\label{Turing.FinTM.f2_lenAction_apply}
\leanok
\uses{Turing.Action.apply, Turing.Cfg, Turing.FinTM, Turing.FinTM.f2_lenAction, Turing.FinTM.f2_lenCfg, Turing.FinTM.f2_pairCountTM, Turing.moveInputPos}
\difficulty{4}
Let $M$ be a finite Turing machine over the Boolean alphabet, let $x$ be a Boolean
input, and let $c$ be a configuration of $M$ on $x$. Choose optional states $q,q'$,
natural numbers $i,j,r,s$ with $i,j\le |x|$, signs $m,d\in\{-1,0,1\}$, and an optional
bit $b$. Suppose moving the input head from position $i+1$ by $m$ gives position $j+1$,
and $(r-1)+d=s-1$. Then applying the length-counting action that moves the input head by
$m$, moves the extra work-tape head by $d$, emits $b$ if present, and enters state $q'$,
to the length-counting configuration built from $c$ at input position $i$, state $q$,
and extra-tape position $r$, gives the corresponding configuration at input position
$j$, state $q'$, and extra-tape position $s$, with output equal to the one-bit word $b$
if present and empty otherwise.
\end{lemma}

\begin{lemma}[Halting after checking a suffix length]
\lean{Turing.FinTM.f2_lenSuffix_run}
\label{Turing.FinTM.f2_lenSuffix_run}
\leanok
\uses{Turing.Action.apply, Turing.Cfg, Turing.Cfg.workTapeSymbols, Turing.FinTM, Turing.FinTM.bufferTape, Turing.FinTM.bufferTape_left, Turing.FinTM.f2_lenAction, Turing.FinTM.f2_lenAction_apply, Turing.FinTM.f2_lenCfg, Turing.FinTM.f2_lenCfg_read, Turing.FinTM.f2_pairCountTM, Turing.MultiTapeTM.runFrom, Turing.MultiTapeTM.runFrom_add, Turing.MultiTapeTM.runFrom_succ_eq_step, Turing.MultiTapeTM.runFrom_zero, Turing.MultiTapeTM.step, Turing.captureCfg, Turing.moveInputPos_pos_of_ne_right}
\difficulty{6}
Let $M$ be a finite Turing machine over the Boolean alphabet, let $x$ and $w$ be Boolean
words, and let $c$ be a configuration of $M$ on input $x$. For every decomposition
$x=uv$ and every natural number $r$ with $r\le |c.output|$, the augmented machine
started from the configuration determined by $c$, $u$, and $r$, in its state numbered
$3$, halts within at most $|v|+1$ steps and has output consisting of the single Boolean
value that is true exactly when $|v|\le r$.
\end{lemma}

\begin{lemma}[Reading the first bit of a block]
\lean{Turing.FinTM.f2_lenParse_first}
\label{Turing.FinTM.f2_lenParse_first}
\leanok
\uses{Turing.Action.apply, Turing.Cfg, Turing.FinTM, Turing.FinTM.f2_lenAction_apply, Turing.FinTM.f2_lenCfg, Turing.FinTM.f2_lenCfg_read, Turing.FinTM.f2_pairCountTM, Turing.MultiTapeTM.step, Turing.moveInputPos_pos_of_ne_right}
\difficulty{3}
Let $x=\mathrm{pre}\,b\,\mathrm{rest}$ be a Boolean word, with $b$ a bit, and let $c$ be
any configuration of a Boolean finite Turing machine $M$. One step of the augmented
machine, started in its pair-reading state with the input head at position
$|\mathrm{pre}|+1$, records $b$ in its parser state and moves the input head to position
$|\mathrm{pre}|+2$.
\end{lemma}

\begin{lemma}[Parsing one aligned bit pair]
\lean{Turing.FinTM.f2_lenParse_block}
\label{Turing.FinTM.f2_lenParse_block}
\leanok
\uses{Turing.Action.apply, Turing.Cfg, Turing.FinTM, Turing.FinTM.f2_lenAction_apply, Turing.FinTM.f2_lenCfg, Turing.FinTM.f2_lenCfg_read, Turing.FinTM.f2_lenParse_first, Turing.FinTM.f2_pairCountTM, Turing.MultiTapeTM.runFrom, Turing.MultiTapeTM.step, Turing.moveInputPos, Turing.moveInputPos_pos_of_ne_right, Turing.moveInputPos_zero}
\difficulty{5}
Let $x=\mathrm{pre}\,b\,d\,\mathrm{rest}$ be a Boolean word, where $b,d$ are bits, and
let $c$ be any configuration of a Boolean finite Turing machine $M$ and $r$ a natural
number. Starting in the parser state at position $|\mathrm{pre}|$, after two steps the
augmented machine is in the parser state at position $|\mathrm{pre}|+2$ if $b=d$; if
$b\ne d$ and $b=1$, it is halted at position $|\mathrm{pre}|+1$ with output $[0]$; and
if $b\ne d$ and $b=0$, it is in state $3$ at position $|\mathrm{pre}|+2$.
\end{lemma}

\begin{lemma}[Bounded parsing and length check]
\lean{Turing.FinTM.f2_lenParse_run}
\label{Turing.FinTM.f2_lenParse_run}
\leanok
\uses{Turing.Action.apply, Turing.Cfg, Turing.FinTM, Turing.FinTM.f2_lenAction, Turing.FinTM.f2_lenCfg, Turing.FinTM.f2_lenCfg_read, Turing.FinTM.f2_lenParse_block, Turing.FinTM.f2_lenParse_first, Turing.FinTM.f2_lenSuffix_run, Turing.FinTM.f2_pairCountTM, Turing.MultiTapeTM.runFrom, Turing.MultiTapeTM.runFrom_add, Turing.MultiTapeTM.runFrom_succ_eq_step, Turing.MultiTapeTM.runFrom_zero, Turing.MultiTapeTM.step, Turing.captureCfg, Turing.pairDecode}
\difficulty{6}
Let $M$ be a Boolean finite Turing machine, let $x$ be a Boolean input, and let $c$ be a
configuration of $M$ on $x$. For every pair of Boolean words
$\mathrm{pre},\mathrm{rest}$ with $x=\mathrm{pre}\,\mathrm{rest}$ and every natural
number $r\le |c.output|$, the machine obtained by augmenting $M$ with the pair parser,
started on the first symbol of $\mathrm{rest}$ with the extra tape containing
$\mathrm{pre}\,c.output$ and positioned at cell $r-1$, halts within at most
$|\mathrm{rest}|+1$ steps and outputs the single bit that is $1$ exactly when the pair
parser decodes $\mathrm{rest}$ as $(a,b)$ for some $a,b$ with $|b|\le r$, and is $0$
otherwise.
\end{lemma}

\begin{lemma}[Bounded input-head rewind]
\lean{Turing.FinTM.f2_catalogRewind}
\label{Turing.FinTM.f2_catalogRewind}
\leanok
\uses{Turing.Cfg, Turing.FinTM.controlAction, Turing.FinTM.controlAction_apply, Turing.FinTM.moveInputPos_neg_val, Turing.FinTM.rewind_scan, Turing.MultiTapeTM, Turing.MultiTapeTM.runFrom, Turing.MultiTapeTM.runFrom_add, Turing.MultiTapeTM.step, Turing.moveInputPos}
\difficulty{3}
Let $x$ be a Boolean input, let $M$ be a machine with $k$ work tapes and state set $S$,
and let $s,t \in S$ and $d \in \mathrm{Option}(S)$. Suppose that from state $s$, on
every input and work-tape reading, $M$ moves the input head left and enters state $t$
without changing any work tape or output; and from state $t$, it moves left and stays in
$t$ whenever it reads a symbol, while on a blank input cell it moves right and enters
state $d$. For every configuration $c$ on input $x$ whose state is $s$, there is an $r
\le c$'s input-head position plus $2$ such that after $r$ steps the configuration has
state $d$ and input-head position $1$, with all other configuration components
unchanged.
\end{lemma}

\begin{lemma}[Initialization of the pair-scanning computation]
\lean{Turing.FinTM.f2_lenStart}
\label{Turing.FinTM.f2_lenStart}
\leanok
\uses{Turing.Action.apply, Turing.Cfg, Turing.Cfg.init, Turing.FinTM, Turing.FinTM.ComputesInTime, Turing.FinTM.ComputesInTime.output_unique, Turing.FinTM.computesInTime_iff, Turing.FinTM.f2_catalogRewind, Turing.FinTM.f2_lenAction, Turing.FinTM.f2_lenCfg, Turing.FinTM.f2_pairCountTM, Turing.MultiTapeTM.initCfg, Turing.MultiTapeTM.runFrom, Turing.MultiTapeTM.runFrom_add, Turing.MultiTapeTM.runFrom_succ_eq_step, Turing.MultiTapeTM.runFrom_zero, Turing.MultiTapeTM.step, Turing.captureCfg, Turing.capture_run, Turing.moveInputPos_zero}
\difficulty{6}
Let $M$ be a finite Turing machine over the Boolean alphabet, let $x,w$ be Boolean
words, and let $T$ be a natural number. If $M$ computes output $w$ from input $x$ within
$T$ steps, then there are a natural number $t\le T+|x|+4$ and a configuration $c$ of $M$
on input $x$ whose output is $w$, such that the machine obtained by augmenting $M$ with
one work tape reaches after $t$ steps a configuration with the control state in the
auxiliary state component (specifically, its nested option component is empty), input
head at the first input position, original work tapes and their head positions as in
$c$, and the extra work tape holding $w$ with its head at position $|w|-1$; its output
tape is empty.
\end{lemma}

\begin{lemma}[Pair count computation]
\lean{Turing.FinTM.f2_pairCount_computes}
\label{Turing.FinTM.f2_pairCount_computes}
\leanok
\uses{Turing.FinTM, Turing.FinTM.ComputesFunInTime, Turing.FinTM.ComputesInTime, Turing.FinTM.ComputesInTime.mono, Turing.FinTM.computesInTime_iff, Turing.FinTM.f2_lenParse_run, Turing.FinTM.f2_lenStart, Turing.FinTM.f2_pairCountTM, Turing.MultiTapeTM.runFrom_add, Turing.pairDecode}
\difficulty{3}
Let $M$ be a finite Turing machine over $\\{0,1\\}$, let $g$ map finite bit strings to
finite bit strings, and let $T : \\mathbb{N} \\to \\mathbb{N}$. Suppose that on every
bit string $x$, $M$ computes $g(x)$ within $T(|x|)$ steps. Then the constructed machine
computes, on every $x$, the one-bit output that is $1$ exactly when the pair parser
succeeds with second component $b$ satisfying $|b| \\le |g(x)|$, and is $0$ otherwise,
within $T(|x|)+2|x|+5$ steps.
\end{lemma}

\begin{lemma}[Pair encoding and decoding]
\lean{Turing.FinTM.f2_catalogPair_inverse}
\label{Turing.FinTM.f2_catalogPair_inverse}
\leanok
\uses{Turing.pairDecode, Turing.pairEncode}
\difficulty{4}
Let $x$ be a finite Boolean word. For any finite Boolean words $a$ and $v$, if decoding
$x$ returns the pair $(a,v)$, then $x$ is the encoding of $(a,v)$.
\end{lemma}

\begin{lemma}[Space bounded by computation time]
\lean{Turing.FinTM.f2_space_of_time}
\label{Turing.FinTM.f2_space_of_time}
\leanok
\uses{Turing.FinTM, Turing.FinTM.ComputesInTime, Turing.FinTM.computesInTime_iff, Turing.MultiTapeTM.initCfg, Turing.MultiTapeTM.runFrom, Turing.MultiTapeTM.runFrom_add, Turing.MultiTapeTM.runFrom_of_halt, Turing.MultiTapeTM.spaceUsed, Turing.MultiTapeTM.visitedByTapeHead}
\difficulty{5}
Let $M$ be a finite Turing machine over Boolean symbols, let $x,y$ be Boolean words, and
let $T$ be a natural number. If $M$ halts after exactly $T$ steps on input $x$, outputs
$y$, and uses some space amount by step $T$, then for every $t\\in\\mathbb N$ its space
usage from the initial configuration on $x$ through step $t$ is at most $k(T+1)$, where
$k$ is the number of work tapes of $M$.
\end{lemma}

\begin{theorem}[Identity computation with bounded space]
\lean{Turing.FinTM.computesFunInTime_id_spaceUsed}
\label{Turing.FinTM.computesFunInTime_id_spaceUsed}
\leanok
\uses{Turing.Action.apply, Turing.Cfg.inputSymbol, Turing.Cfg.workTapeSymbols, Turing.FinTM, Turing.FinTM.ComputesFunInTime, Turing.FinTM.ComputesInTime, Turing.FinTM.ComputesInTime.mono, Turing.FinTM.f2_idTM, Turing.FinTM.f2_idTM_run, Turing.MultiTapeTM.initCfg, Turing.MultiTapeTM.runFrom, Turing.MultiTapeTM.runFrom_succ_eq_step', Turing.MultiTapeTM.spaceUsed, Turing.MultiTapeTM.spaceUsed_zero_tapes_eq_zero, Turing.MultiTapeTM.step}
\difficulty{5}
There are a finite Turing machine $M$ over the Boolean alphabet and a constant
$c\in\mathbb{N}$ such that, on every Boolean input $x$, $M$ outputs $x$ and halts within
$c(|x|+1)$ steps. For every input $x$ and every $t\in\mathbb{N}$, the space used by $M$
by time $t$ is at most $c$.
\end{theorem}

\begin{theorem}[Constant-word computation with bounded space]
\lean{Turing.FinTM.computesFunInTime_const_spaceUsed}
\label{Turing.FinTM.computesFunInTime_const_spaceUsed}
\leanok
\uses{Turing.Cfg.init, Turing.FinTM, Turing.FinTM.ComputesFunInTime, Turing.FinTM.ComputesInTime, Turing.FinTM.ComputesInTime.mono, Turing.FinTM.emit_halts, Turing.FinTM.f2_constTM, Turing.MultiTapeTM.initCfg, Turing.MultiTapeTM.spaceUsed, Turing.MultiTapeTM.spaceUsed_zero_tapes_eq_zero}
\difficulty{3}
For every fixed Boolean word $w$, there are a finite Turing machine $M$ over the Boolean
alphabet and a constant $c\in\mathbb{N}$ such that, on every Boolean input $x$, $M$
outputs $w$ and halts within $c(|x|+1)$ steps. For every input $x$ and every
$t\in\mathbb{N}$, the space used by $M$ by time $t$ is at most $c$.
\end{theorem}

\begin{theorem}[Linear-time prefixing with bounded space]
\lean{Turing.FinTM.computesFunInTime_prepend_spaceUsed}
\label{Turing.FinTM.computesFunInTime_prepend_spaceUsed}
\leanok
\uses{Turing.FinTM, Turing.FinTM.ComputesFunInTime, Turing.FinTM.ComputesInTime.mono, Turing.FinTM.f2_catalogPrefixTM, Turing.FinTM.f2_catalogPrefixTM_computes, Turing.MultiTapeTM.initCfg, Turing.MultiTapeTM.spaceUsed, Turing.MultiTapeTM.spaceUsed_zero_tapes_eq_zero}
\difficulty{3}
For every Boolean word $w$, there exist a finite Turing machine $M$ and a natural number
$c$ such that, on every Boolean input $x$, $M$ halts within $c(|x|+1)$ steps and outputs
the concatenation $w x$. Moreover, for every Boolean word $x$ and every time $t$, the
computation of $M$ on $x$ up to time $t$ uses at most $c$ work-tape cells.
\end{theorem}

\begin{theorem}[Binary Length with a Space Bound]
\lean{Turing.FinTM.computesFunInTime_lengthBits_spaceUsed}
\label{Turing.FinTM.computesFunInTime_lengthBits_spaceUsed}
\leanok
\uses{Turing.FinTM, Turing.FinTM.ComputesFunInTime, Turing.FinTM.f2_counterTM, Turing.FinTM.f2_counter_computes, Turing.FinTM.f2_counter_space, Turing.MultiTapeTM.initCfg, Turing.MultiTapeTM.spaceUsed}
\difficulty{1}
There exist a finite Turing machine $M$ over the Boolean alphabet and a constant
$c\in\mathbb{N}$ such that, on every Boolean word $x$, $M$ halts within $c(|x|+1)$ steps
and outputs the little-endian binary representation of $|x|$. For every such $x$ and
every $t\in\mathbb{N}$, its computation through step $t$ uses at most
$c\bigl(\operatorname{size}(|x|)+1\bigr)$ work-tape cells, where $\operatorname{size}$
is the natural-number size function.
\end{theorem}

\begin{theorem}[Unary polynomial output with a space bound]
\lean{Turing.FinTM.computesFunInTime_polyUnary_spaceUsed}
\label{Turing.FinTM.computesFunInTime_polyUnary_spaceUsed}
\leanok
\uses{Turing.FinTM, Turing.FinTM.ComputesFunInTime, Turing.FinTM.ComputesInTime.mono, Turing.FinTM.computesFunInTime_const_spaceUsed, Turing.FinTM.f2_catalogPolyUnaryTM, Turing.FinTM.f2_catalogPoly_unary_computes, Turing.FinTM.f2_poly_space, Turing.MultiTapeTM.initCfg, Turing.MultiTapeTM.spaceUsed}
\difficulty{4}
For every $C,e\in\mathbb{N}$, there exist a finite Turing machine $M$ over the Boolean
alphabet and a constant $c\in\mathbb{N}$ such that, on every input $x$, $M$ outputs a
string of $C(|x|+1)^e$ true bits within $c(|x|+1)^{e+1}$ steps. For every $x$ and every
$t\in\mathbb{N}$, the total work space used through step $t$ is at most $c(|x|+1)$.
\end{theorem}

\begin{theorem}[Polynomial binary output with a space bound]
\lean{Turing.FinTM.computesFunInTime_polyBits_spaceUsed}
\label{Turing.FinTM.computesFunInTime_polyBits_spaceUsed}
\leanok
\uses{Turing.FinTM, Turing.FinTM.ComputesFunInTime, Turing.FinTM.ComputesInTime.mono, Turing.FinTM.bufferedCompTM, Turing.FinTM.bufferedCompTM_computesInTime, Turing.FinTM.computesFunInTime_const_spaceUsed, Turing.FinTM.f2_catalogPolyUnaryTM, Turing.FinTM.f2_catalogPoly_unary_computes, Turing.FinTM.f2_counterTM, Turing.FinTM.f2_counter_computes, Turing.FinTM.f2_space_of_time, Turing.MultiTapeTM.initCfg, Turing.MultiTapeTM.spaceUsed}
\difficulty{6}
For every $C,e\in\mathbb{N}$, there exist a finite Turing machine $M$ over the Boolean
alphabet and a constant $c\in\mathbb{N}$ such that, on every input $x$, $M$ outputs the
binary representation of $C(|x|+1)^e$, with bits listed from least significant to most
significant, within $c(|x|+1)^{e+1}$ steps. For every $x$ and every $t\in\mathbb{N}$,
the total work space used through step $t$ is at most $c\bigl(C(|x|+1)^e+1\bigr)$.
\end{theorem}

\begin{theorem}[Fixed-prefix pairing in linear time and constant space]
\lean{Turing.FinTM.computesFunInTime_pairEncodeFixed_spaceUsed}
\label{Turing.FinTM.computesFunInTime_pairEncodeFixed_spaceUsed}
\leanok
\uses{Turing.FinTM, Turing.FinTM.ComputesFunInTime, Turing.FinTM.computesFunInTime_prepend_spaceUsed, Turing.MultiTapeTM.initCfg, Turing.MultiTapeTM.spaceUsed, Turing.pairEncode}
\difficulty{1}
For every fixed Boolean word $\alpha$, there are a finite Turing machine $M$ and a
constant $c\in\mathbb{N}$ such that, on every Boolean input $x$, $M$ halts within
$c(|x|+1)$ steps and outputs the pairing of $\alpha$ with $x$. At every time
$t\in\mathbb{N}$, the computation of $M$ on any input uses at most $c$ work-tape cells.
\end{theorem}

\begin{theorem}[First-component extraction with linear space]
\lean{Turing.FinTM.computesFunInTime_pairFst_spaceUsed}
\label{Turing.FinTM.computesFunInTime_pairFst_spaceUsed}
\leanok
\uses{Turing.FinTM, Turing.FinTM.ComputesFunInTime, Turing.FinTM.ComputesInTime.mono, Turing.FinTM.f2_pairExtractTM, Turing.FinTM.f2_pairExtract_computes, Turing.FinTM.f2_space_of_time, Turing.MultiTapeTM.initCfg, Turing.MultiTapeTM.spaceUsed, Turing.pairDecode}
\difficulty{4}
There are a finite Turing machine $M$ and a constant $c\in\mathbb{N}$ such that, on
every Boolean input $x$, $M$ halts within $c(|x|+1)$ steps and outputs the first
component parsed from $x$ by the aligned pair parser, or the empty word if parsing
fails. At every time $t\in\mathbb{N}$, its computation on $x$ uses at most $c(|x|+1)$
work-tape cells.
\end{theorem}

\begin{theorem}[Pair decoding with linear time and space]
\lean{Turing.FinTM.computesFunInTime_pairSnd_spaceUsed}
\label{Turing.FinTM.computesFunInTime_pairSnd_spaceUsed}
\leanok
\uses{Turing.FinTM, Turing.FinTM.ComputesFunInTime, Turing.FinTM.ComputesInTime.mono, Turing.FinTM.f2_pairExtractTM, Turing.FinTM.f2_pairExtract_computes, Turing.FinTM.f2_space_of_time, Turing.MultiTapeTM.initCfg, Turing.MultiTapeTM.spaceUsed, Turing.pairDecode}
\difficulty{4}
There exist a finite Turing machine $M$ over the Boolean alphabet and a constant
$c\in\mathbb{N}$ such that, on every Boolean string $x$, $M$ halts with output equal to
the second component returned by the pair parser when parsing succeeds, and the empty
string when parsing fails, within $c(|x|+1)$ steps. For every $x$ and every
$t\in\mathbb{N}$, the total work space used by $M$ through step $t$ is at most
$c(|x|+1)$.
\end{theorem}

\begin{theorem}[Pair-code validity with bounded space]
\lean{Turing.FinTM.computesFunInTime_pairValid_spaceUsed}
\label{Turing.FinTM.computesFunInTime_pairValid_spaceUsed}
\leanok
\uses{Turing.FinTM, Turing.FinTM.ComputesFunInTime, Turing.FinTM.f2_pairValidTM, Turing.FinTM.f2_pairValid_computes, Turing.MultiTapeTM.initCfg, Turing.MultiTapeTM.spaceUsed, Turing.MultiTapeTM.spaceUsed_zero_tapes_eq_zero, Turing.pairDecode}
\difficulty{3}
There exist a finite Turing machine $M$ over the Boolean alphabet and a constant
$c\in\mathbb{N}$ such that, on every Boolean string $x$, $M$ halts within $c(|x|+1)$
steps and outputs the one-bit string whose bit is true exactly when the pair parser
succeeds on $x$. For every $x$ and every $t\in\mathbb{N}$, the total work space used by
$M$ through step $t$ is at most $c$.
\end{theorem}

\begin{theorem}[Pair Concatenation with a Space Bound]
\lean{Turing.FinTM.computesFunInTime_pairConcat_spaceUsed}
\label{Turing.FinTM.computesFunInTime_pairConcat_spaceUsed}
\leanok
\uses{Turing.FinTM, Turing.FinTM.ComputesFunInTime, Turing.FinTM.ComputesInTime.mono, Turing.FinTM.f2_pairExtractTM, Turing.FinTM.f2_pairExtract_computes, Turing.FinTM.f2_space_of_time, Turing.MultiTapeTM.initCfg, Turing.MultiTapeTM.spaceUsed, Turing.pairDecode}
\difficulty{4}
There exist a finite Turing machine $M$ over the Boolean alphabet and a constant
$c\in\mathbb{N}$ such that, on every Boolean word $z$, $M$ halts within $c(|z|+1)$ steps
and outputs the concatenation of the two components returned by the aligned pair parser
when it succeeds, and outputs the empty word when it fails. For every Boolean word $z$
and every $t\in\mathbb{N}$, the computation through step $t$ uses at most $c(|z|+1)$
work-tape cells.
\end{theorem}

\begin{theorem}[Pair Duplication with a Space Bound]
\lean{Turing.FinTM.computesFunInTime_pairDup_spaceUsed}
\label{Turing.FinTM.computesFunInTime_pairDup_spaceUsed}
\leanok
\uses{Turing.FinTM, Turing.FinTM.ComputesFunInTime, Turing.FinTM.f2_pairDupTM, Turing.FinTM.f2_pairDup_computes, Turing.MultiTapeTM.initCfg, Turing.MultiTapeTM.spaceUsed, Turing.MultiTapeTM.spaceUsed_zero_tapes_eq_zero, Turing.pairEncode}
\difficulty{3}
There exist a finite Turing machine $M$ over the Boolean alphabet and a constant
$c\in\mathbb{N}$ such that, on every Boolean word $x$, $M$ halts within $c(|x|+1)$ steps
and outputs the pair encoding of $(x,x)$: each bit of $x$ doubled, followed by the
separator $01$, followed by $x$. For every Boolean word $x$ and every $t\in\mathbb{N}$,
the computation through step $t$ uses at most $c$ work-tape cells.
\end{theorem}

\begin{lemma}[Unary Output in Polynomial Time]
\lean{Turing.FinTM.f2_unary_sharp}
\label{Turing.FinTM.f2_unary_sharp}
\leanok
\uses{Turing.FinTM, Turing.FinTM.ComputesFunInTime, Turing.FinTM.ComputesInTime.mono, Turing.FinTM.computesFunInTime_const_spaceUsed, Turing.FinTM.f2_catalogPolyUnaryTM, Turing.FinTM.f2_catalogPoly_unary_computes}
\difficulty{4}
For every pair of natural numbers $C$ and $e$, there are a finite Turing machine $M$
over the Boolean alphabet and a natural number $a$ such that, on every natural number
$n$, $M$ computes the word consisting of $C(n+1)^e$ copies of true from an input of
length $n$ within $a\bigl((n+1)^e+n+1\bigr)$ steps.
\end{lemma}

\begin{lemma}[Length of the decoded first component]
\lean{Turing.FinTM.f2_first_length}
\label{Turing.FinTM.f2_first_length}
\leanok
\uses{Turing.FinTM.f2_catalogPair_inverse, Turing.length_pairEncode, Turing.pairDecode}
\difficulty{4}
For every Boolean word $x$, the length of the first component returned by pair decoding
of $x$, taken to be the empty word if decoding is undefined, is at most $|x|$.
\end{lemma}

\begin{theorem}[Pair-length check with a space bound]
\lean{Turing.FinTM.computesFunInTime_pairLenCheck_spaceUsed}
\label{Turing.FinTM.computesFunInTime_pairLenCheck_spaceUsed}
\leanok
\uses{Turing.FinTM, Turing.FinTM.ComputesFunInTime, Turing.FinTM.ComputesInTime.mono, Turing.FinTM.bufferedCompTM, Turing.FinTM.bufferedCompTM_computesInTime, Turing.FinTM.f2_first_length, Turing.FinTM.f2_pairCountTM, Turing.FinTM.f2_pairCount_computes, Turing.FinTM.f2_pairExtractTM, Turing.FinTM.f2_pairExtract_computes, Turing.FinTM.f2_space_of_time, Turing.FinTM.f2_unary_sharp, Turing.MultiTapeTM.initCfg, Turing.MultiTapeTM.spaceUsed, Turing.pairDecode}
\difficulty{6}
For all natural numbers $C$ and $e$, there are a finite Turing machine $M$ and a
constant $c\in\mathbb{N}$ such that, on every Boolean input $x$, $M$ halts within
$c(|x|+1)^{e+1}$ steps and outputs a single bit. The bit is true exactly when the
aligned pair parser returns components $(a,b)$ with $|b|\le C(|a|+1)^e$, and false if
parsing fails. At every time $t\in\mathbb{N}$, the computation on $x$ uses at most
$c\big((|x|+1)^e+|x|+1\big)$ work-tape cells.
\end{theorem}

\begin{definition}[Raw strip machine]
\lean{Turing.FinTM.f2_rawStripTM}
\label{Turing.FinTM.f2_rawStripTM}
\leanok
\uses{Turing.FinTM}
This is a finite Turing machine over the Boolean alphabet with one work tape and four
states, starting in state $0$. In state $0$, it moves right across the input, copying
each scanned bit to the work tape and remaining in state $0$; at the right blank it
moves the work head left and enters state $1$. In state $1$, it halts at a blank work
cell, or erases the scanned bit, moves left, and enters state $2$ if that bit was true
and state $1$ if it was false. In state $2$, it moves left across nonblank work cells;
at a blank cell it moves right and enters state $3$. In state $3$, it outputs each
scanned work bit while moving right, and halts at a blank cell.
\end{definition}

\begin{definition}[Strip configuration]
\lean{Turing.FinTM.f2_stripCfg}
\label{Turing.FinTM.f2_stripCfg}
\leanok
\uses{Turing.Cfg, Turing.FinTM.bufferTape}
For a Boolean input list $x$, a state $q$ that is either one of the four states or the
halting state, an input-head index $i$ with $i \leq |x|$, a Boolean list $w$, an integer
$h$, and a Boolean output list, this is the configuration whose input-head position is
$i+1$. Its work tape contains the buffer represented by $w$, its work-head position is
$h$, and its output tape contains the given output list.
\end{definition}

\begin{lemma}[Erasing the final buffered bit]
\lean{Turing.FinTM.f2_catalogBuffer_erase}
\label{Turing.FinTM.f2_catalogBuffer_erase}
\leanok
\uses{Turing.FinTM.bufferTape, Turing.FinTM.bufferTape_append}
\difficulty{3}
For every Boolean word $w$ and bit $b$, erasing the cell at position $|w|$ from the
buffer tape containing $w$ followed by $b$ yields exactly the buffer tape containing
$w$.
\end{lemma}

\begin{lemma}[Copying an input prefix]
\lean{Turing.FinTM.f2_rawStrip_copy}
\label{Turing.FinTM.f2_rawStrip_copy}
\leanok
\uses{Turing.Action.apply, Turing.Cfg.inputSymbol, Turing.FinTM.bufferTape, Turing.FinTM.bufferTape_append, Turing.FinTM.f2_rawStripTM, Turing.FinTM.f2_stripCfg, Turing.MultiTapeTM.initCfg, Turing.MultiTapeTM.runFrom, Turing.MultiTapeTM.runFrom_succ_eq_step', Turing.MultiTapeTM.step, Turing.inputSymbolInner, Turing.moveInputPos_pos_of_ne_right}
\difficulty{6}
For every Boolean word $x$, every natural number $j$, and every $j\le |x|$, after $j$
steps the raw strip machine started on $x$ is in state $0$, its input head is at
position $j+1$, its work tape contains the prefix of $x$ of length $j$, its work head is
at position $j$, and its output is empty.
\end{lemma}

\begin{lemma}[Rewinding the work-tape head]
\lean{Turing.FinTM.f2_rawStrip_rewind}
\label{Turing.FinTM.f2_rawStrip_rewind}
\leanok
\uses{Turing.Action.apply, Turing.Cfg.workTapeSymbols, Turing.FinTM.bufferTape_left, Turing.FinTM.bufferTape_nat, Turing.FinTM.f2_rawStripTM, Turing.FinTM.f2_stripCfg, Turing.MultiTapeTM.runFrom, Turing.MultiTapeTM.runFrom_succ_eq_step, Turing.MultiTapeTM.runFrom_zero, Turing.MultiTapeTM.step, Turing.moveInputPos_zero}
\difficulty{4}
For every pair of Boolean words $x,a$ and every natural number $i\le |x|$, for every
natural number $j\le |a|$, the raw strip machine started in state $2$ on input $x$, with
work tape containing $a$, work head at position $j-1$, and empty output, reaches state
$3$ with work head at position $0$ and unchanged tape and output after $j+1$ steps.
\end{lemma}

\begin{lemma}[Replaying a word onto the output]
\lean{Turing.FinTM.f2_rawStrip_replay}
\label{Turing.FinTM.f2_rawStrip_replay}
\leanok
\uses{Turing.Action.apply, Turing.Cfg.workTapeSymbols, Turing.FinTM.bufferTape_nat, Turing.FinTM.f2_rawStripTM, Turing.FinTM.f2_stripCfg, Turing.MultiTapeTM.runFrom, Turing.MultiTapeTM.runFrom_succ_eq_step', Turing.MultiTapeTM.step, Turing.moveInputPos_zero}
\difficulty{4}
For every pair of Boolean words $x,a$ and every natural number $i\le |x|$, for each
natural number $j\le |a|$, the raw strip machine started in state $3$ on input $x$, with
work tape containing $a$, work head at position $0$, and empty output, is after $j$
steps in state $3$, with work head at position $j$ and output equal to the prefix of $a$
of length $j$.
\end{lemma}

\begin{lemma}[Finishing the raw strip output]
\lean{Turing.FinTM.f2_rawStrip_finish}
\label{Turing.FinTM.f2_rawStrip_finish}
\leanok
\uses{Turing.Action.apply, Turing.Cfg.workTapeSymbols, Turing.FinTM.f2_rawStripTM, Turing.FinTM.f2_rawStrip_replay, Turing.FinTM.f2_rawStrip_rewind, Turing.FinTM.f2_stripCfg, Turing.MultiTapeTM.runFrom, Turing.MultiTapeTM.runFrom_add, Turing.MultiTapeTM.runFrom_succ_eq_step', Turing.MultiTapeTM.step}
\difficulty{3}
For every pair of Boolean words $x,a$ and every natural number $i\le |x|$, the raw strip
machine started in state $2$ on input $x$, with work tape containing $a$, work head at
position $|a|-1$, and empty output, is halted after $2(|a|+1)$ steps and has output $a$.
\end{lemma}

\begin{lemma}[Erasing the final buffer symbol]
\lean{Turing.FinTM.f2_rawStrip_erase}
\label{Turing.FinTM.f2_rawStrip_erase}
\leanok
\uses{Turing.Action.apply, Turing.Cfg.workTapeSymbols, Turing.FinTM.bufferTape_nat, Turing.FinTM.f2_catalogBuffer_erase, Turing.FinTM.f2_rawStripTM, Turing.FinTM.f2_stripCfg, Turing.MultiTapeTM.step, Turing.moveInputPos_zero}
\difficulty{3}
Let $x,w$ be Boolean words, let $i$ be a natural number with $i\le |x|$, and let $b$ be
a Boolean value. One step of the raw strip machine from the configuration in state $1$,
with input head at position $i+1$, work tape containing $w$ followed by $b$, work head
at position $|w++[b]|-1$, and empty output, yields the configuration in state $2$ if $b$
is true and state $1$ otherwise, with input head still at position $i+1$, work tape
containing $w$, work head at position $|w|-1$, and empty output.
\end{lemma}

\begin{lemma}[Bounded trimming computation]
\lean{Turing.FinTM.f2_rawStrip_trim}
\label{Turing.FinTM.f2_rawStrip_trim}
\leanok
\uses{Turing.Action.apply, Turing.Cfg.workTapeSymbols, Turing.FinTM.f2_rawStripTM, Turing.FinTM.f2_rawStrip_erase, Turing.FinTM.f2_rawStrip_finish, Turing.FinTM.f2_stripCfg, Turing.MultiTapeTM.runFrom, Turing.MultiTapeTM.runFrom_succ_eq_step, Turing.MultiTapeTM.step, Turing.splitAtLastTrue}
\difficulty{5}
For every pair of Boolean words $x,w$ and every position $i$ with $i\le |x|$, the
designated Boolean Turing machine, started in state $1$ with input head at position
$i+1$, work tape holding $w$ and head at $|w|-1$, and empty output, halts within
$3(|w|+1)$ steps and has output equal to the prefix of $w$ strictly before its last true
entry, or the empty word if $w$ has no true entry.
\end{lemma}

\begin{lemma}[Running time of the raw strip machine]
\lean{Turing.FinTM.f2_rawStrip_computes}
\label{Turing.FinTM.f2_rawStrip_computes}
\leanok
\uses{Turing.Action.apply, Turing.Cfg.inputSymbol, Turing.FinTM.ComputesFunInTime, Turing.FinTM.ComputesInTime, Turing.FinTM.ComputesInTime.mono, Turing.FinTM.computesInTime_iff, Turing.FinTM.f2_rawStripTM, Turing.FinTM.f2_rawStrip_copy, Turing.FinTM.f2_rawStrip_trim, Turing.FinTM.f2_stripCfg, Turing.MultiTapeTM.initCfg, Turing.MultiTapeTM.runFrom, Turing.MultiTapeTM.runFrom_add, Turing.MultiTapeTM.runFrom_succ_eq_step', Turing.MultiTapeTM.step, Turing.splitAtLastTrue}
\difficulty{4}
For every Boolean word $x$, the raw strip machine computes the function that returns the
prefix strictly before the last true in $x$, or the empty word if $x$ has no true, in
time $4(|x|+1)$.
\end{lemma}

\begin{definition}[Any true machine]
\lean{Turing.FinTM.f2_anyTrueTM}
\label{Turing.FinTM.f2_anyTrueTM}
\leanok
\uses{Turing.FinTM}
This is a finite Turing machine over the Boolean alphabet with no work tapes and a
single state, which is its initial state. On reading false, it moves the input head
right and continues; on reading true, it outputs true and halts; and on reading a blank
input cell, it outputs false and halts.
\end{definition}

\begin{lemma}[Any-true machine run]
\lean{Turing.FinTM.f2_anyTrue_run}
\label{Turing.FinTM.f2_anyTrue_run}
\leanok
\uses{Turing.Action.apply, Turing.FinTM.f2_anyTrueTM, Turing.FinTM.f2_scanCfg, Turing.FinTM.f2_scanCfg_read, Turing.FinTM.f2_scanStep_right, Turing.MultiTapeTM.runFrom, Turing.MultiTapeTM.runFrom_succ_eq_step, Turing.MultiTapeTM.runFrom_zero, Turing.MultiTapeTM.step}
\difficulty{5}
Let $x$ and $r$ be Boolean words, and write $x=p r$ for a split into a prefix $p$ and
suffix $r$. There is a time $t \le |r|+1$ such that the designated any-true machine,
started in its scanning state with input head just past $p$ and empty output, is halted
and has output consisting of the single bit that is true exactly when $r$ contains a
true bit.
\end{lemma}

\begin{lemma}[Any-true Boolean computation]
\lean{Turing.FinTM.f2_anyTrue_computes}
\label{Turing.FinTM.f2_anyTrue_computes}
\leanok
\uses{Turing.Cfg.ext_zero_tapes, Turing.FinTM.ComputesFunInTime, Turing.FinTM.ComputesInTime, Turing.FinTM.ComputesInTime.mono, Turing.FinTM.computesInTime_iff, Turing.FinTM.f2_anyTrueTM, Turing.FinTM.f2_anyTrue_run, Turing.FinTM.f2_scanCfg, Turing.MultiTapeTM.initCfg}
\difficulty{4}
The designated finite Turing machine over the Boolean alphabet computes the function
sending each Boolean word $x$ to the one-symbol word whose bit is true exactly when at
least one bit of $x$ is true, in time $|x|+1$.
\end{lemma}

\begin{lemma}[Cases for the last true bit]
\lean{Turing.FinTM.f2_catalogMarker_cases}
\label{Turing.FinTM.f2_catalogMarker_cases}
\leanok
\uses{Turing.splitAtLastTrue}
\difficulty{5}
For every finite Boolean word $v$, either $v$ contains no true bit and its cut at the
last true bit is undefined, or $v$ contains a true bit and there is a word $u$ equal to
that cut such that, for every prefix $pre$, the cut of $pre$ followed by $v$ is defined
and equals $pre$ followed by $u$.
\end{lemma}

\begin{lemma}[Linear-time decoding and stripping]
\lean{Turing.FinTM.f2_strip_linear}
\label{Turing.FinTM.f2_strip_linear}
\leanok
\uses{Turing.FinTM, Turing.FinTM.ComputesFunInTime, Turing.FinTM.ComputesInTime, Turing.FinTM.ComputesInTime.mono, Turing.FinTM.computesFunInTime_comp, Turing.FinTM.computesFunInTime_cond, Turing.FinTM.computesFunInTime_const_spaceUsed, Turing.FinTM.computesFunInTime_pairSnd_spaceUsed, Turing.FinTM.f2_anyTrue_computes, Turing.FinTM.f2_catalogMarker_cases, Turing.FinTM.f2_catalogPair_inverse, Turing.FinTM.f2_rawStrip_computes, Turing.pairDecode, Turing.pairEncode, Turing.splitAtLastTrue}
\difficulty{6}
There exist a finite Turing machine $M$ over the Boolean alphabet and a natural number
$c$ such that, for every Boolean word $x$, $M$ computes the following output in at most
$c(|x|+1)$ steps: if $x$ decodes as a pair $(a,v)$ and $v$ has a last occurrence of
$\mathrm{true}$, return the encoding of $a$ paired with the prefix of $v$ strictly
before that last occurrence; if either operation is undefined, return the empty word.
\end{lemma}

\begin{theorem}[Last-marker stripping with bounded space]
\lean{Turing.FinTM.computesFunInTime_stripLast_spaceUsed}
\label{Turing.FinTM.computesFunInTime_stripLast_spaceUsed}
\leanok
\uses{Turing.FinTM, Turing.FinTM.ComputesFunInTime, Turing.FinTM.ComputesInTime.mono, Turing.FinTM.f2_space_of_time, Turing.FinTM.f2_strip_linear, Turing.MultiTapeTM.initCfg, Turing.MultiTapeTM.spaceUsed, Turing.pairDecode, Turing.pairEncode, Turing.splitAtLastTrue}
\difficulty{4}
There exist a finite Turing machine $M$ over the Boolean alphabet and a natural number
$c$ such that, on every Boolean input $x$ of length $n$, $M$ halts within $c(n+1)^2$
steps. It outputs the pair encoding of $a$ and the prefix of $v$ strictly before its
last $1$ when $x$ parses into $(a,v)$ by the aligned-pair parser and $v$ contains a $1$;
otherwise it outputs the empty word. For every Boolean word $x$ and every time $t$, the
computation of $M$ on $x$ up to time $t$ uses at most $c(|x|+1)$ work-tape cells.
\end{theorem}

\begin{theorem}[Fixed-Width Increment in Constant Space]
\lean{Turing.FinTM.computesFunInTime_incFixed_spaceUsed}
\label{Turing.FinTM.computesFunInTime_incFixed_spaceUsed}
\leanok
\uses{Turing.FinTM, Turing.FinTM.ComputesFunInTime, Turing.FinTM.f2_incFixedTM, Turing.FinTM.f2_incFixed_computes, Turing.MultiTapeTM.initCfg, Turing.MultiTapeTM.spaceUsed, Turing.MultiTapeTM.spaceUsed_zero_tapes_eq_zero, Turing.incFixed}
\difficulty{3}
There exist a finite Turing machine $M$ over the Boolean alphabet and a constant
$c\in\mathbb{N}$ such that, on every Boolean word $x$, $M$ halts within $c(|x|+1)$ steps
and outputs the fixed-width increment of $x$, using the empty word when incrementing
overflows. For every Boolean word $x$ and every time $t\in\mathbb{N}$, the computation
of $M$ from its initial configuration on $x$ uses at most $c$ work-tape cells through
step $t$.
\end{theorem}

\begin{definition}[Mapping machine states]
\lean{Turing.FinTM.a2_MapState}
\label{Turing.FinTM.a2_MapState}
\leanok
For a type $Q$, this inductive type has states for parsing with an optional pending
Boolean, copying the second input, moving back across the second or first work area,
emitting from the first work area, emitting again with a specified Boolean, and writing
a separator. It also has a running state carrying a state $q\in Q$ and a Boolean tag.
\end{definition}

\begin{definition}[Action for the mapping machine]
\lean{Turing.FinTM.a2_mapAct}
\label{Turing.FinTM.a2_mapAct}
\leanok
\uses{Turing.Action, Turing.FinTM, Turing.FinTM.a2_MapState, Turing.FinTM.tapeBlocks}
Let $M$ be a finite Turing machine over the Boolean alphabet, with $k$ work tapes, and
let $Q_M$ be its state set. Given an input-head movement, two work-tape instructions, an
optional output bit, and an optional successor state of the mapping machine whose states
encode parsing, copying, returning, output, separation, or a state of $M$, the
associated action has $1+(1+k)$ work tapes. It applies the first instruction to the
first tape, the second instruction to the next tape, and leaves the remaining $k$ tapes
unchanged; its input-head movement, output, and successor state are the supplied ones.
\end{definition}

\begin{definition}[Boolean input mapping machine]
\lean{Turing.FinTM.a2_mapTM}
\label{Turing.FinTM.a2_mapTM}
\leanok
\uses{Turing.FinTM, Turing.FinTM.a2_MapState, Turing.FinTM.a2_mapAct, Turing.FinTM.controlAction, Turing.FinTM.tapeBlocks, Turing.FinTM.virtualMove, Turing.FinTM.virtualNextTag}
Given a finite Boolean Turing machine $M$ and a Boolean flag $f$, this defines a finite
Boolean Turing machine with $1+(1+M.k)$ work tapes. Its initial phase reads the input in
pairs: it copies equal pairs to its first work tape, and when it encounters a pair
consisting of false followed by true, it copies the remaining input to its second work
tape. It then returns to the beginnings of the copied portions, emits the first portion
with false symbols interleaved between its symbols, emits a true separator, and begins a
simulation phase. If $f$ is true, that phase uses a transformed transition of $M$ with
the second work tape and the remaining work tapes; if $f$ is false, it halts in that
phase.
\end{definition}

\begin{definition}[Mapping machine configuration]
\lean{Turing.FinTM.a2_mapCfg}
\label{Turing.FinTM.a2_mapCfg}
\leanok
\uses{Turing.Cfg, Turing.FinTM, Turing.FinTM.a2_MapState, Turing.FinTM.bufferTape, Turing.FinTM.tapeBlocks}
Given a finite Boolean Turing machine $M$, a Boolean input list $x$, a state or halting
marker of its associated mapping machine, an input-head position in the range of
positions for $x$, two Boolean lists $a,b$, two integer head positions $h_a,h_b$, and an
output list, this defines a configuration on input $x$. Its work tapes contain the
buffer for $a$ on the first tape, the buffer for $b$ on the next tape, and blank cells
on the remaining tapes; the corresponding head positions are $h_a$ and $h_b$, with all
remaining work heads at position zero.
\end{definition}

\begin{definition}[Virtual simulation configuration]
\lean{Turing.FinTM.a2_mapVirtual}
\label{Turing.FinTM.a2_mapVirtual}
\leanok
\uses{Turing.Cfg, Turing.FinTM, Turing.FinTM.a2_MapState, Turing.FinTM.bufferTape, Turing.FinTM.tapeBlocks}
Given a finite Boolean Turing machine $M$, Boolean lists $x,y$, a configuration $c$ of
$M$ on input $y$, a Boolean tag, a valid input-head position $p$ for $x$, and Boolean
lists $a$ and $pre$, this defines a configuration for the mapping machine on input $x$.
Its state is the state of $c$ wrapped as a running state with the given tag, its
input-head position is $p$, its work tapes contain the buffer for $a$, the list $y$, and
then the work tapes of $c$, and its output is $pre$ followed by the output of $c$. The
work-head positions are $|a|$, the input-head position of $c$ minus one, and then the
work-head positions of $c$.
\end{definition}

\begin{lemma}[Input symbol at an indexed position]
\lean{Turing.FinTM.a2_mapCfg_read}
\label{Turing.FinTM.a2_mapCfg_read}
\leanok
\uses{Turing.Cfg.inputSymbol, Turing.FinTM, Turing.FinTM.a2_MapState, Turing.FinTM.a2_mapCfg, Turing.FinTM.inputSymbol_at}
\difficulty{1}
Let $M$ be a finite Boolean Turing machine, let $x$ be a Boolean word, let $q$ be an
optional mapping-machine state, let $i\le |x|$, let $a,b$ be Boolean words, let
$h_a,h_b$ be integers, and let $out$ be a Boolean word. In the configuration with input
head at position $i+1$ and these tape words, head positions, state, and output, the
symbol under the input head is the $i$-th symbol of $x$, or blank if $i=|x|$.
\end{lemma}

\begin{lemma}[Effect of a structured simulator action]
\lean{Turing.FinTM.a2_map_move}
\label{Turing.FinTM.a2_map_move}
\leanok
\uses{Turing.Action.apply, Turing.FinTM, Turing.FinTM.a2_MapState, Turing.FinTM.a2_mapAct, Turing.FinTM.a2_mapCfg, Turing.moveInputPos}
\difficulty{4}
Let $M$ be a finite Turing machine over Booleans, let $x,a,b$ be Boolean words, let
$q,q'$ be simulator states, let $p$ be an input-head position for $x$, let $h_a,h_b$ be
integers, and let $o$ be an output word. For movements $m_i,m_a,m_b\in\{-1,0,1\}$ and an
optional output bit $e$, applying the action that moves the input head by $m_i$, makes
no writes, moves the first two work heads by $m_a,m_b$, emits $e$, and enters $q'$ gives
the configuration with state $q'$, input position moved by $m_i$, buffers $a,b$, head
positions $h_a+m_a,h_b+m_b$, and output $o$ followed by $e$ when present.
\end{lemma}

\begin{lemma}[Reading the first input bit]
\lean{Turing.FinTM.a2_map_first}
\label{Turing.FinTM.a2_map_first}
\leanok
\uses{Turing.Action.apply, Turing.FinTM, Turing.FinTM.a2_mapAct, Turing.FinTM.a2_mapCfg, Turing.FinTM.a2_mapCfg_read, Turing.FinTM.a2_mapTM, Turing.FinTM.a2_map_move, Turing.MultiTapeTM.step, Turing.moveInputPos_pos_of_ne_right}
\difficulty{4}
Let $M$ be a finite Turing machine over Booleans, let $x,pre,rest,a$ be Boolean words,
let $b$ be a Boolean, and suppose $x=pre\mathbin{++}[b]\mathbin{++}rest$. One step of
the constructed machine from the parsing state at input position $|pre|+1$, with first
buffer $a$, empty second buffer, head positions $|a|$ and $0$, and empty output,
advances the input head to position $|pre|+2$ and enters the parsing state with $b$
pending; the buffers, head positions, and output remain unchanged.
\end{lemma}

\begin{lemma}[Parsing an aligned pair of input bits]
\lean{Turing.FinTM.a2_map_block}
\label{Turing.FinTM.a2_map_block}
\leanok
\uses{Turing.Action.apply, Turing.FinTM, Turing.FinTM.a2_mapAct, Turing.FinTM.a2_mapCfg, Turing.FinTM.a2_mapCfg_read, Turing.FinTM.a2_mapTM, Turing.FinTM.a2_map_first, Turing.FinTM.bufferTape_append, Turing.FinTM.tapeBlocks_left, Turing.MultiTapeTM.runFrom, Turing.MultiTapeTM.step, Turing.moveInputPos, Turing.moveInputPos_pos_of_ne_right}
\difficulty{5}
Let $M$ be a finite Turing machine over Booleans, let $x,pre,rest$ be Boolean words, let
$b,c$ be Booleans, and suppose $x=pre\mathbin{++}[b,c]\mathbin{++}rest$. Starting with
the parsing state at input position $|pre|+1$, first buffer $a$, empty second buffer,
first-buffer head at $|a|$, second-buffer head at $0$, and empty output, after two steps
the machine is in the parsing state with the pair appended to the first buffer if $b=c$;
if $b\ne c$ and $b$ is true, it has halted with the first buffer unchanged; and if $b\ne
c$ and $b$ is false, it is in the copy state with the first-buffer head at $|a|-1$. In
every case the input head is at position $|pre|+3$, the second buffer remains empty, and
the output remains empty.
\end{lemma}

\begin{lemma}[Copying an input suffix to a buffer]
\lean{Turing.FinTM.a2_map_suffix}
\label{Turing.FinTM.a2_map_suffix}
\leanok
\uses{Turing.Action.apply, Turing.FinTM, Turing.FinTM.a2_mapAct, Turing.FinTM.a2_mapCfg, Turing.FinTM.a2_mapCfg_read, Turing.FinTM.a2_mapTM, Turing.FinTM.a2_map_move, Turing.FinTM.bufferTape_append, Turing.FinTM.tapeBlocks_buffer, Turing.MultiTapeTM.runFrom, Turing.MultiTapeTM.runFrom_succ_eq_step, Turing.MultiTapeTM.runFrom_zero, Turing.MultiTapeTM.step, Turing.moveInputPos_pos_of_ne_right}
\difficulty{5}
Let $M$ be a finite Turing machine over Booleans, and let $x,r,a$ be Boolean words. For
every Boolean prefix $p$ and buffer word $b$ such that $x=pr$, starting the non-forward
simulator in its copying state at input position $|p|+1$, with buffers $a,b$, head
positions $|a|-1,|b|$, and empty output, and running it for $|r|+1$ steps yields the
configuration with state the return-to-buffer phase, input position $|x|+1$, buffers
$a,br$, head positions $|a|-1,|br|-1$, and empty output.
\end{lemma}

\begin{lemma}[Returning from the second buffer]
\lean{Turing.FinTM.a2_map_backB}
\label{Turing.FinTM.a2_map_backB}
\leanok
\uses{Turing.Action.apply, Turing.Cfg.workTapeSymbols, Turing.FinTM, Turing.FinTM.a2_mapAct, Turing.FinTM.a2_mapCfg, Turing.FinTM.a2_mapTM, Turing.FinTM.a2_map_move, Turing.FinTM.bufferTape, Turing.FinTM.bufferTape_left, Turing.FinTM.bufferTape_nat, Turing.FinTM.tapeBlocks_buffer, Turing.MultiTapeTM.runFrom, Turing.MultiTapeTM.runFrom_succ_eq_step, Turing.MultiTapeTM.runFrom_zero, Turing.MultiTapeTM.step}
\difficulty{4}
Let $M$ be a finite Turing machine over Booleans, let $x,a,b$ be Boolean words, and let
$p$ be a valid input-head position on $x$. For every natural number $j$ with $j\le |b|$,
running the constructed machine in its backward-through-the-second-buffer state for
$j+1$ steps, with buffers $a$ and $b$, their head positions $|a|-1$ and $j-1$, and empty
output, yields the configuration in the backward-through-the-first-buffer state with the
same input position and buffers, head positions $|a|-1$ and $0$, and empty output.
\end{lemma}

\begin{lemma}[Returning from the second buffer]
\lean{Turing.FinTM.a2_map_backA}
\label{Turing.FinTM.a2_map_backA}
\leanok
\uses{Turing.Action.apply, Turing.Cfg.workTapeSymbols, Turing.FinTM, Turing.FinTM.a2_mapAct, Turing.FinTM.a2_mapCfg, Turing.FinTM.a2_mapTM, Turing.FinTM.a2_map_move, Turing.FinTM.bufferTape, Turing.FinTM.bufferTape_left, Turing.FinTM.bufferTape_nat, Turing.FinTM.tapeBlocks_left, Turing.MultiTapeTM.runFrom, Turing.MultiTapeTM.runFrom_succ_eq_step, Turing.MultiTapeTM.runFrom_zero, Turing.MultiTapeTM.step}
\difficulty{4}
Let $M$ be a finite Turing machine over Booleans, let $x,a,b$ be Boolean words, and let
$p$ be a valid input-head position on $x$. For every natural number $j$ with $j\le |a|$,
running the constructed machine in its backward-through-the-first-buffer state for $j+1$
steps, with the first and second buffers holding $a$ and $b$, their head positions $j-1$
and $0$, and empty output, yields the configuration in the emit state with the same
input position and buffers, both buffer heads at $0$, and empty output.
\end{lemma}

\begin{lemma}[Emitting the encoded first word]
\lean{Turing.FinTM.a2_map_emit}
\label{Turing.FinTM.a2_map_emit}
\leanok
\uses{Turing.Action.apply, Turing.Cfg.workTapeSymbols, Turing.FinTM, Turing.FinTM.a2_mapAct, Turing.FinTM.a2_mapCfg, Turing.FinTM.a2_mapTM, Turing.FinTM.a2_map_move, Turing.FinTM.bufferTape, Turing.FinTM.bufferTape_nat, Turing.FinTM.tapeBlocks_left, Turing.MultiTapeTM.runFrom, Turing.MultiTapeTM.runFrom_add, Turing.MultiTapeTM.step}
\difficulty{5}
Let $M$ be a finite Turing machine over Booleans, let $x,a,b,rest$ be Boolean words, and
let $p$ be a valid input-head position on $x$. For every Boolean words $pre$ and $out$
such that $a=pre\mathbin{++}rest$, running the constructed machine from the emit state
with first buffer $a$, second buffer $b$, head positions $|pre|$ and $0$, and output
$out$ for $2|rest|+2$ steps yields the run state at $M$'s initial state with tag true.
The first buffer head is then at $|a|$, the second remains at $0$, and the output is
$out$ followed by each bit of $rest$ twice in order, then $0,1$.
\end{lemma}

\begin{lemma}[Completing the input-copying phase]
\lean{Turing.FinTM.a2_map_finish}
\label{Turing.FinTM.a2_map_finish}
\leanok
\uses{Turing.Cfg.init, Turing.FinTM, Turing.FinTM.a2_mapCfg, Turing.FinTM.a2_mapTM, Turing.FinTM.a2_mapVirtual, Turing.FinTM.a2_map_backA, Turing.FinTM.a2_map_backB, Turing.FinTM.a2_map_emit, Turing.FinTM.a2_map_suffix, Turing.MultiTapeTM.initCfg, Turing.MultiTapeTM.runFrom, Turing.MultiTapeTM.runFrom_add, Turing.pairEncode}
\difficulty{4}
Let $M$ be a finite Turing machine over Booleans, let $x,pre,a,b$ be Boolean words, and
suppose $x=pre\mathbin{++}b$. Starting in the copy state at input position $|pre|+1$,
with first buffer $a$, empty second buffer, head positions $|a|-1$ and $0$, and empty
output, after $3|a|+2|b|+5$ steps the constructed machine is in the run state at $M$'s
initial state with tag true. Its input head is at the right boundary position $|x|+1$,
its first buffer is still $a$, and its output is the encoding of the pair
$(a,\epsilon)$, where $\epsilon$ is the empty word.
\end{lemma}

\begin{lemma}[Parsing a remaining input suffix]
\lean{Turing.FinTM.a2_map_parse}
\label{Turing.FinTM.a2_map_parse}
\leanok
\uses{Turing.Action.apply, Turing.FinTM, Turing.FinTM.a2_mapAct, Turing.FinTM.a2_mapCfg, Turing.FinTM.a2_mapCfg_read, Turing.FinTM.a2_mapTM, Turing.FinTM.a2_mapVirtual, Turing.FinTM.a2_map_block, Turing.FinTM.a2_map_finish, Turing.FinTM.a2_map_first, Turing.MultiTapeTM.initCfg, Turing.MultiTapeTM.runFrom, Turing.MultiTapeTM.runFrom_add, Turing.MultiTapeTM.runFrom_succ_eq_step, Turing.MultiTapeTM.runFrom_zero, Turing.MultiTapeTM.step, Turing.pairDecode, Turing.pairEncode}
\difficulty{6}
Let $M$ be a finite Turing machine over Booleans, and let $x$ and $r$ be Boolean words.
For every Boolean prefix $p$ and buffer word $a$ such that $x= p r$, there is a time
$t\le 3|r|+3|a|+5$ with the following property. If decoding $r$ yields $(d,b)$, then the
non-forward simulator started in its parsing state at input position $|p|+1$, with
buffer $a$, empty second buffer, head positions $|a|$ and $0$, and empty output, reaches
the virtual configuration for $M$ initialized on $b$, with tag true, input position
$|x|+1$, buffer $a$ followed by $d$, and output the encoding of $(a d,[])$. If decoding
$r$ is undefined, that configuration at time $t$ is halted and has empty output.
\end{lemma}

\begin{lemma}[Setup of the non-forward simulator]
\lean{Turing.FinTM.a2_map_setup}
\label{Turing.FinTM.a2_map_setup}
\leanok
\uses{Turing.Cfg.init, Turing.FinTM, Turing.FinTM.a2_mapCfg, Turing.FinTM.a2_mapTM, Turing.FinTM.a2_mapVirtual, Turing.FinTM.a2_map_parse, Turing.MultiTapeTM.initCfg, Turing.MultiTapeTM.runFrom, Turing.pairDecode, Turing.pairEncode}
\difficulty{4}
Let $M$ be a finite Turing machine over Booleans and let $x$ be a Boolean word. There is
a time $t\le 5(|x|+1)$ such that, if decoding $x$ yields $(a,b)$, the non-forward
simulator reaches the virtual configuration for $M$ initialized on $b$, with tag true,
input position $|x|+1$, buffer $a$, and output the encoding of $(a,[])$. If decoding $x$
is undefined, the simulator at time $t$ is halted with empty output.
\end{lemma}

\begin{lemma}[One-step simulation of a mapped configuration]
\lean{Turing.FinTM.a2_mapVirtual_step}
\label{Turing.FinTM.a2_mapVirtual_step}
\leanok
\uses{Turing.Action.apply, Turing.Cfg, Turing.Cfg.inputSymbol, Turing.Cfg.workTapeSymbols, Turing.FinTM, Turing.FinTM.VirtualTag, Turing.FinTM.a2_mapTM, Turing.FinTM.a2_mapVirtual, Turing.FinTM.bufferTape_inputSymbol, Turing.FinTM.tapeBlocks_buffer, Turing.FinTM.virtualMove, Turing.FinTM.virtualMove_correct, Turing.FinTM.virtualNextTag, Turing.MultiTapeTM.step, Turing.MultiTapeTM.step_of_halt, Turing.moveInputPos_zero}
\difficulty{6}
Let $M$ be a finite Turing machine over the Boolean alphabet, let $x$ and $y$ be Boolean
words, and let $c$ be a configuration of $M$ on input $y$. Suppose a Boolean tag
correctly marks the boundary of the input-head position of $c$. For any input-head
position $p$ on $x$ and Boolean words $a$ and $pre$, there is a tag correctly marking
the boundary of the input-head position after one step of $M$ from $c$, such that one
step of the constructed machine from the corresponding virtual configuration equals the
virtual configuration associated with that successor of $c$, with $p$, $a$, and $pre$
unchanged.
\end{lemma}

\begin{lemma}[Simulation of a mapped computation]
\lean{Turing.FinTM.a2_mapVirtual_run}
\label{Turing.FinTM.a2_mapVirtual_run}
\leanok
\uses{Turing.Cfg, Turing.FinTM, Turing.FinTM.VirtualTag, Turing.FinTM.a2_mapTM, Turing.FinTM.a2_mapVirtual, Turing.FinTM.a2_mapVirtual_step, Turing.MultiTapeTM.runFrom, Turing.MultiTapeTM.runFrom_succ_eq_step'}
\difficulty{4}
Let $M$ be a finite Turing machine over the Boolean alphabet, let $x$ and $y$ be Boolean
words, and let $c$ be a configuration of $M$ on input $y$. Suppose a Boolean tag
correctly marks the boundary of the input-head position of $c$. For any input-head
position $p$ on $x$, Boolean word $a$, output prefix $pre$, and time $t\in\mathbb{N}$,
there is a tag for the input-head position of the configuration reached by running $M$
from $c$ for $t$ steps, correctly marking its boundary, such that the run for $t$ steps
of the constructed machine on the corresponding virtual configuration equals the virtual
configuration associated with that reached configuration, with $p$, $a$, and $pre$
unchanged.
\end{lemma}

\begin{definition}[Configurations in the run state]
\lean{Turing.FinTM.a2_mapEntered}
\label{Turing.FinTM.a2_mapEntered}
\leanok
\uses{Turing.Cfg, Turing.FinTM, Turing.FinTM.a2_MapState}
Given a finite Turing machine $M$ over the Boolean alphabet, a Boolean input list $x$,
and a configuration $c$ for the auxiliary machine with state set built from the states
of $M$, say that $c$ has entered a run state when its state is $(q,b)$ for some state
$q$ of $M$ and some Boolean tag $b$.
\end{definition}

\begin{lemma}[Non-forward simulation states are stationary]
\lean{Turing.FinTM.a2_mapSetup_stationary}
\label{Turing.FinTM.a2_mapSetup_stationary}
\leanok
\uses{Turing.Cfg, Turing.FinTM, Turing.FinTM.a2_MapState, Turing.FinTM.a2_mapEntered, Turing.FinTM.a2_mapTM, Turing.FinTM.controlAction_apply, Turing.MultiTapeTM.runFrom, Turing.MultiTapeTM.runFrom_succ_eq_step', Turing.MultiTapeTM.step, Turing.moveInputPos_zero}
\difficulty{4}
Let $M$ be a finite Boolean Turing machine, let $x$ be a Boolean word, and let $c$ be a
configuration of the mapping machine on input $x$ whose state represents a simulated
state of $M$. In non-forward mode, the configuration reached from $c$ after any natural
number $t$ of steps is exactly $c$.
\end{lemma}

\begin{lemma}[Setup steps agree across modes]
\lean{Turing.FinTM.a2_mapSetup_step}
\label{Turing.FinTM.a2_mapSetup_step}
\leanok
\uses{Turing.Cfg, Turing.FinTM, Turing.FinTM.a2_MapState, Turing.FinTM.a2_mapEntered, Turing.FinTM.a2_mapTM, Turing.MultiTapeTM.step}
\difficulty{4}
Let $M$ be a finite Boolean Turing machine, let $x$ be a Boolean word, and let $c$ be
any configuration of the mapping machine on input $x$ that is not in a state
representing a simulated state of $M$. The one-step successor of $c$ is the same in
forward and non-forward modes.
\end{lemma}

\begin{lemma}[Agreement before entering the simulation phase]
\lean{Turing.FinTM.a2_mapSetup_run}
\label{Turing.FinTM.a2_mapSetup_run}
\leanok
\uses{Turing.FinTM, Turing.FinTM.a2_mapEntered, Turing.FinTM.a2_mapSetup_step, Turing.FinTM.a2_mapTM, Turing.MultiTapeTM.initCfg, Turing.MultiTapeTM.runFrom, Turing.MultiTapeTM.runFrom_succ_eq_step'}
\difficulty{3}
Let $M$ be a finite Boolean Turing machine, let $x$ be a Boolean word, and let $t$ be a
natural number. If the run of the mapping machine in non-forward setup mode, started on
$x$, has not entered a state representing a simulated state of $M$ at any time $u<t$,
then its configuration at time $t$ is identical to the configuration at time $t$ of the
mapping machine in forward mode started on $x$.
\end{lemma}

\begin{lemma}[Setup preserves simulated tape heads]
\lean{Turing.FinTM.a2_mapSetup_head_step}
\label{Turing.FinTM.a2_mapSetup_head_step}
\leanok
\uses{Turing.Action.apply, Turing.Cfg, Turing.FinTM, Turing.FinTM.a2_MapState, Turing.FinTM.a2_mapAct, Turing.FinTM.a2_mapTM, Turing.FinTM.tapeBlocks_right, Turing.MultiTapeTM.step}
\difficulty{5}
Let $M$ be a finite Boolean Turing machine, let $x$ be a Boolean word, let $c$ be any
configuration of the mapping machine in its non-forward setup mode, and let $i$ be one
of the simulated work tapes of $M$. After one step, the head position on the tape
corresponding to $i$ is unchanged.
\end{lemma}

\begin{lemma}[Setup leaves simulated heads at the origin]
\lean{Turing.FinTM.a2_mapSetup_heads}
\label{Turing.FinTM.a2_mapSetup_heads}
\leanok
\uses{Turing.FinTM, Turing.FinTM.a2_mapSetup_head_step, Turing.FinTM.a2_mapTM, Turing.MultiTapeTM.initCfg, Turing.MultiTapeTM.runFrom, Turing.MultiTapeTM.runFrom_succ_eq_step'}
\difficulty{3}
Let $M$ be a finite Boolean Turing machine, let $x$ be a Boolean word, let $t$ be a
natural number, and let $i$ be one of the simulated work tapes of $M$. After $t$ steps
of the mapping machine in non-forward setup mode, started from its initial configuration
on $x$, the head on the tape corresponding to $i$ is at position $0$.
\end{lemma}

\begin{lemma}[Launch after decoding a paired input]
\lean{Turing.FinTM.a2_map_launch}
\label{Turing.FinTM.a2_map_launch}
\leanok
\uses{Turing.FinTM, Turing.FinTM.a2_mapEntered, Turing.FinTM.a2_mapSetup_run, Turing.FinTM.a2_mapSetup_stationary, Turing.FinTM.a2_mapTM, Turing.FinTM.a2_mapVirtual, Turing.FinTM.a2_map_setup, Turing.MultiTapeTM.initCfg, Turing.MultiTapeTM.runFrom, Turing.MultiTapeTM.runFrom_add, Turing.pairDecode, Turing.pairEncode}
\difficulty{5}
Let $M$ be a finite Turing machine over the Boolean alphabet, and let $x,a,b$ be Boolean
words such that decoding $x$ gives the pair $(a,b)$. Then there is a time $u \le
5(|x|+1)$ at which the forward simulator on input $x$ is in the virtual configuration
representing $M$ initialized on $b$, with boundary tag true, input-head position
$|x|+1$, first buffer $a$, and output prefix the encoding of $(a,[])$. Moreover, at
every time $v \le u$, the forward and non-forward simulators on $x$ have identical
configurations.
\end{lemma}

\begin{lemma}[Rejection of an undecodable input]
\lean{Turing.FinTM.a2_map_reject}
\label{Turing.FinTM.a2_map_reject}
\leanok
\uses{Turing.FinTM, Turing.FinTM.a2_mapEntered, Turing.FinTM.a2_mapSetup_run, Turing.FinTM.a2_mapSetup_stationary, Turing.FinTM.a2_mapTM, Turing.FinTM.a2_map_setup, Turing.MultiTapeTM.initCfg, Turing.MultiTapeTM.runFrom, Turing.MultiTapeTM.runFrom_add, Turing.MultiTapeTM.runFrom_of_halt, Turing.pairDecode}
\difficulty{5}
Let $M$ be a finite Turing machine over Booleans and let $x$ be a Boolean word that has
no valid paired decoding. There is a time $u\le 5(|x|+1)$ at which the forward simulator
on $x$ is halted with empty output. At every time $v$, the forward and non-forward
simulators on $x$ have identical configurations.
\end{lemma}

\begin{definition}[Finite sum equivalence]
\lean{Turing.FinTM.a2_mapSumEquiv}
\label{Turing.FinTM.a2_mapSumEquiv}
\leanok
For natural numbers $a$ and $b$, the disjoint union of the first $a$ natural numbers and
the first $b$ natural numbers is in bijection with the first $a+b$ natural numbers, by
placing the second block after the first.
\end{definition}

\begin{lemma}[Splitting a finite sum]
\lean{Turing.FinTM.a2_map_sum}
\label{Turing.FinTM.a2_map_sum}
\leanok
\uses{Turing.FinTM.a2_mapSumEquiv}
\difficulty{3}
For natural numbers $a,b$ and any function $f:\mathrm{Fin}(a+b)\to\mathbb{N}$, the sum
of $f$ over the first $a+b$ indices equals the sum over the first $a$ indices plus the
sum over the last $b$ indices: $\sum_{i\in\mathrm{Fin}(a+b)}
f(i)=\sum_{i\in\mathrm{Fin}(a)} f(i+b)+\sum_{i\in\mathrm{Fin}(b)} f(a+i)$, where the
terms are indexed by their corresponding positions in the full range.
\end{lemma}

\begin{lemma}[Space bound for the mapped computation]
\lean{Turing.FinTM.a2_map_space}
\label{Turing.FinTM.a2_map_space}
\leanok
\uses{Turing.FinTM, Turing.FinTM.a2_mapTM, Turing.FinTM.a2_map_sum, Turing.MultiTapeTM.initCfg, Turing.MultiTapeTM.runFrom, Turing.MultiTapeTM.spaceUsed, Turing.MultiTapeTM.spaceUsedByTape, Turing.MultiTapeTM.visitedByTapeHead}
\difficulty{4}
Let $M$ be a finite Turing machine over Booleans, let $x$ and $y$ be Boolean words, and
let $t,D,S$ be natural numbers. Suppose that at every time $u\le t$, the positions of
each of the simulator’s first two work-tape heads lie in $[-D,D]$. Suppose also that,
for every work tape of $M$, the cells visited by the corresponding simulator tape
through time $t$ are among the cells visited by that tape of $M$ when initialized on
$y$. If the computation of $M$ initialized on $y$ uses at most $S$ cells through time
$t$, then the simulator initialized on $x$ uses at most $S+2(2D+1)$ cells through time
$t$.
\end{lemma}

\begin{theorem}[Pair mapping with a space bound]
\lean{Turing.FinTM.computesFunInTime_pairMapSnd_spaceUsed}
\label{Turing.FinTM.computesFunInTime_pairMapSnd_spaceUsed}
\leanok
\uses{Turing.Cfg.init, Turing.FinTM, Turing.FinTM.ComputesFunInTime, Turing.FinTM.ComputesInTime, Turing.FinTM.ComputesInTime.mono, Turing.FinTM.VirtualTag, Turing.FinTM.a2_mapSetup_heads, Turing.FinTM.a2_mapTM, Turing.FinTM.a2_mapVirtual, Turing.FinTM.a2_mapVirtual_run, Turing.FinTM.a2_map_launch, Turing.FinTM.a2_map_reject, Turing.FinTM.a2_map_space, Turing.FinTM.f2_head_steps, Turing.FinTM.tapeBlocks_buffer, Turing.FinTM.tapeBlocks_left, Turing.FinTM.tapeBlocks_right, Turing.MultiTapeTM.initCfg, Turing.MultiTapeTM.runFrom, Turing.MultiTapeTM.runFrom_add, Turing.MultiTapeTM.runFrom_of_halt, Turing.MultiTapeTM.spaceUsed, Turing.MultiTapeTM.visitedByTapeHead, Turing.eq_pairEncode_of_pairDecode, Turing.length_pairEncode, Turing.pairDecode, Turing.pairEncode}
\difficulty{6}
Let $g$ be a function from Boolean words to Boolean words, and let
$T_g,S_g:\mathbb{N}\to\mathbb{N}$ be monotone. Suppose a finite Turing machine $M_g$
computes $g$ within time $T_g$, and for every Boolean word $y$ and every time
$t\in\mathbb{N}$ its computation on $y$ uses at most $S_g(|y|)$ work-tape cells. Then
there are a finite Turing machine $M$ and a constant $c\in\mathbb{N}$ such that, on
every Boolean input $z$, $M$ halts within $c(|z|+1+T_g(|z|))$ steps and outputs the
pairing of $a$ with $g(b)$ if the aligned pair parser returns $(a,b)$ on $z$, and the
empty word if parsing fails. At every time $t\in\mathbb{N}$, the computation of $M$ on
$z$ uses at most $S_g(|z|)+c(|z|+1)$ work-tape cells.
\end{theorem}

\begin{lemma}[Nonhalting throughout a run prefix]
\lean{Turing.FinTM.f2_loop_live_prefix}
\label{Turing.FinTM.f2_loop_live_prefix}
\leanok
\uses{Turing.Cfg, Turing.MultiTapeTM, Turing.MultiTapeTM.runFrom, Turing.MultiTapeTM.runFrom_add, Turing.MultiTapeTM.runFrom_of_halt}
\difficulty{3}
Let $k\in\mathbb{N}$, let $S$ be a state type, let $x$ be a Boolean input word, and let
$M$ be a $k$-work-tape Turing machine with initial configuration $c$. If the machine is
still in a state after $t$ steps, then it is still in a state after every number $u$ of
steps with $u\le t$.
\end{lemma}

\begin{lemma}[Empty output throughout a run prefix]
\lean{Turing.FinTM.f2_loop_silent_prefix}
\label{Turing.FinTM.f2_loop_silent_prefix}
\leanok
\uses{Turing.Cfg, Turing.MultiTapeTM, Turing.MultiTapeTM.output_prefix, Turing.MultiTapeTM.runFrom}
\difficulty{3}
Let $k\in\mathbb{N}$, let $S$ be a state type, let $x$ be a Boolean input word, and let
$M$ be a $k$-work-tape Turing machine with initial configuration $c$. If its output
after $t$ steps is empty, then its output is empty after every number $u$ of steps with
$u\le t$.
\end{lemma}

\begin{lemma}[First Halting Time]
\lean{Turing.FinTM.f2_loop_first_halt}
\label{Turing.FinTM.f2_loop_first_halt}
\leanok
\uses{Turing.Cfg, Turing.Cfg.Halted, Turing.MultiTapeTM, Turing.MultiTapeTM.runFrom, Turing.MultiTapeTM.runFrom_add, Turing.MultiTapeTM.runFrom_of_halt, Turing.MultiTapeTM.runFrom_zero}
\difficulty{5}
Let $M$ be a multi-tape Turing machine with $k$ work tapes over Boolean symbols, let $c$
be a configuration on a fixed Boolean input, and let $t$ be a natural number. If $c$ is
not halted but the run from $c$ is halted at time $t$, then there is a time $u$ with
$0<u\le t$ such that the run is not halted at every earlier time, is halted at time $u$,
and has the same configuration at time $u$ as at time $t$.
\end{lemma}

\begin{lemma}[Invariant Under Iterated Transitions]
\lean{Turing.FinTM.f2_loop_orbit_inv}
\label{Turing.FinTM.f2_loop_orbit_inv}
\leanok
\difficulty{4}
Let $R$ be a relation on finite Boolean words, let $T$ assign a finite Boolean word to
each pair of finite Boolean words, and let $u$ assign a finite Boolean word to each
finite Boolean word. Suppose $R(x,u(x))$ holds for every finite Boolean word $x$, and
$R(x,s)$ implies $R(x,T(x,s))$ for all finite Boolean words $x,s$. Then, for every
finite Boolean word $x$ and every $i\in\mathbb{N}$, $R(x,s_i)$ holds, where $s_i$ is
obtained by applying $s\mapsto T(x,s)$ exactly $i$ times to $u(x)$.
\end{lemma}

\begin{lemma}[Bit-string output length bound]
\lean{Turing.FinTM.f2_loop_fuel_width}
\label{Turing.FinTM.f2_loop_fuel_width}
\leanok
\uses{Turing.FinTM, Turing.FinTM.ComputesFunInTime, Turing.MultiTapeTM.output_length_le}
\difficulty{3}
Let $F$ be a finite-state Turing machine over the Boolean alphabet, and let
$R,T\colon\mathbb{N}\to\mathbb{N}$. Suppose that on every input word $y$, $F$ halts
within $T(|y|)$ steps with output the binary expansion of $R(|y|)$. Then for every
Boolean word $x$, the length of the binary expansion of $R(|x|)$ is at most $T(|x|)$.
\end{lemma}

\begin{lemma}[One-step input-head movement bound]
\lean{Turing.FinTM.f2_loop_input_move_le}
\label{Turing.FinTM.f2_loop_input_move_le}
\leanok
\uses{Turing.FinTM.moveInputPos_neg_val, Turing.moveInputPos, Turing.moveInputPos_pos_of_ne_right, Turing.moveInputPos_rightBoundary}
\difficulty{4}
Let $n\in\mathbb{N}$, let $p$ be an input-head position among the $n+2$ positions, and
let $m\in\{-1,0,1\}$ be a movement sign. After the clamped movement by $m$, the new
position is at most $p+1$.
\end{lemma}

\begin{lemma}[Input-head position bound after a run]
\lean{Turing.FinTM.f2_loop_input_run_le}
\label{Turing.FinTM.f2_loop_input_run_le}
\leanok
\uses{Turing.Action.apply, Turing.Cfg, Turing.Cfg.inputSymbol, Turing.Cfg.workTapeSymbols, Turing.FinTM.f2_loop_input_move_le, Turing.MultiTapeTM, Turing.MultiTapeTM.runFrom, Turing.MultiTapeTM.runFrom_succ_eq_step', Turing.MultiTapeTM.step}
\difficulty{4}
Let $k\in\mathbb{N}$, let $S$ be a state type, let $x$ be a Boolean input word, and let
$M$ be a $k$-work-tape Turing machine with a configuration $c$ on input $x$. After $t$
steps from $c$, the input-head position is at most its initial position plus $t$.
\end{lemma}

\begin{lemma}[Output length growth bound]
\lean{Turing.FinTM.f2_loop_output_length_le}
\label{Turing.FinTM.f2_loop_output_length_le}
\leanok
\uses{Turing.Cfg, Turing.MultiTapeTM, Turing.MultiTapeTM.outputSymbol, Turing.MultiTapeTM.runFrom, Turing.MultiTapeTM.runFrom_succ_eq_step', Turing.MultiTapeTM.step, Turing.MultiTapeTM.step_output}
\difficulty{4}
Let $k\in\mathbb{N}$, let $S$ be a state type, let $x$ be a Boolean input word, and let
$M$ be a $k$-work-tape Turing machine with configuration $c$ on input $x$. After $t$
steps, the output length is at most the initial output length plus $t$.
\end{lemma}

\begin{lemma}[Bounded rewind to the initial input position]
\lean{Turing.FinTM.f2_loop_rewind_bounded}
\label{Turing.FinTM.f2_loop_rewind_bounded}
\leanok
\uses{Turing.Cfg, Turing.FinTM.controlAction, Turing.FinTM.controlAction_apply, Turing.FinTM.moveInputPos_neg_val, Turing.FinTM.rewind_scan, Turing.MultiTapeTM, Turing.MultiTapeTM.runFrom, Turing.MultiTapeTM.runFrom_add, Turing.MultiTapeTM.step, Turing.moveInputPos}
\difficulty{5}
Let $k\in\mathbb{N}$, let $S$ be a state type, let $x$ be a Boolean input word, and let
$M$ be a $k$-work-tape Turing machine with states in $S$. Choose states $a,b\in S$ and a
possible destination $d\in\mathrm{Option}(S)$. Suppose that in state $a$, for every
scanned input and work-tape symbols, the transition is the control action that moves the
input head left and enters state $b$. Suppose also that in state $b$, it moves left and
stays in $b$ when the scanned input symbol is present, and moves left and enters $d$
when that symbol is absent; in both cases these are control actions. If a configuration
$c$ on input $x$ has state $a$, then there is a number of steps $t\le c$'s input-head
position $+2$ after which the machine's configuration has state $d$ and input-head
position $1$, with all other components equal to those of $c$.
\end{lemma}

\begin{definition}[Binary decrement with borrow]
\lean{Turing.FinTM.f2_loopDebit}
\label{Turing.FinTM.f2_loopDebit}
\leanok
For a finite Boolean list, define a transformation recursively: the empty list maps to
itself with borrow flag false; a list beginning with true maps to the list with that
first bit changed to false and borrow flag true; and a list beginning with false maps to
the list with that bit changed to true, followed by the recursively transformed tail,
with the tail's borrow flag.
\end{definition}

\begin{definition}[Initial zero-run length]
\lean{Turing.FinTM.f2_loopBorrowPos}
\label{Turing.FinTM.f2_loopBorrowPos}
\leanok
For a finite Boolean list, define its initial zero-run length to be one plus the initial
zero-run length of the tail when the list begins with false, and zero otherwise.
\end{definition}

\begin{lemma}[Bound on the initial false segment]
\lean{Turing.FinTM.f2_loopBorrowPos_le}
\label{Turing.FinTM.f2_loopBorrowPos_le}
\leanok
\uses{Turing.FinTM.f2_loopBorrowPos}
\difficulty{4}
For every finite Boolean word $u$, the length of its initial segment consisting of false
entries is at most the length of $u$.
\end{lemma}

\begin{lemma}[Length preservation by the F2 loop debit]
\lean{Turing.FinTM.f2_loopDebit_length}
\label{Turing.FinTM.f2_loopDebit_length}
\leanok
\uses{Turing.FinTM.f2_loopDebit}
\difficulty{4}
For every finite Boolean list $u$, the list component of the F2 loop debit applied to
$u$ has the same length as $u$.
\end{lemma}

\begin{definition}[Value of a Boolean list]
\lean{Turing.FinTM.f2_loopValue}
\label{Turing.FinTM.f2_loopValue}
\leanok
For a finite Boolean list, define its value recursively by assigning the empty list
value zero and assigning $2v+b$ to a list with first bit $b$ and tail value $v$, where
the Boolean bit $b$ is identified with $0$ or $1$.
\end{definition}

\begin{lemma}[Value of the Binary Bit List]
\lean{Turing.FinTM.f2_loopValue_bits}
\label{Turing.FinTM.f2_loopValue_bits}
\leanok
\uses{Turing.FinTM.f2_loopValue}
\difficulty{4}
For every natural number $n$, evaluating the recursively defined value of the bit list
$n.\mathrm{bits}$ gives $n$.
\end{lemma}

\begin{lemma}[Value after a loop debit]
\lean{Turing.FinTM.f2_loopDebit_value}
\label{Turing.FinTM.f2_loopDebit_value}
\leanok
\uses{Turing.FinTM.f2_loopDebit, Turing.FinTM.f2_loopValue}
\difficulty{4}
For every Boolean word $u$, apply the loop debit to obtain a new word and a Boolean
flag. If the flag is true, the value of the new word plus one equals the value of $u$;
if the flag is false, the value of $u$ is zero.
\end{lemma}

\begin{lemma}[Debit success criterion]
\lean{Turing.FinTM.f2_loopDebit_success}
\label{Turing.FinTM.f2_loopDebit_success}
\leanok
\uses{Turing.FinTM.f2_loopDebit, Turing.FinTM.f2_loopDebit_value, Turing.FinTM.f2_loopValue}
\difficulty{4}
For every Boolean word $u$, the debit operation reports success exactly when the value
of $u$ is positive.
\end{lemma}

\begin{lemma}[Length preservation under repeated word updates]
\lean{Turing.FinTM.f2_loopDebit_iterate_length}
\label{Turing.FinTM.f2_loopDebit_iterate_length}
\leanok
\uses{Turing.FinTM.f2_loopDebit, Turing.FinTM.f2_loopDebit_length}
\difficulty{4}
Define a map $T$ on finite Boolean words recursively by $T([]) = []$,
$T(\mathrm{true}::v) = \mathrm{false}::v$, and $T(\mathrm{false}::v) =
\mathrm{true}::T(v)$. For every finite Boolean word $u$ and every natural number $i$,
applying $T$ repeatedly $i$ times to $u$ preserves its length.
\end{lemma}

\begin{lemma}[Value after repeated debit]
\lean{Turing.FinTM.f2_loopDebit_iterate_value}
\label{Turing.FinTM.f2_loopDebit_iterate_value}
\leanok
\uses{Turing.FinTM.f2_loopDebit, Turing.FinTM.f2_loopDebit_success, Turing.FinTM.f2_loopDebit_value, Turing.FinTM.f2_loopValue, Turing.FinTM.f2_loopValue_bits}
\difficulty{5}
For natural numbers $R$ and $i$ with $i\le R$, start with the binary expansion of $R$
and apply the word-update part of the debit operation $i$ times. The value of the
resulting word is $R-i$.
\end{lemma}

\begin{lemma}[Reading the First Loop Buffer Symbol]
\lean{Turing.FinTM.f2_loopBuffer_read}
\label{Turing.FinTM.f2_loopBuffer_read}
\leanok
\uses{Turing.FinTM.bufferTape, Turing.FinTM.bufferTape_nat}
\difficulty{4}
For any two Boolean words $\mathrm{pre}$ and $\mathrm{bs}$, the buffer tape storing
their concatenation, when read at cell $|\mathrm{pre}|$, contains the first symbol of
$\mathrm{bs}$, if one exists.
\end{lemma}

\begin{lemma}[Updating the First Loop Buffer Symbol]
\lean{Turing.FinTM.f2_loopBuffer_write}
\label{Turing.FinTM.f2_loopBuffer_write}
\leanok
\uses{Turing.FinTM.bufferTape}
\difficulty{5}
For any Boolean words $\mathrm{pre}$ and $\mathrm{bs}$ and Boolean values $\mathrm{old}$
and $\mathrm{new}$, changing the buffer tape of
$\mathrm{pre}\,\mathrm{old}\,\mathrm{bs}$ at cell $|\mathrm{pre}|$ to contain
$\mathrm{new}$ gives exactly the buffer tape of
$\mathrm{pre}\,\mathrm{new}\,\mathrm{bs}$.
\end{lemma}

\begin{definition}[Binary debit machine]
\lean{Turing.FinTM.f2_loopDebitTM}
\label{Turing.FinTM.f2_loopDebitTM}
\leanok
\uses{Turing.FinTM, Turing.FinTM.controlAction}
There is a finite Turing machine with one work tape, Boolean tape symbols, and states
consisting of either an optional Boolean or a Boolean. It starts in the state containing
no pending Boolean; in that state, a scanned false is replaced by true while the head
moves right, a scanned true is replaced by false while the head moves left and records
true as pending, and a blank is left unchanged while the head moves left and records
false as pending. With a pending Boolean, a scanned symbol is left unchanged while the
head moves left; at a blank, the head moves right and the machine enters the state
carrying that Boolean. In a state carrying a Boolean in the second component, the
transition is the specified control action.
\end{definition}

\begin{definition}[Configuration for the debit loop]
\lean{Turing.FinTM.f2_loopDebitCfg}
\label{Turing.FinTM.f2_loopDebitCfg}
\leanok
\uses{Turing.Cfg, Turing.FinTM.bufferTape, Turing.FinTM.f2_loopDebitTM}
Given a Boolean input list $x$, a valid input-head position $p$ among its $|x|+2$
positions, a state $q$ that is either an optional Boolean or a Boolean, an integer $z$,
and a Boolean list $u$, define a configuration for the one-tape debit machine with state
$q$, input-head position $p$, work tape initialized from $u$ by the machine's buffer
encoding, work-head position $z$, and empty output.
\end{definition}

\begin{lemma}[One step of the borrow computation]
\lean{Turing.FinTM.f2_loopBorrow_step}
\label{Turing.FinTM.f2_loopBorrow_step}
\leanok
\uses{Turing.Action.apply, Turing.Cfg.workTapeSymbols, Turing.FinTM.f2_loopBuffer_read, Turing.FinTM.f2_loopBuffer_write, Turing.FinTM.f2_loopDebitCfg, Turing.FinTM.f2_loopDebitTM, Turing.MultiTapeTM.step, Turing.moveInputPos_zero}
\difficulty{5}
Let $x$ be a Boolean input, let $p$ be an allowed input-head position, and let $v$ and
$s$ be finite Boolean words. If the one-tape debit machine is in its initial search
state, its head is at position $|v|$, and its tape contains $v s$, then one step yields
the following configuration: if $s$ is empty, it enters the borrow-result state carrying
false, moves its head to $|v|-1$, and leaves the tape as $v$; if $s$ begins with true,
it enters the borrow-result state carrying true, moves its head to $|v|-1$, and changes
that first symbol to false; if $s$ begins with false, it remains in the initial search
state, moves its head to $|v|+1$, and changes that first symbol to true.
\end{lemma}

\begin{lemma}[Run of the borrow computation]
\lean{Turing.FinTM.f2_loopBorrow_run}
\label{Turing.FinTM.f2_loopBorrow_run}
\leanok
\uses{Turing.FinTM.f2_loopBorrowPos, Turing.FinTM.f2_loopBorrow_step, Turing.FinTM.f2_loopDebit, Turing.FinTM.f2_loopDebitCfg, Turing.FinTM.f2_loopDebitTM, Turing.MultiTapeTM.runFrom, Turing.MultiTapeTM.runFrom_succ_eq_step}
\difficulty{4}
Let $x$ be a Boolean input, let $p$ be an allowed input-head position, and let $u$ and
$v$ be finite Boolean words. Starting the one-tape debit machine in its initial search
state, with its head at position $|v|$ and its tape containing the concatenation $v u$,
after $b+1$ steps it is in the borrow-result state carrying the Boolean outcome of the
borrow operation on $u$, its head is at position $|v|+b-1$, and its tape contains $v$
followed by the resulting word, where $b$ is the length of the initial segment of $u$
consisting of false entries.
\end{lemma}

\begin{lemma}[Rewinding after a borrow]
\lean{Turing.FinTM.f2_loopBorrow_rewind}
\label{Turing.FinTM.f2_loopBorrow_rewind}
\leanok
\uses{Turing.Action.apply, Turing.Cfg.workTapeSymbols, Turing.FinTM.bufferTape_left, Turing.FinTM.bufferTape_nat, Turing.FinTM.f2_loopDebitCfg, Turing.FinTM.f2_loopDebitTM, Turing.MultiTapeTM.runFrom, Turing.MultiTapeTM.runFrom_succ_eq_step, Turing.MultiTapeTM.runFrom_zero, Turing.MultiTapeTM.step, Turing.moveInputPos_zero}
\difficulty{5}
Let $x$ be a Boolean input, let $p$ be an allowed input-head position, let $u$ be a
finite Boolean word, and let $b$ be a Boolean. For every integer $j$ with $0\le j\le
|u|$, the one-tape debit machine started in the borrow-result state carrying $b$, with
its head at position $j-1$ and tape containing $u$, reaches after $j+1$ steps the
terminal result state carrying $b$, with its head at zero and its tape still containing
$u$.
\end{lemma}

\begin{lemma}[Correctness of the borrow computation]
\lean{Turing.FinTM.f2_loopBorrow_correct}
\label{Turing.FinTM.f2_loopBorrow_correct}
\leanok
\uses{Turing.FinTM.f2_loopBorrowPos, Turing.FinTM.f2_loopBorrowPos_le, Turing.FinTM.f2_loopBorrow_rewind, Turing.FinTM.f2_loopBorrow_run, Turing.FinTM.f2_loopDebit, Turing.FinTM.f2_loopDebitCfg, Turing.FinTM.f2_loopDebitTM, Turing.FinTM.f2_loopDebit_length, Turing.MultiTapeTM.runFrom, Turing.MultiTapeTM.runFrom_add}
\difficulty{4}
Let $x$ be a Boolean input, let $p$ be one of its allowed input-head positions, and let
$u$ be a finite Boolean word. If $b$ is the length of the initial segment of $u$
consisting of false entries, then $2b+2\le 2|u|+2$. Moreover, the one-tape debit
machine, started in its initial search state with its head at cell zero and the tape
containing $u$, reaches after $2b+2$ steps the configuration with its result state equal
to the Boolean outcome of the borrow operation on $u$, its head at cell zero, and its
tape containing the resulting word.
\end{lemma}

\begin{definition}[Loop body machine construction]
\lean{Turing.FinTM.f2_loopBodyTM}
\label{Turing.FinTM.f2_loopBodyTM}
\leanok
\uses{Turing.FinTM}
Given a finite Boolean-alphabet Turing machine $B$ and a state $a$ of $B$, construct a
machine with one additional work tape and states $(q,b)$, where $q$ is a state of $B$
and $b$ is Boolean. Its initial state is $(q_0,\mathrm{false})$. When the state is
$(a,\mathrm{false})$, it halts, leaves the input head stationary, leaves the original
work tapes unchanged, and writes false on the added tape without moving its head. In
every other state, it performs $B$'s transition using the original work tapes; it
preserves that action on the input and original work tapes, emits the same output, and,
if $B$ halts on that transition, writes true on the added tape while leaving its head
stationary. If $B$ continues in state $s$, the new state is $(s,\mathrm{false})$, and
the added tape is left unchanged.
\end{definition}

\begin{definition}[Loop Body Configuration Extension]
\lean{Turing.FinTM.f2_loopBodyCfg}
\label{Turing.FinTM.f2_loopBodyCfg}
\leanok
\uses{Turing.Cfg, Turing.FinTM, Turing.FinTM.f2_loopBodyTM}
Given a finite Boolean Turing machine $B$, a state $a$ of $B$, an input $x$, a
configuration $c$ of $B$ on $x$, a Boolean release value $r$, and an optional Boolean
flag value $f$, form a configuration for the loop-body machine with state obtained by
pairing the state of $c$ with $r$. It keeps the input-head position, output, and all
work tapes and head positions of $c$, and adds one work tape whose head is at position
$0$ and whose only possibly nonblank cell is at position $0$, with contents $f$.
\end{definition}

\begin{lemma}[Stopping at the anchor state]
\lean{Turing.FinTM.f2_loopBody_stop}
\label{Turing.FinTM.f2_loopBody_stop}
\leanok
\uses{Turing.Action.apply, Turing.Cfg, Turing.FinTM, Turing.FinTM.f2_loopBodyCfg, Turing.FinTM.f2_loopBodyTM, Turing.MultiTapeTM.step, Turing.moveInputPos_zero}
\difficulty{4}
Let $B$ be a finite Turing machine over the Boolean alphabet, let $a$ be one of its
states, and let $c$ be a configuration on a Boolean input whose state is $a$. For any
optional Boolean flag on the added tape, one step of the machine with the extra tape and
release flag set to false yields the lifted configuration with the base machine halted,
release flag false, and the added tape's cell zero set to false.
\end{lemma}

\begin{lemma}[One step of the loop-body machine]
\lean{Turing.FinTM.f2_loopBody_step}
\label{Turing.FinTM.f2_loopBody_step}
\leanok
\uses{Turing.Action, Turing.Action.apply, Turing.Cfg, Turing.Cfg.inputSymbol, Turing.Cfg.workTapeSymbols, Turing.FinTM, Turing.FinTM.f2_loopBodyCfg, Turing.FinTM.f2_loopBodyTM, Turing.MultiTapeTM.step}
\difficulty{6}
Let $B$ be a finite Turing machine over the Boolean alphabet, let $a$ and $q$ be states
of $B$, and let $c$ be a configuration of $B$ on a Boolean input with state $q$. Let $r$
be a Boolean and let the extra tape's cell zero contain an optional Boolean flag. If $q$
is not the anchor state with release flag false, then one step of the machine formed by
adding the release flag and extra tape yields the lift of one step of $B$ from $c$, with
release flag false. The extra tape's cell zero contains true if the base machine's next
configuration is halted, and otherwise retains its previous contents.
\end{lemma}

\begin{lemma}[Run of the loop-body machine]
\lean{Turing.FinTM.f2_loopBody_run}
\label{Turing.FinTM.f2_loopBody_run}
\leanok
\uses{Turing.Cfg, Turing.FinTM, Turing.FinTM.f2_loopBodyCfg, Turing.FinTM.f2_loopBodyTM, Turing.FinTM.f2_loopBody_step, Turing.MultiTapeTM.runFrom, Turing.MultiTapeTM.runFrom_succ_eq_step'}
\difficulty{5}
Let $B$ be a finite Turing machine over the Boolean alphabet, let $a$ be a state of $B$,
and let $x$ be a Boolean input. Let $c$ be a configuration of $B$ on $x$ whose state is
present, let $r$ be a Boolean, and let $t\in\mathbb{N}$. Suppose that for every $u<t$,
the configuration reached from $c$ after $u$ steps has a present state, and that its
state is not $a$ except possibly when $u=0$ and $r$ is true. Then after $t$ steps, the
machine formed by adding one work tape and a Boolean release flag to the control state
has the configuration obtained from the base machine's configuration after $t$ steps:
the release flag is $r$ if $t=0$ and false otherwise, and the added tape has true at
cell zero exactly when the base machine has halted by time $t$, and is blank there
otherwise.
\end{lemma}

\begin{lemma}[Captured loop-body computation]
\lean{Turing.FinTM.f2_loopBody_capture}
\label{Turing.FinTM.f2_loopBody_capture}
\leanok
\uses{Turing.Cfg, Turing.Cfg.Halted, Turing.FinTM, Turing.FinTM.f2_loopBodyCfg, Turing.FinTM.f2_loopBodyTM, Turing.FinTM.f2_loopBody_run, Turing.MultiTapeTM, Turing.MultiTapeTM.runFrom, Turing.captureAction, Turing.captureCfg, Turing.capture_run}
\difficulty{4}
Let $\mathrm{body}$ be a Boolean finite Turing machine with anchor state $a$, and let
$H$ be a state set with a machine whose transitions agree with the capture transform of
the loop-body machine's transitions after applying an embedding $e$ and return state
$r$. Suppose $c$ is a configuration of $\mathrm{body}$ on input $x$, its state is
present, and for every $u<t$ the run from $c$ is still active and has not reached $a$.
Suppose also that for every $u<t$, either $u=0$ and the release flag is true, or the run
from $c$ is not in state $a$. Then after $t$ steps the host machine from the captured
loop-body configuration agrees with the capture of the body machine's configuration
after $t$ steps; the release flag in that configuration is the original flag if $t=0$
and false otherwise, and its extra tape contains the halt marker exactly when the body
computation has halted.
\end{lemma}

\begin{definition}[Loop Host State Space]
\lean{Turing.FinTM.f2_LoopHostState}
\label{Turing.FinTM.f2_LoopHostState}
\leanok
\uses{Turing.FinTM}
For finite Boolean Turing machines $B$ and $F$, the loop host state space is the
disjoint union of the states of $F$, the states consisting of a startup bit together
with a state of $B$ and a release bit, and a set of fourteen control states.
\end{definition}

\begin{definition}[Loop fuel source machine]
\lean{Turing.FinTM.f2_loopFuelSource}
\label{Turing.FinTM.f2_loopFuelSource}
\leanok
\uses{Turing.FinTM, Turing.FinTM.rightAction, Turing.MultiTapeTM}
Given finite Turing machines $B$ and $F$ over the Boolean alphabet, there is a
multi-tape machine with $B.k+1+(1+F.k)$ work tapes and the state set of $F$. Its initial
state is that of $F$, and its transition uses the transition of $F$ on the same state
and input symbol, with the symbols on the final $F.k$ work tapes as its work-tape input;
the action is extended to account for the additional work tapes.
\end{definition}

\begin{definition}[Loop Body Source Machine]
\lean{Turing.FinTM.f2_loopBodySource}
\label{Turing.FinTM.f2_loopBodySource}
\leanok
\uses{Turing.FinTM, Turing.FinTM.f2_loopBodyTM, Turing.FinTM.leftAction, Turing.MultiTapeTM}
Let $B$ and $F$ be finite Turing machines over the Boolean alphabet, with $k_B$ and
$k_F$ work tapes, respectively, and let $a$ be a state of $B$. The loop body source
machine has $k_B+1+(1+k_F)$ work tapes and states in
$Q_B\times\{\mathrm{false},\mathrm{true}\}$; its initial state is
$(q_0^B,\mathrm{false})$. On a transition, it applies the transition of the loop body
machine to the input and the first $k_B+1$ work tapes, using the corresponding initial
segment of the larger machine's work-tape contents, and extends that action by the
specified left-action operation over the remaining $1+k_F$ tapes.
\end{definition}

\begin{definition}[Loop control action]
\lean{Turing.FinTM.f2_loopControlAction}
\label{Turing.FinTM.f2_loopControlAction}
\leanok
\uses{Turing.Action, Turing.FinTM, Turing.FinTM.f2_LoopHostState}
Let $B$ and $F$ be finite Turing machines over the Boolean alphabet. Given an input-head
movement, a flag-tape instruction, counter-tape and payload-tape instructions, an
optional output bit, and an optional next state in $F$'s state set or a state of the
form $(b,(q,c))$ with $b,c\in\{\mathrm{false},\mathrm{true}\}$ and $q$ a state of $B$,
define an action on $B.k+1+(1+F.k)+1$ work tapes. The flag instruction is assigned to
tape $B.k$, the counter instruction to tape $B.k+1$, and the payload instruction to tape
$B.k+1+(1+F.k)$; every other work tape receives no write and no movement. The specified
input-head movement, output bit, and next state are used unchanged.
\end{definition}

\begin{definition}[Boolean loop host machine]
\lean{Turing.FinTM.f2_loopHost}
\label{Turing.FinTM.f2_loopHost}
\leanok
\uses{Turing.FinTM, Turing.FinTM.f2_LoopHostState, Turing.FinTM.f2_loopBodySource, Turing.FinTM.f2_loopControlAction, Turing.FinTM.f2_loopFuelSource, Turing.captureAction}
Given finite Boolean Turing machines $B$ and $F$, a state $a$ of $B$, and a Boolean mode
$m$, define a finite Boolean Turing machine with $B.k+1+(1+F.k)+1$ work tapes. Its
initial state is the initial state of $F$. Its transition function combines a fuel phase
and calls to the loop body, using additional tapes to track a flag, a counter, and a
payload; depending on the mode and control phase, it emits Boolean symbols or halts.
\end{definition}

\begin{lemma}[Capture of a live body run]
\lean{Turing.FinTM.f2_loopHost_body_capture}
\label{Turing.FinTM.f2_loopHost_body_capture}
\leanok
\uses{Turing.Cfg, Turing.Cfg.Halted, Turing.FinTM, Turing.FinTM.f2_loopBodySource, Turing.FinTM.f2_loopHost, Turing.MultiTapeTM.runFrom, Turing.captureCfg, Turing.capture_run}
\difficulty{1}
Let $B$ and $F$ be finite Boolean Turing machines, let $a$ be a state of $B$, and choose
find and startup modes. For an input word $x$, configuration $c$ of the padded body
machine on $x$, and time $t$, suppose the body-source run from $c$ is not halted at
every time $u<t$. Then the host machine, started with the capture of $c$ using the
corresponding body-call state, reaches after $t$ steps exactly the capture of the
body-source configuration reached from $c$ after $t$ steps.
\end{lemma}

\begin{lemma}[Fuel simulation under capture]
\lean{Turing.FinTM.f2_loopHost_fuel_capture}
\label{Turing.FinTM.f2_loopHost_fuel_capture}
\leanok
\uses{Turing.Cfg, Turing.Cfg.Halted, Turing.FinTM, Turing.FinTM.f2_loopFuelSource, Turing.FinTM.f2_loopHost, Turing.MultiTapeTM.runFrom, Turing.captureCfg, Turing.capture_run}
\difficulty{1}
Let $B$ and $F$ be finite Boolean Turing machines, let $a$ be a state of $B$, and let
$m$ be either Boolean mode. For any input $x$, fuel-source configuration $c$, and time
$t$, suppose the fuel-source run from $c$ is not halted at any time $u<t$. Then the
loop-host run for $t$ steps from the capture configuration of $c$ agrees exactly with
the capture configuration of the fuel-source run after $t$ steps.
\end{lemma}

\begin{lemma}[Initial configuration of the loop host]
\lean{Turing.FinTM.f2_loopHost_init}
\label{Turing.FinTM.f2_loopHost_init}
\leanok
\uses{Turing.Cfg.ofWords, Turing.FinTM, Turing.FinTM.f2_loopFuelSource, Turing.FinTM.f2_loopHost, Turing.MultiTapeTM.initCfg, Turing.captureCfg, Turing.initCfg_ofWords}
\difficulty{3}
For Boolean Turing machines $B$ and $F$, a state $a$ of $B$, a mode bit, and an input
word $x$, the loop host started on $x$ is in the configuration obtained by capturing the
initial configuration of its fuel source, with the state embedding from $F$ into the
host state space, empty prefix and output, and control return state $0$.
\end{lemma}

\begin{lemma}[Idle Loop Control Action]
\lean{Turing.FinTM.f2_loopControl_idle}
\label{Turing.FinTM.f2_loopControl_idle}
\leanok
\uses{Turing.FinTM, Turing.FinTM.controlAction, Turing.FinTM.f2_LoopHostState, Turing.FinTM.f2_loopControlAction}
\difficulty{1}
For finite Boolean Turing machines $\mathrm{body}$ and $F$, any input-head movement
$\mathrm{inp}$, and any optional next loop-host state, the loop-control action that
makes no writes or movements on its work tapes and emits no output is the control action
that makes the same input-head movement and passes to that next state.
\end{lemma}

\begin{lemma}[Rewinding the input head]
\lean{Turing.FinTM.f2_loopHost_input_rewind}
\label{Turing.FinTM.f2_loopHost_input_rewind}
\leanok
\uses{Turing.Cfg, Turing.FinTM, Turing.FinTM.f2_loopControl_idle, Turing.FinTM.f2_loopHost, Turing.FinTM.f2_loop_rewind_bounded, Turing.MultiTapeTM.runFrom}
\difficulty{4}
Let $B$ and $F$ be Boolean finite Turing machines, let $a$ be a state of $B$, choose a
mode bit and an input word $x$, and let $c$ be a configuration of the corresponding loop
host on $x$. If $c$ is in control phase $4$, then after some $t\leq c$'s input-head
position plus $2$ steps, the host is in the startup body-call state with the body in its
initial state and its input head at position $1$; all other configuration data are
unchanged.
\end{lemma}

\begin{definition}[Loop frame configuration]
\lean{Turing.FinTM.f2_loopFrame}
\label{Turing.FinTM.f2_loopFrame}
\leanok
\uses{Turing.Cfg, Turing.FinTM, Turing.FinTM.f2_LoopHostState}
Let $B$ and $F$ be finite Turing machines over the Boolean alphabet, let $x$ be a
Boolean input, and let a base configuration on $x$ have $B.k+1+(1+F.k)+1$ work tapes and
states from $F$’s state set or from
$(	ext{Bool}	imes(B.	ext{State}	imes	ext{Bool}))	imes	ext{Fin}(14)$. Choose an optional
state $q$ from this combined state set, an input-head position $p$ among the $|x|+2$
positions, three work-tape contents $f,c,u:br	o	ext{Option}(	ext{Bool})$, head
positions $h_c,h_u\inbr$, and an output string $o$ of Booleans. The resulting
configuration has state $q$, input-head position $p$, and output $o$; its tapes at
indices $B.k$, $B.k+1$, and $B.k+1+(1+F.k)$ have contents $f,c,u$ and head positions
$0,h_c,h_u$, respectively, while every other tape and head position agrees with the base
configuration.
\end{definition}

\begin{definition}[Conditional Tape Write]
\lean{Turing.FinTM.f2_loopWrite}
\label{Turing.FinTM.f2_loopWrite}
\leanok
For a Boolean tape $t$ and head position $h$, a write instruction either leaves the tape
unchanged or replaces its value at $h$ with the specified optional Boolean value, where
an absent Boolean denotes a blank cell.
\end{definition}

\begin{lemma}[Loop Control Action on a Framed Configuration]
\lean{Turing.FinTM.f2_loopControl_apply}
\label{Turing.FinTM.f2_loopControl_apply}
\leanok
\uses{Turing.Action.apply, Turing.Cfg, Turing.FinTM, Turing.FinTM.f2_LoopHostState, Turing.FinTM.f2_loopControlAction, Turing.FinTM.f2_loopFrame, Turing.FinTM.f2_loopWrite, Turing.moveInputPos}
\difficulty{5}
Let $\mathrm{body}$ and $F$ be finite Boolean Turing machines, let $x$ be a Boolean
input word, and choose a base configuration, optional control state $q$, input-head
position $p$, three tape contents $\mathrm{flag},\mathrm{counter},\mathrm{payload}$,
head positions $\mathrm{ch},\mathrm{ph}$, output word $\mathrm{out}$, input movement
$\mathrm{inp}$, optional write values $\mathrm{fw},\mathrm{ca},\mathrm{pa}$ for the
flag, counter, and payload tapes, optional emitted Boolean symbol $\mathrm{emit}$, and
optional next state. Applying the loop-control action to the frame configuration
produces the frame with next state, input position moved from $p$ by $\mathrm{inp}$, the
flag tape updated at cell $0$ according to $\mathrm{fw}$, the counter and payload tapes
updated at their respective head positions according to the write components of
$\mathrm{ca}$ and $\mathrm{pa}$, head positions advanced by their movement components,
and output $\mathrm{out}$ followed by $\mathrm{emit}$ if present; all remaining frame
data are unchanged.
\end{lemma}

\begin{definition}[Boolean loop replay machine]
\lean{Turing.FinTM.f2_loopReplayTM}
\label{Turing.FinTM.f2_loopReplayTM}
\leanok
\uses{Turing.FinTM}
Consider a Turing machine with one work tape, Boolean symbols, and a single state. On
every transition it leaves the input head in place and makes no change or movement on
the work tape. If the scanned work-tape cell contains a Boolean value, the machine
outputs that value and remains in its state; if the cell is blank, it outputs nothing
and halts.
\end{definition}

\begin{definition}[Loop replay configuration]
\lean{Turing.FinTM.f2_loopReplayCfg}
\label{Turing.FinTM.f2_loopReplayCfg}
\leanok
\uses{Turing.Cfg, Turing.FinTM.bufferTape}
Given a Boolean input string $x$, an input-head position $p$ among its $|x|+2$
positions, an optional state $q$, an integer $z$, and Boolean lists $w$ and $o$, form a
configuration with state $q$, input-head position $p$, output string $o$, and its single
work tape equal to the list-buffer encoding of $w$, with its head at position $z$.
\end{definition}

\begin{lemma}[One Replay Step]
\lean{Turing.FinTM.f2_loopReplay_step}
\label{Turing.FinTM.f2_loopReplay_step}
\leanok
\uses{Turing.Action.apply, Turing.Cfg.workTapeSymbols, Turing.FinTM.f2_loopBuffer_read, Turing.FinTM.f2_loopReplayCfg, Turing.FinTM.f2_loopReplayTM, Turing.MultiTapeTM.step, Turing.moveInputPos_zero}
\difficulty{4}
Let $x$ be a Boolean input word, let $p$ be an input-head position, and let $u$, $v$,
and $o$ be Boolean words. If the replay machine is active, its work tape contains $uv$,
and its head is at position $|u|$, then one step halts with the head still at $|u|$, the
tape containing $u$, and output unchanged when $v$ is empty. If $v=b::v'$ for a bit $b$,
one step leaves the machine active, advances the head to $|u|+1$, keeps the tape equal
to $uv$, and appends $b$ to the output.
\end{lemma}

\begin{lemma}[Replay run on a Boolean word]
\lean{Turing.FinTM.f2_loopReplay_run}
\label{Turing.FinTM.f2_loopReplay_run}
\leanok
\uses{Turing.FinTM.f2_loopReplayCfg, Turing.FinTM.f2_loopReplayTM, Turing.FinTM.f2_loopReplay_step, Turing.MultiTapeTM.runFrom, Turing.MultiTapeTM.runFrom_succ_eq_step}
\difficulty{4}
Let $x$ and $r$ be Boolean words, and let $p$ be a valid input-head position for $x$.
For all Boolean words $u,v$ and output words $o$, the replay machine started on input
$x$ with input-head position $p$, active state, work tape containing $uv$ with its head
at position $|u|$, and output $o$ reaches after $|v|+1$ steps a halted configuration
with input-head position $p$, work tape still containing $uv$ with its head at position
$|uv|$, and output $ov$.
\end{lemma}

\begin{lemma}[Loop Control Payload Action]
\lean{Turing.FinTM.f2_loopControl_payload}
\label{Turing.FinTM.f2_loopControl_payload}
\leanok
\uses{Turing.Action, Turing.FinTM, Turing.FinTM.f2_LoopHostState, Turing.FinTM.f2_loopControlAction, Turing.FinTM.rightAction}
\difficulty{4}
For finite Boolean Turing machines $\mathrm{body}$ and $F$, any work-head movement $d$,
optional emitted Boolean symbol, and optional next loop-host state, the loop-control
action with no input-head movement and no actions on the flag and counter tapes, and
with payload-head movement $d$ but no payload write, equals the right-block extension of
the one-tape action that has no input-head movement, no write and movement $d$ on its
sole work tape, and the specified output and next state.
\end{lemma}

\begin{lemma}[Replaying a stored word]
\lean{Turing.FinTM.f2_loopHost_replay}
\label{Turing.FinTM.f2_loopHost_replay}
\leanok
\uses{Turing.FinTM, Turing.FinTM.f2_loopControlAction, Turing.FinTM.f2_loopControl_payload, Turing.FinTM.f2_loopHost, Turing.FinTM.f2_loopReplayCfg, Turing.FinTM.f2_loopReplayTM, Turing.FinTM.f2_loopReplay_run, Turing.FinTM.rightAction, Turing.FinTM.rightCfg, Turing.FinTM.rightCfg_run, Turing.MultiTapeTM.runFrom}
\difficulty{5}
Let $B$ and $F$ be Boolean finite Turing machines, let $a$ be a state of $B$, choose a
mode bit and an input word $x$, an input-head position $p$, words $w$ and $o$, and
arbitrary contents and head positions for the other work tapes. Starting with the replay
machine's sole tape storing $w$ and its head at $0$, control phase $13$, and output $o$,
after $|w|+1$ steps the loop host remains in phase $13$, its tape still stores $w$ with
head at position $|w|$, and its output is $o$ followed by $w$; all other specified data
are preserved.
\end{lemma}

\begin{definition}[Loop fuel configuration]
\lean{Turing.FinTM.f2_loopFuelCfg}
\label{Turing.FinTM.f2_loopFuelCfg}
\leanok
\uses{Turing.Cfg, Turing.FinTM, Turing.FinTM.rightCfg}
Let $body$ and $F$ be finite Turing machines over the Boolean alphabet, let $x$ be a
finite Boolean input, and let $c$ be a configuration of $F$ on input $x$. The
construction defines a configuration on the same input with $body.k+1+(1+F.k)$ work
tapes by adjoining $1+(body.k+1)$ blank tapes to the work tapes of $c$, with each added
tape head at position zero; the state, input-head position, and output are preserved.
\end{definition}

\begin{lemma}[Run of the loop fuel machine]
\lean{Turing.FinTM.f2_loopFuel_run}
\label{Turing.FinTM.f2_loopFuel_run}
\leanok
\uses{Turing.Cfg, Turing.FinTM, Turing.FinTM.f2_loopFuelCfg, Turing.FinTM.f2_loopFuelSource, Turing.FinTM.rightAction, Turing.FinTM.rightCfg_run, Turing.MultiTapeTM, Turing.MultiTapeTM.runFrom}
\difficulty{3}
Let $B$ and $F$ be finite Boolean Turing machines, let $x$ be a Boolean input, let $c$
be a configuration of $F$ on $x$, and let $t$ be a natural number. Running the loop fuel
source machine for $t$ steps from the loop fuel configuration built from $c$ yields the
loop fuel configuration built from the configuration reached by running $F$ for $t$
steps from $c$.
\end{lemma}

\begin{lemma}[Initial configuration of the loop fuel machine]
\lean{Turing.FinTM.f2_loopFuel_init}
\label{Turing.FinTM.f2_loopFuel_init}
\leanok
\uses{Turing.Cfg.init, Turing.FinTM, Turing.FinTM.f2_loopFuelCfg, Turing.FinTM.f2_loopFuelSource, Turing.FinTM.rightCfg, Turing.MultiTapeTM.initCfg}
\difficulty{4}
For finite Boolean machines $B$ and $F$ and any Boolean input $x$, the initial
configuration of the loop fuel source machine equals the loop fuel configuration built
from the initial configuration of $F$ on $x$.
\end{lemma}

\begin{lemma}[Payload tape in a loop frame]
\lean{Turing.FinTM.f2_loopFrame_payload}
\label{Turing.FinTM.f2_loopFrame_payload}
\leanok
\uses{Turing.Cfg, Turing.Cfg.workTapeSymbols, Turing.FinTM, Turing.FinTM.f2_LoopHostState, Turing.FinTM.f2_loopFrame}
\difficulty{3}
Let $B$ and $F$ be finite Turing machines over the Boolean alphabet, with state sets
$Q_B,Q_F$ and respectively $b,f$ work tapes. Use the state set $Q_F \sqcup
(\{0,1\}\times(Q_B\times\{0,1\})\sqcup\{0,1,\ldots,13\})$ for a host machine, and let
$x$ be a finite Boolean input. Take a base configuration on $b+1+(1+f)+1$ work tapes
with this host state set, and choose any optional host state, valid input-head position,
three tape-content functions from $\mathbb{Z}$ to optional Boolean symbols (for the
flag, counter, and payload tapes), integer counter and payload head positions, and
Boolean output word. Form the loop-frame configuration by placing the flag tape at index
$b$, the counter tape at index $b+1$, and the payload tape at index $b+1+(1+f)$, with
the flag head at $0$ and the other two heads at their chosen positions, and using the
chosen state, input-head position, and output. The symbol scanned by the last work tape
of this configuration is the payload tape's content at the chosen payload-head position.
\end{lemma}

\begin{lemma}[Counter tape in an F2 loop frame]
\lean{Turing.FinTM.f2_loopFrame_counter}
\label{Turing.FinTM.f2_loopFrame_counter}
\leanok
\uses{Turing.Cfg, Turing.Cfg.workTapeSymbols, Turing.FinTM, Turing.FinTM.f2_LoopHostState, Turing.FinTM.f2_loopFrame}
\difficulty{1}
Let $B$ and $F$ be finite Turing machines over the Boolean alphabet, and let $x$ be a
finite Boolean input. For any configuration with $B.k+1+(1+F.k)+1$ work tapes and states
in $F.State \oplus ((\{0,1\}\times(B.State\times\{0,1\}))\oplus\mathrm{Fin}(14))$, any
optional state and valid input-head position, any tape contents
$\mathrm{flag},\mathrm{counter},\mathrm{payload}:\mathbb{Z}\to\mathrm{Option}(\{0,1\})$,
head positions $ch,ph\in\mathbb{Z}$, and output word, the symbol scanned on work tape
$B.k+1$ in the resulting loop-frame configuration is $\mathrm{counter}(ch)$.
\end{lemma}

\begin{lemma}[Fuel tape rewind]
\lean{Turing.FinTM.f2_loopHost_fuel_rewind}
\label{Turing.FinTM.f2_loopHost_fuel_rewind}
\leanok
\uses{Turing.Action.apply, Turing.Cfg, Turing.Cfg.workTapeSymbols, Turing.FinTM, Turing.FinTM.bufferTape, Turing.FinTM.bufferTape_left, Turing.FinTM.bufferTape_nat, Turing.FinTM.f2_LoopHostState, Turing.FinTM.f2_loopControlAction, Turing.FinTM.f2_loopControl_apply, Turing.FinTM.f2_loopFrame, Turing.FinTM.f2_loopFrame_payload, Turing.FinTM.f2_loopHost, Turing.FinTM.f2_loopWrite, Turing.MultiTapeTM.runFrom, Turing.MultiTapeTM.runFrom_succ_eq_step, Turing.MultiTapeTM.runFrom_zero, Turing.MultiTapeTM.step}
\difficulty{5}
Let $B$ and $F$ be finite Boolean Turing machines, let $a$ be a state of $B$, and let
$m$ be either Boolean mode. Choose a host base configuration, an input-head position,
flag and counter tapes, an integer counter-head position, a word $w$, and an output
word. For every integer $j$ with $0\le j\le |w|$, a host configuration in phase $1$ with
payload tape storing $w$, payload head at $j-1$, and the specified counter head reaches
phase $2$ after $j+1$ steps, with the payload tape still storing $w$ and its head at
zero; the other specified components are unchanged.
\end{lemma}

\begin{definition}[Loop Copy Tape Lookup]
\lean{Turing.FinTM.f2_loopCopyTape}
\label{Turing.FinTM.f2_loopCopyTape}
\leanok
\uses{Turing.FinTM.bufferTape}
For Boolean lists $pre$ and $rest$ and an integer position $z$, let $B(L,z)$ denote the
list-buffer lookup at position $z$ in the list $L$. Define the tape lookup to be blank
when $z<|pre|$, and otherwise to be $B(pre ++ rest,z)$.
\end{definition}

\begin{lemma}[Reading the next copy symbol]
\lean{Turing.FinTM.f2_loopCopy_read}
\label{Turing.FinTM.f2_loopCopy_read}
\leanok
\uses{Turing.FinTM.f2_loopBuffer_read, Turing.FinTM.f2_loopCopyTape}
\difficulty{1}
For any Boolean words $p$ and $r$, the tape $T(p,r)$ contains the first symbol of $r$ at
cell $|p|$, with a blank there when $r$ is empty.
\end{lemma}

\begin{lemma}[Erasing the copied tape symbol]
\lean{Turing.FinTM.f2_loopCopy_erase}
\label{Turing.FinTM.f2_loopCopy_erase}
\leanok
\uses{Turing.FinTM.f2_loopCopyTape}
\difficulty{4}
Let $\mathrm{buffer}(w)$ denote the tape containing the word $w$ from cell zero onward,
and let $T(p,r)$ be the tape that is blank before cell $|p|$ and otherwise contains $p$
followed by $r$. For Boolean words $p,r$ and a Boolean symbol $b$, erasing cell $|p|$
from $T(p,b::r)$ gives exactly $T(p++[b],r)$.
\end{lemma}

\begin{lemma}[Initial copy-loop tape]
\lean{Turing.FinTM.f2_loopCopy_initial}
\label{Turing.FinTM.f2_loopCopy_initial}
\leanok
\uses{Turing.FinTM.bufferTape, Turing.FinTM.f2_loopCopyTape}
\difficulty{4}
For every Boolean word $w$, the tape $T([],w)$ is $\mathrm{buffer}(w)$.
\end{lemma}

\begin{lemma}[Final copy-loop tape]
\lean{Turing.FinTM.f2_loopCopy_final}
\label{Turing.FinTM.f2_loopCopy_final}
\leanok
\uses{Turing.FinTM.bufferTape, Turing.FinTM.f2_loopCopyTape}
\difficulty{4}
For every Boolean word $w$, the tape $T(w,[])$ is the blank tape $\mathrm{buffer}([])$.
\end{lemma}

\begin{lemma}[Fuel tape copying]
\lean{Turing.FinTM.f2_loopHost_fuel_copy}
\label{Turing.FinTM.f2_loopHost_fuel_copy}
\leanok
\uses{Turing.Action.apply, Turing.Cfg, Turing.Cfg.workTapeSymbols, Turing.FinTM, Turing.FinTM.bufferTape, Turing.FinTM.bufferTape_append, Turing.FinTM.f2_LoopHostState, Turing.FinTM.f2_loopControlAction, Turing.FinTM.f2_loopControl_apply, Turing.FinTM.f2_loopCopyTape, Turing.FinTM.f2_loopCopy_erase, Turing.FinTM.f2_loopCopy_final, Turing.FinTM.f2_loopCopy_read, Turing.FinTM.f2_loopFrame, Turing.FinTM.f2_loopFrame_payload, Turing.FinTM.f2_loopHost, Turing.FinTM.f2_loopWrite, Turing.MultiTapeTM.runFrom, Turing.MultiTapeTM.runFrom_succ_eq_step, Turing.MultiTapeTM.runFrom_zero, Turing.MultiTapeTM.step}
\difficulty{6}
Let $B$ and $F$ be finite Boolean Turing machines, let $a$ be a state of $B$, and let
$m$ be either Boolean mode. Choose a host base configuration, an input-head position, a
flag tape, an output word, and a remainder word. For every prefix word $p$, start the
host in control phase $2$, with the counter and payload tapes and head positions
arranged as specified by $p$ and the tape that stores the concatenation of $p$ with the
remainder. After $|\mathrm{rest}|+1$ steps, the host is in phase $3$, its counter tape
stores $p\mathbin{+
egthinspace+}\mathrm{rest}$ with the prescribed head positions, its payload tape is
blank with head at zero, and its other specified components are unchanged.
\end{lemma}

\begin{lemma}[Fuel tape return]
\lean{Turing.FinTM.f2_loopHost_fuel_return}
\label{Turing.FinTM.f2_loopHost_fuel_return}
\leanok
\uses{Turing.Action.apply, Turing.Cfg, Turing.Cfg.workTapeSymbols, Turing.FinTM, Turing.FinTM.bufferTape, Turing.FinTM.bufferTape_left, Turing.FinTM.bufferTape_nat, Turing.FinTM.f2_LoopHostState, Turing.FinTM.f2_loopControlAction, Turing.FinTM.f2_loopControl_apply, Turing.FinTM.f2_loopFrame, Turing.FinTM.f2_loopFrame_counter, Turing.FinTM.f2_loopHost, Turing.FinTM.f2_loopWrite, Turing.MultiTapeTM.runFrom, Turing.MultiTapeTM.runFrom_succ_eq_step, Turing.MultiTapeTM.runFrom_zero, Turing.MultiTapeTM.step}
\difficulty{6}
Let $B$ and $F$ be finite Boolean Turing machines, let $a$ be a state of $B$, and let
$m$ be either Boolean mode. Choose a host base configuration, an input-head position, a
flag tape, a counter word, and an output word. For every integer $j$ with $0\le j\le
|w|$, a host configuration in phase $3$ whose counter tape stores $w$, whose payload
tape is blank, and whose counter and payload heads are both at $j-1$ reaches, after
$j+1$ steps, phase $4$ with the same counter word and output, a blank payload tape, and
both heads at zero.
\end{lemma}

\begin{lemma}[Fuel tape setup]
\lean{Turing.FinTM.f2_loopHost_fuel_setup}
\label{Turing.FinTM.f2_loopHost_fuel_setup}
\leanok
\uses{Turing.Action.apply, Turing.Cfg, Turing.FinTM, Turing.FinTM.bufferTape, Turing.FinTM.f2_LoopHostState, Turing.FinTM.f2_loopControlAction, Turing.FinTM.f2_loopControl_apply, Turing.FinTM.f2_loopCopy_initial, Turing.FinTM.f2_loopFrame, Turing.FinTM.f2_loopHost, Turing.FinTM.f2_loopHost_fuel_copy, Turing.FinTM.f2_loopHost_fuel_return, Turing.FinTM.f2_loopHost_fuel_rewind, Turing.FinTM.f2_loopWrite, Turing.MultiTapeTM.runFrom, Turing.MultiTapeTM.runFrom_add, Turing.MultiTapeTM.runFrom_succ_eq_step, Turing.MultiTapeTM.step}
\difficulty{5}
Let $B$ and $F$ be finite Boolean Turing machines, let $a$ be a state of $B$, and let
$m$ be either Boolean mode. Choose a host base configuration, an input-head position, a
flag tape, a word $w$, and an output word. Starting in phase $0$ with an empty counter
tape, payload tape storing $w$, counter head at zero, and payload head at $|w|$, the
loop host reaches phase $4$ after $3|w|+4$ steps, with the counter tape storing $w$, the
payload tape blank, and both heads at zero; the other specified components are
unchanged.
\end{lemma}

\begin{definition}[Captured loop fuel configuration]
\lean{Turing.FinTM.f2_loopFuelCaptured}
\label{Turing.FinTM.f2_loopFuelCaptured}
\leanok
\uses{Turing.Cfg, Turing.FinTM, Turing.FinTM.f2_LoopHostState, Turing.FinTM.f2_loopFuelCfg, Turing.captureCfg}
For Boolean machines $body$ and $F$, an input string $x$, and a configuration $c$ of $F$
on $x$, define a configuration of $body.k+1+(1+F.k)+1$ work tapes whose state space is
$F.State\oplus((\{\mathrm{false},\mathrm{true}\}\times(body.State\times\{\mathrm{false},\mathrm{true}\}))\oplus
Fin(14))$. It is obtained by capturing the loop-fuel configuration associated with $c$,
using the left state summand and the element $0$ of the final $Fin(14)$ summand.
\end{definition}

\begin{definition}[Ready configuration for the F2 loop]
\lean{Turing.FinTM.f2_loopReady}
\label{Turing.FinTM.f2_loopReady}
\leanok
\uses{Turing.Cfg, Turing.FinTM, Turing.FinTM.bufferTape, Turing.FinTM.f2_LoopHostState, Turing.FinTM.f2_loopFrame, Turing.FinTM.f2_loopFuelCaptured}
Let $body$ and $F$ be finite Turing machines over the Boolean alphabet, let $x$ be a
Boolean input string, and let $c$ be a configuration of $F$ on input $x$. The ready
configuration is obtained by capturing the fuel of $c$ and applying the loop-frame
construction with initial host state in the $body$ branch, body state $body.tm.q₀$,
Boolean flag true, fuel parameter $1$, the specified buffer tapes for the empty string,
the output of $c$, and the empty string, head-position parameters $0,0$, and empty
output; it has $body.k+1+(1+F.k)+1$ work tapes and state set
$F.State\oplus((Bool\times(body.State\times Bool))\oplus Fin(14))$.
\end{definition}

\begin{lemma}[Captured Halting Configuration Frame]
\lean{Turing.FinTM.f2_loopFuelCaptured_frame}
\label{Turing.FinTM.f2_loopFuelCaptured_frame}
\leanok
\uses{Turing.Cfg, Turing.FinTM, Turing.FinTM.bufferTape, Turing.FinTM.f2_loopFrame, Turing.FinTM.f2_loopFuelCaptured, Turing.FinTM.f2_loopFuelCfg, Turing.FinTM.rightCfg, Turing.captureCfg}
\difficulty{6}
Let $B$ and $F$ be finite Turing machines over the Boolean alphabet, with $k_B$ and
$k_F$ work tapes, respectively. For any Boolean input $x$ and configuration $c$ of $F$
on $x$ whose state is absent, the configuration obtained by embedding $c$ into the loop
construction for $B$ and $F$ equals the loop frame based on that embedded configuration,
with the designated return state, input-head position inherited from $c$, blank flag and
counter tapes with heads at cell $0$, a payload tape containing $c$'s output with its
head at the output length, and an empty output word.
\end{lemma}

\begin{lemma}[Preparing the loop host]
\lean{Turing.FinTM.f2_loopHost_prepare}
\label{Turing.FinTM.f2_loopHost_prepare}
\leanok
\uses{Turing.Cfg, Turing.Cfg.Halted, Turing.Cfg.init, Turing.Cfg.ofWords, Turing.FinTM, Turing.FinTM.ComputesFunInTime, Turing.FinTM.bufferTape, Turing.FinTM.f2_head_steps, Turing.FinTM.f2_loopFrame, Turing.FinTM.f2_loopFuelCaptured, Turing.FinTM.f2_loopFuelCaptured_frame, Turing.FinTM.f2_loopFuelCfg, Turing.FinTM.f2_loopFuel_init, Turing.FinTM.f2_loopFuel_run, Turing.FinTM.f2_loopHost, Turing.FinTM.f2_loopHost_fuel_capture, Turing.FinTM.f2_loopHost_fuel_setup, Turing.FinTM.f2_loopHost_init, Turing.FinTM.f2_loopHost_input_rewind, Turing.FinTM.f2_loopReady, Turing.FinTM.f2_loop_first_halt, Turing.FinTM.f2_loop_fuel_width, Turing.FinTM.f2_loop_input_run_le, Turing.FinTM.rightCfg, Turing.MultiTapeTM.initCfg, Turing.MultiTapeTM.runFrom, Turing.MultiTapeTM.runFrom_add}
\difficulty{5}
Let $B$ and $F$ be Boolean finite Turing machines, let $a$ be a state of $B$, choose a
mode bit, functions $R,T:\mathbb{N}\to\mathbb{N}$, and an input word $x$. Assume that
$F$ computes the function sending each input $y$ to the binary expansion of $R(|y|)$ in
time $T(|y|)$. Then there are a configuration $c$ of $F$ on $x$ and a time $t$ such that
$c$ is halted, its output is the binary expansion of $R(|x|)$, $t\leq 5T(|x|)+7$, and
the loop host reaches its ready configuration built from $c$ after $t$ steps; moreover,
every work-tape head of $c$ lies between $-T(|x|)$ and $T(|x|)$, inclusive.
\end{lemma}

\begin{definition}[Padded loop body configuration]
\lean{Turing.FinTM.f2_loopBodyPadded}
\label{Turing.FinTM.f2_loopBodyPadded}
\leanok
\uses{Turing.Cfg, Turing.FinTM, Turing.FinTM.bufferTape, Turing.FinTM.f2_loopBodyCfg, Turing.FinTM.leftCfg}
For Boolean finite Turing machines $body$ and $F$, a state $anchor$ of $body$, an input
word $x$, a configuration $c$ of $body$ on $x$, a Boolean $release$, an optional Boolean
$flag$, a Boolean word $word$, and a configuration $fuel$ of $F$ on $x$, the
construction that applies the loop-body configuration operation to $c$ and the given
control data, then appends the buffer encoding of $word$ and the work tapes of $fuel$,
yields a configuration on $x$ with $body.k+1+(1+F.k)$ work tapes and states in
$body.State\times\{\mathrm{false},\mathrm{true}\}$. The heads on the added tapes are at
position zero.
\end{definition}

\begin{definition}[Loop-call configuration]
\lean{Turing.FinTM.f2_loopCall}
\label{Turing.FinTM.f2_loopCall}
\leanok
\uses{Turing.Cfg, Turing.FinTM, Turing.FinTM.f2_LoopHostState, Turing.FinTM.f2_loopBodyPadded, Turing.captureCfg}
Given Boolean finite Turing machines $B$ and $F$, a state $a$ of $B$, an input word $x$,
a Boolean $s$, a configuration $c$ of $B$ on $x$, a Boolean $r$, an optional Boolean
$f$, a Boolean word $w$, and a configuration $g$ of $F$ on $x$, the construction returns
a configuration on $x$ with $B.k+F.k+3$ work tapes and state set
$F.State\oplus((\{0,1\}\times(B.State\times\{0,1\}))\oplus\{0,1,\ldots,13\})$. Its state
is encoded from $s$ and the state of the padded loop-body configuration; the additional
tag is $6$ if $s$ is true and $7$ otherwise.
\end{definition}

\begin{lemma}[Run of the extended loop-body machine]
\lean{Turing.FinTM.f2_loopBodySource_run}
\label{Turing.FinTM.f2_loopBodySource_run}
\leanok
\uses{Turing.Cfg, Turing.FinTM, Turing.FinTM.f2_loopBodySource, Turing.FinTM.f2_loopBodyTM, Turing.FinTM.leftCfg, Turing.FinTM.leftCfg_run, Turing.MultiTapeTM.runFrom}
\difficulty{1}
Let $B$ and $F$ be finite Turing machines over the Boolean alphabet, let $a$ be a state
of $B$, and let $x$ be a Boolean word. Form the loop-body machine from $B$ and $a$, and
form its source machine by adjoining $1+F.k$ work tapes. For every configuration $c$ of
the loop-body machine on input $x$, every choice of contents and head positions for
those additional tapes, and every $t\in\mathbb N$, running the source machine for $t$
steps from the configuration obtained by adjoining the chosen tapes to $c$ gives the
configuration obtained by adjoining those same tapes and head positions to the loop-body
machine's configuration after $t$ steps from $c$.
\end{lemma}

\begin{lemma}[Return from an anchor state]
\lean{Turing.FinTM.f2_loopHost_anchor_return}
\label{Turing.FinTM.f2_loopHost_anchor_return}
\leanok
\uses{Turing.Cfg, Turing.Cfg.Halted, Turing.FinTM, Turing.FinTM.f2_loopBodyCfg, Turing.FinTM.f2_loopBodyPadded, Turing.FinTM.f2_loopBodySource_run, Turing.FinTM.f2_loopBodyTM, Turing.FinTM.f2_loopBody_run, Turing.FinTM.f2_loopBody_stop, Turing.FinTM.f2_loopCall, Turing.FinTM.f2_loopHost, Turing.FinTM.f2_loopHost_body_capture, Turing.FinTM.f2_loop_live_prefix, Turing.FinTM.leftCfg, Turing.MultiTapeTM.runFrom, Turing.MultiTapeTM.runFrom_succ_eq_step'}
\difficulty{5}
Let $B$ and $F$ be finite Boolean Turing machines, let $a$ be a state of $B$, and choose
find and startup modes. For an input word $x$, configuration $c$ of $B$ on $x$, release
bit $r$, time $t$, word $w$, and configuration $f$ of $F$ on $x$, suppose the run of $B$
from $c$ is in state $a$ at time $t$, that $r$ is false if $t=0$, and that at every
earlier time $u<t$ either $u=0$ and $r$ is true or the run is not in state $a$. Then
after $t+1$ steps the host machine, started with the corresponding body-call
configuration, is exactly in the call configuration formed from the time-$t$ body
configuration with its state set to halted, release set to false, and flag set to false;
the word, fuel configuration, and other supplied data are preserved.
\end{lemma}

\begin{lemma}[Ready Configuration as Startup Call]
\lean{Turing.FinTM.f2_loopReady_call}
\label{Turing.FinTM.f2_loopReady_call}
\leanok
\uses{Turing.Cfg, Turing.Cfg.init, Turing.FinTM, Turing.FinTM.f2_loopBodyCfg, Turing.FinTM.f2_loopBodyPadded, Turing.FinTM.f2_loopCall, Turing.FinTM.f2_loopFrame, Turing.FinTM.f2_loopFuelCaptured, Turing.FinTM.f2_loopFuelCfg, Turing.FinTM.f2_loopReady, Turing.FinTM.leftCfg, Turing.FinTM.rightCfg, Turing.MultiTapeTM.initCfg, Turing.captureCfg}
\difficulty{5}
Let $B$ and $F$ be finite Boolean Turing machines, let $a$ be a state of $B$, and let
$x$ be a Boolean input word. For every configuration $c$ of $F$ on $x$, the ready
configuration determined by $c$ equals the call configuration for the initial
configuration of $B$ on $x$, with startup set to true, release set to false, no flag,
and the fuel word and fuel configuration inherited from $c$.
\end{lemma}

\begin{lemma}[Releasing the loop body]
\lean{Turing.FinTM.f2_loopHost_release}
\label{Turing.FinTM.f2_loopHost_release}
\leanok
\uses{Turing.Action.apply, Turing.Cfg, Turing.FinTM, Turing.FinTM.f2_loopBodyCfg, Turing.FinTM.f2_loopBodyPadded, Turing.FinTM.f2_loopCall, Turing.FinTM.f2_loopControlAction, Turing.FinTM.f2_loopHost, Turing.FinTM.leftCfg, Turing.MultiTapeTM.step, Turing.captureCfg, Turing.moveInputPos_zero}
\difficulty{5}
Let $B$ and $F$ be Boolean finite Turing machines, let $a$ be a state of $B$, choose a
mode bit and an input word $x$, and let $c$ be a configuration of $B$ on $x$, $w$ a
Boolean word, and $f$ a configuration of $F$ on $x$. One step of the loop host from the
call configuration with startup enabled, body configuration $c$ halted, release
disabled, and flag false produces the call configuration with startup disabled, body
state set to $a$, release enabled, and blank flag; the word, fuel configuration, and
remaining configuration data are preserved.
\end{lemma}

\begin{lemma}[F2 Loop Host Startup]
\lean{Turing.FinTM.f2_loopHost_start}
\label{Turing.FinTM.f2_loopHost_start}
\leanok
\uses{Turing.Cfg, Turing.Cfg.ofWords, Turing.FinTM, Turing.FinTM.f2_loopCall, Turing.FinTM.f2_loopHost, Turing.FinTM.f2_loopHost_anchor_return, Turing.FinTM.f2_loopHost_release, Turing.FinTM.f2_loopReady, Turing.FinTM.f2_loopReady_call, Turing.MultiTapeTM.initCfg, Turing.MultiTapeTM.runFrom, Turing.MultiTapeTM.runFrom_succ_eq_step', Turing.stateWord}
\difficulty{3}
Let $B$ and $F$ be finite Boolean Turing machines, let $a$ be a state of $B$, let $x$
and $s$ be Boolean words, and let $c$ be a configuration of $F$ on input $x$. Suppose
that the run of $B$ from its initial configuration on $x$ does not visit $a$ at any time
$u<t$, and that at time $t$ it is exactly the configuration with state $a$ and tape
contents given by placing $s$ on tape $0$ and leaving the other tapes blank. Then after
exactly $t+2$ steps, the loop host in the chosen mode, started in its ready
configuration determined by $c$, is the call configuration for the body state and
contents just described, with startup set to false, release set to true, no flag, and
both the fuel word and fuel configuration inherited from $c$.
\end{lemma}

\begin{lemma}[Return after a body halts]
\lean{Turing.FinTM.f2_loopHost_halt_return}
\label{Turing.FinTM.f2_loopHost_halt_return}
\leanok
\uses{Turing.Cfg, Turing.Cfg.Halted, Turing.FinTM, Turing.FinTM.f2_loopBodyCfg, Turing.FinTM.f2_loopBodyPadded, Turing.FinTM.f2_loopBodySource_run, Turing.FinTM.f2_loopBody_run, Turing.FinTM.f2_loopCall, Turing.FinTM.f2_loopHost, Turing.FinTM.f2_loopHost_body_capture, Turing.FinTM.leftCfg, Turing.MultiTapeTM.runFrom}
\difficulty{5}
Let $B$ and $F$ be finite Boolean Turing machines, let $a$ be a state of $B$, and let
$m$ be either Boolean mode. Let $c$ be a configuration of $B$ on input $x$, let $t$ be a
positive integer, and choose a word $w$ and a fuel-machine configuration. Suppose the
run of $B$ from $c$ does not halt at any time $u<t$, does not visit state $a$ at any
strictly intermediate positive time, and is halted at time $t$. Then the loop host,
started in the call configuration corresponding to $c$ with release flag set and no flag
bit, reaches after $t$ steps the call configuration corresponding to the run of $B$ at
time $t$, with release flag unset and the flag bit set; the word, fuel configuration,
input, and other call data are retained.
\end{lemma}

\begin{lemma}[Loop Call Frame Representation]
\lean{Turing.FinTM.f2_loopCall_frame}
\label{Turing.FinTM.f2_loopCall_frame}
\leanok
\uses{Turing.Cfg, Turing.FinTM, Turing.FinTM.bufferTape, Turing.FinTM.f2_loopBodyCfg, Turing.FinTM.f2_loopBodyPadded, Turing.FinTM.f2_loopCall, Turing.FinTM.f2_loopFrame, Turing.FinTM.leftCfg, Turing.captureCfg}
\difficulty{4}
Let $\mathrm{body}$ and $F$ be finite Turing machines over the Boolean alphabet, let
$\mathrm{anchor}$ be a state of $\mathrm{body}$, and let $x$ be a Boolean input word.
For any startup flag, body configuration $c$, release flag, optional Boolean flag value,
word, and configuration $\mathrm{fuel}$ of $F$ on $x$, the loop-call configuration
equals its frame representation: the frame uses the call's state, the input-head
position of $c$, the tape containing the optional flag at cell $0$, the buffer tapes for
$\mathrm{word}$ and $c$'s output, head positions $0$ and $|c.\mathrm{output}|$, and
empty output.
\end{lemma}

\begin{lemma}[Reframing a loop-call configuration]
\lean{Turing.FinTM.f2_loopCall_reframe}
\label{Turing.FinTM.f2_loopCall_reframe}
\leanok
\uses{Turing.Cfg, Turing.FinTM, Turing.FinTM.bufferTape, Turing.FinTM.f2_loopBodyCfg, Turing.FinTM.f2_loopBodyPadded, Turing.FinTM.f2_loopCall, Turing.FinTM.f2_loopFrame, Turing.FinTM.leftCfg, Turing.captureCfg}
\difficulty{5}
Let $B$ and $F$ be finite Boolean Turing machines, let $a$ be a state of $B$, let $x$ be
a Boolean input word, and let $c$ be a configuration of $B$ on $x$. Choose Boolean
values $s,s',r,r'$, optional Boolean symbols $u,u'$, Boolean words $w,w'$, a
configuration $f$ of $F$ on $x$, and an optional state $q$ of $B$. The configuration
obtained by reframing the loop-call configuration built from $B,F,a,s,c,r,u,w,f$, using
the state of the loop-call configuration built from $B,F,a,s',c[q],r',u',w',f$, the
input-head position of $c$, tape contents $u'$ only at cell $0$, buffer tapes storing
$w'$ and the output of $c$, head positions $0$ and the length of $c$'s output, and empty
output, is exactly the loop-call configuration built from $B,F,a,s',c[q],r',u',w',f$,
where $c[q]$ denotes $c$ with its state replaced by $q$.
\end{lemma}

\begin{lemma}[Single step of the borrow operation]
\lean{Turing.FinTM.f2_loopHost_borrow_step}
\label{Turing.FinTM.f2_loopHost_borrow_step}
\leanok
\uses{Turing.Action.apply, Turing.Cfg, Turing.Cfg.workTapeSymbols, Turing.FinTM, Turing.FinTM.bufferTape, Turing.FinTM.f2_LoopHostState, Turing.FinTM.f2_loopBuffer_read, Turing.FinTM.f2_loopBuffer_write, Turing.FinTM.f2_loopControlAction, Turing.FinTM.f2_loopControl_apply, Turing.FinTM.f2_loopFrame, Turing.FinTM.f2_loopFrame_counter, Turing.FinTM.f2_loopHost, Turing.FinTM.f2_loopWrite, Turing.MultiTapeTM.step}
\difficulty{4}
Let $B$ and $F$ be finite Boolean Turing machines, let $a$ be a state of $B$, choose
find mode, and let $x$ be an input word. For any base host configuration, input-head
position, and words $v$ and $u$, one step from control phase 8 with the counter tape
storing $vu$, its head at $|v|$, and the other designated tapes blank has the following
effect: if $u$ is empty, the machine enters phase 10, the counter tape stores $v$, and
its head moves to $|v|-1$; if $u$ begins with true, it enters phase 9, changes that
first bit to false, and moves the head to $|v|-1$; if $u$ begins with false, it stays in
phase 8, changes that first bit to true, and moves the head to $|v|+1$. In each case the
input position and output are preserved.
\end{lemma}

\begin{lemma}[Run of the borrow operation]
\lean{Turing.FinTM.f2_loopHost_borrow_run}
\label{Turing.FinTM.f2_loopHost_borrow_run}
\leanok
\uses{Turing.Cfg, Turing.FinTM, Turing.FinTM.bufferTape, Turing.FinTM.f2_LoopHostState, Turing.FinTM.f2_loopBorrowPos, Turing.FinTM.f2_loopDebit, Turing.FinTM.f2_loopFrame, Turing.FinTM.f2_loopHost, Turing.FinTM.f2_loopHost_borrow_step, Turing.MultiTapeTM.runFrom, Turing.MultiTapeTM.runFrom_succ_eq_step}
\difficulty{4}
Let $B$ and $F$ be finite Boolean Turing machines, let $a$ be a state of $B$, choose
find mode, and let $x$ be an input word. For any base host configuration, input-head
position, and Boolean word $w$, and for every prefix word $v$, start in control phase 8
with the counter tape storing the concatenation $vw$, its head at position $|v|$, and
the other designated tapes blank. After $b(w)+1$ steps, the counter tape stores $v$
followed by the debited word obtained from $w$; its head is at $|v|+b(w)-1$, and the
control phase is 9 if the debit reports a carry and 10 otherwise. The remaining
designated frame data are preserved as specified.
\end{lemma}

\begin{lemma}[Rewind after borrowing]
\lean{Turing.FinTM.f2_loopHost_borrow_rewind}
\label{Turing.FinTM.f2_loopHost_borrow_rewind}
\leanok
\uses{Turing.Action.apply, Turing.Cfg, Turing.Cfg.workTapeSymbols, Turing.FinTM, Turing.FinTM.bufferTape, Turing.FinTM.bufferTape_left, Turing.FinTM.bufferTape_nat, Turing.FinTM.f2_LoopHostState, Turing.FinTM.f2_loopControlAction, Turing.FinTM.f2_loopControl_apply, Turing.FinTM.f2_loopFrame, Turing.FinTM.f2_loopFrame_counter, Turing.FinTM.f2_loopHost, Turing.FinTM.f2_loopWrite, Turing.MultiTapeTM.runFrom, Turing.MultiTapeTM.runFrom_succ_eq_step, Turing.MultiTapeTM.runFrom_zero, Turing.MultiTapeTM.step}
\difficulty{5}
Let $B$ and $F$ be finite Boolean Turing machines, let $a$ be a state of $B$, choose
find mode, and let $x$ be an input word. For any base host configuration, input-head
position, Boolean word $w$, and success bit $s$, for every natural number $j\le |w|$,
start in control phase 9 if $s$ is true and phase 10 otherwise, with the counter tape
storing $w$ and its head at $j-1$, and with the other designated tapes blank. After
$j+1$ steps the frame has the same counter contents, the counter head is reset to zero,
the other designated heads, input position, and output are preserved, and the control
state is the return call state with anchor and release flag true when $s$ is true, or
control phase 11 otherwise.
\end{lemma}

\begin{lemma}[Borrow operation and running-time bound]
\lean{Turing.FinTM.f2_loopHost_borrow}
\label{Turing.FinTM.f2_loopHost_borrow}
\leanok
\uses{Turing.Cfg, Turing.FinTM, Turing.FinTM.bufferTape, Turing.FinTM.f2_LoopHostState, Turing.FinTM.f2_loopBorrowPos, Turing.FinTM.f2_loopBorrowPos_le, Turing.FinTM.f2_loopDebit, Turing.FinTM.f2_loopDebit_length, Turing.FinTM.f2_loopFrame, Turing.FinTM.f2_loopHost, Turing.FinTM.f2_loopHost_borrow_rewind, Turing.FinTM.f2_loopHost_borrow_run, Turing.MultiTapeTM.runFrom, Turing.MultiTapeTM.runFrom_add}
\difficulty{4}
Let $B$ and $F$ be finite Boolean Turing machines, let $a$ be a state of $B$, choose
find mode, and let $x$ be an input word. For any host configuration used as a base,
input-head position, and Boolean word $w$, the borrow position $b(w)$ satisfies
$2b(w)+2\le 2|w|+2$. Starting from the specified host frame in control phase 8, with the
counter tape storing $w$ and the other designated tapes blank, after $2b(w)+2$ steps the
frame is in the return call state with the anchor and release flag true if the debit
operation on $w$ reports a carry; otherwise it is in control phase 11. The counter tape
then stores the debited word, and the designated head positions and output are reset as
specified by the frame.
\end{lemma}

\begin{lemma}[Flag tape in a loop frame]
\lean{Turing.FinTM.f2_loopFrame_flag}
\label{Turing.FinTM.f2_loopFrame_flag}
\leanok
\uses{Turing.Cfg, Turing.Cfg.workTapeSymbols, Turing.FinTM, Turing.FinTM.f2_LoopHostState, Turing.FinTM.f2_loopFrame}
\difficulty{1}
Let $B$ and $F$ be finite Turing machines over the Boolean alphabet, and let $x$ be a
finite Boolean input. In a configuration with $B.k+1+(1+F.k)+1$ work tapes, form a new
configuration by placing the flag tape on tape $B.k$ with head at position $0$, the
counter tape on tape $B.k+1$ with head at position $ch$, and the payload tape on tape
$B.k+1+(1+F.k)$ with head at position $ph$, while retaining the base configuration's
other tapes; use the specified state, input-head position, and output word. The symbol
scanned by tape $B.k$ in this new configuration is the flag tape's content at position
$0$.
\end{lemma}

\begin{lemma}[Clearing the loop flag]
\lean{Turing.FinTM.f2_loopFlag_clear}
\label{Turing.FinTM.f2_loopFlag_clear}
\leanok
\uses{Turing.FinTM.bufferTape, Turing.FinTM.f2_loopWrite}
\difficulty{4}
For every optional Boolean flag, start with a tape whose cell at the origin contains
that flag and whose other cells are blank. Writing a blank at the origin makes the tape
blank everywhere, equal to the tape storing the empty word.
\end{lemma}

\begin{lemma}[Rejecting or debiting a word]
\lean{Turing.FinTM.f2_loopHost_reject}
\label{Turing.FinTM.f2_loopHost_reject}
\leanok
\uses{Turing.Action.apply, Turing.Cfg, Turing.Cfg.workTapeSymbols, Turing.FinTM, Turing.FinTM.bufferTape, Turing.FinTM.f2_loopBodyCfg, Turing.FinTM.f2_loopBodyPadded, Turing.FinTM.f2_loopBorrowPos, Turing.FinTM.f2_loopBorrowPos_le, Turing.FinTM.f2_loopCall, Turing.FinTM.f2_loopCall_frame, Turing.FinTM.f2_loopCall_reframe, Turing.FinTM.f2_loopControlAction, Turing.FinTM.f2_loopControl_apply, Turing.FinTM.f2_loopDebit, Turing.FinTM.f2_loopFlag_clear, Turing.FinTM.f2_loopFrame, Turing.FinTM.f2_loopFrame_flag, Turing.FinTM.f2_loopHost, Turing.FinTM.f2_loopHost_borrow, Turing.FinTM.f2_loopWrite, Turing.FinTM.leftCfg, Turing.MultiTapeTM.runFrom, Turing.MultiTapeTM.runFrom_succ_eq_step, Turing.MultiTapeTM.runFrom_succ_eq_step', Turing.MultiTapeTM.step, Turing.captureCfg}
\difficulty{5}
Let $B$ and $F$ be Boolean finite Turing machines, let $a$ be a state of $B$, choose a
mode bit, an input word $x$, a configuration $c$ of $B$ on $x$, a word $w$, and a
configuration of $F$ on $x$. Assume $c$ is halted and has empty output. For some $t\leq
2|w|+4$, if the debit operation on $w$ succeeds, the loop host run from the call
configuration with $c$, flag false, and payload $w$ reaches the call configuration with
$c$ restored to state $a$, release enabled, flag blank, and the debited word as payload,
with the fuel configuration unchanged. If the debit operation does not succeed, the host
instead halts and its output is empty in find mode and the one-bit word
$[\mathrm{false}]$ otherwise.
\end{lemma}

\begin{lemma}[Rewinding the payload tape]
\lean{Turing.FinTM.f2_loopHost_payload_rewind}
\label{Turing.FinTM.f2_loopHost_payload_rewind}
\leanok
\uses{Turing.Action.apply, Turing.Cfg, Turing.Cfg.workTapeSymbols, Turing.FinTM, Turing.FinTM.bufferTape, Turing.FinTM.bufferTape_left, Turing.FinTM.bufferTape_nat, Turing.FinTM.f2_LoopHostState, Turing.FinTM.f2_loopControlAction, Turing.FinTM.f2_loopControl_apply, Turing.FinTM.f2_loopFrame, Turing.FinTM.f2_loopFrame_payload, Turing.FinTM.f2_loopHost, Turing.FinTM.f2_loopWrite, Turing.MultiTapeTM.runFrom, Turing.MultiTapeTM.runFrom_succ_eq_step, Turing.MultiTapeTM.runFrom_zero, Turing.MultiTapeTM.step}
\difficulty{4}
Let $B$ and $F$ be Boolean finite Turing machines, let $a$ be a state of $B$, choose a
mode bit and an input word $x$, and fix a base configuration, an input-head position,
flag and counter tape contents, counter-head position, payload word $w$, and output
word. For every $j$ with $j\leq |w|$, starting in control phase $12$ with the payload
tape storing $w$ and its head at position $j-1$, after $j+1$ steps the host is in phase
$13$ with the same payload and other specified data, and the payload head is at position
$0$.
\end{lemma}

\begin{lemma}[Payload replay reaches halt]
\lean{Turing.FinTM.f2_loopHost_frame_replay}
\label{Turing.FinTM.f2_loopHost_frame_replay}
\leanok
\uses{Turing.Cfg, Turing.FinTM, Turing.FinTM.bufferTape, Turing.FinTM.f2_LoopHostState, Turing.FinTM.f2_loopFrame, Turing.FinTM.f2_loopHost, Turing.FinTM.f2_loopHost_replay, Turing.FinTM.f2_loopReplayCfg, Turing.FinTM.rightCfg, Turing.MultiTapeTM.runFrom}
\difficulty{5}
Let $B$ and $F$ be finite Boolean Turing machines, let $a$ be a state of $B$, and choose
a host configuration, an input-head position, flag and counter tapes, an integer
counter-head position, and a Boolean word $w$. Form a host configuration in control
phase $13$ whose payload tape stores $w$ from cell zero, whose payload head is at zero,
whose output is empty, and whose other components are those specified by the chosen host
configuration, position, flag, and counter data. After $|w|+1$ steps of the loop host,
the resulting configuration is halted and its output is exactly $w$.
\end{lemma}

\begin{lemma}[Halting bound for the loop host]
\lean{Turing.FinTM.f2_loopHost_accept}
\label{Turing.FinTM.f2_loopHost_accept}
\leanok
\uses{Turing.Action.apply, Turing.Cfg, Turing.Cfg.workTapeSymbols, Turing.FinTM, Turing.FinTM.bufferTape, Turing.FinTM.f2_loopBodyCfg, Turing.FinTM.f2_loopBodyPadded, Turing.FinTM.f2_loopCall, Turing.FinTM.f2_loopCall_frame, Turing.FinTM.f2_loopControlAction, Turing.FinTM.f2_loopControl_apply, Turing.FinTM.f2_loopFrame, Turing.FinTM.f2_loopFrame_flag, Turing.FinTM.f2_loopHost, Turing.FinTM.f2_loopHost_frame_replay, Turing.FinTM.f2_loopHost_payload_rewind, Turing.FinTM.f2_loopWrite, Turing.FinTM.leftCfg, Turing.MultiTapeTM.runFrom, Turing.MultiTapeTM.runFrom_add, Turing.MultiTapeTM.runFrom_succ_eq_step, Turing.MultiTapeTM.step, Turing.captureCfg}
\difficulty{6}
Let $B$ and $F$ be finite Turing machines over the Boolean alphabet, let $a$ be a state
of $B$, and choose a Boolean mode. For any input word $x$, any configuration $c$ of $B$
on $x$ with no current state, any Boolean word $w$, and any configuration $f$ of $F$ on
$x$, there is a time $t \le 2|u|+3$, where $u$ is the output word of $c$, such that the
loop host built from $B$, $F$, $a$, and the chosen mode, started from the configuration
formed using $c$, the empty prefix and initial output, the flag set to true, $w$, and
$f$, has halted and has output $u$ in find mode and the one-letter word
$[\mathrm{true}]$ otherwise.
\end{lemma}

\begin{lemma}[F2 Loop Round Behavior]
\lean{Turing.FinTM.f2_loopHost_round}
\label{Turing.FinTM.f2_loopHost_round}
\leanok
\uses{Turing.Cfg, Turing.Cfg.ofWords, Turing.FinTM, Turing.FinTM.f2_loopCall, Turing.FinTM.f2_loopDebit, Turing.FinTM.f2_loopHost, Turing.FinTM.f2_loopHost_accept, Turing.FinTM.f2_loopHost_anchor_return, Turing.FinTM.f2_loopHost_halt_return, Turing.FinTM.f2_loopHost_reject, Turing.FinTM.f2_loop_first_halt, Turing.FinTM.f2_loop_output_length_le, Turing.MultiTapeTM.runFrom, Turing.MultiTapeTM.runFrom_add, Turing.stateWord}
\difficulty{5}
Let $B$ and $F$ be finite Boolean Turing machines, let $a$ be a state of $B$, and choose
a mode, an input word, words $s$, $s'$, $p$, and $w$, a Boolean acceptance value, a
configuration of $F$, and a positive integer $t$. Suppose that, starting from the
configuration with state $a$ and tape contents given by placing $s$ on tape $0$ and
leaving the other tapes blank, the run of $B$ never visits state $a$ at any time
strictly between $0$ and $t$. Suppose also that at time $t$ the run has halted with
output $p$ if it is accepted, and otherwise is exactly the configuration with state $a$
and tape contents given by placing $s'$ on tape $0$ and leaving the other tapes blank.
Then the loop host, started from the corresponding call configuration with the given
word $w$ and fuel configuration, reaches within $3t+2|w|+5$ steps either a halted
configuration whose output is $p$ in find mode and $[\mathrm{true}]$ otherwise (in the
accepted case); the next call configuration with word equal to the debit of $w$ (if the
run was not accepted and the debit reports success); or a halted configuration with
output $[]$ in find mode and $[\mathrm{false}]$ otherwise (if the run was not accepted
and the debit reports failure).
\end{lemma}

\begin{lemma}[Loop Call Head Position Bound]
\lean{Turing.FinTM.f2_loopCall_heads}
\label{Turing.FinTM.f2_loopCall_heads}
\leanok
\uses{Turing.Cfg, Turing.Cfg.ofWords, Turing.FinTM, Turing.FinTM.f2_loopBodyCfg, Turing.FinTM.f2_loopBodyPadded, Turing.FinTM.f2_loopCall, Turing.FinTM.leftCfg, Turing.captureCfg, Turing.stateWord}
\difficulty{4}
Let $\mathrm{body}$ and $F$ be finite Boolean Turing machines, let $\mathrm{anchor}$ be
a state of $\mathrm{body}$, and let $x$, $s$, and $\mathrm{word}$ be Boolean words. For
a configuration $\mathrm{fuel}$ of $F$ on $x$ and a natural number $B$, suppose every
work-tape head of $\mathrm{fuel}$ lies in the integer interval $[-B,B]$. Then every
work-tape head of the loop-call configuration built with startup flag false, the seam
configuration at state $\mathrm{anchor}$ from the tape words assigning $s$ to tape zero
and empty words to the other body tapes, release flag true, no flag symbol,
$\mathrm{word}$, and $\mathrm{fuel}$ also lies in $[-B,B]$.
\end{lemma}

\begin{definition}[Loop host bound]
\lean{Turing.FinTM.f2_loopHost_bound}
\label{Turing.FinTM.f2_loopHost_bound}
\leanok
The loop host bound is the natural number $10$.
\end{definition}

\begin{lemma}[Loop host contracts]
\lean{Turing.FinTM.f2_loopHost_contracts}
\label{Turing.FinTM.f2_loopHost_contracts}
\leanok
\uses{Turing.Cfg, Turing.Cfg.ofWords, Turing.FinTM, Turing.FinTM.ComputesFunInTime, Turing.FinTM.f2_head_steps, Turing.FinTM.f2_loopCall, Turing.FinTM.f2_loopCall_heads, Turing.FinTM.f2_loopDebit, Turing.FinTM.f2_loopDebit_iterate_length, Turing.FinTM.f2_loopDebit_iterate_value, Turing.FinTM.f2_loopDebit_success, Turing.FinTM.f2_loopHost, Turing.FinTM.f2_loopHost_bound, Turing.FinTM.f2_loopHost_prepare, Turing.FinTM.f2_loopHost_round, Turing.FinTM.f2_loopHost_start, Turing.FinTM.f2_loop_fuel_width, Turing.FinTM.f2_loop_orbit_inv, Turing.MultiTapeTM.initCfg, Turing.MultiTapeTM.runFrom, Turing.MultiTapeTM.runFrom_add, Turing.stateWord}
\difficulty{7}
Let $B$ and $F$ be finite Boolean Turing machines, let $a$ be a state of $B$, and let
$I$, $S$, $A$, and $O$ be predicates or functions on Boolean words, with $S(x)$ the
initial auxiliary word. Suppose $F$ computes the binary representation of $R(|x|)$
within $T(|x|)$ steps on every input $x$; $I(x,S(x))$ holds, and $I(x,s)$ implies
$I(x,S_1(x,s))$ for every $x,s$, where $S_1$ is the state update. Suppose also that each
input has a run of $B$ from its initial configuration to the seam configuration in state
$a$ with tape words given by $S(x)$, taking at most $T(|x|)$ steps and never visiting
state $a$ earlier. For every $x,s$ with $I(x,s)$, suppose there is a positive time at
most $T(|x|)$ at which a run from the seam configuration in state $a$ either halts with
output $O(x,s)$ when $A(x,s)$ holds, or returns to the seam configuration with updated
word $S_1(x,s)$ when $A(x,s)$ fails; in either case it does not visit state $a$ at any
strictly intermediate positive time. Then there is a constant $c$ such that, for every
input $x$, the loop host has a computation sequence $(C_i)_{i\in\mathbb N}$ and a
startup time at most $c(T(|x|)+1)$ with the following properties: the host run from its
initial configuration reaches $C_0$ at that startup time, and $C_i$ has empty output for
every $i\le R(|x|)$. The configuration $C_{R(|x|)+1}$ is halted and its output is empty
in find mode and the one-bit word $[\mathrm{false}]$ otherwise. For each $i\le R(|x|)$
there is a time at most $c(T(|x|)+1)$ such that, if $A(x,S_1^{[i]}(x))$ holds, the host
run from $C_i$ halts with output $O(x,S_1^{[i]}(x))$ in find mode and $[\mathrm{true}]$
otherwise; if it fails, that run reaches exactly $C_{i+1}$. At each $C_i$ for $i\le
R(|x|)$, every work-tape head lies between $-T(|x|)$ and $T(|x|)$; at the final
configuration, every work-tape head lies between $-(T(|x|)+c(T(|x|)+1))$ and
$T(|x|)+c(T(|x|)+1)$. Here $S_1^{[i]}$ denotes $i$-fold iteration of the state update.
\end{lemma}

\begin{lemma}[Bounded Loop Search]
\lean{Turing.FinTM.f2_loop_find_run}
\label{Turing.FinTM.f2_loop_find_run}
\leanok
\uses{Turing.Cfg, Turing.MultiTapeTM, Turing.MultiTapeTM.runFrom, Turing.MultiTapeTM.runFrom_add}
\difficulty{4}
Let $M$ be a multi-tape Turing machine with $k$ work tapes over Boolean symbols, let
$(c_j)_{j\in\mathbb N}$ be a sequence of configurations on a fixed Boolean input, let
$A(j)$ be a Boolean-valued acceptance test, and let $P(j)$ be a Boolean word for each
$j$. For natural numbers $B,N$, suppose the configuration $c_N$ is halted with empty
output, and for every $j<N$ there is a time at most $B$ such that either the run from
$c_j$ halts with output $P(j)$ when $A(j)$ holds, or reaches exactly $c_{j+1}$ when
$A(j)$ fails. Then the run from $c_0$ halts within $NB$ steps, with output $P(i)$ for
the first index $i<N$ satisfying $A(i)$, if such an index exists, and with empty output
if none exists.
\end{lemma}

\begin{lemma}[Head positions along bounded computation segments]
\lean{Turing.FinTM.f2_segment_heads}
\label{Turing.FinTM.f2_segment_heads}
\leanok
\uses{Turing.Cfg, Turing.FinTM.f2_head_steps, Turing.MultiTapeTM, Turing.MultiTapeTM.runFrom, Turing.MultiTapeTM.runFrom_add, Turing.MultiTapeTM.runFrom_of_halt}
\difficulty{4}
Let $M$ be a multi-tape Turing machine with $k$ work tapes over Boolean symbols, let $x$
be a Boolean input word, and let $(c_j)_{j\\in\\mathbb N}$ be configurations of $M$ on
$x$. Suppose $N,B,H$ are natural numbers, every tape head in $c_j$ lies in $[-H,H]$ for
each $j<N$, $c_N$ is halted, every tape head in $c_N$ lies in $[-(H+B),H+B]$, and for
each $j<N$ there is a $u\\le B$ such that the run from $c_j$ is halted after $u$ steps
or reaches $c_{j+1}$ after $u$ steps. Then, at every time $t$ and on every work tape,
the head position in the run from $c_0$ lies in $[-(H+B),H+B]$.
\end{lemma}

\begin{lemma}[Space bound from a head-position radius]
\lean{Turing.FinTM.f2_space_radius}
\label{Turing.FinTM.f2_space_radius}
\leanok
\uses{Turing.FinTM, Turing.MultiTapeTM.initCfg, Turing.MultiTapeTM.runFrom, Turing.MultiTapeTM.spaceUsed, Turing.MultiTapeTM.spaceUsedByTape, Turing.MultiTapeTM.visitedByTapeHead}
\difficulty{4}
Let $M$ be a finite Turing machine over Boolean symbols, let $x$ be a Boolean input
word, and let $B$ be a natural number. Suppose that at every time and on every work
tape, the head position in the run of $M$ from its initial configuration on $x$ lies
between $-B$ and $B$. Then, for every $t\\in\\mathbb N$, the total space used through
step $t$ is at most $k(2B+1)$, where $k$ is the number of work tapes of $M$.
\end{lemma}

\begin{lemma}[Space bound from bounded computation segments]
\lean{Turing.FinTM.f2_seamed_space}
\label{Turing.FinTM.f2_seamed_space}
\leanok
\uses{Turing.Cfg, Turing.FinTM, Turing.FinTM.f2_head_steps, Turing.FinTM.f2_segment_heads, Turing.FinTM.f2_space_radius, Turing.MultiTapeTM.initCfg, Turing.MultiTapeTM.runFrom, Turing.MultiTapeTM.runFrom_add, Turing.MultiTapeTM.spaceUsed}
\difficulty{5}
Let $M$ be a finite Turing machine over $\\{0,1\\}$ with $k$ work tapes, let $x$ be a
Boolean input word, and let $(c_j)_{j\\in\\mathbb N}$ be configurations of $M$ on $x$.
Suppose $N,B,H,$ and $r$ are natural numbers with $r\\le B$, the computation from the
initial configuration reaches $c_0$ after $r$ steps, and for every $j<N$ and every tape
the head position in $c_j$ lies between $-H$ and $H$. Suppose also that $c_N$ is halted,
its tape-head positions lie between $-(H+B)$ and $H+B$, and for each $j<N$ there is a
$u\\le B$ such that the run from $c_j$ is halted after $u$ steps or reaches $c_{j+1}$
after $u$ steps. Then, for every $t\\in\\mathbb N$, the total space used from the
initial configuration through step $t$ is at most $k(2(H+B)+1)$.
\end{lemma}

\begin{lemma}[Space Bound for a Bounded Loop Search]
\lean{Turing.FinTM.f2_exists_loopFind_space}
\label{Turing.FinTM.f2_exists_loopFind_space}
\leanok
\uses{Turing.Cfg.ofWords, Turing.FinTM, Turing.FinTM.ComputesFunInTime, Turing.FinTM.ComputesInTime, Turing.FinTM.ComputesInTime.mono, Turing.FinTM.f2_loopHost, Turing.FinTM.f2_loopHost_contracts, Turing.FinTM.f2_loop_find_run, Turing.FinTM.f2_seamed_space, Turing.MultiTapeTM.initCfg, Turing.MultiTapeTM.runFrom, Turing.MultiTapeTM.runFrom_add, Turing.MultiTapeTM.spaceUsed, Turing.stateWord}
\difficulty{6}
Let $B$ and $F$ be finite Turing machines over the Boolean alphabet, let $a$ be a state
of $B$, and let $I$ be a predicate on pairs of Boolean words. Let $g$ be a state-update
function, $A$ an acceptance test, $o$ an output function, $s_0$ an initial-state
function, and $R,T:\mathbb{N}\to\mathbb{N}$. Assume that $F$ computes the function
sending an input $x$ to the binary expansion of $R(|x|)$ within time $T(|x|)$; that
$I(x,s_0(x))$ holds and is preserved by $g$; and that, from the initial configuration of
$B$ on $x$, some time $t\le T(|x|)$ reaches the seam configuration with state $a$ and
tape contents given by $s_0(x)$, with no earlier time having state $a$. Assume also that
for every $x,s$ with $I(x,s)$, there is a time $t$ with $0<t\le T(|x|)$ such that no
time strictly between $0$ and $t$ has state $a$ when started from the seam configuration
for $(x,s)$; at time $t$, if $A(x,s)$ holds, the machine halts with output $o(x,s)$, and
otherwise it reaches the seam configuration for $(x,g(x,s))$. Then there is a finite
Turing machine $E$ and a natural number $c$ such that, on every input $x$, $E$ computes
the result of searching the iterates $g^i(x,s_0(x))$ for the least $i\le R(|x|)$
satisfying $A(x,g^i(x,s_0(x)))$, returning $o(x,g^i(x,s_0(x)))$ if one is found and the
empty word otherwise, within time $c(T(|x|)+1)(R(|x|)+2)$. Moreover, for every input $x$
and every time $t$, the space used by $E$ by time $t$ is at most $c(T(|x|)+1)$.
\end{lemma}

\begin{definition}[Conditional split step]
\lean{Turing.FinTM.f2_splitStep}
\label{Turing.FinTM.f2_splitStep}
\leanok
For Boolean lists $w$ and $s$, the split step returns $s$ with a final `true` appended
when the length of $s$ is at most the length of $w$; otherwise it returns $s$ unchanged.
\end{definition}

\begin{definition}[Split acceptance condition]
\lean{Turing.FinTM.f2_splitAccept}
\label{Turing.FinTM.f2_splitAccept}
\leanok
For natural numbers $C,e$ and Boolean lists $w,s$, the split acceptance predicate is
true exactly when $|s|+C(|s|+1)^e=|w|$.
\end{definition}

\begin{lemma}[Length bound for a split step]
\lean{Turing.FinTM.f2_splitStep_inv}
\label{Turing.FinTM.f2_splitStep_inv}
\leanok
\uses{Turing.FinTM.f2_splitStep}
\difficulty{3}
Let $w$ and $s$ be finite Boolean sequences, and suppose that the length of $s$ is at
most the length of $w$ plus one. Form a new sequence by appending $	exttt{true}$ to $s$
if the length of $s$ is at most the length of $w$, and otherwise leaving $s$ unchanged.
The resulting sequence has length at most the length of $w$ plus one.
\end{lemma}

\begin{lemma}[Orbit of the split step]
\lean{Turing.FinTM.f2_splitStep_orbit}
\label{Turing.FinTM.f2_splitStep_orbit}
\leanok
\uses{Turing.FinTM.f2_splitStep}
\difficulty{4}
For every Boolean word $w$, every natural number $i$ with $i\le |w|+1$ satisfies:
iterating the split step determined by $w$ exactly $i$ times from the empty word yields
the word consisting of $i$ copies of $\mathrm{true}$.
\end{lemma}

\begin{lemma}[Congruence of list search]
\lean{Turing.FinTM.f2_catalogFind_congr}
\label{Turing.FinTM.f2_catalogFind_congr}
\leanok
\difficulty{4}
Let $\alpha$ be a type, let $xs$ be a finite list of elements of $\alpha$, and let $p,q
: \alpha \to \{\mathrm{false},\mathrm{true}\}$. If $p(a)=q(a)$ for every $a$ in $xs$,
then searching $xs$ for the first element satisfying $p$ gives the same result as
searching for the first element satisfying $q$.
\end{lemma}

\begin{lemma}[Iterated split search agrees with polynomial search]
\lean{Turing.FinTM.f2_splitFind_eq}
\label{Turing.FinTM.f2_splitFind_eq}
\leanok
\uses{Turing.FinTM.f2_catalogFind_congr, Turing.FinTM.f2_splitAccept, Turing.FinTM.f2_splitStep, Turing.FinTM.f2_splitStep_orbit, Turing.solveSplit}
\difficulty{4}
For natural numbers $C,e$ and a Boolean word $w$, searching the indices $0,\ldots,|w|$
for the first one at which the acceptance test holds after iterating the split step from
the empty word gives the same optional index as solving $i+C(i+1)^e=|w|$ by searching
the same range.
\end{lemma}

\begin{lemma}[No split solution and failed acceptance tests]
\lean{Turing.FinTM.f2_splitFind_none}
\label{Turing.FinTM.f2_splitFind_none}
\leanok
\uses{Turing.FinTM.f2_splitAccept, Turing.FinTM.f2_splitFind_eq, Turing.FinTM.f2_splitStep, Turing.solveSplit}
\difficulty{3}
For natural numbers $C,e$ and a Boolean word $w$, there is no index in $0,\ldots,|w|$
satisfying $i+C(i+1)^e=|w|$ if and only if, for every $i\le |w|$, the acceptance test
after iterating the split step from the empty word $i$ times is false.
\end{lemma}

\begin{lemma}[Agreement of split procedures]
\lean{Turing.FinTM.f2_splitLoop_result}
\label{Turing.FinTM.f2_splitLoop_result}
\leanok
\uses{Turing.FinTM.f2_splitAccept, Turing.FinTM.f2_splitFind_eq, Turing.FinTM.f2_splitStep, Turing.FinTM.f2_splitStep_orbit, Turing.pairEncode, Turing.solveSplit}
\difficulty{4}
For natural numbers $C,e$ and a Boolean word $w$, the following two results agree:
search the integers from $0$ through $|w|$ for the first index $i$ satisfying
$i+C(i+1)^e=|w|$, and, if found, encode the prefix of $w$ of length equal to the loop
counter after $i$ iterations of the rule that appends a true bit whenever the counter's
length is at most $|w|$; otherwise return the empty word. This equals the encoding of
the prefix of length $i$ and the corresponding suffix when the equation has a solution
in that range, and is empty when it does not.
\end{lemma}

\begin{lemma}[Polynomial loop bound]
\lean{Turing.FinTM.f2_splitLoop_bound}
\label{Turing.FinTM.f2_splitLoop_bound}
\leanok
\difficulty{4}
For all natural numbers $c,A,e,n$, $c\bigl(A(n+1)^{e+1}+1\bigr)(n+2)\le
2c(A+1)(n+1)^{e+2}$.
\end{lemma}

\begin{definition}[Split input position]
\lean{Turing.FinTM.f2_splitPos}
\label{Turing.FinTM.f2_splitPos}
\leanok
For a Boolean list $w$ and a natural number $j$, the split input position is the
position indexed by $\min(j,|w|)+1$ among the $|w|+2$ positions.
\end{definition}

\begin{lemma}[Input symbol at a split position]
\lean{Turing.FinTM.f2_splitPos_read}
\label{Turing.FinTM.f2_splitPos_read}
\leanok
\uses{Turing.Cfg, Turing.Cfg.inputSymbol, Turing.FinTM.f2_splitPos, Turing.inputSymbolInner}
\difficulty{4}
Let $k$ be a natural number, let $S$ be a type, let $w$ be a Boolean word, and let
$\mathrm{cfg}$ be a configuration of a $k$-work-tape machine on input $w$. For any
natural number $j$, if the input head of $\mathrm{cfg}$ is at position $1+\min(j,|w|)$,
then the symbol read is the $j$-th bit of $w$ when $j<|w|$, and no symbol is read
otherwise.
\end{lemma}

\begin{lemma}[Successor split position]
\lean{Turing.FinTM.f2_splitPos_succ}
\label{Turing.FinTM.f2_splitPos_succ}
\leanok
\uses{Turing.FinTM.f2_splitPos, Turing.moveInputPos, Turing.moveInputPos_pos_of_ne_right, Turing.moveInputPos_rightBoundary}
\difficulty{5}
For every Boolean word $w$ and natural number $j$, moving the input head one position to
the right from position $1+\min(j,|w|)$ gives position $1+\min(j+1,|w|)$, with positions
clamped at the right boundary.
\end{lemma}

\begin{definition}[Split scratch tape contents]
\lean{Turing.FinTM.f2_splitScratch}
\label{Turing.FinTM.f2_splitScratch}
\leanok
For natural numbers $q,j$ and an integer position $z$, the split scratch tape contains
`true` at $z$ exactly when $j\le z<q$, interpreting $j$ and $q$ as integers; it is blank
otherwise.
\end{definition}

\begin{lemma}[Erasing a Scratch Tape Cell]
\lean{Turing.FinTM.f2_splitScratch_erase}
\label{Turing.FinTM.f2_splitScratch_erase}
\leanok
\uses{Turing.FinTM.f2_splitScratch}
\difficulty{4}
For natural numbers $q$ and $j$, define a Boolean-valued partial tape that contains true
exactly at integer positions $z$ satisfying $j\le z<q$. Updating this tape at position
$j$ to blank gives the tape whose true cells are exactly the positions $z$ satisfying
$j+1\le z<q$.
\end{lemma}

\begin{definition}[Split restoration machine]
\lean{Turing.FinTM.f2_splitRestoreTM}
\label{Turing.FinTM.f2_splitRestoreTM}
\leanok
\uses{Turing.FinTM, Turing.FinTM.controlAction}
For each natural number $k$, define a Boolean-alphabet Turing machine with $k+1$ work
tapes, state set $	ext{Fin}(5)\times\{\mathrm{false},\mathrm{true}\}$, and initial state
$(0,\mathrm{false})$. Its transition function is specified as follows. In a state whose
first coordinate is $0$, if work tape $0$ scans a symbol, the input head moves right,
work tape $0$ moves right, and each other work tape moves right while its scanned cell
is cleared; the successor state has first coordinate $0$ and Boolean coordinate equal to
the disjunction of the current Boolean coordinate and the condition that the input
symbol is blank. If work tape $0$ scans blank in that state, the input head stays put,
all work heads move left, and tape $0$ is written with true unless the Boolean
coordinate is already true; the successor state is $(1,\mathrm{false})$. In state $1$,
if tape $0$ scans a symbol, all work heads move left without writing and the machine
remains in state $(1,\mathrm{false})$; if it scans blank, all work heads move right
without writing and the machine enters $(2,\mathrm{false})$. In state $2$, the input
head moves left and the machine enters $(3,\mathrm{false})$. In state $3$, if the input
symbol is present, the input head moves left and the machine remains in
$(3,\mathrm{false})$; if the input symbol is blank, the input head moves right and the
machine enters $(4,\mathrm{false})$. In every other state, all heads stay put and the
machine enters $(4,\mathrm{false})$. No transition emits an output symbol.
\end{definition}

\begin{definition}[F2 split restoration scan configuration]
\lean{Turing.FinTM.f2_splitRestoreScan}
\label{Turing.FinTM.f2_splitRestoreScan}
\leanok
\uses{Turing.Cfg, Turing.FinTM.bufferTape, Turing.FinTM.f2_splitPos, Turing.FinTM.f2_splitRestoreTM, Turing.FinTM.f2_splitScratch}
For a natural number $k$, Boolean lists $w$ and $s$, and a natural number $j$, define a
configuration of the specified machine on input $w$. Its state is $(0, [|w|<j])$, its
input-head position is the position assigned to $w$ and $j$ by the split-position
construction, and its output is empty. Every work-head position is $j$; the first work
tape is the buffer encoding of $s$, and the remaining work tapes are given by the
scratch-tape construction with parameters $|s|+1$ and $j$.
\end{definition}

\begin{lemma}[Split restoration scan]
\lean{Turing.FinTM.f2_splitRestore_scan}
\label{Turing.FinTM.f2_splitRestore_scan}
\leanok
\uses{Turing.Action.apply, Turing.Cfg.inputSymbol, Turing.Cfg.workTapeSymbols, Turing.FinTM.f2_splitPos_read, Turing.FinTM.f2_splitPos_succ, Turing.FinTM.f2_splitRestoreScan, Turing.FinTM.f2_splitRestoreTM, Turing.FinTM.f2_splitScratch_erase, Turing.MultiTapeTM.runFrom, Turing.MultiTapeTM.runFrom_succ_eq_step', Turing.MultiTapeTM.step}
\difficulty{6}
For every natural number $k$, Boolean words $w,s$, and natural number $j\le |s|$,
running the split-restoration machine for $j$ steps from its scan configuration at
position $0$ yields the scan configuration at position $j$. There the state is $(0,
|w|<j)$, the input head is at position $\min(j,|w|)+1$, the first work tape stores $s$,
each other work tape contains true in cells $j\le z<|s|+1$ and is blank elsewhere, every
work-tape head is at $j$, and the output is empty.
\end{lemma}

\begin{definition}[Split restoration clean configuration]
\lean{Turing.FinTM.f2_splitRestoreClean}
\label{Turing.FinTM.f2_splitRestoreClean}
\leanok
\uses{Turing.Cfg, Turing.FinTM.bufferTape, Turing.FinTM.f2_splitRestoreTM}
For natural number $k$, Boolean lists $w,s$, a state $q$ of the specified restoration
machine, an input-head position among the $|w|+2$ positions, and an integer $h$, the
clean configuration on input $w$ has state $q$, that input-head position, empty output,
every work-head position equal to $h$, first work tape given by the buffer encoding of
$s$, and all remaining work tapes blank at every position.
\end{definition}

\begin{lemma}[Split restoration append step]
\lean{Turing.FinTM.f2_splitRestore_append}
\label{Turing.FinTM.f2_splitRestore_append}
\leanok
\uses{Turing.Action.apply, Turing.Cfg.workTapeSymbols, Turing.FinTM.bufferTape_append, Turing.FinTM.f2_splitPos, Turing.FinTM.f2_splitRestoreClean, Turing.FinTM.f2_splitRestoreScan, Turing.FinTM.f2_splitRestoreTM, Turing.FinTM.f2_splitScratch, Turing.FinTM.f2_splitScratch_erase, Turing.FinTM.f2_splitStep, Turing.MultiTapeTM.step, Turing.moveInputPos_zero}
\difficulty{5}
For every natural number $k$ and Boolean words $w,s$, one step of the split-restoration
machine from its scan configuration at position $|s|$ yields the clean configuration
with state $(1,\text{false})$, input head at position $\min(|s|,|w|)+1$, first work tape
storing $s$, every other work tape blank, all work-tape heads at $|s|-1$, and empty
output. The word stored on the first work tape in this resulting configuration is $s$
with true appended if $|s|\le |w|$, and is $s$ otherwise.
\end{lemma}

\begin{lemma}[Split restoration rewind]
\lean{Turing.FinTM.f2_splitRestore_rewind}
\label{Turing.FinTM.f2_splitRestore_rewind}
\leanok
\uses{Turing.Action.apply, Turing.Cfg.workTapeSymbols, Turing.FinTM.bufferTape_left, Turing.FinTM.bufferTape_nat, Turing.FinTM.f2_splitRestoreClean, Turing.FinTM.f2_splitRestoreScan, Turing.FinTM.f2_splitRestoreTM, Turing.MultiTapeTM.runFrom, Turing.MultiTapeTM.runFrom_succ_eq_step, Turing.MultiTapeTM.runFrom_zero, Turing.MultiTapeTM.step, Turing.moveInputPos_zero}
\difficulty{5}
For every natural number $k$, Boolean words $w,s$, input-head position $p$ among the
$|w|+2$ positions, and natural number $j\le |s|$, running the split-restoration machine
for $j+1$ steps from the clean configuration with state $(1,\text{false})$, input-head
position $p$, first work tape storing $s$, all other work tapes blank, and every
work-tape head at $j-1$ yields the clean configuration with state $(2,\text{false})$,
the same input-head position and tape contents, all other work tapes blank, all
work-tape heads at $0$, and empty output.
\end{lemma}

\begin{lemma}[Split restoration run]
\lean{Turing.FinTM.f2_splitRestore_run}
\label{Turing.FinTM.f2_splitRestore_run}
\leanok
\uses{Turing.Cfg.ofWords, Turing.FinTM.f2_catalogRewind, Turing.FinTM.f2_splitPos, Turing.FinTM.f2_splitRestoreClean, Turing.FinTM.f2_splitRestoreScan, Turing.FinTM.f2_splitRestoreTM, Turing.FinTM.f2_splitRestore_append, Turing.FinTM.f2_splitRestore_rewind, Turing.FinTM.f2_splitRestore_scan, Turing.FinTM.f2_splitStep, Turing.MultiTapeTM.runFrom, Turing.MultiTapeTM.runFrom_add, Turing.MultiTapeTM.runFrom_succ_eq_step', Turing.stateWord}
\difficulty{5}
For every natural number $k$ and Boolean words $w,s$, there is a natural number $t\le
2|s|+|w|+5$ such that the split-restoration machine, run from its scan configuration at
position $0$ for $t$ steps, reaches the configuration with state $(4,\text{false})$,
input head at position $1$, first work tape storing $s$ with true appended when $|s|\le
|w|$ and otherwise storing $s$, all other work tapes blank, all work-tape heads at $0$,
and empty output.
\end{lemma}

\begin{definition}[Output-aware action transformation]
\lean{Turing.FinTM.f2_splitCountAction}
\label{Turing.FinTM.f2_splitCountAction}
\leanok
\uses{Turing.Action}
Let $k$ be a nonnegative integer, let $S$ and $H$ be sets, and let
$e:S\times\{\mathrm{false},\mathrm{true}\}\to H$ and
$r:\{\mathrm{false},\mathrm{true}\}\to H$ be functions. Given a Boolean flag $o$, an
optional input bit $i$, and a one-step action $a$ of a $k$-work-tape machine over the
Boolean alphabet with states $S$, form an action with $k+1$ work tapes and states $H$ as
follows: its input head moves right if $a$ emits a bit and stays otherwise; its first
work tape is unchanged and its remaining work-tape actions are those of $a$; it emits no
output; and its successor state is always present, equal to $e(q,o')$ if $a$ has
successor state $q$, and to $r(o')$ otherwise, where $o'=o\lor(a\text{ emits a bit and
}i\text{ is absent})$.
\end{definition}

\begin{definition}[Split-count configuration construction]
\lean{Turing.FinTM.f2_splitCountCfg}
\label{Turing.FinTM.f2_splitCountCfg}
\leanok
\uses{Turing.Cfg, Turing.FinTM.bufferTape, Turing.FinTM.f2_splitPos}
Let $k$ be a nonnegative integer, let $S$ and $H$ be types, and let
$e:S\times\{\mathrm{false},\mathrm{true}\}\to H$ and
$r:\{\mathrm{false},\mathrm{true}\}\to H$ be functions. For finite Boolean strings $w$
and $s$, and a configuration $c$ with $k$ work tapes, Boolean symbols, state set $S$,
and empty input, define $o$ to be the truth value of $|w|<|s|+|c.\mathrm{output}|$. The
resulting configuration has $k+1$ work tapes, state set $H$, input $w$, state $e(q,o)$
if $c$ has state $q$ and $r(o)$ if $c$ has no state, input-head position given by the
split-position construction for $w$ and $|s|+|c.\mathrm{output}|$, first work tape given
by the buffer construction for $s$, remaining work tapes and their head positions
inherited from $c$, new first work-tape head at $0$, and empty output.
\end{definition}

\begin{lemma}[Updating the overflow flag]
\lean{Turing.FinTM.f2_splitCount_over}
\label{Turing.FinTM.f2_splitCount_over}
\leanok
\uses{Turing.Cfg, Turing.Cfg.inputSymbol, Turing.FinTM.f2_splitPos, Turing.FinTM.f2_splitPos_read}
\difficulty{4}
Let $k$ be a natural number, let $S$ be a type, let $w$ be a Boolean word, let $cfg$ be
a configuration over $w$, and let $j$ be a natural number. If the input-head position of
$cfg$ is the position associated with $j$, then the disjunction of the conditions
$|w|<j$ and that the scanned input symbol is blank is equivalent to $|w|<j+1$.
\end{lemma}

\begin{lemma}[One-step compatibility of the split counter]
\lean{Turing.FinTM.f2_splitCount_apply}
\label{Turing.FinTM.f2_splitCount_apply}
\leanok
\uses{Turing.Action, Turing.Action.apply, Turing.Cfg, Turing.Cfg.inputSymbol, Turing.FinTM.f2_splitCountAction, Turing.FinTM.f2_splitCountCfg, Turing.FinTM.f2_splitCount_over, Turing.FinTM.f2_splitPos_succ}
\difficulty{4}
Let $k$ be a natural number, let $S,H$ be types, let $emb:S\to H$ and
$ret:\mathrm{Bool}\to H$, let $w,s$ be Boolean words, let $c$ be a configuration over
the empty input with $k$ work tapes, and let $a$ be a one-step action for that
configuration. Applying the lifted action to the associated counter configuration gives
exactly the counter configuration associated with the result of applying $a$ to $c$.
\end{lemma}

\begin{lemma}[Simulation of a live counter run]
\lean{Turing.FinTM.f2_splitCount_run}
\label{Turing.FinTM.f2_splitCount_run}
\leanok
\uses{Turing.Cfg, Turing.Cfg.Halted, Turing.Cfg.inputSymbol, Turing.Cfg.workTapeSymbols, Turing.FinTM.f2_splitCountAction, Turing.FinTM.f2_splitCountCfg, Turing.FinTM.f2_splitCount_apply, Turing.MultiTapeTM, Turing.MultiTapeTM.runFrom, Turing.MultiTapeTM.runFrom_succ_eq_step', Turing.MultiTapeTM.step}
\difficulty{5}
Let $k$ be a natural number, let $S,H$ be types, and let $tm$ and $host$ be
Boolean-alphabet machines with respectively $k$ and $k+1$ work tapes. Let $emb:S\to H$
and $ret:\mathrm{Bool}\to H$, and suppose that for every state $q$, Boolean flag $over$,
input symbol option $inp$, and tuple $work$ of work-tape symbols, the transition of
$host$ from state $emb(q,over)$ agrees with the lifted action formed from the transition
of $tm$ on blank input and the last $k$ work tapes. For Boolean words $w,s$, a
configuration $c$ of $tm$ on empty input, and a natural number $t$, suppose that the run
from $c$ is not halted at any time $j<t$. Then after $t$ steps, $host$ started from the
counter configuration associated with $c$ is exactly the counter configuration
associated with the configuration reached by $tm$ from $c$ after $t$ steps.
\end{lemma}

\begin{lemma}[Polynomial loop termination configuration]
\lean{Turing.FinTM.f2_splitPoly_loop_end}
\label{Turing.FinTM.f2_splitPoly_loop_end}
\leanok
\uses{Turing.Action.apply, Turing.FinTM.f2_catalogPolyCfg, Turing.FinTM.f2_catalogPolyCost, Turing.FinTM.f2_catalogPolyUnaryTM, Turing.FinTM.f2_catalogPoly_loop, Turing.MultiTapeTM.runFrom, Turing.MultiTapeTM.runFrom_succ_eq_step', Turing.MultiTapeTM.step}
\difficulty{4}
For natural numbers $c,C,q$ with $q>0$, run the specified unary Boolean Turing machine
on the empty input from the configuration with loop control at index $c$, all work heads
at zero, and empty output. After $\mathrm{cost}(q,C,c+1)+1$ steps, the machine is halted
in the configuration with advance control at its last index, all work heads at zero, and
output consisting of $Cq^{c+1}$ true bits; its tape contents and input-head position are
those of the specified catalog configuration.
\end{lemma}

\begin{lemma}[First visit to a state]
\lean{Turing.FinTM.f2_catalogFirstEntry}
\label{Turing.FinTM.f2_catalogFirstEntry}
\leanok
\uses{Turing.Cfg, Turing.MultiTapeTM, Turing.MultiTapeTM.runFrom, Turing.MultiTapeTM.runFrom_add, Turing.MultiTapeTM.step}
\difficulty{4}
Let $M$ be a multi-tape Turing machine, let $q$ be a state, and let $c,d$ be
configurations on a fixed Boolean input. Suppose every configuration whose state is $q$
is fixed by one step of $M$. If $d$ has state $q$ and running $M$ from $c$ for $T$ steps
reaches $d$, then there is a time $t\le T$ at which the run from $c$ reaches $d$, and at
every earlier time $j<t$ its state is different from $q$.
\end{lemma}

\begin{lemma}[First split restoration phase]
\lean{Turing.FinTM.f2_splitRestore_first}
\label{Turing.FinTM.f2_splitRestore_first}
\leanok
\uses{Turing.Action.apply, Turing.Cfg, Turing.Cfg.ofWords, Turing.FinTM.controlAction, Turing.FinTM.controlAction_apply, Turing.FinTM.f2_catalogFirstEntry, Turing.FinTM.f2_splitRestoreScan, Turing.FinTM.f2_splitRestoreTM, Turing.FinTM.f2_splitRestore_run, Turing.FinTM.f2_splitStep, Turing.MultiTapeTM.runFrom, Turing.MultiTapeTM.runFrom_zero, Turing.MultiTapeTM.step, Turing.moveInputPos_zero, Turing.stateWord}
\difficulty{5}
For every natural number $k$ and Boolean words $w,s$, there is a positive integer $t\le
2|s|+|w|+5$ such that, for every $j<t$, the split-restoration machine run from its scan
configuration at position $0$ is not in state $(4,\text{false})$ after $j$ steps, and
after $t$ steps it is in the configuration with state $(4,\text{false})$, input head at
position $1$, first work tape storing the word obtained from $s$ by appending true when
$|s|\le |w|$ and leaving $s$ unchanged otherwise, all other work tapes blank, all
work-tape heads at $0$, and empty output.
\end{lemma}

\begin{lemma}[Acceptance condition for the split counter]
\lean{Turing.FinTM.f2_splitCount_accept}
\label{Turing.FinTM.f2_splitCount_accept}
\leanok
\uses{Turing.Cfg, Turing.Cfg.inputSymbol, Turing.FinTM.f2_splitCountCfg, Turing.FinTM.f2_splitPos_read}
\difficulty{4}
Let $k$ be a natural number, let $S$ and $H$ be types, let $\mathrm{emb}:S\to H$ and
$\mathrm{ret}:\{\mathrm{false},\mathrm{true}\}\to H$, and let $w$ and $s$ be finite
Boolean words. For a configuration $c$ of a $k$-tape Turing machine with Boolean
work-tape symbols and empty input, the conjunction of the condition $|w|\ge
|s|+|c.\mathrm{output}|$ and the condition that the input head of the associated
$(k+1)$-tape configuration reads no symbol is equivalent to
$|s|+|c.\mathrm{output}|=|w|$.
\end{lemma}

\begin{lemma}[Preserving a halted computation]
\lean{Turing.FinTM.f2_splitCount_firstHalt}
\label{Turing.FinTM.f2_splitCount_firstHalt}
\leanok
\uses{Turing.Cfg, Turing.FinTM.f2_splitCountAction, Turing.FinTM.f2_splitCountCfg, Turing.FinTM.f2_splitCount_run, Turing.MultiTapeTM, Turing.MultiTapeTM.runFrom, Turing.MultiTapeTM.runFrom_add, Turing.MultiTapeTM.runFrom_of_halt}
\difficulty{4}
Let $k$ be a natural number, let $S,H$ be types, and let $tm$ and $host$ be
Boolean-alphabet machines with respectively $k$ and $k+1$ work tapes. Let $emb:S\to H$
and $ret:\mathrm{Bool}\to H$. Suppose that, for every state $q$, Boolean flag $over$,
input symbol option $inp$, and tuple of work-tape symbols $work$, the transition of
$host$ from state $emb(q,over)$ agrees with the lifted action formed from the transition
of $tm$ on blank input and the last $k$ work tapes. For Boolean words $w,s$,
configurations $c,d$ of $tm$ on empty input, and a natural number $T$, if $d$ is halted
and $tm$ reaches $d$ from $c$ in $T$ steps, then there is a time $t\le T$ at which
$host$, started from the counter configuration associated with $c$, reaches the counter
configuration associated with $d$.
\end{lemma}

\begin{definition}[Split-preparation machine]
\lean{Turing.FinTM.f2_splitPrepareTM}
\label{Turing.FinTM.f2_splitPrepareTM}
\leanok
\uses{Turing.FinTM, Turing.FinTM.controlAction}
For each nonnegative integer $k$, define a finite Turing machine over the Boolean
alphabet with $k+1$ work tapes and state set
$\{0,1,2\}\times\{\mathrm{false},\mathrm{true}\}$, initially in $(0,\mathrm{false})$. In
phase $0$, if the first work-tape cell is nonblank, it moves the input head right, moves
every work head right, writes the bit $1$ on each work tape other than the first, and
stays in phase $0$, setting the Boolean component to true if it was already true or the
input cell is blank. If that cell is blank, it keeps the input head still, moves every
work head left, writes $1$ on each tape other than the first, and enters phase $1$
without changing the Boolean component. In phase $1$, it moves the input head still and
moves all work heads left while the first work-tape cell is nonblank; when that cell is
blank, it moves all work heads right and enters phase $2$. In phase $2$ it stays in
place and remains in phase $2$, preserving the Boolean component.
\end{definition}

\begin{definition}[Split preparation scan configuration]
\lean{Turing.FinTM.f2_splitPrepareScan}
\label{Turing.FinTM.f2_splitPrepareScan}
\leanok
\uses{Turing.Cfg, Turing.FinTM.bufferTape, Turing.FinTM.f2_catalogPolyTape, Turing.FinTM.f2_splitPos, Turing.FinTM.f2_splitPrepareTM}
For $k\in\mathbb N$, Boolean strings $w,s$, and $j\in\mathbb N$, let $p(w,j)$ denote the
input-head position assigned to $w$ and $j$ by the split-position construction, let
$B_s$ be the buffer tape formed from $s$, and let $P_j$ be the catalog polynomial tape
associated with $j$. The split-preparation machine with $k+1$ work tapes has a
configuration on input $w$ whose state is $(0,\mathbf
1_{|w|>j\,\text{or}\,|w|<j})=(0,\mathbf 1_{|w|<j})$, whose input-head position is
$p(w,j)$, whose first work tape is $B_s$ and whose remaining work tapes are each $P_j$,
with every work-tape head at position $j$ and empty output.
\end{definition}

\begin{lemma}[Split preparation scan]
\lean{Turing.FinTM.f2_splitPrepare_scan}
\label{Turing.FinTM.f2_splitPrepare_scan}
\leanok
\uses{Turing.Action.apply, Turing.Cfg.workTapeSymbols, Turing.FinTM.f2_catalogPolyTape_write, Turing.FinTM.f2_splitCount_over, Turing.FinTM.f2_splitPos_succ, Turing.FinTM.f2_splitPrepareScan, Turing.FinTM.f2_splitPrepareTM, Turing.MultiTapeTM.runFrom, Turing.MultiTapeTM.runFrom_succ_eq_step', Turing.MultiTapeTM.step}
\difficulty{5}
For every natural number $k$, Boolean words $w,s$, and natural number $j\le |s|$,
running the split-preparation machine for $j$ steps from the scan configuration at
position $0$ yields the scan configuration at position $j$. At position $j$, the state
is $(0, |w|<j)$, the input head is at $\min(j,|w|)+1$, the first work tape stores $s$,
each other work tape contains true exactly in cells $j,\ldots,|s|$, every work-tape head
is at $j$, and the output is empty.
\end{lemma}

\begin{definition}[Split-preparation ready configuration]
\lean{Turing.FinTM.f2_splitPrepareReady}
\label{Turing.FinTM.f2_splitPrepareReady}
\leanok
\uses{Turing.Cfg, Turing.FinTM.bufferTape, Turing.FinTM.f2_catalogPolyTape, Turing.FinTM.f2_splitPos, Turing.FinTM.f2_splitPrepareTM}
For each $k\in\mathbb{N}$, let $w$ and $s$ be finite Boolean strings, let
$q\in\{0,1,2\}$, and let $h\in\mathbb{Z}$. The split-preparation construction assigns a
configuration on input $w$ with state $(q,\,|w|<|s|)$, input-head position $P(w,|s|)$,
first work tape $B(s)$, and each remaining work tape $C_{|s|+1}$; every work-tape head
is at $h$, and the output is empty.
\end{definition}

\begin{lemma}[First preparation step]
\lean{Turing.FinTM.f2_splitPrepare_extra}
\label{Turing.FinTM.f2_splitPrepare_extra}
\leanok
\uses{Turing.Action.apply, Turing.Cfg.workTapeSymbols, Turing.FinTM.f2_catalogPolyTape_write, Turing.FinTM.f2_splitPrepareReady, Turing.FinTM.f2_splitPrepareScan, Turing.FinTM.f2_splitPrepareTM, Turing.MultiTapeTM.step, Turing.moveInputPos_zero}
\difficulty{4}
For every natural number $k$ and Boolean words $w,s$, one step of the specified
preparation machine from its scan configuration at position $|s|$ gives its ready
configuration in state $1$, with the flag recording whether $|w|<|s|$ and the first
work-tape head at $|s|-1$.
\end{lemma}

\begin{lemma}[Preparation rewind completion]
\lean{Turing.FinTM.f2_splitPrepare_rewind}
\label{Turing.FinTM.f2_splitPrepare_rewind}
\leanok
\uses{Turing.Action.apply, Turing.Cfg.workTapeSymbols, Turing.FinTM.bufferTape_left, Turing.FinTM.bufferTape_nat, Turing.FinTM.f2_splitPrepareReady, Turing.FinTM.f2_splitPrepareTM, Turing.MultiTapeTM.runFrom, Turing.MultiTapeTM.runFrom_succ_eq_step, Turing.MultiTapeTM.runFrom_zero, Turing.MultiTapeTM.step, Turing.moveInputPos_zero}
\difficulty{5}
For every natural number $k$, Boolean words $w,s$, and natural number $j\le |s|$,
starting from the specified ready configuration in state $1$ with first work-tape head
at $j-1$, the preparation machine reaches the specified ready configuration in state $2$
with all work-tape heads at zero after $j+1$ steps.
\end{lemma}

\begin{lemma}[Split preparation run]
\lean{Turing.FinTM.f2_splitPrepare_run}
\label{Turing.FinTM.f2_splitPrepare_run}
\leanok
\uses{Turing.Cfg.ofWords, Turing.FinTM.f2_catalogPolyTape, Turing.FinTM.f2_splitPos, Turing.FinTM.f2_splitPrepareReady, Turing.FinTM.f2_splitPrepareScan, Turing.FinTM.f2_splitPrepareTM, Turing.FinTM.f2_splitPrepare_extra, Turing.FinTM.f2_splitPrepare_rewind, Turing.FinTM.f2_splitPrepare_scan, Turing.MultiTapeTM.runFrom, Turing.MultiTapeTM.runFrom_add, Turing.MultiTapeTM.runFrom_succ_eq_step', Turing.stateWord}
\difficulty{4}
For every natural number $k$ and Boolean words $w,s$, running the split-preparation
machine on input $w$, from its initial configuration with state $(0,	ext{false})$ and
with $s$ on the first work tape and all other work tapes empty, for $2(|s|+1)$ steps
yields the configuration with state $(2, |w|<|s|)$, input head at position
$\min(|s|,|w|)+1$, first work tape storing $s$, every other work tape containing true
exactly in cells $0,\ldots,|s|$, all work-tape heads at $0$, and empty output.
\end{lemma}

\begin{lemma}[Preparation reaches ready state]
\lean{Turing.FinTM.f2_splitPrepare_first}
\label{Turing.FinTM.f2_splitPrepare_first}
\leanok
\uses{Turing.Action.apply, Turing.Cfg.ofWords, Turing.FinTM.controlAction, Turing.FinTM.controlAction_apply, Turing.FinTM.f2_catalogFirstEntry, Turing.FinTM.f2_splitPrepareReady, Turing.FinTM.f2_splitPrepareTM, Turing.FinTM.f2_splitPrepare_run, Turing.MultiTapeTM.runFrom, Turing.MultiTapeTM.step, Turing.moveInputPos_zero, Turing.stateWord}
\difficulty{3}
For every natural number $k$ and Boolean words $w,s$, there is a time $t\le 2(|s|+1)$
such that, from the specified initial configuration on input $w$ with $s$ on the first
work tape and all other work tapes empty, the preparation machine is never in state
$(2,\,[|w|<|s|])$ at any time $j<t$, and at time $t$ it is in the specified ready
configuration with that state, input-head position $1+\min(|s|,|w|)$, and all work-tape
heads at zero.
\end{lemma}

\begin{definition}[Avoiding an anchor state]
\lean{Turing.FinTM.f2_splitSafe}
\label{Turing.FinTM.f2_splitSafe}
\leanok
\uses{Turing.Cfg, Turing.MultiTapeTM, Turing.MultiTapeTM.runFrom}
Let $k$ be a nonnegative integer, let $S$ be a set, let $w$ be a finite Boolean string,
let $M$ be a deterministic machine with $k$ work tapes over Boolean symbols and states
$S$, let $a\in S$, let $c$ be a configuration of $M$ on input $w$, and let $t\in\mathbb
N$. The configuration reached from $c$ after every number $j\le t$ of steps has state
different from $a$.
\end{definition}

\begin{lemma}[Concatenating safe computation intervals]
\lean{Turing.FinTM.f2_splitSafe_add}
\label{Turing.FinTM.f2_splitSafe_add}
\leanok
\uses{Turing.Cfg, Turing.FinTM.f2_splitSafe, Turing.MultiTapeTM, Turing.MultiTapeTM.runFrom, Turing.MultiTapeTM.runFrom_add}
\difficulty{4}
Let $k$ be a natural number, $S$ a state set, $w$ a Boolean input word, $M$ a
$k$-work-tape machine over Booleans, $a\in S$ an anchor state, and $c$ a configuration
of $M$ on input $w$. If the run from $c$ avoids state $a$ at every time $j\le u$, and
the run of length $v$ starting from the configuration reached after $u$ steps from $c$
also avoids state $a$ at every time $j\le v$, then the run from $c$ avoids state $a$ at
every time $j\le u+v$.
\end{lemma}

\begin{lemma}[Safe embedded run to a stopped configuration]
\lean{Turing.FinTM.f2_splitEmbed_cut}
\label{Turing.FinTM.f2_splitEmbed_cut}
\leanok
\uses{Turing.Action.mapState, Turing.Cfg, Turing.Cfg.mapState, Turing.FinTM.f2_splitSafe, Turing.MultiTapeTM, Turing.MultiTapeTM.runFrom, Turing.MultiTapeTM.runFrom_add, Turing.MultiTapeTM.runFrom_succ_eq_step', Turing.MultiTapeTM.step}
\difficulty{6}
Let $k$ be a natural number, let $S$ and $H$ be state sets, let $w$ be a Boolean input
word, and let $tm$ and $host$ be machines with $k$ work tapes over the Boolean alphabet.
Suppose $emb:S\to H$ and $anchor\in H$ satisfy $emb(q)\ne anchor$ for every $q\in S$;
suppose also that every configuration of $tm$ whose state is a stopped state is fixed by
one step, and that for each non-stopped state $q$, the transition of $host$ at $emb(q)$
is the state-renaming of the transition of $tm$ at $q$, for every input symbol and tuple
of work-tape symbols. If a configuration $d$ has a stopped state and is reached from a
configuration $c$ of $tm$ after $T$ steps, then there is some $t\le T$ such that running
$host$ from the state-renaming of $c$ for $t$ steps gives the state-renaming of $d$, and
at every time $j\le t$ the state of that host computation is different from $anchor$.
\end{lemma}

\begin{definition}[Input-head rewind machine]
\lean{Turing.FinTM.f2_splitRewindTM}
\label{Turing.FinTM.f2_splitRewindTM}
\leanok
\uses{Turing.FinTM, Turing.FinTM.controlAction}
For each nonnegative integer $k$, define a finite Turing machine over the Boolean
alphabet with $k+1$ work tapes and state set $\{0,1,2\}$, initially in state $0$. In
state $0$ it moves the input head left and enters state $1$. In state $1$, it moves left
and stays in state $1$ when the scanned input cell contains a bit; when that cell is
blank, it moves right and enters state $2$. In state $2$ it stays in place and remains
in state $2$. In all cases it leaves every work tape unchanged, emits no output, and has
a successor state.
\end{definition}

\begin{definition}[Split-emission machine]
\lean{Turing.FinTM.f2_splitEmitTM}
\label{Turing.FinTM.f2_splitEmitTM}
\leanok
\uses{Turing.FinTM, Turing.FinTM.controlAction}
For each nonnegative integer $k$, define a finite Turing machine over the Boolean
alphabet with $k+1$ work tapes and state set $\{0,1,2,3\}$, initially in state $0$. In
state $0$, if the first work-tape cell is blank, it stays in place, emits $0$, and
enters state $2$; otherwise it stays in place, emits the scanned input bit if present,
and enters state $1$. In state $1$ it moves the input head right, moves the first work
head right, leaves the other work heads still, emits the scanned input bit if present,
and returns to state $0$. In state $2$ it stays in place, emits $1$, and enters state
$3$. In state $3$, if the scanned input cell contains a bit, it moves the input head
right, emits that bit, and remains in state $3$; if the cell is blank, it stays in
place, emits nothing, and halts. All work tapes are otherwise unchanged.
\end{definition}

\begin{definition}[Split-body control states]
\lean{Turing.FinTM.f2_SplitBodyState}
\label{Turing.FinTM.f2_SplitBodyState}
\leanok
For a type $S$, define a control-state type with an anchor state; preparation states
indexed by $	ext{Fin}(3)\times\{\mathrm{false},\mathrm{true}\}$; counting states indexed
by $S$ and a Boolean; checking states indexed by a Boolean; rewind states indexed by a
Boolean and $	ext{Fin}(3)$; restore states indexed by
$	ext{Fin}(5)\times\{\mathrm{false},\mathrm{true}\}$; and emit states indexed by
$	ext{Fin}(4)$.
\end{definition}

\begin{definition}[Split-body machine]
\lean{Turing.FinTM.f2_splitBodyTM}
\label{Turing.FinTM.f2_splitBodyTM}
\leanok
\uses{Turing.Action.mapState, Turing.FinTM, Turing.FinTM.controlAction, Turing.FinTM.f2_SplitBodyState, Turing.FinTM.f2_splitCountAction, Turing.FinTM.f2_splitEmitTM, Turing.FinTM.f2_splitPrepareTM, Turing.FinTM.f2_splitRestoreTM, Turing.FinTM.f2_splitRewindTM}
Given a finite Turing machine $M$ over the Boolean alphabet and a start state of $M$,
define a finite Turing machine with one more work tape and control states for the
anchor, preparation, counting, checking, rewinding, restoring, and emitting stages. Its
initial state is the anchor; its transitions implement these stages, using $M$ during
counting and the corresponding preparation, rewind, restore, and emit procedures in
their stages.
\end{definition}

\begin{definition}[Split-bank configuration]
\lean{Turing.FinTM.f2_splitBank}
\label{Turing.FinTM.f2_splitBank}
\leanok
\uses{Turing.Cfg, Turing.FinTM, Turing.FinTM.f2_catalogPolyTape}
Given a finite Turing machine $M$ over the Boolean alphabet, a Boolean list $s$, an
optional state of $M$, and a Boolean output list, define a configuration for $M$ on
empty input with that state and output, input-head position $1$, every work tape equal
to the catalog polynomial tape of length parameter $|s|+1$, and every work-head position
$0$.
\end{definition}

\begin{lemma}[Initial split-body configuration]
\lean{Turing.FinTM.f2_splitBody_start}
\label{Turing.FinTM.f2_splitBody_start}
\leanok
\uses{Turing.Cfg.ofWords, Turing.FinTM, Turing.FinTM.f2_splitBodyTM, Turing.MultiTapeTM.initCfg, Turing.initCfg_ofWords, Turing.stateWord}
\difficulty{3}
Let $M$ be a finite Turing machine over the Boolean alphabet, let $start$ be one of its
states, and let $w$ be a Boolean input word. The initial configuration of the split-body
machine on $w$ has the anchor state, its input head at position $1$, blank work tapes
with all heads at position $0$, and empty output.
\end{lemma}

\begin{lemma}[State renaming preserves computation]
\lean{Turing.FinTM.f2_splitEmbed_run}
\label{Turing.FinTM.f2_splitEmbed_run}
\leanok
\uses{Turing.Action.mapState, Turing.Cfg, Turing.Cfg.mapState, Turing.MultiTapeTM, Turing.MultiTapeTM.runFrom, Turing.MultiTapeTM.runFrom_comm_of_step, Turing.MultiTapeTM.step}
\difficulty{4}
Let $k$ be a natural number, let $S$ and $H$ be state sets, let $w$ be a Boolean input
word, and let $tm$ and $host$ be machines with $k$ work tapes over Booleans. Suppose
$emb:S\to H$ is a state map such that, for every state and every input and work-tape
symbol tuple, the transition of $host$ at the embedded state is the state-renaming of
the transition of $tm$. Then for every configuration $c$ and every natural number $t$,
running $host$ from the configuration obtained by renaming the state of $c$ for $t$
steps gives the state-renaming of the configuration reached by running $tm$ from $c$ for
$t$ steps.
\end{lemma}

\begin{lemma}[Preparation reaches the counting configuration]
\lean{Turing.FinTM.f2_splitBody_prepare}
\label{Turing.FinTM.f2_splitBody_prepare}
\leanok
\uses{Turing.Action.apply, Turing.Cfg.mapState, Turing.Cfg.ofWords, Turing.FinTM, Turing.FinTM.controlAction, Turing.FinTM.controlAction_apply, Turing.FinTM.f2_splitBank, Turing.FinTM.f2_splitBodyTM, Turing.FinTM.f2_splitCountCfg, Turing.FinTM.f2_splitEmbed_cut, Turing.FinTM.f2_splitPrepareReady, Turing.FinTM.f2_splitPrepareTM, Turing.FinTM.f2_splitPrepare_run, Turing.FinTM.f2_splitSafe, Turing.MultiTapeTM.runFrom, Turing.MultiTapeTM.runFrom_succ_eq_step', Turing.MultiTapeTM.step, Turing.moveInputPos_zero, Turing.stateWord}
\difficulty{6}
Let $M$ be a finite Turing machine over the Boolean alphabet, let $a$ be a state of $M$,
and let $w$ and $s$ be Boolean words. There exist a natural number $t\le 2(|s|+1)$ such
that, after $t+1$ steps, the split-body machine started in its preparation state with
$s$ on its first work tape and all other work tapes blank, on input $w$, is in the
counting configuration initialized with the simulated machine in state $a$, the
auxiliary tape storing $s$, and empty output. At every time from $0$ through $t+1$, its
state is not the anchor state.
\end{lemma}

\begin{lemma}[A bounded safe split-body simulation]
\lean{Turing.FinTM.f2_splitBody_count}
\label{Turing.FinTM.f2_splitBody_count}
\leanok
\uses{Turing.FinTM, Turing.FinTM.f2_splitBank, Turing.FinTM.f2_splitBodyTM, Turing.FinTM.f2_splitCountCfg, Turing.FinTM.f2_splitCount_run, Turing.FinTM.f2_splitSafe, Turing.MultiTapeTM.runFrom, Turing.MultiTapeTM.runFrom_add, Turing.MultiTapeTM.runFrom_of_halt}
\difficulty{5}
Let $M$ be a finite Turing machine over Boolean symbols, let $a$ be one of its states,
let $w,s$ be Boolean words, and let $o$ be a Boolean output word. Suppose that the run
of $M$ on the empty input, started in state $a$ with the bank tape initialized from $s$
and empty output, halts after exactly $T$ steps with output $o$. Then the constructed
split-body machine, started in its count configuration encoding this initial run
together with $w$ and $s$, reaches the corresponding count configuration encoding the
halted run and output $o$ after some $t\\le T$ steps; throughout steps $0$ through $t$,
it never enters its anchor state.
\end{lemma}

\begin{lemma}[Rewind to the first input position]
\lean{Turing.FinTM.f2_splitBody_rewind}
\label{Turing.FinTM.f2_splitBody_rewind}
\leanok
\uses{Turing.Action.apply, Turing.Cfg, Turing.Cfg.mapState, Turing.FinTM, Turing.FinTM.controlAction, Turing.FinTM.controlAction_apply, Turing.FinTM.f2_catalogRewind, Turing.FinTM.f2_splitBodyTM, Turing.FinTM.f2_splitEmbed_cut, Turing.FinTM.f2_splitRewindTM, Turing.FinTM.f2_splitSafe, Turing.MultiTapeTM.runFrom, Turing.MultiTapeTM.step, Turing.moveInputPos_zero}
\difficulty{4}
Let $M$ be a finite Turing machine over the Boolean alphabet, let $start$ be a state of
$M$, and let $w$ be a Boolean input word. Suppose $c$ is a configuration on input $w$
with one more work tape than $M$, with control state $0$ among three numbered states.
For either Boolean value $ok$, if $c$ is in control state $0$, then starting from the
same configuration with its state renamed to the rewind phase carrying $ok$, the
split-body machine reaches, after some $t \le r+2$ steps, a configuration with the same
tapes, work-head positions, and output as $c$, input-head position $1$, and control
state $2$, where $r$ is the initial input-head position of $c$. At every time $j \le t$
in this run, the machine is not in its anchor state.
\end{lemma}

\begin{lemma}[Restoring the split-body configuration]
\lean{Turing.FinTM.f2_splitBody_restore}
\label{Turing.FinTM.f2_splitBody_restore}
\leanok
\uses{Turing.Action.apply, Turing.Cfg.mapState, Turing.Cfg.ofWords, Turing.FinTM, Turing.FinTM.controlAction, Turing.FinTM.controlAction_apply, Turing.FinTM.f2_splitBodyTM, Turing.FinTM.f2_splitEmbed_cut, Turing.FinTM.f2_splitRestoreScan, Turing.FinTM.f2_splitRestoreTM, Turing.FinTM.f2_splitRestore_run, Turing.FinTM.f2_splitSafe, Turing.FinTM.f2_splitStep, Turing.MultiTapeTM.runFrom, Turing.MultiTapeTM.step, Turing.moveInputPos_zero, Turing.stateWord}
\difficulty{3}
Let $M$ be a finite Turing machine over Boolean symbols, let $a$ be one of its states,
and let $w,s$ be Boolean words. Starting from the constructed restore-scan configuration
indexed by zero, the split-body machine reaches the seam configuration with restore
state $(4,\mathrm{false})$ and the updated tape word obtained by appending $1$ to $s$
when $|s|\\le |w|$, and otherwise leaving $s$ unchanged, after some $t\\le 2|s|+|w|+5$
steps. Throughout these steps it avoids its anchor state.
\end{lemma}

\begin{definition}[Split-emission configuration]
\lean{Turing.FinTM.f2_splitEmitCfg}
\label{Turing.FinTM.f2_splitEmitCfg}
\leanok
\uses{Turing.Cfg, Turing.FinTM.bufferTape, Turing.FinTM.f2_catalogPolyTape, Turing.FinTM.f2_splitPos}
Given a natural number $k$, Boolean lists $w$ and $s$, an optional state from
$	ext{Fin}(4)$, natural numbers $j$ and $h$, and a Boolean output list, define a
configuration with $k+1$ work tapes on input $w$. Its input-head position is the split
position determined by $w$ and $j$; its first work tape is the buffer tape determined by
$s$, each remaining work tape is the catalog polynomial tape with length parameter
$|s|+1$, the first work-head position is $h$, and the remaining work-head positions are
$0$; its state and output are the specified ones.
\end{definition}

\begin{lemma}[Output after paired input steps]
\lean{Turing.FinTM.f2_splitEmit_double}
\label{Turing.FinTM.f2_splitEmit_double}
\leanok
\uses{Turing.Action.apply, Turing.Cfg.inputSymbol, Turing.Cfg.workTapeSymbols, Turing.FinTM.f2_splitEmitCfg, Turing.FinTM.f2_splitEmitTM, Turing.FinTM.f2_splitPos_read, Turing.FinTM.f2_splitPos_succ, Turing.MultiTapeTM.runFrom, Turing.MultiTapeTM.runFrom_succ_eq_step', Turing.MultiTapeTM.step, Turing.moveInputPos_zero}
\difficulty{5}
For every natural number $k$, Boolean words $w,s$ with $|s|\le |w|$, and natural number
$j\le |s|$, running the specified split-emitting machine from its initial configuration
on $w$ with buffer $s$ for $2j$ steps yields the configuration with state $0$, input
head at position $j$, auxiliary head at position $j$, and output
$(w\mathbin{\upharpoonright}j)$ with each bit repeated twice in order. Here
$w\mathbin{\upharpoonright}j$ denotes the prefix of $w$ of length $j$.
\end{lemma}

\begin{lemma}[Separator emitted in two steps]
\lean{Turing.FinTM.f2_splitEmit_separator}
\label{Turing.FinTM.f2_splitEmit_separator}
\leanok
\uses{Turing.Action.apply, Turing.Cfg.workTapeSymbols, Turing.FinTM.f2_splitEmitCfg, Turing.FinTM.f2_splitEmitTM, Turing.MultiTapeTM.runFrom, Turing.MultiTapeTM.runFrom_succ_eq_step, Turing.MultiTapeTM.runFrom_zero, Turing.MultiTapeTM.step, Turing.moveInputPos_zero}
\difficulty{3}
For every natural number $k$, Boolean words $w,s$, and output word $out$, running the
specified split-emitting machine from the configuration in state $0$ with both heads at
position $|s|$ and output $out$ for two steps gives the configuration in state $3$, with
both heads still at position $|s|$ and output $out$ followed by $01$.
\end{lemma}

\begin{lemma}[Suffix emission and halting]
\lean{Turing.FinTM.f2_splitEmit_suffix}
\label{Turing.FinTM.f2_splitEmit_suffix}
\leanok
\uses{Turing.Action.apply, Turing.Cfg.inputSymbol, Turing.FinTM.controlAction, Turing.FinTM.f2_splitEmitCfg, Turing.FinTM.f2_splitEmitTM, Turing.FinTM.f2_splitPos_read, Turing.FinTM.f2_splitPos_succ, Turing.MultiTapeTM.runFrom, Turing.MultiTapeTM.runFrom_succ_eq_step, Turing.MultiTapeTM.runFrom_zero, Turing.MultiTapeTM.step}
\difficulty{5}
For every natural number $k$ and Boolean words $w,s,rest$, for every prefix $pre$,
output word $out$, and natural number $h$, if $w=pre\mathbin{++}rest$, then running the
specified split-emitting machine from state $3$, with input head at the position
determined by $|pre|$, auxiliary head at position $h$, and output $out$, for $|rest|+1$
steps yields the halted configuration with input head at position $|w|$, auxiliary head
still at $h$, and output $out$ followed by $rest$.
\end{lemma}

\begin{lemma}[Final split-emitting computation]
\lean{Turing.FinTM.f2_splitEmit_run}
\label{Turing.FinTM.f2_splitEmit_run}
\leanok
\uses{Turing.FinTM.f2_splitEmitCfg, Turing.FinTM.f2_splitEmitTM, Turing.FinTM.f2_splitEmit_double, Turing.FinTM.f2_splitEmit_separator, Turing.FinTM.f2_splitEmit_suffix, Turing.MultiTapeTM.runFrom, Turing.MultiTapeTM.runFrom_add, Turing.pairEncode}
\difficulty{3}
For every natural number $k$ and Boolean words $w,s$ with $|s|\le |w|$, running the
specified split-emitting machine from its initial configuration on input $w$ with buffer
$s$ for $|s|+|w|+3$ steps yields the halted configuration whose input head is at
position $|w|$, whose auxiliary head is at position $|s|$, and whose output is the
encoding formed by doubling each bit of the prefix of $w$ of length $|s|$, then
appending $01$, then appending the suffix of $w$ after that prefix.
\end{lemma}

\begin{lemma}[One-Step Safe Segment]
\lean{Turing.FinTM.f2_splitSafe_one}
\label{Turing.FinTM.f2_splitSafe_one}
\leanok
\uses{Turing.Cfg, Turing.FinTM.f2_splitSafe, Turing.MultiTapeTM, Turing.MultiTapeTM.runFrom, Turing.MultiTapeTM.step}
\difficulty{4}
Let $M$ be a deterministic multi-tape Turing machine with Boolean work-tape symbols, let
$a$ be a state, and let $c,d$ be configurations on the same input. If one step from $c$
produces $d$, and neither $c$ nor $d$ has state $a$, then the configuration after one
step from $c$ is $d$, and the run from $c$ avoids state $a$ at steps $0$ and $1$.
\end{lemma}

\begin{lemma}[Joining Safe Computation Segments]
\lean{Turing.FinTM.f2_splitSafe_join}
\label{Turing.FinTM.f2_splitSafe_join}
\leanok
\uses{Turing.Cfg, Turing.FinTM.f2_splitSafe, Turing.FinTM.f2_splitSafe_add, Turing.MultiTapeTM, Turing.MultiTapeTM.runFrom, Turing.MultiTapeTM.runFrom_add}
\difficulty{3}
Let $M$ be a deterministic multi-tape Turing machine with Boolean work-tape symbols, let
$a$ be a state, and let $c,d,f$ be configurations on the same input. Suppose that after
$u$ steps from $c$ the machine reaches $d$, and that during steps $0$ through $u$ its
state is never $a$. Suppose also that after $v$ steps from $d$ it reaches $f$, and
during steps $0$ through $v$ its state is never $a$. Then after $u+v$ steps from $c$ it
reaches $f$, and its state is never $a$ during any step from $0$ through $u+v$.
\end{lemma}

\begin{lemma}[One round of the split-body machine]
\lean{Turing.FinTM.f2_splitBody_round}
\label{Turing.FinTM.f2_splitBody_round}
\leanok
\uses{Turing.Action.apply, Turing.Cfg, Turing.Cfg.inputSymbol, Turing.Cfg.mapState, Turing.Cfg.ofWords, Turing.FinTM, Turing.FinTM.bufferTape, Turing.FinTM.controlAction, Turing.FinTM.controlAction_apply, Turing.FinTM.f2_SplitBodyState, Turing.FinTM.f2_catalogPolyTape, Turing.FinTM.f2_splitBank, Turing.FinTM.f2_splitBodyTM, Turing.FinTM.f2_splitBody_count, Turing.FinTM.f2_splitBody_prepare, Turing.FinTM.f2_splitBody_restore, Turing.FinTM.f2_splitBody_rewind, Turing.FinTM.f2_splitCountCfg, Turing.FinTM.f2_splitCount_accept, Turing.FinTM.f2_splitEmbed_run, Turing.FinTM.f2_splitEmitCfg, Turing.FinTM.f2_splitEmitTM, Turing.FinTM.f2_splitEmit_run, Turing.FinTM.f2_splitPos, Turing.FinTM.f2_splitRestoreScan, Turing.FinTM.f2_splitSafe, Turing.FinTM.f2_splitSafe_join, Turing.FinTM.f2_splitSafe_one, Turing.FinTM.f2_splitScratch, Turing.FinTM.f2_splitStep, Turing.MultiTapeTM.runFrom, Turing.MultiTapeTM.runFrom_add, Turing.MultiTapeTM.runFrom_succ_eq_step', Turing.MultiTapeTM.step, Turing.moveInputPos_zero, Turing.pairEncode, Turing.stateWord}
\difficulty{7}
Let $M$ be a finite Turing machine over the Boolean alphabet, let $start$ be one of its
states, and let $w,s,out$ be Boolean words. Suppose that from the bank configuration
storing $s$, the machine $M$ on empty input halts after $T$ steps with output $out$.
Then the split-body machine started on input $w$ in the anchor state with $s$ on its
first work tape and all other tapes blank has a time $t$ with $0<t\le T+5|s|+3|w|+20$,
and it is not in the anchor state at any intermediate time $j$ with $0<j<t$. If
$|s|+|out|=|w|$, its configuration at time $t$ is halted and its output is the encoding
of the first $|s|$ bits of $w$ paired with the remaining bits. Otherwise, its
configuration at time $t$ is the anchor seam configuration with the first tape word
obtained by appending one true bit to $s$ when $|s|\le |w|$, and equal to $s$ otherwise;
all other tape words are empty.
\end{lemma}

\begin{lemma}[Constructing a Polynomial-Time Split Solver]
\lean{Turing.FinTM.f2_splitSolve_of_body}
\label{Turing.FinTM.f2_splitSolve_of_body}
\leanok
\uses{Turing.Cfg.ofWords, Turing.FinTM, Turing.FinTM.ComputesFunInTime, Turing.FinTM.ComputesInTime.mono, Turing.FinTM.computesFunInTime_lengthBits_spaceUsed, Turing.FinTM.f2_exists_loopFind_space, Turing.FinTM.f2_splitAccept, Turing.FinTM.f2_splitLoop_bound, Turing.FinTM.f2_splitLoop_result, Turing.FinTM.f2_splitStep, Turing.FinTM.f2_splitStep_inv, Turing.MultiTapeTM.initCfg, Turing.MultiTapeTM.runFrom, Turing.MultiTapeTM.spaceUsed, Turing.pairEncode, Turing.solveSplit, Turing.stateWord}
\difficulty{5}
Let $C,e,A$ be natural numbers, let $B$ be a finite Turing machine over the Boolean
alphabet, and choose a state $a$ of $B$. Assume that for every Boolean word $w$ there is
a time $t\le A(|w|+1)^{e+1}$ such that, at every earlier time $t'<t$, the run of $B$
from its initial configuration on $w$ is not in state $a$, and at time $t$ it is in the
seam configuration with state $a$ and all work-tape words empty. Assume also that for
every pair of Boolean words $w,s$ with $|s|\le |w|+1$, there is a positive time $t\le
A(|w|+1)^{e+1}$ such that at every intermediate time $t'$ with $0<t'<t$ the run from the
seam configuration with state $a$, input $w$, and tape words assigned from $s$ is not in
state $a$. If $|s|+C(|s|+1)^e=|w|$, then at time $t$ this run is halted and has output
equal to the encoding of the prefix of $w$ of length $|s|$ together with the remaining
suffix; otherwise, at time $t$ it is the seam configuration with state $a$ and tape
words assigned from the word obtained by appending true to $s$ when $|s|\le |w|$, and
from $s$ unchanged otherwise. Under these assumptions, there exist a finite Boolean
Turing machine $M$ and a natural-number constant $c$ that compute, on each input $w$,
the encoded split specified by the equation $i+C(i+1)^e=|w|$ (and the empty word if no
such $i$ exists) within $c(|w|+1)^{e+2}$ steps, while for every input $x$ and every time
$t$ using at most $c(|x|+1)^{e+1}$ space.
\end{lemma}

\begin{lemma}[Polynomial Source Output]
\lean{Turing.FinTM.f2_splitSource_poly}
\label{Turing.FinTM.f2_splitSource_poly}
\leanok
\uses{Turing.FinTM.f2_catalogPolyCfg, Turing.FinTM.f2_catalogPolyCost, Turing.FinTM.f2_catalogPolyUnaryTM, Turing.FinTM.f2_splitBank, Turing.FinTM.f2_splitPoly_loop_end, Turing.MultiTapeTM.runFrom}
\difficulty{1}
For natural numbers $c,C$ and a Boolean word $s$, consider the specified unary-output
machine with $c+1$ work tapes and its polynomial-cost bound. Starting from the bank
configuration on input $s$ in the loop state indexed by $c$, after
$\mathrm{cost}(|s|+1,C,c+1)+1$ steps it is halted and its output consists of exactly
$C(|s|+1)^{c+1}$ copies of true.
\end{lemma}

\begin{lemma}[Constant-Source Prefix Output]
\lean{Turing.FinTM.f2_splitSource_constant}
\label{Turing.FinTM.f2_splitSource_constant}
\leanok
\uses{Turing.Action.apply, Turing.Cfg.ext_zero_tapes, Turing.Cfg.inputSymbol, Turing.FinTM.f2_catalogPrefixCfg, Turing.FinTM.f2_catalogPrefixTM, Turing.FinTM.f2_catalogPrefixTM_emit, Turing.FinTM.f2_splitBank, Turing.MultiTapeTM.initCfg, Turing.MultiTapeTM.runFrom, Turing.MultiTapeTM.runFrom_succ_eq_step', Turing.MultiTapeTM.step}
\difficulty{4}
For every natural number $C$ and Boolean word $s$, consider the machine that copies its
Boolean input to the output and then emits one true symbol for each of $C$ additional
steps. Starting from the specified bank configuration with this machine, input $s$, its
initial state, and empty output, after $C+1$ steps the machine is halted and the output
consists of exactly $C$ copies of true.
\end{lemma}

\begin{lemma}[Polynomial bound for the split-body running time]
\lean{Turing.FinTM.f2_splitBody_envelope}
\label{Turing.FinTM.f2_splitBody_envelope}
\leanok
\difficulty{4}
Let $C,e,l,n,T$ be natural numbers. If $l\\le n+1$ and $T\\le (C+1+5e)(l+1)^e+1$, then
$T+5l+3n+20\\le ((C+1+5e)2^e+40)(n+1)^{e+1}$.
\end{lemma}

\begin{lemma}[Polynomial-Time Split Solver from a Source Machine]
\lean{Turing.FinTM.f2_splitSolve_source}
\label{Turing.FinTM.f2_splitSolve_source}
\leanok
\uses{Turing.FinTM, Turing.FinTM.ComputesFunInTime, Turing.FinTM.f2_splitAccept, Turing.FinTM.f2_splitBank, Turing.FinTM.f2_splitBodyTM, Turing.FinTM.f2_splitBody_envelope, Turing.FinTM.f2_splitBody_round, Turing.FinTM.f2_splitBody_start, Turing.FinTM.f2_splitSolve_of_body, Turing.MultiTapeTM.initCfg, Turing.MultiTapeTM.runFrom, Turing.MultiTapeTM.spaceUsed, Turing.pairEncode, Turing.solveSplit}
\difficulty{4}
Let $C,e$ be natural numbers, let $M$ be a finite Turing machine over the Boolean
alphabet, choose a state $a$ of $M$, and let $B$ be a function from natural numbers to
natural numbers. Suppose that for every Boolean word $s$, running $M$ for $B(|s|)$ steps
from the specified bank configuration with state $a$, input empty, and empty output
yields the bank configuration with no state and output consisting of $C(|s|+1)^e$ copies
of true. Suppose also that $B(\ell)\le (C+1+5e)(\ell+1)^e+1$ for every natural number
$\ell$. Then there exist a finite Boolean Turing machine $N$ and a natural-number
constant $c$ that compute, on every input $w$, the encoded split determined by an $i$
satisfying $i+C(i+1)^e=|w|$, outputting the empty word if no such $i$ exists, within
$c(|w|+1)^{e+2}$ steps; moreover, for every input $x$ and every time $t$, the total
space used by $N$ through step $t$ is at most $c(|x|+1)^{e+1}$.
\end{lemma}

\begin{lemma}[Polynomial-Time Split Solver]
\lean{Turing.FinTM.f2_splitSolve_closed}
\label{Turing.FinTM.f2_splitSolve_closed}
\leanok
\uses{Turing.FinTM, Turing.FinTM.ComputesFunInTime, Turing.FinTM.f2_catalogPolyCost, Turing.FinTM.f2_catalogPolyCost_le, Turing.FinTM.f2_catalogPolyUnaryTM, Turing.FinTM.f2_catalogPrefixTM, Turing.FinTM.f2_splitSolve_source, Turing.FinTM.f2_splitSource_constant, Turing.FinTM.f2_splitSource_poly, Turing.MultiTapeTM.initCfg, Turing.MultiTapeTM.spaceUsed, Turing.pairEncode, Turing.solveSplit}
\difficulty{4}
For all natural numbers $C$ and $e$, there exist a finite Turing machine $M$ over the
Boolean alphabet and a natural-number constant $c$ such that, on every input word $w$,
$M$ halts within $c(|w|+1)^{e+2}$ steps and outputs the following word: if there is an
$i$ with $0\le i\le |w|$ and $i+C(i+1)^e=|w|$, it outputs the self-delimiting encoding
of the prefix of $w$ of length $i$ together with the remaining suffix; otherwise it
outputs the empty word. Moreover, for every input $x$ and every number of steps $t$, the
total space used by $M$ through step $t$ is at most $c(|x|+1)^{e+1}$.
\end{lemma}

\begin{theorem}[Polynomial-time split search with bounded space]
\lean{Turing.FinTM.computesFunInTime_splitSolve_spaceUsed}
\label{Turing.FinTM.computesFunInTime_splitSolve_spaceUsed}
\leanok
\uses{Turing.FinTM, Turing.FinTM.ComputesFunInTime, Turing.FinTM.f2_splitSolve_closed, Turing.MultiTapeTM.initCfg, Turing.MultiTapeTM.spaceUsed, Turing.pairEncode, Turing.solveSplit}
\difficulty{1}
For all natural numbers $C$ and $e$, there exist a finite Turing machine $M$ and a
natural number $c$ such that, on every Boolean input $w$ of length $n$, $M$ halts within
$c(n+1)^{e+2}$ steps and outputs the encoding that doubles each bit of the prefix of $w$
of length $i$, appends $01$, and then appends the remaining suffix, where $i$ is the
least integer in $0,\ldots,n$ satisfying $i+C(i+1)^e=n$; if no such $i$ exists, its
output is the empty word. For every Boolean word $x$ and every time $t$, the computation
of $M$ on $x$ up to time $t$ uses at most $c(|x|+1)^{e+1}$ work-tape cells.
\end{theorem}

\begin{definition}[Catalog redirection state]
\lean{Turing.FinTM.catalog_redirectState}
\label{Turing.FinTM.catalog_redirectState}
\leanok
For a type $S$, a Boolean $b$, an optional state $q$ of $S$, and an optional Boolean
$r$, define a redirected state as follows: if $q$ is present, it is the pair $(q,r)$ in
the left summand; if $q$ is absent, it is absent when $r=b$, and otherwise it is the
unique element of the right summand.
\end{definition}

\begin{definition}[Catalog redirection action]
\lean{Turing.FinTM.catalog_redirectAction}
\label{Turing.FinTM.catalog_redirectAction}
\leanok
\uses{Turing.Action, Turing.FinTM.catalog_redirectState}
Given a Boolean $b$, a single-step action on $k$ work tapes with Boolean symbols and
states in $S$, and an optional Boolean $r$, define an action with states in
$(S\times\mathrm{Option}(\mathrm{Bool}))\sqcup\mathrm{Unit}$. It preserves the
input-head movement and all work-tape actions, emits no symbol, and sets its successor
state to the catalog redirected state determined by $b$, the original successor state,
and the optional value obtained by taking the action's output when present and otherwise
using $r$.
\end{definition}

\begin{definition}[Catalog redirected configuration]
\lean{Turing.FinTM.catalog_redirectCfg}
\label{Turing.FinTM.catalog_redirectCfg}
\leanok
\uses{Turing.Cfg, Turing.FinTM, Turing.FinTM.catalog_redirectState, Turing.FinTM.redirectTM}
Given a finite Turing machine $M$ over the Boolean alphabet, a Boolean $b$, an input
list $x$, and a configuration $c$ of $M$ on $x$, define a configuration of the
redirected machine on the same input. Its state is the catalog redirected state
determined by $b$, the state of $c$, and the last output symbol of $c$ if one exists; it
preserves the input-head position, work tapes, and work-head positions, and has empty
output.
\end{definition}

\begin{lemma}[Redirect-state loop invariance]
\lean{Turing.FinTM.catalog_redirect_loop}
\label{Turing.FinTM.catalog_redirect_loop}
\leanok
\uses{Turing.Action.apply, Turing.Cfg, Turing.FinTM, Turing.FinTM.redirectTM, Turing.MultiTapeTM.runFrom, Turing.MultiTapeTM.runFrom_succ_eq_step', Turing.MultiTapeTM.step}
\difficulty{4}
Let $M$ be a finite-state machine over the Boolean alphabet, let $b$ be a Boolean value,
and let $x$ be a Boolean input word. Every configuration of the redirected machine on
$x$ whose state is the added loop state remains exactly unchanged after any number $t$
of steps.
\end{lemma}

\begin{lemma}[Redirected action correspondence]
\lean{Turing.FinTM.catalog_redirect_apply}
\label{Turing.FinTM.catalog_redirect_apply}
\leanok
\uses{Turing.Action, Turing.Action.apply, Turing.Cfg, Turing.FinTM, Turing.FinTM.catalog_redirectAction, Turing.FinTM.catalog_redirectCfg}
\difficulty{4}
Let $M$ be a finite-state machine over the Boolean alphabet with $k$ work tapes, let $b$
be a Boolean value, let $x$ be a Boolean input word, and let $c$ be a configuration of
$M$ on $x$. For any action $a$ of $M$ and $c$, applying the corresponding redirected
action to the redirected configuration of $c$ gives the redirected configuration of the
result of applying $a$ to $c$.
\end{lemma}

\begin{lemma}[Redirected step correspondence]
\lean{Turing.FinTM.catalog_redirect_step}
\label{Turing.FinTM.catalog_redirect_step}
\leanok
\uses{Turing.Cfg, Turing.Cfg.inputSymbol, Turing.Cfg.workTapeSymbols, Turing.FinTM, Turing.FinTM.catalog_redirectAction, Turing.FinTM.catalog_redirectCfg, Turing.FinTM.catalog_redirectState, Turing.FinTM.catalog_redirect_apply, Turing.FinTM.catalog_redirect_loop, Turing.FinTM.redirectTM, Turing.MultiTapeTM.step, Turing.MultiTapeTM.step_of_halt}
\difficulty{5}
Let $M$ be a finite-state machine over the Boolean alphabet, let $b$ be a Boolean value,
let $x$ be a Boolean input word, and let $c$ be any configuration of $M$ on $x$. Taking
one step from the redirected configuration of $c$ yields the redirected configuration of
the result of taking one step from $c$ in $M$.
\end{lemma}

\begin{lemma}[Redirected run correspondence]
\lean{Turing.FinTM.catalog_redirect_run}
\label{Turing.FinTM.catalog_redirect_run}
\leanok
\uses{Turing.FinTM, Turing.FinTM.catalog_redirectCfg, Turing.FinTM.catalog_redirect_step, Turing.FinTM.redirectTM, Turing.MultiTapeTM.initCfg, Turing.MultiTapeTM.runFrom, Turing.MultiTapeTM.runFrom_comm_of_step}
\difficulty{3}
Let $M$ be a finite-state machine over the Boolean alphabet, let $b$ be a Boolean value,
let $x$ be a Boolean input word, and let $t$ be a nonnegative integer. After $t$ steps
from its initial configuration on $x$, the redirected machine is in the redirected state
corresponding to $M$'s state and last output bit at time $t$, with the same input-head
position and work tapes as $M$; its output tape is empty.
\end{lemma}

\begin{theorem}[Redirection preserves per-tape space]
\lean{Turing.FinTM.redirectTM_spaceUsedByTape}
\label{Turing.FinTM.redirectTM_spaceUsedByTape}
\leanok
\uses{Turing.FinTM, Turing.FinTM.catalog_redirect_run, Turing.FinTM.redirectTM, Turing.MultiTapeTM.initCfg, Turing.MultiTapeTM.spaceUsedByTape, Turing.MultiTapeTM.visitedByTapeHead}
\difficulty{3}
Let $M$ be a finite-state Turing machine over the Boolean alphabet, let $b$ be a bit,
let $x$ be a Boolean word, let $t \in \mathbb{N}$, and let $i$ be one of its work-tape
indices. The number of distinct cells visited by tape $i$ in the redirected machine's
computation from its initial configuration on $x$ through step $t$ equals the number
visited by tape $i$ in the corresponding computation of $M$.
\end{theorem}

\begin{definition}[Finite-machine padding construction]
\lean{Turing.FinTM.f2_timedPadTM}
\label{Turing.FinTM.f2_timedPadTM}
\leanok
\uses{Turing.FinTM, Turing.FinTM.leftAction, Turing.MultiTapeTM}
Given a finite Turing machine $D$ over the Boolean alphabet and a natural number $r$,
define a multi-tape Turing machine with $D$’s initial state and state set, and with
$D.k+r$ work tapes. Its transition in state $q$, on input symbol $a$ and work-tape
symbols $w$, is obtained by applying the left-action transformation with parameter $r$
and the identity map to the action of $D$ in state $q$ on $a$ and the first $D.k$
work-tape symbols.
\end{definition}

\begin{definition}[Timed conditional machine]
\lean{Turing.FinTM.f2_timedCondTM}
\label{Turing.FinTM.f2_timedCondTM}
\leanok
\uses{Turing.FinTM, Turing.FinTM.branchTM, Turing.FinTM.controlAction, Turing.FinTM.f2_timedPadTM, Turing.FinTM.leftAction, Turing.FinTM.rightAction, Turing.captureAction}
Given three finite Turing machines $D$, $M_1$, and $M_2$ over the Boolean alphabet,
define a finite machine with $k_D+k_1+k_2+1$ work tapes and state set
$Q_D\sqcup(\{0,1\}\sqcup((\{0,1\}\times\{0,1\})\sqcup(Q_1\sqcup Q_2)))$. Its initial
state is the initial state of $D$, tagged in the left summand. In states tagged by
$Q_D$, it simulates $D$ on its tapes, padded by the other tapes, while recording any
output bit on the extra tape; the remaining states control the subsequent scan and
selection phases, then simulate the selected branch machine on its tape block, with the
transition table on branch states given by the branch union started in its false
setting.
\end{definition}

\begin{definition}[Timed branch configuration]
\lean{Turing.FinTM.f2_timedBranchCfg}
\label{Turing.FinTM.f2_timedBranchCfg}
\leanok
\uses{Turing.Cfg, Turing.FinTM, Turing.FinTM.bufferTape, Turing.FinTM.f2_timedCondTM, Turing.FinTM.leftCfg, Turing.FinTM.rightCfg}
Given finite Boolean Turing machines $D$, $M_1$, and $M_2$, an input string $x$, a
configuration $c$ of the branch union of $M_1$ and $M_2$ on $x$, work-tape contents $T$
for the $D$ tapes, work-head positions $h$ for those tapes, and a bit $b$, define a
configuration of the timed conditional machine on $x$. It places $c$ in the branch
machine's tape and state blocks, places $T$ and $h$ in the $D$ tape block, and
initializes the extra tape with the one-bit string $b$ and its head at position zero.
\end{definition}

\begin{definition}[Timed control configuration]
\lean{Turing.FinTM.f2_timedControlCfg}
\label{Turing.FinTM.f2_timedControlCfg}
\leanok
\uses{Turing.Cfg, Turing.FinTM, Turing.FinTM.f2_timedCondTM, Turing.FinTM.leftCfg, Turing.captureCfg}
Given finite Boolean Turing machines $D$, $M_1$, and $M_2$, an input string $x$, and a
configuration $c$ of $D$ on $x$, define a configuration of the timed conditional machine
on $x$. It embeds the state and configuration of $D$ into the machine's $D$ block, pads
the work tapes belonging to $M_1$ and $M_2$ with blank tapes whose heads are at zero,
and captures the output of $c$ on the extra tape; the resulting control state is the
designated false control state.
\end{definition}

\begin{lemma}[Captured run before halting]
\lean{Turing.FinTM.f2_timed_capture}
\label{Turing.FinTM.f2_timed_capture}
\leanok
\uses{Turing.Cfg, Turing.Cfg.Halted, Turing.FinTM, Turing.FinTM.f2_timedCondTM, Turing.FinTM.f2_timedControlCfg, Turing.FinTM.f2_timedPadTM, Turing.FinTM.leftCfg, Turing.FinTM.leftCfg_run, Turing.MultiTapeTM.runFrom, Turing.capture_run}
\difficulty{4}
Let $D,M_1,M_2$ be finite Turing machines over the Boolean alphabet, let $c$ be a
configuration of $D$ on input $x$, and let $t\in\mathbb{N}$. If the run of $D$ from $c$
is not halted at any step $s<t$, then running the timed conditional machine for $t$
steps from the control configuration associated with $c$ yields the control
configuration associated with the configuration reached by $D$ after $t$ steps.
\end{lemma}

\begin{lemma}[Initial control configuration]
\lean{Turing.FinTM.f2_timed_control_init}
\label{Turing.FinTM.f2_timed_control_init}
\leanok
\uses{Turing.Cfg.init, Turing.FinTM, Turing.FinTM.f2_timedCondTM, Turing.FinTM.f2_timedControlCfg, Turing.FinTM.leftCfg, Turing.MultiTapeTM.initCfg, Turing.captureCfg}
\difficulty{4}
For finite Turing machines $D,M_1,M_2$ over the Boolean alphabet and any Boolean input
word $x$, the initial configuration of the timed conditional machine on $x$ is the
control configuration associated with the initial configuration of $D$ on $x$.
\end{lemma}

\begin{lemma}[Timed Simulation of a Selected Branch]
\lean{Turing.FinTM.f2_timed_branch_run}
\label{Turing.FinTM.f2_timed_branch_run}
\leanok
\uses{Turing.Action.apply, Turing.Cfg, Turing.Cfg.inputSymbol, Turing.Cfg.workTapeSymbols, Turing.FinTM, Turing.FinTM.branchTM, Turing.FinTM.bufferTape, Turing.FinTM.f2_timedBranchCfg, Turing.FinTM.f2_timedCondTM, Turing.FinTM.leftAction, Turing.FinTM.leftCfg, Turing.FinTM.leftCfg_apply, Turing.FinTM.rightAction, Turing.FinTM.rightCfg, Turing.FinTM.rightCfg_apply, Turing.MultiTapeTM.runFrom, Turing.MultiTapeTM.runFrom_comm_of_step, Turing.MultiTapeTM.step}
\difficulty{5}
Let $D,M_1,M_2$ be finite-state Boolean Turing machines, let $x$ be a Boolean input, and
let $c$ be a configuration of the branch machine combining $M_1$ and $M_2$ on $x$. For
any contents and head positions on $D$’s work tapes, any Boolean $b$, and any
$t\in\mathbb N$, running the timed conditional machine for $t$ steps from the
configuration embedding $c$, those tapes and head positions, and the bit $b$ yields
exactly the embedded configuration obtained by running the branch machine, started in
branch $b$, for $t$ steps.
\end{lemma}

\begin{definition}[Timed ready configuration]
\lean{Turing.FinTM.f2_timedReadyCfg}
\label{Turing.FinTM.f2_timedReadyCfg}
\leanok
\uses{Turing.Cfg, Turing.FinTM, Turing.FinTM.branchTM, Turing.FinTM.f2_timedBranchCfg, Turing.FinTM.f2_timedCondTM, Turing.MultiTapeTM.initCfg}
Given finite Boolean Turing machines $D$, $M_1$, and $M_2$, an input string $x$, a
configuration $c$ of $D$ on $x$, and a bit $b$, define a configuration of the timed
conditional machine on $x$. Its branch block is initialized from the initial
configuration of the branch union of $M_1$ and $M_2$ selected by $b$, and its $D$-tape
contents and head positions come from $c$; the extra tape records $b$. Its state is the
control state carrying $(b,\mathrm{false})$, and its input-head position is that of $c$.
\end{definition}

\begin{lemma}[Reading the selected branch]
\lean{Turing.FinTM.f2_timed_read}
\label{Turing.FinTM.f2_timed_read}
\leanok
\uses{Turing.Action, Turing.Action.apply, Turing.Cfg, Turing.Cfg.workTapeSymbols, Turing.FinTM, Turing.FinTM.branchTM, Turing.FinTM.bufferTape, Turing.FinTM.controlAction, Turing.FinTM.controlAction_apply, Turing.FinTM.f2_timedBranchCfg, Turing.FinTM.f2_timedCondTM, Turing.FinTM.f2_timedControlCfg, Turing.FinTM.f2_timedReadyCfg, Turing.FinTM.leftCfg, Turing.FinTM.rightCfg, Turing.MultiTapeTM.initCfg, Turing.MultiTapeTM.runFrom, Turing.MultiTapeTM.step, Turing.captureCfg, Turing.moveInputPos_zero}
\difficulty{6}
Let $D,M_1,M_2$ be finite Turing machines over the Boolean alphabet, and let $c$ be a
configuration of $D$ on input $x$. If $c$ is halted and its output is the one-bit word
$[b]$, then after two steps the timed conditional machine, started from the control
configuration associated with $c$, is in the ready configuration corresponding to the
branch selected by $b$, with $c$'s input-head position and work tapes and heads
retained.
\end{lemma}

\begin{lemma}[Timed Branch Initialization]
\lean{Turing.FinTM.f2_timed_start}
\label{Turing.FinTM.f2_timed_start}
\leanok
\uses{Turing.Cfg.init, Turing.FinTM, Turing.FinTM.ComputesInTime, Turing.FinTM.ComputesInTime.output_unique, Turing.FinTM.branchTM, Turing.FinTM.computesInTime_iff, Turing.FinTM.f2_timedBranchCfg, Turing.FinTM.f2_timedCondTM, Turing.FinTM.f2_timedControlCfg, Turing.FinTM.f2_timedReadyCfg, Turing.FinTM.f2_timed_capture, Turing.FinTM.f2_timed_control_init, Turing.FinTM.f2_timed_read, Turing.FinTM.timed_rewind, Turing.MultiTapeTM.initCfg, Turing.MultiTapeTM.runFrom, Turing.MultiTapeTM.runFrom_add, Turing.MultiTapeTM.timed_input_bound}
\difficulty{5}
Let $D$, $M_1$, and $M_2$ be finite Turing machines over the Boolean alphabet, let $x$
be a Boolean input, let $b$ be a Boolean, and let $T$ be a natural number. If $D$ halts
on $x$ within $T$ steps with output the one-symbol word $[b]$, then there are a natural
number $a \le 2T+5$, work-tape contents for the $D$-tape block, and head positions for
that block such that after $a$ steps the timed conditional machine reaches the
configuration obtained by placing the initial configuration of the branch machine
selected by $b$ in its branch-tape block, with those contents and head positions on the
$D$-tape block and with $b$ stored on the extra tape.
\end{lemma}

\begin{lemma}[Conditional computation time]
\lean{Turing.FinTM.f2_cond_time}
\label{Turing.FinTM.f2_cond_time}
\leanok
\uses{Turing.FinTM, Turing.FinTM.ComputesFunInTime, Turing.FinTM.ComputesInTime, Turing.FinTM.ComputesInTime.mono, Turing.FinTM.branchTM, Turing.FinTM.branchTM_computes, Turing.FinTM.computesInTime_iff, Turing.FinTM.f2_timedBranchCfg, Turing.FinTM.f2_timedCondTM, Turing.FinTM.f2_timed_branch_run, Turing.FinTM.f2_timed_start, Turing.FinTM.leftCfg, Turing.FinTM.rightCfg, Turing.MultiTapeTM.runFrom_add}
\difficulty{4}
Let $D,M_1,M_2$ be finite Turing machines over the Boolean alphabet, let $p$ map Boolean
words to Booleans, and let $f_1,f_2$ map Boolean words to Boolean words. Suppose $D$
computes $[p(x)]$ on every input $x$ in exactly $T_0(|x|)$ steps, and $M_1,M_2$ compute
$f_1(x),f_2(x)$ in exactly $T_1(|x|),T_2(|x|)$ steps, respectively. Then the conditional
machine computes $f_1(x)$ when $p(x)$ is true and $f_2(x)$ when $p(x)$ is false, on
every input $x$, in exactly $5\bigl(T_0(|x|)+\max(T_1(|x|),T_2(|x|))+1\bigr)$ steps.
\end{lemma}

\begin{definition}[Finite-sum index equivalence]
\lean{Turing.FinTM.f2_finSumEquiv}
\label{Turing.FinTM.f2_finSumEquiv}
\leanok
For natural numbers $a$ and $b$, the disjoint union of the index sets of sizes $a$ and
$b$ is equivalent to the index set of size $a+b$, by sending the first block to the
initial indices and the second block to the following indices.
\end{definition}

\begin{lemma}[Addition of finite sums]
\lean{Turing.FinTM.f2_sum_add}
\label{Turing.FinTM.f2_sum_add}
\leanok
\uses{Turing.FinTM.f2_finSumEquiv}
\difficulty{3}
For natural numbers $a,b$ and a function $f$ from the first $a+b$ natural-number indices
to $\mathbb{N}$, the sum of $f$ over all these indices is the sum over the first $a$
indices plus the sum over the remaining $b$ indices.
\end{lemma}

\begin{lemma}[Space used by the branch machine]
\lean{Turing.FinTM.f2_branch_space}
\label{Turing.FinTM.f2_branch_space}
\leanok
\uses{Turing.FinTM, Turing.FinTM.branchTM, Turing.FinTM.f2_sum_add, Turing.FinTM.leftCfg, Turing.FinTM.leftCfg_run, Turing.FinTM.rightCfg, Turing.FinTM.rightCfg_run, Turing.MultiTapeTM.initCfg, Turing.MultiTapeTM.spaceUsed, Turing.MultiTapeTM.spaceUsedByTape, Turing.MultiTapeTM.visitedByTapeHead}
\difficulty{5}
Let $M_1$ and $M_2$ be finite Turing machines over the Boolean alphabet, let $b$ be a
Boolean, let $x$ be a Boolean word, and let $t$ be a nonnegative integer. The total
number of work-tape cells visited by the branch machine through time $t$, starting from
its initial configuration on $x$, is the space used by $M_1$ through time $t$ plus the
number of work tapes of $M_2$ when $b$ is true; when $b$ is false, it is the space used
by $M_2$ through time $t$ plus the number of work tapes of $M_1$.
\end{lemma}

\begin{definition}[Conditional machine head positions]
\lean{Turing.FinTM.f2_condHeads}
\label{Turing.FinTM.f2_condHeads}
\leanok
\uses{Turing.FinTM}
Given finite Boolean Turing machines $D$, $M_1$, and $M_2$, head-position assignments
$d$ for the $D$ tapes and $b$ for the combined $M_1$ and $M_2$ tapes, and an integer
$z$, define head positions for all tapes of the timed conditional machine by using $d$
on the first block, $b$ on the next block, and $z$ on the final extra tape.
\end{definition}

\begin{lemma}[Conditional control head positions]
\lean{Turing.FinTM.f2_control_heads}
\label{Turing.FinTM.f2_control_heads}
\leanok
\uses{Turing.Cfg, Turing.FinTM, Turing.FinTM.f2_condHeads, Turing.FinTM.f2_timedControlCfg, Turing.FinTM.leftCfg, Turing.captureCfg}
\difficulty{3}
Let $D,M_1,M_2$ be finite Turing machines over the Boolean alphabet, let $x$ be a
Boolean word, and let $c$ be a configuration of $D$ on input $x$. The vector of
work-tape head positions in the conditional machine's control configuration constructed
from $c$ equals the vector consisting of $D$'s head positions in $c$, zero positions for
the tapes of $M_1$ and $M_2$, and the length of $c$'s output as the final coordinate.
\end{lemma}

\begin{lemma}[Branch configuration head positions]
\lean{Turing.FinTM.f2_branch_heads}
\label{Turing.FinTM.f2_branch_heads}
\leanok
\uses{Turing.Cfg, Turing.FinTM, Turing.FinTM.f2_condHeads, Turing.FinTM.f2_timedBranchCfg}
\difficulty{1}
Let $D$, $M_1$, and $M_2$ be finite Turing machines over the Boolean alphabet, let $x$
be a Boolean word, let $c$ be a configuration of the branch machine on $x$, and let $d$
and $h$ be integer head-position vectors for the tapes of $D$ and the branch machine,
respectively. For either Boolean $b$, the work-head positions in the timed branch
configuration formed from $c$, the supplied tapes and head positions, and the one-cell
buffer containing $b$ are exactly the concatenation of $d$, $h$, and a final position
$0$.
\end{lemma}

\begin{lemma}[Work head positions after control steps]
\lean{Turing.FinTM.f2_read_heads}
\label{Turing.FinTM.f2_read_heads}
\leanok
\uses{Turing.Action.apply, Turing.Cfg, Turing.FinTM, Turing.FinTM.f2_condHeads, Turing.FinTM.f2_control_heads, Turing.FinTM.f2_timedCondTM, Turing.FinTM.f2_timedControlCfg, Turing.FinTM.f2_timed_read, Turing.FinTM.leftCfg, Turing.MultiTapeTM.runFrom, Turing.MultiTapeTM.step, Turing.captureCfg}
\difficulty{5}
Let $D,M_1,M_2$ be finite Boolean Turing machines and let $c$ be a configuration of $D$
on an input $x$. If $c$ is halted and its output is the one-bit word $[b]$, then for
every $u\le 2$ there is an integer $z$ with $0\le z\le 1$ such that, after $u$ steps of
the timed conditional machine started from the timed control configuration built from
$c$, the work-tape head positions agree with the following: the first block has the head
positions of $c$, the next block has all head positions zero, and the final head
position is $z$.
\end{lemma}

\begin{lemma}[Conditional machine head ledger]
\lean{Turing.FinTM.f2_cond_ledger}
\label{Turing.FinTM.f2_cond_ledger}
\leanok
\uses{Turing.FinTM, Turing.FinTM.ComputesInTime, Turing.FinTM.ComputesInTime.output_unique, Turing.FinTM.branchTM, Turing.FinTM.computesInTime_iff, Turing.FinTM.f2_branch_heads, Turing.FinTM.f2_condHeads, Turing.FinTM.f2_control_heads, Turing.FinTM.f2_read_heads, Turing.FinTM.f2_timedBranchCfg, Turing.FinTM.f2_timedCondTM, Turing.FinTM.f2_timedControlCfg, Turing.FinTM.f2_timedReadyCfg, Turing.FinTM.f2_timed_branch_run, Turing.FinTM.f2_timed_capture, Turing.FinTM.f2_timed_control_init, Turing.FinTM.f2_timed_read, Turing.FinTM.rewind_heads, Turing.MultiTapeTM.initCfg, Turing.MultiTapeTM.output_prefix, Turing.MultiTapeTM.runFrom, Turing.MultiTapeTM.runFrom_add}
\difficulty{6}
Let $D,M_1,M_2$ be finite Turing machines over the Boolean alphabet, let $x$ be a
Boolean word, let $b$ be a Boolean, and let $T$ be a natural number. Suppose $D$
computes the one-symbol output $[b]$ on input $x$ in exactly $T$ steps. For every
natural number $u$, there are $v \le T$, $w \le u$, and an integer $z$ with $0 \le z \le
1$ such that the vector of work-tape head positions of the conditional machine after $u$
steps equals the vector formed from the head positions of $D$ after $v$ steps, those of
the branch machine selected by $b$ after $w$ steps, and the final coordinate $z$.
\end{lemma}

\begin{lemma}[Conditional machine space bound]
\lean{Turing.FinTM.f2_cond_space}
\label{Turing.FinTM.f2_cond_space}
\leanok
\uses{Turing.FinTM, Turing.FinTM.ComputesInTime, Turing.FinTM.branchTM, Turing.FinTM.f2_condHeads, Turing.FinTM.f2_cond_ledger, Turing.FinTM.f2_sum_add, Turing.FinTM.f2_timedCondTM, Turing.MultiTapeTM.initCfg, Turing.MultiTapeTM.runFrom, Turing.MultiTapeTM.spaceUsed, Turing.MultiTapeTM.spaceUsedByTape, Turing.MultiTapeTM.visitedByTapeHead}
\difficulty{5}
Let $D,M_1,M_2$ be finite Turing machines over the Boolean alphabet, let $x$ be a
Boolean word, let $b$ be a Boolean, and let $T,t$ be natural numbers. Suppose $D$
computes the one-symbol output $[b]$ on input $x$ in exactly $T$ steps. The space used
by the conditional machine on input $x$ through step $t$ is at most the sum of the space
used by $D$ through step $T$, the space used by the branch machine selected by $b$
through step $t$, and $2$.
\end{lemma}

\begin{theorem}[Space-bounded timed conditional composition]
\lean{Turing.FinTM.computesFunInTime_cond_spaceUsed}
\label{Turing.FinTM.computesFunInTime_cond_spaceUsed}
\leanok
\uses{Turing.FinTM, Turing.FinTM.ComputesFunInTime, Turing.FinTM.ComputesInTime.mono, Turing.FinTM.f2_branch_space, Turing.FinTM.f2_cond_space, Turing.FinTM.f2_cond_time, Turing.FinTM.f2_timedCondTM, Turing.MultiTapeTM.initCfg, Turing.MultiTapeTM.spaceUsed}
\difficulty{4}
Let $D$, $M_1$, and $M_2$ be finite Turing machines over the Boolean alphabet; let $p$
be a predicate on Boolean words, let $f_1,f_2$ map Boolean words to Boolean words, let
$T_0,T_1,T_2$ map natural numbers to natural numbers, and let $s_D,s_1,s_2$ map natural
numbers to natural numbers. Suppose $D$ computes $[p(x)]$ in time $T_0(|x|)$, $M_i$
computes $f_i(x)$ in time $T_i(|x|)$ for $i=1,2$, and for every Boolean word $x$ and
every $t\in\mathbb{N}$, the space used by each machine $D,M_1,M_2$ by time $t$ is at
most $s_D(|x|),s_1(|x|),s_2(|x|)$, respectively. Then there are a finite Turing machine
$M$ over the Boolean alphabet and $c\in\mathbb{N}$ such that $M$ computes the function
equal to $f_1(x)$ when $p(x)$ is true and $f_2(x)$ otherwise in time
$c\bigl(T_0(|x|)+\max(T_1(|x|),T_2(|x|))+1\bigr)$, and for every $x$ and every
$t\in\mathbb{N}$ its space use by time $t$ is at most
$s_D(|x|)+\max(s_1(|x|),s_2(|x|))+c$.
\end{theorem}

\begin{lemma}[Space bounds the work head radius]
\lean{Turing.FinTM.a2_source_radius}
\label{Turing.FinTM.a2_source_radius}
\leanok
\uses{Turing.Cfg, Turing.MultiTapeTM, Turing.MultiTapeTM.runFrom, Turing.MultiTapeTM.runFrom_succ_eq_step', Turing.MultiTapeTM.runFrom_zero, Turing.MultiTapeTM.spaceUsed, Turing.MultiTapeTM.spaceUsedByTape_le_spaceUsed, Turing.MultiTapeTM.visitedByTapeHead, Turing.MultiTapeTM.workTapePos_step_le}
\difficulty{6}
Let $M$ be a $k$-tape machine over the Boolean alphabet and let $c$ be a configuration
on a Boolean input. Suppose every work-tape head starts at position $0$, and the total
number of distinct work-tape cells visited through time $t$ is at most $B$. Then, for
every work tape $i$, its head position at time $t$ lies in the interval $[-B,B]$.
\end{lemma}

\begin{definition}[Bounded work-head positions]
\lean{Turing.FinTM.a2_heads}
\label{Turing.FinTM.a2_heads}
\leanok
\uses{Turing.Cfg}
For a configuration $c$ of a Boolean Turing machine on an input string and a natural
number $B$, say that its work-head positions are bounded by $B$ when every work-head
position lies in the interval $[-B,B]$.
\end{definition}

\begin{lemma}[Monotonicity of head bounds]
\lean{Turing.FinTM.a2_heads_mono}
\label{Turing.FinTM.a2_heads_mono}
\leanok
\uses{Turing.Cfg, Turing.FinTM.a2_heads}
\difficulty{3}
For a Boolean Turing-machine configuration $c$ with $k$ work tapes and natural numbers
$A\le B$, if every work-tape head of $c$ lies in $[-A,A]$, then every work-tape head of
$c$ lies in $[-B,B]$.
\end{lemma}

\begin{lemma}[Head bounds over a bounded number of steps]
\lean{Turing.FinTM.a2_heads_steps}
\label{Turing.FinTM.a2_heads_steps}
\leanok
\uses{Turing.Cfg, Turing.FinTM.a2_heads, Turing.FinTM.f2_head_steps, Turing.MultiTapeTM, Turing.MultiTapeTM.runFrom}
\difficulty{3}
Let $M$ be a Boolean Turing machine with $k$ work tapes, and let $c$ be a configuration
on input $x$. If every work-tape head of $c$ lies in $[-A,A]$ and $t\le B$, then every
work-tape head after $t$ steps lies in $[-(A+B),A+B]$.
\end{lemma}

\begin{lemma}[Combining head bounds across runs]
\lean{Turing.FinTM.a2_heads_join}
\label{Turing.FinTM.a2_heads_join}
\leanok
\uses{Turing.Cfg, Turing.FinTM.a2_heads, Turing.MultiTapeTM, Turing.MultiTapeTM.runFrom, Turing.MultiTapeTM.runFrom_add}
\difficulty{4}
Let $M$ be a Boolean Turing machine with $k$ work tapes, let $c$ be a configuration on
input $x$, and let $a,b,B$ be natural numbers. Suppose every configuration reached from
$c$ by time $a$ has all work-tape heads in $[-B,B]$, and every configuration reached in
at most $b$ steps from the configuration at time $a$ has the same bound. Then every
configuration reached from $c$ by time $a+b$ has all work-tape heads in $[-B,B]$.
\end{lemma}

\begin{lemma}[Head bounds after halting]
\lean{Turing.FinTM.a2_heads_halted}
\label{Turing.FinTM.a2_heads_halted}
\leanok
\uses{Turing.Cfg, Turing.FinTM.a2_heads, Turing.MultiTapeTM, Turing.MultiTapeTM.runFrom, Turing.MultiTapeTM.runFrom_add, Turing.MultiTapeTM.runFrom_of_halt}
\difficulty{4}
Let $M$ be a Boolean Turing machine with $k$ work tapes, let $c$ be a configuration on
input $x$, and let $a,B$ be natural numbers. If the run from $c$ has halted by time $a$,
and every configuration reached at a time $u\le a$ has all work-tape heads in $[-B,B]$,
then every configuration reached from $c$ at any natural time has all work-tape heads in
$[-B,B]$.
\end{lemma}

\begin{lemma}[Head positions in a captured call]
\lean{Turing.FinTM.a2_call_heads}
\label{Turing.FinTM.a2_call_heads}
\leanok
\uses{Turing.Cfg, Turing.FinTM, Turing.FinTM.a2_heads, Turing.FinTM.f2_loopBodyCfg, Turing.FinTM.f2_loopBodyPadded, Turing.FinTM.f2_loopCall, Turing.FinTM.leftCfg, Turing.captureCfg}
\difficulty{5}
Let $B$ be a natural number, let $\mathrm{body}$ and $F$ be finite Turing machines over
the Boolean alphabet, let $a$ be a state of $\mathrm{body}$, and let $x$ be a Boolean
word. Suppose $c$ is a configuration of $\mathrm{body}$ on $x$, $\mathrm{fuel}$ is a
configuration of $F$ on $x$, and $\mathrm{startup}$ and $\mathrm{release}$ are Boolean
values. If every work-tape head of $c$ and of $\mathrm{fuel}$ lies between $-B$ and $B$,
and the output of $c$ has length at most $B$, then every work-tape head in the
configuration formed by capturing this call also lies between $-B$ and $B$.
\end{lemma}

\begin{lemma}[Head bounds after fuel capture]
\lean{Turing.FinTM.a2_fuel_heads}
\label{Turing.FinTM.a2_fuel_heads}
\leanok
\uses{Turing.Cfg, Turing.FinTM, Turing.FinTM.a2_heads, Turing.FinTM.f2_loopFuelCaptured, Turing.FinTM.f2_loopFuelCfg, Turing.FinTM.rightCfg, Turing.captureCfg}
\difficulty{5}
Let $B$ and $F$ be finite Boolean Turing machines, let $x$ be an input word, and let $c$
be a configuration of $F$ whose work-tape heads all lie between $-N$ and $N$, inclusive.
If the output of $c$ has length at most $N$, then in the captured fuel configuration for
the constructed machine with body $B$, every work-tape head also lies between $-N$ and
$N$, inclusive.
\end{lemma}

\begin{lemma}[Simulation of a loop-body call]
\lean{Turing.FinTM.a2_call_run}
\label{Turing.FinTM.a2_call_run}
\leanok
\uses{Turing.Cfg, Turing.Cfg.Halted, Turing.FinTM, Turing.FinTM.f2_loopBodyCfg, Turing.FinTM.f2_loopBodyPadded, Turing.FinTM.f2_loopBodySource_run, Turing.FinTM.f2_loopBody_run, Turing.FinTM.f2_loopCall, Turing.FinTM.f2_loopHost, Turing.FinTM.f2_loopHost_body_capture, Turing.FinTM.leftCfg, Turing.MultiTapeTM.runFrom}
\difficulty{4}
Let $B$ and $F$ be finite Boolean Turing machines, let $a$ be a state of $B$, and choose
Boolean parameters for find mode and startup. For an input word $x$, an initial
configuration $c$ of $B$, a release flag, a time $t$, a Boolean word $w$, and a
configuration of $F$ used as fuel, assume that $c$ is not halted, every configuration
reached from $c$ in fewer than $t$ steps is not halted, and none of those configurations
is in state $a$ except possibly at time $0$ when the release flag is true. Then after
$t$ steps the constructed host machine, started in the encoded call configuration, is in
the encoded call configuration for the body machine after $t$ steps; its release flag is
the original one exactly when $t=0$, and its additional flag is true exactly when the
body configuration at time $t$ is halted. The word and fuel configuration are unchanged.
\end{lemma}

\begin{lemma}[Head bounds during a captured loop prefix]
\lean{Turing.FinTM.a2_call_prefix}
\label{Turing.FinTM.a2_call_prefix}
\leanok
\uses{Turing.Cfg, Turing.FinTM, Turing.FinTM.a2_call_heads, Turing.FinTM.a2_call_run, Turing.FinTM.a2_heads, Turing.FinTM.a2_source_radius, Turing.FinTM.f2_loopCall, Turing.FinTM.f2_loopHost, Turing.MultiTapeTM.output_prefix, Turing.MultiTapeTM.runFrom, Turing.MultiTapeTM.spaceUsed}
\difficulty{3}
Let $B,t$ be natural numbers, let $\mathrm{body}$ and $F$ be finite Turing machines over
the Boolean alphabet, let $a$ be a state of $\mathrm{body}$, and let $x$ be a Boolean
word. Take a Boolean value $\mathrm{startup}$, a configuration $c$ of $\mathrm{body}$ on
$x$, a Boolean value $\mathrm{release}$, a Boolean word $\mathrm{word}$, and a
configuration $\mathrm{fuel}$ of $F$ on $x$. Assume that $c$ is not halted, all its
work-tape heads start at position zero, the computation from $c$ is not halted at any
time $u<t$, and for every $u<t$ either $u=0$ and $\mathrm{release}$ is true or the state
at time $u$ is not $a$. Assume also that the computation from $c$ uses at most $B$
work-tape cells through every time $u\le t$, all work-tape heads of $\mathrm{fuel}$ lie
in $[-B,B]$, and the output at time $t$ has length at most $B$. Then throughout every
time $u\le t$, every work-tape head in the loop host computation started from the
captured call with startup value $\mathrm{startup}$, release value $\mathrm{release}$,
no flag, word $\mathrm{word}$, and fuel configuration $\mathrm{fuel}$ lies in $[-B,B]$.
\end{lemma}

\begin{lemma}[Preparation for bounded loop simulation]
\lean{Turing.FinTM.a2_loop_prepare}
\label{Turing.FinTM.a2_loop_prepare}
\leanok
\uses{Turing.Cfg, Turing.Cfg.Halted, Turing.Cfg.init, Turing.FinTM, Turing.FinTM.ComputesFunInTime, Turing.FinTM.a2_fuel_heads, Turing.FinTM.a2_heads, Turing.FinTM.a2_heads_join, Turing.FinTM.a2_heads_mono, Turing.FinTM.a2_heads_steps, Turing.FinTM.a2_source_radius, Turing.FinTM.bufferTape, Turing.FinTM.f2_loopControl_idle, Turing.FinTM.f2_loopFrame, Turing.FinTM.f2_loopFuelCaptured, Turing.FinTM.f2_loopFuelCaptured_frame, Turing.FinTM.f2_loopFuelCfg, Turing.FinTM.f2_loopFuel_init, Turing.FinTM.f2_loopFuel_run, Turing.FinTM.f2_loopHost, Turing.FinTM.f2_loopHost_fuel_capture, Turing.FinTM.f2_loopHost_fuel_setup, Turing.FinTM.f2_loopHost_init, Turing.FinTM.f2_loopReady, Turing.FinTM.f2_loop_first_halt, Turing.FinTM.f2_loop_fuel_width, Turing.FinTM.f2_loop_input_run_le, Turing.FinTM.rewind_heads, Turing.FinTM.rightCfg, Turing.MultiTapeTM.initCfg, Turing.MultiTapeTM.output_prefix, Turing.MultiTapeTM.runFrom, Turing.MultiTapeTM.runFrom_add, Turing.MultiTapeTM.spaceUsed}
\difficulty{6}
Let $B$ and $F$ be finite Boolean Turing machines, let $a$ be a state of $B$, and let
$R,T,S$ be functions from natural numbers to natural numbers. Suppose that $F$, on every
input $x$, halts within $T(|x|)$ steps and outputs the binary expansion of $R(|x|)$, and
that its space use at every time is at most $S(|x|)$. Then for every Boolean input $x$
there are a configuration $c$ of $F$ on $x$ and a time $t$ such that $c$ is halted with
output the binary expansion of $R(|x|)$, $t\le 5T(|x|)+7$, and the constructed loop host
machine reaches its ready configuration for $c$ at time $t$. The work-tape heads of $c$
lie in $[-S(|x|),S(|x|)]$, and at every time up to $t$ the host machine's work-tape
heads lie in the interval with radius $S(|x|)+4|\operatorname{bits}(R(|x|))|+4$.
\end{lemma}

\begin{lemma}[Bounded heads during loop startup]
\lean{Turing.FinTM.a2_loop_start_prefix}
\label{Turing.FinTM.a2_loop_start_prefix}
\leanok
\uses{Turing.Cfg, Turing.Cfg.init, Turing.Cfg.ofWords, Turing.FinTM, Turing.FinTM.a2_call_prefix, Turing.FinTM.a2_heads, Turing.FinTM.a2_heads_join, Turing.FinTM.a2_heads_mono, Turing.FinTM.a2_heads_steps, Turing.FinTM.f2_loopHost, Turing.FinTM.f2_loopReady, Turing.FinTM.f2_loopReady_call, Turing.FinTM.f2_loop_live_prefix, Turing.MultiTapeTM.initCfg, Turing.MultiTapeTM.runFrom, Turing.MultiTapeTM.spaceUsed, Turing.stateWord}
\difficulty{4}
Let $body$ and $F$ be finite Boolean Turing machines, let $anchor$ be a state of $body$,
let $x$ and $s$ be Boolean words, let $t$ and $B$ be natural numbers, and let $fuel$ be
a configuration of $F$ on input $x$. Suppose that the run of $body$ from its initial
configuration on $x$ is never in state $anchor$ at any time $u<t$, that at time $t$ it
is exactly the seam configuration with state $anchor$ and word $s$, that its space use
is at most $B$ at every time $u\le t$, and that every work-tape head of $fuel$ lies
between $-B$ and $B$. Then, for the loop-host machine in non-search mode started in its
ready configuration with fuel $fuel$, every work-tape head lies between $-(B+2)$ and
$B+2$ at each time $u\le t+2$.
\end{lemma}

\begin{lemma}[A loop round with bounded heads]
\lean{Turing.FinTM.a2_loop_round}
\label{Turing.FinTM.a2_loop_round}
\leanok
\uses{Turing.Cfg, Turing.Cfg.ofWords, Turing.FinTM, Turing.FinTM.a2_call_prefix, Turing.FinTM.a2_heads, Turing.FinTM.a2_heads_join, Turing.FinTM.a2_heads_mono, Turing.FinTM.a2_heads_steps, Turing.FinTM.f2_loopCall, Turing.FinTM.f2_loopDebit, Turing.FinTM.f2_loopHost, Turing.FinTM.f2_loopHost_accept, Turing.FinTM.f2_loopHost_anchor_return, Turing.FinTM.f2_loopHost_halt_return, Turing.FinTM.f2_loopHost_reject, Turing.FinTM.f2_loop_first_halt, Turing.FinTM.f2_loop_live_prefix, Turing.MultiTapeTM.runFrom, Turing.MultiTapeTM.runFrom_add, Turing.MultiTapeTM.spaceUsed, Turing.stateWord}
\difficulty{6}
Let $B$ and $t$ be natural numbers with $t>0$. Let $body$ and $F$ be finite Boolean
Turing machines, let $anchor$ be a state of $body$, let $x,s,next,word$ be Boolean
words, let $fuel$ be a configuration of $F$ on input $x$, and let $accepted$ be a
Boolean. Start $body$ from the seam configuration with state $anchor$ and word $s$ on
tape $0$ (all other work tapes blank). Assume that at every time $u$ with $0<u<t$ its
state is not $anchor$; if $accepted$ is true, assume that at time $t$ it has halted and
output exactly the one-bit word $[true]$, and otherwise assume that at time $t$ it is
exactly the seam configuration with state $anchor$ and word $next$. Assume also that
throughout times $u\le t$ this run uses at most $B$ work-tape cells in total, and that
every work-tape head of $fuel$ lies between $-B$ and $B$. Then, for the loop-host
machine in non-search mode, started in its ready configuration with this body seam, the
given word as payload, and fuel $fuel$, there is a time $v$ such that through every time
$u\le v$ all work-tape heads lie between $-(B+2|word|+6)$ and $B+2|word|+6$, and either
the host is halted at time $v$, or the word admits a successful one-unit debit and at
some time $v$ the host configuration is exactly its ready configuration for the seam
with state $anchor$ and word $next$, with the debited word as payload and the same fuel
configuration. Here a successful debit means the recursive operation that changes the
first true bit to false, changing each preceding false bit to true; it fails if the word
contains no true bit.
\end{lemma}

\begin{lemma}[Work head bounds across computation segments]
\lean{Turing.FinTM.a2_segments}
\label{Turing.FinTM.a2_segments}
\leanok
\uses{Turing.Cfg, Turing.FinTM.a2_heads, Turing.MultiTapeTM, Turing.MultiTapeTM.runFrom, Turing.MultiTapeTM.runFrom_add, Turing.MultiTapeTM.runFrom_of_halt}
\difficulty{5}
Let $k,N,B$ be natural numbers with $N>0$, let $Q$ be a state type, let $x$ be a Boolean
input, let $M$ be a $k$-tape machine over the Boolean alphabet, and let
$(c_i)_{i\in\mathbb{N}}$ be a sequence of configurations on $x$. Suppose that for every
$i<N$ there is a time $u$ such that every configuration reached from $c_i$ at a time
$v\le u$ has all work-tape head positions in $[-B,B]$, and at time $u$ the machine is
either halted or $i+1<N$ and its configuration equals $c_{i+1}$. Then, at every time
$t$, the run of $M$ from $c_0$ has all work-tape head positions in $[-B,B]$.
\end{lemma}

\begin{lemma}[Bounded-time loop evaluation]
\lean{Turing.FinTM.a2_loop_halted_run}
\label{Turing.FinTM.a2_loop_halted_run}
\leanok
\uses{Turing.Cfg, Turing.MultiTapeTM, Turing.MultiTapeTM.runFrom, Turing.MultiTapeTM.runFrom_add}
\difficulty{6}
Let $M$ be a Boolean Turing machine with $k$ work tapes, let $C_j$ be configurations on
a fixed input indexed by natural numbers, and let $P(j)$ be a Boolean predicate. Suppose
that at index $N$, $C_N$ is halted with output the one-element word $[\mathrm{false}]$.
Also suppose that for every $j<N$ there is a time $t\le B$ such that, if $P(j)$ holds,
the run from $C_j$ halts with output $[\mathrm{true}]$ after $t$ steps; if $P(j)$ does
not hold, that run reaches $C_{j+1}$ after $t$ steps. Then the run from $C_0$ halts
within $N B$ steps with output the one-element word whose bit is true exactly when
$P(j)$ holds for some $j<N$.
\end{lemma}

\begin{theorem}[Loop machine with space bound]
\lean{Turing.FinTM.exists_loopTM_spaceUsed}
\label{Turing.FinTM.exists_loopTM_spaceUsed}
\leanok
\uses{Turing.Cfg.ofWords, Turing.FinTM, Turing.FinTM.ComputesFunInTime, Turing.FinTM.ComputesInTime, Turing.FinTM.ComputesInTime.mono, Turing.FinTM.a2_heads, Turing.FinTM.a2_heads_join, Turing.FinTM.a2_heads_mono, Turing.FinTM.a2_loop_halted_run, Turing.FinTM.a2_loop_prepare, Turing.FinTM.a2_loop_round, Turing.FinTM.a2_loop_start_prefix, Turing.FinTM.a2_segments, Turing.FinTM.f2_loopCall, Turing.FinTM.f2_loopDebit, Turing.FinTM.f2_loopDebit_iterate_length, Turing.FinTM.f2_loopDebit_iterate_value, Turing.FinTM.f2_loopDebit_success, Turing.FinTM.f2_loopHost, Turing.FinTM.f2_loopHost_contracts, Turing.FinTM.f2_loopHost_start, Turing.FinTM.f2_loop_fuel_width, Turing.FinTM.f2_loop_orbit_inv, Turing.FinTM.f2_space_radius, Turing.MultiTapeTM.initCfg, Turing.MultiTapeTM.runFrom, Turing.MultiTapeTM.runFrom_add, Turing.MultiTapeTM.spaceUsed, Turing.stateWord}
\difficulty{7}
Let $B$ and $F$ be finite Turing machines over Boolean words, let $q$ be a state of $B$,
let $I$ be a predicate on pairs of Boolean words, let $u$ map such pairs to Boolean
words, let $A$ be Boolean-valued on such pairs, let $s_0$ map Boolean words to Boolean
words, and let $R,T,S : \\mathbb{N} \\to \\mathbb{N}$. Suppose $F$ computes the binary
digits of $R(|x|)$ within $T(|x|)$ steps, $I(x,s_0(x))$ holds, and $I(x,s)$ implies
$I(x,u(x,s))$. Suppose startup on each $x$ reaches the clean seam with word $s_0(x)$
within $T(|x|)$ steps without visiting $q$ earlier. For every invariant pair $(x,s)$,
suppose each round takes positive time at most $T(|x|)$ without an intermediate
positive-time visit to $q$; it halts with output $[\\mathrm{true}]$ if $A(x,s)$ and
otherwise reaches the clean seam with word $u(x,s)$. Assume also that for every $x,t$,
$F$ uses at most $S(|x|)$ space through time $t$; startup uses at most $S(|x|)$ space at
every time $t\\le T(|x|)$; and every invariant round uses at most $S(|x|)$ space at
every time $t\\le T(|x|)$. Then there are a finite Turing machine $E$ and a natural
number $c$ such that $E$ computes the one-bit word indicating whether
$A(x,u^i(x,s_0(x)))$ holds for some $i\\le R(|x|)$ within $c(T(n)+1)(R(n)+2)$ steps on
inputs of length $n$, and for every $x,t$ its space use through time $t$ is at most
$c(S(|x|)+T(|x|)+1)$.
\end{theorem}

\begin{definition}[Single-tape transition action]
\lean{Turing.tapeAction}
\label{Turing.tapeAction}
\leanok
\uses{Turing.Action}
Let $i$ be one of the $k$ work tapes, let $u$ be a movement of the input head, let $w$
specify an optional write value for the selected tape (a symbol, blank, or no write),
let $m$ be its head movement, and let $q$ be a state. Define an action that moves the
input head by $u$, applies $w$ and $m$ to tape $i$, leaves every other work tape
unchanged with its head stationary, emits no output, and enters state $q$.
\end{definition}

\begin{definition}[Input movement action]
\lean{Turing.inputMoveAction}
\label{Turing.inputMoveAction}
\leanok
\uses{Turing.Action}
For a machine with $k$ work tapes, a symbol alphabet, and a state set, given an
input-head movement $d\in\{-1,0,1\}$ and a state $q$, the corresponding action attempts
movement $d$ on the input head, leaves every work tape unchanged, emits no output, and
continues in state $q$.
\end{definition}

\begin{theorem}[Step from a tape movement action]
\lean{Turing.step_of_tapeAction_move}
\label{Turing.step_of_tapeAction_move}
\leanok
\uses{Turing.Action.apply, Turing.Cfg, Turing.Cfg.inputSymbol, Turing.Cfg.workTapeSymbols, Turing.MultiTapeTM, Turing.MultiTapeTM.step, Turing.moveInputPos, Turing.tapeAction}
\difficulty{4}
Let $M$ be a $k$-tape Turing machine and $c$ a configuration. Suppose $c$ has state $q$,
and the transition of $M$ at $q$ on the symbols currently read is the action that moves
the input head by $\mathrm{inp}$, moves tape $i$ by $m$, makes no writes or output, and
enters state $q'$. Then the next configuration has state $q'$, input-head position
obtained by moving $c$'s input head by $\mathrm{inp}$ with boundary clamping, and
work-head positions equal to those of $c$ except that tape $i$'s position is increased
by $m$; all other components are unchanged.
\end{theorem}

\begin{theorem}[Step from a tape write action]
\lean{Turing.step_of_tapeAction_write}
\label{Turing.step_of_tapeAction_write}
\leanok
\uses{Turing.Action.apply, Turing.Cfg, Turing.Cfg.inputSymbol, Turing.Cfg.workTapeSymbols, Turing.MultiTapeTM, Turing.MultiTapeTM.step, Turing.moveInputPos, Turing.tapeAction}
\difficulty{4}
Let $M$ be a $k$-tape Turing machine and $c$ a configuration. Suppose $c$ has state $q$,
and the transition of $M$ at $q$ on the symbols currently read is the action that moves
the input head by $\mathrm{inp}$, writes $s$ on the scanned cell of tape $i$, moves that
tape's head by $m$, makes no other writes or output, and enters state $q'$. Then the
next configuration has state $q'$, input-head position obtained by moving $c$'s input
head by $\mathrm{inp}$ with boundary clamping, tape $i$'s scanned cell set to $s$, and
tape $i$'s head position increased by $m$; all other components are unchanged.
\end{theorem}

\begin{theorem}[Step prescribed by an input-only action]
\lean{Turing.step_of_inputMoveAction}
\label{Turing.step_of_inputMoveAction}
\leanok
\uses{Turing.Action.apply, Turing.Cfg, Turing.Cfg.inputSymbol, Turing.Cfg.workTapeSymbols, Turing.MultiTapeTM, Turing.MultiTapeTM.step, Turing.inputMoveAction, Turing.moveInputPos}
\difficulty{3}
Let $M$ be a deterministic multi-tape Turing machine with $k$ work tapes, let $x$ be an
input word, and let $c$ be a configuration whose state is $q$. If the transition of $M$
at state $q$, on the symbols currently scanned on the input and work tapes, is the
action that moves the input head by a sign $m$, leaves all work tapes and their heads
unchanged, emits no output, and enters state $q'$, then one step from $c$ has state $q'$
and input-head position obtained by moving the current position by $m$ with the
prescribed boundary clamping.
\end{theorem}

\begin{definition}[Joining fields with separators]
\lean{Turing.joinFields}
\label{Turing.joinFields}
\leanok
Given a finite list of finite lists of symbols, concatenate the lists in order, placing
one blank separator between each consecutive pair of lists. Each symbol is represented
as a present symbol, and each separator as an absent symbol.
\end{definition}

\begin{definition}[Symbols on the joined tape]
\lean{Turing.fieldsTape}
\label{Turing.fieldsTape}
\leanok
\uses{Turing.joinFields}
Given a finite list of finite lists of symbols and an integer position $z$, the tape
contains the symbol at position $z$ in the joined encoding when $z\ge 0$ and that
position holds a symbol; it is blank when $z<0$, the position is outside the encoding,
or the position is a separator.
\end{definition}

\begin{definition}[Offset of a field]
\lean{Turing.fieldOffset}
\label{Turing.fieldOffset}
\leanok
Given a finite list of finite lists of symbols and a natural number $j$, the offset of
field $j$ is the sum, over the first $j$ fields (or all fields if there are fewer than
$j$), of each field's length plus one.
\end{definition}

\begin{theorem}[Joining the first two fields]
\lean{Turing.joinFields_cons_cons}
\label{Turing.joinFields_cons_cons}
\leanok
\uses{Turing.joinFields}
\difficulty{1}
Let $\Gamma$ be a type, let $u,v$ be finite lists of elements of $\Gamma$, and let $fs$
be a finite list of such lists. Joining the fields in $u::v::fs$ gives the list obtained
by mapping $u$ to optional elements, then appending a blank separator and the joined
fields in $v::fs$.
\end{theorem}

\begin{theorem}[Shift of joined fields]
\lean{Turing.fieldsTape_cons_shift}
\label{Turing.fieldsTape_cons_shift}
\leanok
\uses{Turing.fieldsTape, Turing.joinFields_cons_cons}
\difficulty{4}
For lists $u,v,fs$ of symbols and an integer $z\ge -1$, the symbol stored at position
$|u|+1+z$ in the tape formed by joining $u,v$, and the fields in $fs$ equals the symbol
stored at position $z$ in the tape formed by joining $v$ and the fields in $fs$.
\end{theorem}

\begin{theorem}[Successive field offsets]
\lean{Turing.fieldOffset_succ}
\label{Turing.fieldOffset_succ}
\leanok
\uses{Turing.fieldOffset}
\difficulty{1}
Let $fs$ be a finite list of finite fields and let $j<|fs|$. The offset after field $j$
is the offset before it plus the length of field $j$ and one.
\end{theorem}

\begin{theorem}[Offset after a leading field]
\lean{Turing.fieldOffset_cons_succ}
\label{Turing.fieldOffset_cons_succ}
\leanok
\uses{Turing.fieldOffset}
\difficulty{1}
For a list $u$ and a list of fields $fs$, and every natural number $j$, the offset of
field $j+1$ in the list with $u$ prepended to $fs$ equals $|u|+1$ plus the offset of
field $j$ in $fs$.
\end{theorem}

\begin{theorem}[Length of joined fields]
\lean{Turing.length_joinFields_add_one}
\label{Turing.length_joinFields_add_one}
\leanok
\uses{Turing.fieldOffset, Turing.fieldOffset_cons_succ, Turing.joinFields, Turing.joinFields_cons_cons}
\difficulty{4}
Let $\Gamma$ be a type and let $fs$ be a nonempty finite list of finite lists of
elements of $\Gamma$. Join the fields by inserting one separator between consecutive
fields, and let the offset after all fields be the sum, over each field, of its length
plus one. Then the length of the joined list plus one equals that offset.
\end{theorem}

\begin{theorem}[Field offset lies within the joined fields]
\lean{Turing.fieldOffset_add_length_le}
\label{Turing.fieldOffset_add_length_le}
\leanok
\uses{Turing.fieldOffset, Turing.fieldOffset_succ, Turing.joinFields, Turing.length_joinFields_add_one}
\difficulty{5}
Let $\Gamma$ be a type, let $fs$ be a finite list of finite lists of elements of
$\Gamma$, and let $j<|fs|$. The offset of field $j$ plus the length of that field is at
most the length of the list formed by joining the fields with blank separators.
\end{theorem}

\begin{theorem}[Tape contents within a field]
\lean{Turing.fieldsTape_field}
\label{Turing.fieldsTape_field}
\leanok
\uses{Turing.FinTM.bufferTape, Turing.fieldOffset, Turing.fieldOffset_cons_succ, Turing.fieldsTape, Turing.fieldsTape_cons_shift, Turing.joinFields, Turing.joinFields_cons_cons}
\difficulty{6}
Let $\Gamma$ be a type, let $fs$ be a finite list of finite words over $\Gamma$, and let
$j$ index one of its fields. For every integer $p$ with $-1 \le p \le |fs_j|$, the
joined fields tape at the position obtained by adding $p$ to the offset of field $j$
equals the buffer tape storing $fs_j$ at position $p$.
\end{theorem}

\begin{theorem}[Comparison from a general configuration]
\lean{Turing.compareTM_run_ofCfg}
\label{Turing.compareTM_run_ofCfg}
\leanok
\uses{Turing.Cfg, Turing.FinTM.bufferTape, Turing.FlagPhase, Turing.MultiTapeTM.runFrom, Turing.catalogCompareF, Turing.catalogCompareR, Turing.catalogTrace, Turing.catalog_compare_stop, Turing.catalog_compare_trace, Turing.compareTM}
\difficulty{5}
Let $Sigma$ be a type with decidable equality, let $k$ be a natural number, let $x$ be
an input word over $Sigma$, choose distinct work-tape indices $fst,sndin k$, choose
words $u,v$ over $Sigma$, and let $d$ be a configuration on input $x$ whose state is
the running state. Suppose that, for every integer $p$ with $-1\le p\le |u|$, tape $fst$
at position $d$'s head position plus $p$ contains the symbol prescribed by storing $u$
contiguously from position zero; assume the analogous condition for tape $snd$ and word
$v$. Then there exists a natural number $delta\le	ext{min}(|u|,|v|)$ such that after
$2delta+2$ steps the machine is in configuration $d$ with only its state changed to
completed with the truth value of $u=v$; before that time it has not entered either
completed state. At every time up to and including $2delta+2$, the heads of tapes $fst$
and $snd$ remain between their respective initial positions minus one and initial
positions plus $delta$, inclusive, and every other work-tape head remains at its
initial position.
\end{theorem}

\begin{definition}[Seeking phases]
\lean{Turing.SeekPhase}
\label{Turing.SeekPhase}
\leanok
The seeking procedure has four phases: the main run, the phase immediately after
crossing a blank, the backward return, and the completed phase with a live exit anchor.
\end{definition}

\begin{theorem}[Contract for a Finite Pass]
\lean{Turing.catalog_pass_contract}
\label{Turing.catalog_pass_contract}
\leanok
\uses{Turing.Cfg, Turing.MultiTapeTM, Turing.MultiTapeTM.runFrom, Turing.MultiTapeTM.step, Turing.catalogTrace, Turing.catalog_trace_run}
\difficulty{4}
Let $M$ be a $k$-tape machine, let $F:br	o C$ be a sequence of configurations, let $D$
be a configuration, and let $L$ be a natural number. Suppose each step from $F(r)$ is
$F(r+1)$ for $r<L$, the step from $F(L)$ is $D$, and $D$ is fixed by a step. If
predicates $P$ and $Q$ both hold at every $F(r)$ for $r\le L$, and $Q$ holds at $D$,
then after $L+1$ steps from $F(0)$ the machine is at $D$; $P$ holds at every time
$t<L+1$, and $Q$ holds at every time $t\le L+1$.
\end{theorem}

\begin{definition}[Catalog seek configuration update]
\lean{Turing.catalogSeekCfg}
\label{Turing.catalogSeekCfg}
\leanok
\uses{Turing.Cfg, Turing.SeekPhase}
Given a finite-tape configuration in a seeking phase, a tape index $i$, a seeking phase
$q$, and an integer $z$, form a new configuration by setting its state to $q$ and adding
$z$ to the head position on tape $i$. All other configuration data and all other
tape-head positions remain unchanged.
\end{definition}

\begin{theorem}[A Seek Step Updates the Head Position]
\lean{Turing.catalog_seek_step}
\label{Turing.catalog_seek_step}
\leanok
\uses{Turing.Cfg, Turing.Cfg.inputSymbol, Turing.Cfg.workTapeSymbols, Turing.MultiTapeTM, Turing.MultiTapeTM.step, Turing.SeekPhase, Turing.catalogSeekCfg, Turing.step_of_tapeAction_move, Turing.tapeAction}
\difficulty{3}
Let $M$ be a $k$-tape machine with the four seek phases, let $i$ be a work-tape index,
and let $d$ be a configuration. For phases $q,q'$, integer $z$, and movement
$m\in\{-1,0,1\}$, suppose the transition from the configuration obtained from $d$ by
setting the phase to $q$ and shifting head $i$ by $z$ is the action that leaves the
input head and all tape contents unchanged, moves work-tape head $i$ by $m$ (and every
other work-tape head by zero), emits nothing, and enters phase $q'$. Then one machine
step gives the configuration obtained from $d$ by setting the phase to $q'$ and shifting
head $i$ by $z+m$.
\end{theorem}

\begin{definition}[Forward tape-skipping machine]
\lean{Turing.skipTM}
\label{Turing.skipTM}
\leanok
\uses{Turing.MultiTapeTM, Turing.SeekPhase, Turing.inputMoveAction, Turing.tapeAction}
For a positive number $k$ of work tapes and a selected tape $i<k$, the machine starts in
the run phase. In that phase, if the selected tape's scanned cell contains a symbol, it
leaves the cell unchanged, moves that tape's head one step to the right, and remains in
the run phase; if the cell is blank, it makes the same tape movement and enters the done
phase. In every other phase, it leaves the work tapes unchanged, keeps the input head
still, emits no output, and remains in its current phase.
\end{definition}

\begin{theorem}[Skipping over a word]
\lean{Turing.skipTM_run_ofCfg}
\label{Turing.skipTM_run_ofCfg}
\leanok
\uses{Turing.Cfg, Turing.Cfg.workTapeSymbols, Turing.FinTM.bufferTape, Turing.MultiTapeTM.runFrom, Turing.MultiTapeTM.step, Turing.SeekPhase, Turing.catalogSeekCfg, Turing.catalog_pass_contract, Turing.catalog_seek_step, Turing.moveInputPos_zero, Turing.skipTM, Turing.step_of_inputMoveAction}
\difficulty{5}
Let $k$ be a nonnegative integer, let $i<k$ index a work tape, let $x$ be an input word,
and let $w$ be a finite word. Consider a configuration $d$ in the running phase such
that, for every integer $p$ with $0\le p\le |w|$, tape $i$ at position $d_i+p$ contains
the symbol of $w$ at position $p$ (and is blank where that position has no symbol).
After $|w|+1$ steps, the skipping machine is in the done phase with tape head $i$ at
$d_i+|w|+1$; it is not in the done phase at any earlier time. At every time up to and
including $|w|+1$, the head on tape $i$ lies between its initial position and
$d_i+|w|+1$, inclusive, and every other work-tape head remains at its initial position.
\end{theorem}

\begin{definition}[Backward tape-skipping machine]
\lean{Turing.skipBackTM}
\label{Turing.skipBackTM}
\leanok
\uses{Turing.MultiTapeTM, Turing.SeekPhase, Turing.inputMoveAction, Turing.tapeAction}
For a positive number $k$ of work tapes and a selected tape $i<k$, the machine starts in
the run phase. In the run phase it leaves the selected tape unchanged, moves its head
one step to the left, and enters the gap phase; in the gap phase it makes the same tape
movement and enters the back phase. In the back phase, if the selected tape's scanned
cell contains a symbol, it moves that tape's head one step to the left and remains in
the back phase; if the cell is blank, it moves the head one step to the right and enters
the done phase. In the done phase, it leaves the work tapes unchanged, keeps the input
head still, emits no output, and remains in the done phase.
\end{definition}

\begin{theorem}[Skipping backward over a word]
\lean{Turing.skipBackTM_run_ofCfg}
\label{Turing.skipBackTM_run_ofCfg}
\leanok
\uses{Turing.Cfg, Turing.Cfg.workTapeSymbols, Turing.FinTM.bufferTape, Turing.MultiTapeTM.runFrom, Turing.MultiTapeTM.step, Turing.SeekPhase, Turing.catalogSeekCfg, Turing.catalog_pass_contract, Turing.catalog_seek_step, Turing.moveInputPos_zero, Turing.skipBackTM, Turing.step_of_inputMoveAction}
\difficulty{6}
Let $k$ be a nonnegative integer, let $i<k$ index a work tape, let $x$ be an input word,
and let $u$ be a finite word. Consider a configuration $d$ in the running phase such
that, for every integer $p$ with $-1\le p\le |u|$, tape $i$ at position $d_i-|u|-1+p$
contains the symbol of $u$ at position $p$ (and is blank where that position has no
symbol). After $|u|+3$ steps, the skipping machine is in the done phase with tape head
$i$ at $d_i-|u|-1$; it is not in the done phase at any earlier time. At every time up to
and including $|u|+3$, the head on tape $i$ lies between $d_i-|u|-2$ and its initial
position, inclusive, and every other work-tape head remains at its initial position.
\end{theorem}

\begin{theorem}[No adjacent blank cells between fields]
\lean{Turing.fieldsTape_no_double_blank}
\label{Turing.fieldsTape_no_double_blank}
\leanok
\uses{Turing.FinTM.bufferTape, Turing.fieldOffset, Turing.fieldOffset_cons_succ, Turing.fieldsTape, Turing.fieldsTape_cons_shift, Turing.fieldsTape_field}
\difficulty{6}
Let $\Gamma$ be a type and let $ws$ be a finite list of nonempty finite words over
$\Gamma$. For every integer $z\ge 0$ such that $z+1$ is strictly less than the total
length of the fields together with one separator per field, the joined fields tape is
not blank at both positions $z$ and $z+1$.
\end{theorem}

\begin{definition}[Rightward search for a tape boundary]
\lean{Turing.seekEndTM}
\label{Turing.seekEndTM}
\leanok
\uses{Turing.MultiTapeTM, Turing.SeekPhase, Turing.inputMoveAction, Turing.tapeAction}
For an alphabet of symbols, a number $k$ of work tapes, and a selected tape $i$ among
them, define a machine with four phases: run, gap, back, and done. It starts in run;
while scanning tape $i$, it moves right and remains in run on a nonblank symbol, and
moves right into gap on a blank. In gap it moves right and returns to run on a nonblank
symbol; on a blank it enters done without moving the input head. In back it stays in
place on tape $i$ and enters done, while in done it stays in place; these transitions
make no writes or outputs.
\end{definition}

\begin{theorem}[Seeking the end of a field sequence]
\lean{Turing.seekEndTM_run_ofCfg}
\label{Turing.seekEndTM_run_ofCfg}
\leanok
\uses{Turing.Cfg, Turing.Cfg.inputSymbol, Turing.Cfg.workTapeSymbols, Turing.MultiTapeTM.runFrom, Turing.MultiTapeTM.step, Turing.SeekPhase, Turing.catalogSeekCfg, Turing.catalog_pass_contract, Turing.catalog_seek_step, Turing.fieldOffset, Turing.fieldsTape, Turing.fieldsTape_no_double_blank, Turing.joinFields, Turing.length_joinFields_add_one, Turing.moveInputPos_zero, Turing.seekEndTM, Turing.step_of_inputMoveAction, Turing.tapeAction}
\difficulty{6}
Let $k$ be a nonnegative integer, let $i<k$ index a work tape, and let $x$ be an input
word. Let $ws$ be a nonempty finite list of nonempty words, and let $L$ be the sum, over
the words in $ws$, of their lengths plus one. Consider a configuration $d$ in the
running phase, and suppose that for every integer $p$ with $0\le p\le L$, tape $i$ at
position $d_i+p$ contains the symbol at position $p$ of the concatenation of the words
in $ws$, with a blank separator between consecutive words and no symbol outside that
concatenation. After $L+1$ steps, the seeking machine is in the done phase with tape
head $i$ at $d_i+L$; it is not in the done phase at any earlier time. At every time up
to and including $L+1$, the head on tape $i$ lies between its initial position and
$d_i+L$, inclusive, and every other work-tape head remains at its initial position.
\end{theorem}

\begin{definition}[Leftward search for a tape boundary]
\lean{Turing.seekStartTM}
\label{Turing.seekStartTM}
\leanok
\uses{Turing.MultiTapeTM, Turing.SeekPhase, Turing.inputMoveAction, Turing.tapeAction}
For an alphabet of symbols, a number $k$ of work tapes, and a selected tape $i$ among
them, define a machine with phases run, gap, back, and done. It starts in run; while
scanning tape $i$, it moves left and remains in run on a nonblank symbol, and moves left
into gap on a blank. In gap it moves left and returns to run on a nonblank symbol; on a
blank it enters back without moving the input head. In back it moves right on tape $i$
and enters done; in done it stays in place. These transitions make no writes or outputs.
\end{definition}

\begin{theorem}[Seeking the start of a field sequence]
\lean{Turing.seekStartTM_run_ofCfg}
\label{Turing.seekStartTM_run_ofCfg}
\leanok
\uses{Turing.Cfg, Turing.Cfg.inputSymbol, Turing.Cfg.workTapeSymbols, Turing.FinTM.bufferTape, Turing.MultiTapeTM.runFrom, Turing.MultiTapeTM.step, Turing.SeekPhase, Turing.catalogSeekCfg, Turing.catalog_pass_contract, Turing.catalog_seek_step, Turing.fieldOffset, Turing.fieldsTape, Turing.fieldsTape_field, Turing.fieldsTape_no_double_blank, Turing.moveInputPos_zero, Turing.seekStartTM, Turing.step_of_inputMoveAction, Turing.tapeAction}
\difficulty{7}
Let $k$ be a nonnegative integer, let $i<k$ index a work tape, and let $x$ be an input
word. Let $ws$ be a finite list of nonempty words, and let $L$ be the sum of their
lengths plus one per word. Consider a configuration $d$ in the running phase whose
scanned cell on tape $i$ is nonblank. Suppose that for every integer $p$ with $-2\le
p<L$, tape $i$ at position $d_i-L+p$ contains the symbol at position $p$ of the
concatenation of the words in $ws$, with blank separators between consecutive words and
blanks outside the concatenation. After $L+4$ steps, the seeking machine is in the done
phase with tape head $i$ at $d_i-L$; it is not in the done phase at any earlier time. At
every time up to and including $L+4$, the head on tape $i$ lies between $d_i-L-2$ and
its initial position, inclusive, and every other work-tape head remains at its initial
position.
\end{theorem}

\begin{definition}[Mealy transducer run]
\lean{Turing.mealyRun}
\label{Turing.mealyRun}
\leanok
Given sets $A$ and $\Gamma$ and a transition function $\delta:A\times\Gamma\to
A\times\Gamma$, define the run on an initial state and a finite input list recursively.
On the empty list it returns the initial state and the empty output list; on a nonempty
list it applies $\delta$ to the current state and first symbol, then runs on the
remaining list from the resulting state, returning the final state and the list of
emitted symbols in their original order.
\end{definition}

\begin{theorem}[Concatenation law for a Mealy run]
\lean{Turing.mealyRun_append}
\label{Turing.mealyRun_append}
\leanok
\uses{Turing.mealyRun}
\difficulty{4}
Let $A$ and $\Gamma$ be types, let $\delta:A\times\Gamma\to A\times\Gamma$ be a
transition rule, let $a\in A$, and let $l_1,l_2$ be finite strings over $\Gamma$.
Running the rule on the concatenation $l_1l_2$ yields the same final state and output as
first running it on $l_1$, then running it on $l_2$ from the state thereby obtained; the
two output strings concatenate in that order.
\end{theorem}

\begin{theorem}[Induction for Mealy runs]
\lean{Turing.mealyRun_induction}
\label{Turing.mealyRun_induction}
\leanok
\uses{Turing.mealyRun}
\difficulty{4}
Let $A$ be a state space, let $\Gamma$ be an alphabet, and let $\delta:A\times\Gamma\to
A\times\Gamma$ be a transition and output rule. Suppose a property $P(a,w,r)$ holds for
every state $a$ when the input list is empty and the output pair is $(a,[])$. Suppose
also that whenever $P(\delta(a,\gamma)_1,w,r)$ holds, it holds for the state $a$, the
input list $\gamma::w$, and output pair $(r_1,\delta(a,\gamma)_2::r_2)$. Then
$P(a,w,\operatorname{run}(a,w))$ holds for every state $a$ and finite input list $w$,
where the run repeatedly applies $\delta$ and records its outputs.
\end{theorem}

\begin{theorem}[Length preservation of a Mealy run]
\lean{Turing.mealyRun_length}
\label{Turing.mealyRun_length}
\leanok
\uses{Turing.mealyRun, Turing.mealyRun_induction}
\difficulty{3}
Let $A$ be a state space, let $\Gamma$ be an alphabet, let $\delta:A\times\Gamma\to
A\times\Gamma$ be a transition and output rule, and let $a\in A$. For every finite input
list $w$, the output list produced by running the machine from $a$ on $w$ has length
equal to the length of $w$.
\end{theorem}

\begin{theorem}[State and output at a run position]
\lean{Turing.mealyRun_next}
\label{Turing.mealyRun_next}
\leanok
\uses{Turing.mealyRun, Turing.mealyRun_append, Turing.mealyRun_length}
\difficulty{5}
Let $A$ be a state space, let $\Gamma$ be an alphabet, and let $\delta:A\times\Gamma\to
A\times\Gamma$ be a transition and output rule. For an initial state $a$, a finite input
list $w$, and a natural number $r<|w|$, the state after processing the first $r+1$
symbols is the state obtained by applying $\delta$ to the state after the first $r$
symbols and the symbol $w_r$. The $r$th symbol in the full output list is the output
component of that same transition.
\end{theorem}

\begin{definition}[Forward and backward scan result]
\lean{Turing.scanResult}
\label{Turing.scanResult}
\leanok
\uses{Turing.mealyRun}
Given sets $A$ and $\Gamma$, two transition functions
$\delta_f,\delta_b:A\times\Gamma\to A\times\Gamma$, an initial state $a_0$, and a finite
word $w$ over $\Gamma$, first run $\delta_f$ from $a_0$ across $w$. Then run $\delta_b$
from the resulting state across the reverse of the first output word, and reverse the
second output word. The result is the final state together with that reversed output.
\end{definition}

\begin{definition}[Three phases of a scan]
\lean{Turing.ScanPhase}
\label{Turing.ScanPhase}
\leanok
For a type $A$, define a phase type whose elements are a forward phase carrying a state
in $A$, a backward phase carrying a state in $A$, or a done phase carrying a state in
$A$.
\end{definition}

\begin{definition}[Two-pass scan machine]
\lean{Turing.scanTM}
\label{Turing.scanTM}
\leanok
\uses{Turing.MultiTapeTM, Turing.ScanPhase, Turing.inputMoveAction, Turing.tapeAction}
Let $\Gamma$ be a symbol alphabet, $A$ a state type, $k$ a natural number, and $i$ one
of the $k$ work tapes. Given forward and backward transition functions
$\delta_f,\delta_b:A\times\Gamma\to A\times\Gamma$ and an initial state $a_0$, define a
machine that starts in the forward phase carrying $a_0$. On each nonblank symbol of tape
$i$, the forward phase applies $\delta_f$, writes its output symbol on that tape, moves
right, and carries the new state. At a blank it moves left into the backward phase,
carrying the current state. The backward phase applies $\delta_b$ to each nonblank
symbol, writes the resulting symbol, moves left, and carries the new state; at a blank
it moves right into the done phase. In the done phase it stays in place. The input head
does not move, and no output symbols are emitted.
\end{definition}

\begin{theorem}[Configuration trace of a two-pass scan]
\lean{Turing.scanTM_run_ofCfg}
\label{Turing.scanTM_run_ofCfg}
\leanok
\uses{Turing.Cfg, Turing.Cfg.workTapeSymbols, Turing.FinTM.bufferTape, Turing.FinTM.bufferTape_nat, Turing.MultiTapeTM.runFrom, Turing.MultiTapeTM.step, Turing.ScanPhase, Turing.catalogTape, Turing.catalogTrace, Turing.catalog_trace_run, Turing.catalog_write_middle, Turing.catalog_write_take, Turing.mealyRun, Turing.mealyRun_append, Turing.mealyRun_length, Turing.mealyRun_next, Turing.moveInputPos_zero, Turing.scanResult, Turing.scanTM, Turing.step_of_inputMoveAction, Turing.step_of_tapeAction_move, Turing.step_of_tapeAction_write}
\difficulty{7}
Let $k\in\mathbb N$, let $\Gamma$ be a symbol type, let $A$ be a state type, let $x,w$
be finite words over $\Gamma$, choose a work-tape index $i<k$, two transducers
$\delta_f,\delta_b:A\times\Gamma\to A\times\Gamma$, an initial state $a_0$, and a
configuration $d$ of the scan machine on input $x$. Suppose $d$ is in forward state
$a_0$ and, on tape $i$, every cell at offset $p$ from its head contains the
corresponding symbol of $w$ for $-1\le p\le |w|$, with cells outside the word blank.
After $2|w|+2$ steps, the machine is in the done state carrying the final state of the
two-pass transducer scan, and tape $i$ contains that scan's output in the $|w|$ cells
starting at its original head position; every other tape cell retains its original
content. At every earlier time $t<2|w|+2$ the state is not done, and throughout these
steps tape $i$'s head stays between one cell left of its initial position and the cell
$|w|$ to its right, while all other tape heads remain fixed.
\end{theorem}

\begin{definition}[Conditional scan with carry]
\lean{Turing.scanCarry}
\label{Turing.scanCarry}
\leanok
Given a set $\Gamma$, a predicate
$\operatorname{top}:\Gamma\to\{\mathrm{false},\mathrm{true}\}$, and a function
$\operatorname{next}:\Gamma\to\Gamma$, define a transition on a Boolean carry $c$ and
symbol $\gamma$. If $c$ is true, it returns
$(\operatorname{top}(\gamma),\operatorname{next}(\gamma))$; if $c$ is false, it returns
$(\mathrm{false},\gamma)$.
\end{definition}

\begin{theorem}[Run with an identity transition]
\lean{Turing.mealyRun_keep}
\label{Turing.mealyRun_keep}
\leanok
\uses{Turing.mealyRun, Turing.mealyRun_induction, Turing.scanKeep}
\difficulty{4}
Let $A$ be a state space, let $\Gamma$ be an alphabet, and let $a\in A$ and $w$ be a
finite list over $\Gamma$. If each step leaves the state and symbol unchanged, then
running the machine from $a$ on $w$ returns the pair $(a,w)$.
\end{theorem}

\begin{theorem}[A second keep scan changes nothing]
\lean{Turing.scanResult_keep}
\label{Turing.scanResult_keep}
\leanok
\uses{Turing.mealyRun, Turing.mealyRun_keep, Turing.scanKeep, Turing.scanResult}
\difficulty{1}
Let $A$ be a state set, let $\Gamma$ be a symbol set, let $\delta_f:A\to\Gamma\to
A\times\Gamma$ be a transition that reads and possibly changes a symbol, let $a_0\in A$,
and let $w$ be a finite word over $\Gamma$. A two-pass scan first runs this transition
from left to right over $w$, then runs a transition that leaves both state and symbols
unchanged from right to left over the first output in reverse. The resulting state and
word equal those from the first pass alone.
\end{theorem}

\begin{theorem}[Reversing a left-shifted list update]
\lean{Turing.shiftLeft_list}
\label{Turing.shiftLeft_list}
\leanok
\difficulty{4}
Let $B$ and $\Gamma$ be types, let $g:\Gamma\to B$ and $s:\Gamma\to B\to\Gamma$ be
functions, let $b_0\in B$, and let $w$ be a finite list in $\Gamma$. Updating the
reversed list $w$ elementwise with the corresponding values of $g$, using $b_0$ at the
added final position, and then reversing the result gives the same list as updating $w$
elementwise with the values of $g$ on its tail, using $b_0$ at the final position.
\end{theorem}

\begin{theorem}[A cleared carry leaves a word unchanged]
\lean{Turing.mealyRun_carry_clear}
\label{Turing.mealyRun_carry_clear}
\leanok
\uses{Turing.mealyRun, Turing.scanCarry}
\difficulty{4}
Let $\Gamma$ be a type, let $\mathrm{top}:\Gamma\to\{\mathrm{false},\mathrm{true}\}$ and
$\mathrm{next}:\Gamma\to\Gamma$ be functions, and let $w$ be a finite string over
$\Gamma$. Scanning $w$ with the carry rule that, when the carry is true, outputs
$\mathrm{top}(\gamma)$ and replaces $\gamma$ by $\mathrm{next}(\gamma)$, and when the
carry is false, outputs false and leaves $\gamma$ unchanged, starting with carry false
returns carry false and leaves the entire string unchanged.
\end{theorem}

\begin{theorem}[Fold scan output]
\lean{Turing.scanResult_fold}
\label{Turing.scanResult_fold}
\leanok
\uses{Turing.mealyRun_induction, Turing.scanFold, Turing.scanKeep, Turing.scanResult, Turing.scanResult_keep}
\difficulty{4}
Let $A$ be a state set, let $\Gamma$ be a symbol set, let $\delta:A\to\Gamma\to A$, let
$a_0\in A$, and let $w$ be a finite word over $\Gamma$. In a two-pass scan, the first
pass reads $w$ from left to right, updates the state by $\delta$, and leaves each symbol
unchanged; the second pass reads the resulting word in reverse and also leaves every
symbol unchanged. Starting the first pass at $a_0$ and the second at its final state,
the overall result is the state obtained by folding $\delta$ from left to right over
$w$, together with the unchanged word $w$.
\end{theorem}

\begin{theorem}[Map scan output]
\lean{Turing.scanResult_map}
\label{Turing.scanResult_map}
\leanok
\uses{Turing.mealyRun_induction, Turing.scanKeep, Turing.scanMap, Turing.scanResult, Turing.scanResult_keep}
\difficulty{4}
Let $A$ and $\Gamma$ be sets, let $f:\Gamma\to\Gamma$, let $a_0\in A$, and let $w$ be a
finite word over $\Gamma$. In a two-pass scan, the first pass maps each symbol by $f$
and leaves the state unchanged, while the second pass leaves both state and symbols
unchanged. The overall result is the unchanged state $a_0$ and the word obtained by
applying $f$ to every symbol of $w$.
\end{theorem}

\begin{theorem}[Rightward scan transformation]
\lean{Turing.scanResult_shiftRight}
\label{Turing.scanResult_shiftRight}
\leanok
\uses{Turing.mealyRun, Turing.scanKeep, Turing.scanResult, Turing.scanResult_keep, Turing.scanShift}
\difficulty{4}
Let $B$ and $\Gamma$ be types, let $g:\Gamma\to B$ and $s:\Gamma\to B\to\Gamma$ be
functions, let $b_0\in B$, and let $w$ be a finite word over $\Gamma$. A scan that, at
each symbol, replaces it by the previous state and sets the state to the symbol’s
$g$-value, followed by a scan that leaves the state and symbols unchanged, returns the
state $(g(w))_{\mathrm{last}}$ if $w$ is nonempty and $b_0$ otherwise, together with the
word obtained by pairing each symbol of $w$ with its corresponding entry in
$b_0::(g(w))$ and applying $s$ to each pair.
\end{theorem}

\begin{theorem}[Left-shift scan output]
\lean{Turing.scanResult_shiftLeft}
\label{Turing.scanResult_shiftLeft}
\leanok
\uses{Turing.mealyRun_keep, Turing.scanKeep, Turing.scanResult, Turing.scanResult_keep, Turing.scanResult_shiftRight, Turing.scanShift, Turing.shiftLeft_list}
\difficulty{3}
Let $B$ and $\Gamma$ be sets, let $\operatorname{get}:\Gamma\to B$ and
$\operatorname{set}:\Gamma\to B\to\Gamma$, let $b_0\in B$, and let $w$ be a finite word
over $\Gamma$. A two-pass scan first leaves the state and symbols unchanged from left to
right, then scans from right to left, replacing each symbol $\gamma$ by
$\operatorname{set}(\gamma,b)$ and updating the state to $\operatorname{get}(\gamma)$.
The final state is the first value of $\operatorname{get}$ along $w$, or $b_0$ if $w$ is
empty; the output replaces each symbol of $w$ by setting its stored value to the value
of $\operatorname{get}$ at the next symbol, using $b_0$ for the last symbol.
\end{theorem}

\begin{theorem}[Carry scan output]
\lean{Turing.scanResult_carry}
\label{Turing.scanResult_carry}
\leanok
\uses{Turing.mealyRun, Turing.mealyRun_carry_clear, Turing.scanCarry, Turing.scanKeep, Turing.scanResult, Turing.scanResult_keep}
\difficulty{4}
Let $\Gamma$ be a set, let
$\operatorname{top}:\Gamma\to\{\mathrm{true},\mathrm{false}\}$ and
$\operatorname{next}:\Gamma\to\Gamma$, and let $w$ be a finite word over $\Gamma$.
Consider a two-pass scan that starts with a true carry, scans left to right changing a
symbol to $\operatorname{next}(\gamma)$ while the carry is true and setting the carry to
$\operatorname{top}(\gamma)$, then scans the resulting word right to left without
changing its symbols. Its final state is the conjunction of $\operatorname{top}$ over
$w$, and its output word consists of the symbols in the initial consecutive
$\operatorname{top}$-true prefix changed by $\operatorname{next}$, followed by the first
symbol after that prefix changed by $\operatorname{next}$ if one exists, followed by the
remaining symbols unchanged.
\end{theorem}

%%% added by the blueprint pipeline
\section{Additional declarations}

\begin{definition}[Identity scan transition]
\lean{Turing.scanKeep}
\label{Turing.scanKeep}
\leanok
For any sets $A$ and $\Gamma$, define the transition that takes a state $a\in A$ and
symbol $\gamma\in\Gamma$ to the unchanged pair $(a,\gamma)$.
\end{definition}

\begin{definition}[Scan fold]
\lean{Turing.scanFold}
\label{Turing.scanFold}
\leanok
Given types $A$ and $\Gamma$ and a function $\delta : A \times \Gamma \to A$, define the
scan fold to send each pair $(a,\gamma)$ to $(\delta(a,\gamma),\gamma)$; it updates the
state using the current symbol and returns that symbol unchanged.
\end{definition}

\begin{definition}[Scan map]
\lean{Turing.scanMap}
\label{Turing.scanMap}
\leanok
Given types $A$ and $\Gamma$ and a function $f : \Gamma \to \Gamma$, define a map from
$A \times \Gamma$ to $A \times \Gamma$ by sending each pair $(a,\gamma)$ to
$(a,f(\gamma))$.
\end{definition}

\begin{definition}[Read-and-write scan shift]
\lean{Turing.scanShift}
\label{Turing.scanShift}
\leanok
For types $B$ and $\Gamma$, a read function $g : \Gamma \to B$, and a write function $u
: \Gamma \to B \to \Gamma$, define a map taking $(b,\gamma)$ to
$(g(\gamma),u(\gamma,b))$.
\end{definition}
